\documentclass[type=report,theme=light,language=de]{semio}

%region 🔖️Apparatus
\ExplSyntaxOn
% 🌍 Every user-visible string of this reference in both document languages; there is no default language.
\NewExpandableDocumentCommand \ApiText { m m }
  { \str_case:Vn \l_semio_language_tl { { en } {#1} { de } {#2} } }
\ExplSyntaxOff

% 🔤 A control sequence name, typeset with its backslash.
\NewDocumentCommand \Cs { m } { \texttt{\textbackslash #1} }
% 🔑 A key, an option value or a column name.
\NewDocumentCommand \Key { m } { \texttt{#1} }
% 🧩 A placeholder inside a signature.
\NewDocumentCommand \meta { m } { \texttt{<#1>} }
% 🙂 One emoji of a repository path: the monospaced face has no emoji coverage, so the glyph is taken
% from the emoji face instead of being dropped silently.
\NewDocumentCommand \E { m } { {\SemioEmoji #1} }
% ✍️ The signature line of a public command or environment.
\NewDocumentCommand \ApiSig { m } { \par\medskip\noindent\texttt{#1}\par\medskip }
% 📦 The package a command lives in.
\NewDocumentCommand \ApiPkg { m } { \par\noindent\textbf{\ApiText{Package}{Paket}:}~\texttt{#1}\par }
% 📐 The header cells of a key table; \SemioTableLong splits its header argument on & before any
% expansion, so the cells are written out at every call site instead of hidden behind one macro.
\NewDocumentCommand \ApiKey { } { \ApiText{Key}{Schlüssel} }
\NewDocumentCommand \ApiType { } { \ApiText{Type}{Typ} }
\NewDocumentCommand \ApiDefault { } { \ApiText{Default}{Vorgabe} }
\NewDocumentCommand \ApiMeaning { } { \ApiText{Meaning}{Bedeutung} }
\NewDocumentCommand \ApiValue { } { \ApiText{Value}{Wert} }
%endregion 🔖️Apparatus

\title{\ApiText{The semio visualization API}{Die semio-Visualisierungs-API}}
\subtitle{\ApiText{Reference of every public command, key and default}{Referenz aller öffentlichen Befehle, Schlüssel und Vorgaben}}
\author{Semio}
\date{\today}

\begin{document}
\maketitle
\tableofcontents

\chapter{\ApiText{Reading this reference}{Diese Referenz lesen}}

\ApiText{%
  The semio visualization library is a grammar of graphics for print. Every chart in the catalogue is
  a preset over the same small set of primitives: a table of data, a chain of transforms, a set of
  scales, a coordinate system, marks, guides and a theme. This reference documents that grammar
  command by command, in the order of taxonomy section 79, and then one chapter per package namespace
  of taxonomy section 76.%
}{%
  Die semio-Visualisierungsbibliothek ist eine Grammatik der Grafik für den Druck. Jedes Diagramm des
  Katalogs ist eine Voreinstellung über derselben kleinen Menge von Grundbausteinen: eine Datentabelle,
  eine Kette von Transformationen, eine Menge von Skalen, ein Koordinatensystem, Marken, Hilfslinien
  und ein Thema. Diese Referenz dokumentiert diese Grammatik Befehl für Befehl, in der Reihenfolge von
  Taxonomieabschnitt 79, und danach mit je einem Kapitel pro Paketnamensraum aus Taxonomieabschnitt 76.%
}

\section{\ApiText{Conventions}{Konventionen}}

\ApiText{%
  Every public command is named \Cs{SemioViz\meta{Thing}}; every internal one is
  \Cs{semio\_viz\_\meta{module}\_\meta{operation}} in expl3. Options are \Key{l3keys} in the key
  family \Key{semio / viz / \meta{module}}, so every option has a documented name, type and default,
  and no renderer hard-codes a count, a radius or a colour.%
}{%
  Jeder öffentliche Befehl heißt \Cs{SemioViz\meta{Sache}}; jeder interne heißt
  \Cs{semio\_viz\_\meta{Modul}\_\meta{Operation}} in expl3. Optionen sind \Key{l3keys} in der
  Schlüsselfamilie \Key{semio / viz / \meta{Modul}}, jede Option hat also einen dokumentierten Namen,
  Typ und Vorgabewert, und kein Renderer verdrahtet Anzahl, Radius oder Farbe fest.%
}

\ApiText{%
  Coordinates are plain millimetre numbers without units, measured inside the current figure frame
  from its lower-left corner. Colours are never literal: a renderer asks the theme for a role or a
  palette slot, so the same figure is correct in the light and the dark theme. All user-visible text
  exists in English and German and is selected by the document language; this document is written
  that way itself.%
}{%
  Koordinaten sind schlichte Millimeterzahlen ohne Einheit, gemessen in dem aktuellen Figurenrahmen
  ab dessen unterer linker Ecke. Farben sind nie wörtlich: Ein Renderer fragt beim Thema eine Rolle
  oder einen Palettenplatz an, dieselbe Abbildung stimmt also im hellen wie im dunklen Thema. Alle
  sichtbaren Texte gibt es auf Englisch und Deutsch und werden über die Dokumentsprache gewählt;
  dieses Dokument selbst ist so geschrieben.%
}

\section{\ApiText{The frame: \Cs{begin}\{VizFigure\}}{Der Rahmen: \Cs{begin}\{VizFigure\}}}

\ApiSig{\Cs{begin}\{VizFigure\}[width=, height=, title=]\ \meta{recipe}\ \Cs{end}\{VizFigure\}}
\ApiPkg{semio-viz}

\ApiText{%
  Nothing draws outside a frame. \Key{VizFigure} is a figure window that opens a millimetre TikZ
  canvas; \Key{width} and \Key{height} are millimetres and default to 80 by 40. Every example in this
  reference is a real \Key{VizFigure} compiled from the source printed beside it.%
}{%
  Außerhalb eines Rahmens wird nichts gezeichnet. \Key{VizFigure} ist ein Abbildungsfenster, das eine
  Millimeter-TikZ-Leinwand öffnet; \Key{width} und \Key{height} sind Millimeter und haben die
  Vorgaben 80 und 40. Jedes Beispiel dieser Referenz ist eine echte, aus dem daneben abgedruckten
  Quelltext übersetzte \Key{VizFigure}.%
}

\begin{VizFigure}[title={\ApiText{Marks in a frame}{Marken in einem Rahmen}}, width=80, height=40]
\SemioVizMark{dot}[size=6, fill=semio-primary]{20, 14}
\SemioVizMark{square}[size=5, fill=semio-secondary]{40, 22}
\SemioVizMark{triangle}[size=6, fill=semio-tertiary]{60, 30}
\SemioVizPath{polyline}[points={8,8; 28,26; 52,12; 74,30}, stroke=semio-info]
\end{VizFigure}

\chapter{\ApiText{Data}{Daten}}

\ApiText{%
  Data is declared, never loaded implicitly. A table is a named list of columns and rows; the
  remaining commands mark an existing table as a hierarchy, a graph, a matrix or a geometry, which
  binds the column roles the layout algorithms of that shape expect.%
}{%
  Daten werden deklariert, nie implizit geladen. Eine Tabelle ist eine benannte Liste von Spalten und
  Zeilen; die übrigen Befehle markieren eine vorhandene Tabelle als Hierarchie, Graph, Matrix oder
  Geometrie und binden damit die Spaltenrollen, die die Layoutalgorithmen dieser Gestalt erwarten.%
}

\ApiPkg{semio-viz-data}

\section{\Cs{SemioVizTable}, \Cs{SemioVizRow}}

\ApiSig{\Cs{SemioVizTable}\{\meta{name}\}\{\meta{columns}\}}
\ApiSig{\Cs{SemioVizRow}\{\meta{name}\}\{\meta{values}\}}

\ApiText{%
  \Cs{SemioVizTable} declares a named table and its comma-separated column list; \Cs{SemioVizRow}
  appends one comma-separated row. Column names are the identifiers every encoding, transform and
  family option refers to.%
}{%
  \Cs{SemioVizTable} deklariert eine benannte Tabelle und ihre kommaseparierte Spaltenliste;
  \Cs{SemioVizRow} hängt eine kommaseparierte Zeile an. Spaltennamen sind die Bezeichner, auf die
  sich jede Kodierung, jede Transformation und jede Familienoption bezieht.%
}

\begin{VizFigure}[title={\ApiText{A table drawn as bars}{Eine als Balken gezeichnete Tabelle}}, width=80, height=40]
\SemioVizTable{api}{cat, val}
\SemioVizRow{api}{A, 4}
\SemioVizRow{api}{B, 7}
\SemioVizRow{api}{C, 5}
\SemioVizRow{api}{D, 9}
\SemioVizPlot[data=api, mark=bar, x=cat, y=val, guide={axis=x, axis=y}]
\end{VizFigure}

\section{\Cs{SemioVizTableFromCSV}}

\ApiSig{\Cs{SemioVizTableFromCSV}\{\meta{name}\}\{\meta{file}\}[delimiter=, quote=, header=, columns=]}

\ApiText{%
  Reads an RFC~4180 CSV file from disk. A quoted field may contain the delimiter; a doubled quote is
  one literal quote. With \Key{header=false} the \Key{columns} key supplies the column names.%
}{%
  Liest eine CSV-Datei nach RFC~4180 von der Festplatte. Ein gequotetes Feld darf das Trennzeichen
  enthalten; ein doppeltes Anführungszeichen ist ein literales. Bei \Key{header=false} liefert der
  Schlüssel \Key{columns} die Spaltennamen.%
}

\section{\ApiText{Shaped data}{Geformte Daten}}

\ApiSig{\Cs{SemioVizHierarchy}\{\meta{name}\}[id=, parent=, value=, label=]}
\ApiSig{\Cs{SemioVizGraph}\{\meta{name}\}[source=, target=, weight=, nodes=, edges=]}
\ApiSig{\Cs{SemioVizMatrix}\{\meta{name}\}[row=, col=, value=]}
\ApiSig{\Cs{SemioVizGeometry}\{\meta{name}\}[ring=, lon=, lat=, label=]}
\ApiSig{\Cs{SemioVizFunction}\{\meta{name}\}\{\meta{expression}\}[domain=, samples=]}

\ApiText{%
  Each of these marks an existing table with the role its columns play, so a layout algorithm can
  read it without being told again. \Cs{SemioVizFunction} is the exception: it samples an l3fp
  expression in the variable \Key{x} over \Key{domain} into a fresh table of \Key{samples} rows.%
}{%
  Jeder dieser Befehle markiert eine vorhandene Tabelle mit der Rolle, die ihre Spalten spielen, damit
  ein Layoutalgorithmus sie lesen kann, ohne erneut informiert zu werden. \Cs{SemioVizFunction} ist die
  Ausnahme: Er tastet einen l3fp-Ausdruck in der Variablen \Key{x} über \Key{domain} in eine neue
  Tabelle mit \Key{samples} Zeilen ab.%
}

\begin{VizFigure}[title={\ApiText{A sampled function}{Eine abgetastete Funktion}}, width=80, height=40]
\SemioVizFunction{apisin}{sin(x)}[domain={0, 6.2832}, samples=40]
\SemioVizPlot[data=apisin, mark=line, x=x, y=y, guide={axis=x, axis=y}]
\end{VizFigure}

\section{\ApiText{Demo tables}{Demotabellen}}

\ApiText{%
  The library ships named demo tables so that every chart kind and every example renders without a
  data file. They all begin with \Key{demo-}; the catalogue names the table each chart kind reads.%
}{%
  Die Bibliothek liefert benannte Demotabellen mit, damit jede Diagrammart und jedes Beispiel ohne
  Datendatei zeichnet. Alle beginnen mit \Key{demo-}; der Katalog nennt die Tabelle, die jede
  Diagrammart liest.%
}

\SemioTableLong[text-size=8pt]{\ApiText{Core demo tables}{Kern-Demotabellen}}{0.26,0.74}{\ApiText{Table}{Tabelle} & \ApiText{Content}{Inhalt}}{%
  \SemioTableRow{\Key{demo} & \ApiText{Five categorical rows with two measures and a grouping column.}{Fünf kategoriale Zeilen mit zwei Messgrößen und einer Gruppierungsspalte.}}
  \SemioTableRow{\Key{demo-series} & \ApiText{Three numeric series over a shared ordinal index.}{Drei numerische Reihen über einem gemeinsamen Ordinalindex.}}
  \SemioTableRow{\Key{demo-time} & \ApiText{A daily time series in open-high-low-close-volume shape.}{Eine tägliche Zeitreihe in Eröffnung-Hoch-Tief-Schluss-Volumen-Form.}}
  \SemioTableRow{\Key{demo-distribution} & \ApiText{Repeated observations per group for distribution renderers.}{Wiederholte Beobachtungen je Gruppe für Verteilungsdarstellungen.}}
  \SemioTableRow{\Key{demo-hierarchy} & \ApiText{A three-level hierarchy with leaf weights.}{Eine dreistufige Hierarchie mit Blattgewichten.}}
  \SemioTableRow{\Key{demo-graph} & \ApiText{A small directed weighted edge list.}{Eine kleine gerichtete, gewichtete Kantenliste.}}
  \SemioTableRow{\Key{demo-matrix} & \ApiText{A dense row/column/value matrix.}{Eine dicht besetzte Zeilen-Spalten-Wert-Matrix.}}
  \SemioTableRow{\Key{demo-geo} & \ApiText{Ring and point geometry with one measure per region.}{Ring- und Punktgeometrie mit einer Messgröße je Region.}}
  \SemioTableRow{\Key{demo-points} & \ApiText{Twelve scattered points, the input of triangulation, hull and voronoi layouts.}{Zwölf gestreute Punkte, die Eingabe von Triangulations-, Hüllen- und Voronoi-Layouts.}}
  \SemioTableRow{\Key{demo-grid} & \ApiText{A six by five scalar grid with a smooth bump, for contours and heat surfaces.}{Ein Skalargitter aus sechs mal fünf Zellen mit einer glatten Erhebung, für Höhenlinien und Wärmeflächen.}}
  \SemioTableRow{\Key{demo-flow} & \ApiText{Staged flows between nodes for sankey-like renderers.}{Stufenweise Flüsse zwischen Knoten für Sankey-artige Darstellungen.}}
  \SemioTableRow{\Key{demo-parts} & \ApiText{Parts of one whole, summing to a hundred.}{Teile eines Ganzen, die sich zu hundert summieren.}}
}

\chapter{\ApiText{Transform}{Transformation}}

\ApiPkg{semio-viz-transform}
\ApiSig{\Cs{SemioVizTransform}\{\meta{out}\}\{\meta{in}\}[\meta{operations}]}

\ApiText{%
  One command derives a table from another. Several operations may be given at once; they are applied
  in one fixed, documented order, so the result never depends on the order the keys were written in:
  \Key{filter}, \Key{sort}, \Key{bin}, \Key{group}/\Key{rollup}, \Key{fold}, \Key{pivot}, \Key{join},
  \Key{window}, \Key{stack}, \Key{normalize}, \Key{quantile}, \Key{kde}, \Key{regression}. Writing
  \meta{out} equal to \meta{in} transforms a table in place.%
}{%
  Ein Befehl leitet eine Tabelle aus einer anderen ab. Es dürfen mehrere Operationen gleichzeitig
  angegeben werden; sie werden in einer festen, dokumentierten Reihenfolge angewandt, das Ergebnis
  hängt also nie von der Schreibreihenfolge der Schlüssel ab: \Key{filter}, \Key{sort}, \Key{bin},
  \Key{group}/\Key{rollup}, \Key{fold}, \Key{pivot}, \Key{join}, \Key{window}, \Key{stack},
  \Key{normalize}, \Key{quantile}, \Key{kde}, \Key{regression}. Ist \meta{out} gleich \meta{in}, wird
  die Tabelle an Ort und Stelle transformiert.%
}

\SemioTableLong[text-size=8pt]{\ApiText{Transform operations}{Transformationsoperationen}}{0.17,0.19,0.16,0.48}{\ApiKey & \ApiType & \ApiDefault & \ApiMeaning}{%
  \SemioTableRow{\Key{filter} & \Key{fp predicate} & --- & \ApiText{Keeps the rows for which the predicate over bare column names holds, e.g. \Key{val > 4 \&\& val2 < 6}.}{Behält die Zeilen, für die das Prädikat über blanken Spaltennamen gilt, z.\,B. \Key{val > 4 \&\& val2 < 6}.}}
  \SemioTableRow{\Key{sort} & \Key{column} & --- & \ApiText{Sorts by that column.}{Sortiert nach dieser Spalte.}}
  \SemioTableRow{\Key{descending} & \Key{boolean} & \Key{false} & \ApiText{Reverses the sort order.}{Kehrt die Sortierreihenfolge um.}}
  \SemioTableRow{\Key{group} & \Key{column} & --- & \ApiText{The column whose distinct values become the groups.}{Die Spalte, deren verschiedene Werte die Gruppen bilden.}}
  \SemioTableRow{\Key{rollup} & \Key{agg:column} & --- & \ApiText{Aggregates each group; \Key{agg} is \Key{sum}, \Key{mean}, \Key{min}, \Key{max}, \Key{count} or \Key{median}.}{Aggregiert jede Gruppe; \Key{agg} ist \Key{sum}, \Key{mean}, \Key{min}, \Key{max}, \Key{count} oder \Key{median}.}}
  \SemioTableRow{\Key{pivot} & \Key{key:value} & --- & \ApiText{Turns the distinct key values into columns; the group column stays as the row key.}{Macht die verschiedenen Schlüsselwerte zu Spalten; die Gruppenspalte bleibt der Zeilenschlüssel.}}
  \SemioTableRow{\Key{fold} & \Key{id:c1;c2;…} & --- & \ApiText{Long form: the named columns become the columns \Key{key} and \Key{value}.}{Langform: Die genannten Spalten werden zu den Spalten \Key{key} und \Key{value}.}}
  \SemioTableRow{\Key{join} & \Key{table:left:right} & --- & \ApiText{Inner join on equal keys.}{Innerer Verbund über gleiche Schlüssel.}}
  \SemioTableRow{\Key{window} & \Key{kind:column:size} & --- & \ApiText{\Key{movingAverage}, \Key{cumulativeSum}, \Key{rank} or \Key{lag} over a sliding window.}{\Key{movingAverage}, \Key{cumulativeSum}, \Key{rank} oder \Key{lag} über ein gleitendes Fenster.}}
  \SemioTableRow{\Key{bin} & \Key{column} & --- & \ApiText{Bins the column into the columns \Key{x0}, \Key{x1}, \Key{count}.}{Klassiert die Spalte in die Spalten \Key{x0}, \Key{x1}, \Key{count}.}}
  \SemioTableRow{\Key{thresholds} & \Key{rule} & \Key{sturges} & \ApiText{\Key{sturges}, \Key{count:n}, \Key{width:w} or an explicit \Key{a;b;c} list.}{\Key{sturges}, \Key{count:n}, \Key{width:w} oder eine explizite Liste \Key{a;b;c}.}}
  \SemioTableRow{\Key{stack} & \Key{x:series:value} & --- & \ApiText{Stacks into the columns \Key{x}, \Key{series}, \Key{y0}, \Key{y1}.}{Stapelt in die Spalten \Key{x}, \Key{series}, \Key{y0}, \Key{y1}.}}
  \SemioTableRow{\Key{order} & \Key{name} & \Key{none} & \ApiText{Stack order: \Key{none}, \Key{ascending}, \Key{descending}, \Key{inside-out}, \Key{reverse}, \Key{appearance}.}{Stapelreihenfolge: \Key{none}, \Key{ascending}, \Key{descending}, \Key{inside-out}, \Key{reverse}, \Key{appearance}.}}
  \SemioTableRow{\Key{offset} & \Key{name} & \Key{none} & \ApiText{Stack offset: \Key{none}, \Key{expand}, \Key{diverging}, \Key{silhouette}, \Key{wiggle}.}{Stapelversatz: \Key{none}, \Key{expand}, \Key{diverging}, \Key{silhouette}, \Key{wiggle}.}}
  \SemioTableRow{\Key{normalize} & \Key{column} & --- & \ApiText{Divides the column by the sum of its values.}{Teilt die Spalte durch die Summe ihrer Werte.}}
  \SemioTableRow{\Key{quantile} & \Key{column} & --- & \ApiText{Reduces the column to the columns \Key{p} and \Key{value} (R-7 quantiles).}{Reduziert die Spalte auf die Spalten \Key{p} und \Key{value} (R-7-Quantile).}}
  \SemioTableRow{\Key{probabilities} & \Key{clist} & \Key{0,0.25,0.5,0.75,1} & \ApiText{The probabilities the quantile operation evaluates.}{Die Wahrscheinlichkeiten, die die Quantiloperation auswertet.}}
  \SemioTableRow{\Key{kde} & \Key{column} & --- & \ApiText{Kernel density estimate into the columns \Key{x} and \Key{density}.}{Kerndichteschätzung in die Spalten \Key{x} und \Key{density}.}}
  \SemioTableRow{\Key{bandwidth} & \Key{number} & \Key{0} & \ApiText{Kernel bandwidth; \Key{0} selects Silverman's rule of thumb.}{Kernbandbreite; \Key{0} wählt Silvermans Faustregel.}}
  \SemioTableRow{\Key{kernel} & \Key{name} & \Key{gaussian} & \ApiText{\Key{gaussian} or \Key{epanechnikov}.}{\Key{gaussian} oder \Key{epanechnikov}.}}
  \SemioTableRow{\Key{samples} & \Key{integer} & \Key{50} & \ApiText{Sample count of the kde and function grids.}{Stützstellenzahl der Dichte- und Funktionsgitter.}}
  \SemioTableRow{\Key{regression} & \Key{kind:x:y} & --- & \ApiText{Fits \Key{linear}, \Key{poly}, \Key{exp}, \Key{log} or \Key{pow}.}{Passt \Key{linear}, \Key{poly}, \Key{exp}, \Key{log} oder \Key{pow} an.}}
  \SemioTableRow{\Key{degree} & \Key{integer} & \Key{2} & \ApiText{Polynomial degree of \Key{regression=poly}.}{Polynomgrad von \Key{regression=poly}.}}
}

\begin{VizFigure}[title={\ApiText{Filter and rollup}{Filtern und Aggregieren}}, width=80, height=40]
\SemioVizTransform{apiroll}{demo}[group=grp, rollup=sum:val]
\SemioVizPlot[data=apiroll, mark=bar, x=grp, y=value, guide={axis=x, axis=y}]
\end{VizFigure}

\chapter{\ApiText{Scale}{Skala}}

\ApiPkg{semio-viz-scale}
\ApiSig{\Cs{SemioVizScale}\{\meta{name}\}\{\meta{kind}\}\{\meta{domain}\}\{\meta{range}\}[\meta{options}]}

\ApiText{%
  A scale is a named mapping from a data domain to a drawing range, with the accessors and the tick
  logic of d3-scale. Fifteen kinds are registered; the guides, the plot engine and every family look
  a scale up by name, so an axis always describes exactly the mapping the marks were placed with.%
}{%
  Eine Skala ist eine benannte Abbildung von einem Datenbereich auf einen Zeichenbereich, mit den
  Zugriffsfunktionen und der Teilstrichlogik von d3-scale. Fünfzehn Arten sind registriert; die
  Hilfslinien, die Plot-Maschine und jede Familie schlagen eine Skala über ihren Namen nach, eine
  Achse beschreibt also immer genau die Abbildung, mit der die Marken gesetzt wurden.%
}

\SemioTableLong[text-size=8pt]{\ApiText{Scale kinds}{Skalenarten}}{0.20,0.80}{\ApiValue & \ApiMeaning}{%
  \SemioTableRow{\Key{linear} & \ApiText{Affine mapping of a continuous domain.}{Affine Abbildung eines stetigen Bereichs.}}
  \SemioTableRow{\Key{log} & \ApiText{Logarithmic, base \Key{base}; the domain must not contain zero.}{Logarithmisch zur Basis \Key{base}; der Bereich darf die Null nicht enthalten.}}
  \SemioTableRow{\Key{pow} & \ApiText{Power mapping with exponent \Key{exponent}.}{Potenzabbildung mit dem Exponenten \Key{exponent}.}}
  \SemioTableRow{\Key{sqrt} & \ApiText{Power mapping with the exponent fixed at 0.5, the correct scale for areas.}{Potenzabbildung mit dem festen Exponenten 0,5, die richtige Skala für Flächen.}}
  \SemioTableRow{\Key{symlog} & \ApiText{Symmetric logarithm with a linear region of width \Key{constant} around zero.}{Symmetrischer Logarithmus mit einem linearen Bereich der Breite \Key{constant} um die Null.}}
  \SemioTableRow{\Key{identity} & \ApiText{The value itself; useful when data already is in millimetres.}{Der Wert selbst; nützlich, wenn die Daten bereits in Millimetern vorliegen.}}
  \SemioTableRow{\Key{ordinal} & \ApiText{Discrete domain onto a discrete range, cycling the range.}{Diskreter Bereich auf einen diskreten Wertevorrat, der zyklisch durchlaufen wird.}}
  \SemioTableRow{\Key{band} & \ApiText{Discrete domain onto contiguous bands, with inner and outer padding and an alignment.}{Diskreter Bereich auf zusammenhängende Bänder, mit innerem und äußerem Abstand und einer Ausrichtung.}}
  \SemioTableRow{\Key{point} & \ApiText{Discrete domain onto points, a band scale of zero width.}{Diskreter Bereich auf Punkte, eine Bandskala der Breite null.}}
  \SemioTableRow{\Key{quantile} & \ApiText{Sample values split into as many quantile classes as the range has entries.}{Stichprobenwerte, aufgeteilt in so viele Quantilklassen, wie der Wertevorrat Einträge hat.}}
  \SemioTableRow{\Key{quantize} & \ApiText{Continuous domain cut into equal intervals, one range entry each.}{Stetiger Bereich, in gleich große Intervalle geschnitten, je ein Eintrag des Wertevorrats.}}
  \SemioTableRow{\Key{threshold} & \ApiText{Explicit cut points; the range has one entry more than the domain.}{Explizite Schnittstellen; der Wertevorrat hat einen Eintrag mehr als der Bereich.}}
  \SemioTableRow{\Key{sequential} & \ApiText{Continuous domain onto a colour ramp named by \Key{scheme}.}{Stetiger Bereich auf eine Farbrampe, die \Key{scheme} benennt.}}
  \SemioTableRow{\Key{diverging} & \ApiText{Three-point domain onto a ramp with a neutral midpoint.}{Dreipunktbereich auf eine Rampe mit neutraler Mitte.}}
  \SemioTableRow{\Key{temporal} & \ApiText{Timestamps onto a range, with calendar ticks selected by \Key{interval}.}{Zeitstempel auf einen Wertebereich, mit Kalenderteilstrichen nach \Key{interval}.}}
}

\SemioTableLong[text-size=8pt]{\ApiText{Scale options}{Skalenoptionen}}{0.17,0.16,0.13,0.54}{\ApiKey & \ApiType & \ApiDefault & \ApiMeaning}{%
  \SemioTableRow{\Key{nice} & \Key{integer|true} & \Key{false} & \ApiText{Extends the domain to round tick values.}{Erweitert den Bereich auf runde Teilstrichwerte.}}
  \SemioTableRow{\Key{clamp} & \Key{boolean} & \Key{false} & \ApiText{Clamps inputs to the domain extent.}{Beschneidet Eingaben auf die Bereichsgrenzen.}}
  \SemioTableRow{\Key{padding} & \Key{number} & \Key{0} & \ApiText{Point or band outer padding; on a band scale it sets both paddings.}{Äußerer Abstand bei Punkt- und Bandskalen; bei einer Bandskala setzt er beide Abstände.}}
  \SemioTableRow{\Key{paddingInner} & \Key{number} & \Key{0} & \ApiText{Band inner padding, as a fraction of the step.}{Innerer Bandabstand als Anteil der Schrittweite.}}
  \SemioTableRow{\Key{paddingOuter} & \Key{number} & \Key{0} & \ApiText{Band outer padding, as a fraction of the step.}{Äußerer Bandabstand als Anteil der Schrittweite.}}
  \SemioTableRow{\Key{align} & \Key{number} & \Key{0.5} & \ApiText{How the leftover space of a band or point scale is distributed.}{Wie der Restplatz einer Band- oder Punktskala verteilt wird.}}
  \SemioTableRow{\Key{base} & \Key{number} & \Key{10} & \ApiText{Logarithm base, also used by the tick generator.}{Logarithmusbasis, die auch der Teilstrichgenerator verwendet.}}
  \SemioTableRow{\Key{exponent} & \Key{number} & \Key{1} & \ApiText{Exponent of a power scale; \Key{sqrt} fixes it at 0.5.}{Exponent einer Potenzskala; \Key{sqrt} setzt ihn fest auf 0,5.}}
  \SemioTableRow{\Key{constant} & \Key{number} & \Key{1} & \ApiText{Width of the linear region of a symlog scale.}{Breite des linearen Bereichs einer Symlog-Skala.}}
  \SemioTableRow{\Key{ticks} & \Key{integer} & \Key{10} & \ApiText{Target tick count.}{Angestrebte Anzahl Teilstriche.}}
  \SemioTableRow{\Key{tickValues} & \Key{clist} & --- & \ApiText{Explicit tick values, overriding the generator.}{Explizite Teilstrichwerte, die den Generator übergehen.}}
  \SemioTableRow{\Key{tickFormat} & \Key{spec} & --- & \ApiText{A \Cs{SemioVizFormat} specifier for the tick labels.}{Eine Formatangabe für \Cs{SemioVizFormat} für die Teilstrichbeschriftung.}}
  \SemioTableRow{\Key{interpolator} & \Key{name} & \Key{oklab} & \ApiText{Colour space a sequential or diverging ramp is sampled in.}{Farbraum, in dem eine sequentielle oder divergierende Rampe abgetastet wird.}}
  \SemioTableRow{\Key{scheme} & \Key{name} & --- & \ApiText{Named colour scheme, resolved by \Key{semio-viz-theme}.}{Benanntes Farbschema, das \Key{semio-viz-theme} auflöst.}}
  \SemioTableRow{\Key{reverse} & \Key{boolean} & \Key{false} & \ApiText{Reverses the range.}{Kehrt den Wertebereich um.}}
  \SemioTableRow{\Key{round} & \Key{boolean} & \Key{false} & \ApiText{Rounds band positions to whole millimetres.}{Rundet Bandpositionen auf ganze Millimeter.}}
  \SemioTableRow{\Key{interval} & \Key{name} & \Key{auto} & \ApiText{Temporal tick interval: \Key{auto}, \Key{day}, \Key{week}, \Key{month} or \Key{year}.}{Zeitliches Teilstrichintervall: \Key{auto}, \Key{day}, \Key{week}, \Key{month} oder \Key{year}.}}
}

\begin{VizFigure}[title={\ApiText{Three scales, three axes}{Drei Skalen, drei Achsen}}, width=80, height=44]
\SemioVizScale{apilin}{linear}{0, 100}{10, 70}[nice=5, ticks=5]
\SemioVizScale{apilog}{log}{1, 1000}{10, 70}[base=10, ticks=4]
\SemioVizScale{apisym}{symlog}{-100, 100}{10, 70}[constant=10, ticks=5]
\SemioVizAxis[scale=apilin, orient=bottom, at=8, title={linear}]
\SemioVizAxis[scale=apilog, orient=bottom, at=20, title={log}]
\SemioVizAxis[scale=apisym, orient=bottom, at=32, title={symlog}]
\end{VizFigure}

\ApiSig{\Cs{SemioVizInterpolate}\{\meta{kind}\}\{\meta{a}\}\{\meta{b}\}\{\meta{t}\}}

\ApiText{%
  The interpolators of the scale kernel, reachable without declaring a scale. \Key{number} is the
  affine one and \Key{round} the same value rounded to the nearest integer; \Key{rgb}, \Key{lab},
  \Key{hcl} and \Key{oklab} interpolate two \Key{RRGGBB} colours in that space and return a
  \Key{\#RRGGBB} one; \Key{array} walks two equally long comma lists of numbers element by element.
  Every kind is d3-interpolate's function of the same name, and the colour spaces are the ones the
  sequential and diverging ramps are sampled in.%
}{%
  Die Interpolatoren des Skalenkerns, erreichbar ohne eine Skala zu deklarieren. \Key{number} ist der
  affine, \Key{round} derselbe Wert auf die nächste ganze Zahl gerundet; \Key{rgb}, \Key{lab},
  \Key{hcl} und \Key{oklab} interpolieren zwei \Key{RRGGBB}-Farben in diesem Raum und liefern eine
  \Key{\#RRGGBB}-Farbe; \Key{array} durchläuft zwei gleich lange Zahlenlisten elementweise. Jede Art
  ist die gleichnamige Funktion aus d3-interpolate, und die Farbräume sind jene, in denen die
  sequentiellen und divergierenden Rampen abgetastet werden.%
}

\chapter{\ApiText{Format}{Formatierung}}

\ApiPkg{semio-viz-format}
\ApiSig{\Cs{SemioVizFormat}\{\meta{spec}\}\{\meta{value}\}}
\ApiSig{\Cs{SemioVizTimeFormat}\{\meta{spec}\}\{\meta{ISO timestamp}\}}
\ApiSig{\Cs{SemioVizLocale}\{en|de\}}

\ApiText{%
  Numbers are formatted with the d3-format grammar
  \Key{[[fill]align][sign][symbol][0][width][,][.precision][\textasciitilde][type]}, timestamps with
  the d3-time-format directives. The locale --- decimal separator, group separator, month and weekday
  names --- follows the document language; \Cs{SemioVizLocale} overrides it for a passage.%
}{%
  Zahlen werden mit der d3-format-Grammatik
  \Key{[[fill]align][sign][symbol][0][width][,][.precision][\textasciitilde][type]} formatiert,
  Zeitstempel mit den d3-time-format-Direktiven. Das Gebietsschema --- Dezimaltrenner,
  Gruppentrenner, Monats- und Wochentagsnamen --- folgt der Dokumentsprache;
  \Cs{SemioVizLocale} setzt es für einen Abschnitt außer Kraft.%
}

\SemioTableLong[text-size=8pt]{\ApiText{Number format types}{Zahlenformattypen}}{0.16,0.84}{\ApiValue & \ApiMeaning}{%
  \SemioTableRow{\Key{f} & \ApiText{Fixed point with the given number of decimals.}{Festkomma mit der angegebenen Anzahl Nachkommastellen.}}
  \SemioTableRow{\Key{e} & \ApiText{Exponential notation.}{Exponentialschreibweise.}}
  \SemioTableRow{\Key{g} & \ApiText{Significant digits, switching to exponential where that is shorter.}{Signifikante Stellen, mit Wechsel zur Exponentialform, wo diese kürzer ist.}}
  \SemioTableRow{\Key{r} & \ApiText{Rounded to significant digits, never exponential.}{Auf signifikante Stellen gerundet, nie exponentiell.}}
  \SemioTableRow{\Key{s} & \ApiText{SI prefix chosen automatically.}{Automatisch gewählter SI-Präfix.}}
  \SemioTableRow{\Key{p} & \ApiText{Percentage: multiplied by a hundred and rounded.}{Prozent: mit hundert multipliziert und gerundet.}}
  \SemioTableRow{\Key{d} & \ApiText{Rounded to an integer.}{Auf eine ganze Zahl gerundet.}}
  \SemioTableRow{\Key{b}, \Key{o}, \Key{x}, \Key{X} & \ApiText{Binary, octal, lower-case and upper-case hexadecimal.}{Binär, oktal, klein- und großgeschrieben hexadezimal.}}
}

\SemioTableLong[text-size=8pt]{\ApiText{Time format directives}{Zeitformat-Direktiven}}{0.20,0.80}{\ApiValue & \ApiMeaning}{%
  \SemioTableRow{\Key{\%Y}, \Key{\%y} & \ApiText{Four-digit and two-digit year.}{Vierstelliges und zweistelliges Jahr.}}
  \SemioTableRow{\Key{\%m}, \Key{\%d}, \Key{\%e} & \ApiText{Month, zero-padded day, space-padded day.}{Monat, nullgefüllter Tag, leerzeichengefüllter Tag.}}
  \SemioTableRow{\Key{\%H}, \Key{\%I}, \Key{\%M}, \Key{\%S} & \ApiText{Hour on 24 and on 12, minute, second.}{Stunde im 24- und im 12-Stunden-Zyklus, Minute, Sekunde.}}
  \SemioTableRow{\Key{\%p} & \ApiText{The AM/PM marker of the locale.}{Die Vormittags-/Nachmittagskennung des Gebietsschemas.}}
  \SemioTableRow{\Key{\%a}, \Key{\%A} & \ApiText{Abbreviated and full weekday name.}{Abgekürzter und ausgeschriebener Wochentagsname.}}
  \SemioTableRow{\Key{\%b}, \Key{\%B} & \ApiText{Abbreviated and full month name.}{Abgekürzter und ausgeschriebener Monatsname.}}
  \SemioTableRow{\Key{\%j} & \ApiText{Day of the year.}{Tag des Jahres.}}
  \SemioTableRow{\Key{\%U}, \Key{\%W} & \ApiText{Week of the year, counted from Sunday and from Monday.}{Woche des Jahres, ab Sonntag und ab Montag gezählt.}}
  \SemioTableRow{\Key{\%Z} & \ApiText{Time zone offset.}{Zeitzonenversatz.}}
}

\ApiText{%
  Worked examples, typeset by the library itself: \Key{,.1f} of 1234.56 is
  \SemioVizFormat{,.1f}{1234.56}; \Key{.2s} of 48000 is \SemioVizFormat{.2s}{48000};
  \Key{+.1p} of 0.128 is \SemioVizFormat{+.1p}{0.128}; \Key{\%d.~\%B~\%Y} of 2026-09-05 is
  \SemioVizTimeFormat{\%d.~\%B~\%Y}{2026-09-05}.%
}{%
  Ausgeführte Beispiele, von der Bibliothek selbst gesetzt: \Key{,.1f} von 1234.56 ergibt
  \SemioVizFormat{,.1f}{1234.56}; \Key{.2s} von 48000 ergibt \SemioVizFormat{.2s}{48000};
  \Key{+.1p} von 0.128 ergibt \SemioVizFormat{+.1p}{0.128}; \Key{\%d.~\%B~\%Y} von 2026-09-05 ergibt
  \SemioVizTimeFormat{\%d.~\%B~\%Y}{2026-09-05}.%
}

\chapter{\ApiText{Coordinate systems}{Koordinatensysteme}}

\ApiPkg{semio-viz-coordinate}
\ApiSig{\Cs{SemioVizCoordinate}\{\meta{kind}\}[\meta{options}]}
\ApiSig{\Cs{SemioVizProject}\{\meta{x}\}\{\meta{y}\} \ \Cs{SemioVizProjectedX} \ \Cs{SemioVizProjectedY}}

\ApiText{%
  A coordinate system maps the two data channels of a plot onto the frame. Installing one is the only
  thing a family has to do to be drawn in polar instead of cartesian space: it reads its two channels
  through the coordinate kernel and never learns which system it is in. \Cs{SemioVizProject} maps a
  single point and leaves the millimetres in the two expandable accessors.%
}{%
  Ein Koordinatensystem bildet die beiden Datenkanäle eines Diagramms auf den Rahmen ab. Ein System zu
  installieren ist das Einzige, was eine Familie tun muss, um statt kartesisch polar gezeichnet zu
  werden: Sie liest ihre beiden Kanäle über den Koordinatenkern und erfährt nie, in welchem System sie
  steckt. \Cs{SemioVizProject} bildet einen einzelnen Punkt ab und hinterlässt die Millimeter in den
  beiden expandierbaren Zugriffsbefehlen.%
}

\SemioTableLong[text-size=8pt]{\ApiText{Coordinate kinds}{Koordinatenarten}}{0.20,0.80}{\ApiValue & \ApiMeaning}{%
  \SemioTableRow{\Key{cartesian} & \ApiText{Horizontal and vertical channel, the default.}{Waagerechter und senkrechter Kanal, die Vorgabe.}}
  \SemioTableRow{\Key{polar} & \ApiText{Angle and radius; pies, coxcombs, radar and radial bars live here.}{Winkel und Radius; Torten-, Coxcomb-, Netz- und Radialbalkendiagramme leben hier.}}
  \SemioTableRow{\Key{ternary} & \ApiText{Three-part compositions on an equilateral triangle.}{Dreiteilige Zusammensetzungen auf einem gleichseitigen Dreieck.}}
  \SemioTableRow{\Key{barycentric} & \ApiText{Barycentric weights of the same triangle.}{Baryzentrische Gewichte desselben Dreiecks.}}
  \SemioTableRow{\Key{parallel} & \ApiText{Parallel axes; the x channel is the one-based axis index.}{Parallele Achsen; der x-Kanal ist der ab eins zählende Achsenindex.}}
  \SemioTableRow{\Key{logpolar} & \ApiText{Polar with a logarithmic radius of base \Key{base}.}{Polar mit logarithmischem Radius zur Basis \Key{base}.}}
  \SemioTableRow{\Key{geographic} & \ApiText{Degrees of longitude and latitude through a named projection.}{Grad geographischer Länge und Breite durch eine benannte Projektion.}}
}

\SemioTableLong[text-size=8pt]{\ApiText{Coordinate options}{Koordinatenoptionen}}{0.19,0.14,0.13,0.54}{\ApiKey & \ApiType & \ApiDefault & \ApiMeaning}{%
  \SemioTableRow{\Key{width}, \Key{height} & \Key{mm} & \Key{80}, \Key{40} & \ApiText{The frame the system maps into.}{Der Rahmen, in den das System abbildet.}}
  \SemioTableRow{\Key{originX}, \Key{originY} & \Key{mm} & \Key{0} & \ApiText{Lower-left corner of that frame.}{Untere linke Ecke dieses Rahmens.}}
  \SemioTableRow{\Key{flipY} & \Key{boolean} & \Key{false} & \ApiText{Maps the y channel downwards, for image-like data.}{Bildet den y-Kanal nach unten ab, für bildartige Daten.}}
  \SemioTableRow{\Key{startAngle} & \Key{rad} & \Key{0} & \ApiText{Polar angle of the first datum, clockwise from twelve o'clock.}{Polarwinkel des ersten Datums, im Uhrzeigersinn ab zwölf Uhr.}}
  \SemioTableRow{\Key{endAngle} & \Key{rad} & \Key{2*pi} & \ApiText{Polar angle one full turn later.}{Polarwinkel eine volle Umdrehung später.}}
  \SemioTableRow{\Key{innerRadius} & \Key{mm} & \Key{0} & \ApiText{Radius the domain minimum maps to.}{Radius, auf den das Bereichsminimum abgebildet wird.}}
  \SemioTableRow{\Key{outerRadius} & \Key{mm} & \Key{20} & \ApiText{Radius the domain maximum maps to.}{Radius, auf den das Bereichsmaximum abgebildet wird.}}
  \SemioTableRow{\Key{clockwise} & \Key{boolean} & \Key{true} & \ApiText{Direction the polar angle grows in.}{Richtung, in der der Polarwinkel wächst.}}
  \SemioTableRow{\Key{base} & \Key{number} & \Key{10} & \ApiText{Logarithm base of the logpolar radius.}{Logarithmusbasis des logpolaren Radius.}}
  \SemioTableRow{\Key{axes} & \Key{integer} & \Key{3} & \ApiText{Number of parallel axes.}{Anzahl paralleler Achsen.}}
  \SemioTableRow{\Key{size} & \Key{mm} & \Key{60} & \ApiText{Edge length of the ternary triangle.}{Kantenlänge des ternären Dreiecks.}}
  \SemioTableRow{\Key{domainX}, \Key{domainY} & \Key{clist} & \Key{0,1} & \ApiText{The data range each channel spans before mapping.}{Der Datenbereich, den jeder Kanal vor der Abbildung überspannt.}}
}

\begin{VizFigure}[title={\ApiText{The same bars in polar space}{Dieselben Balken im Polarraum}}, width=80, height=44]
\SemioVizPlot[data=demo, mark=bar, coordinate=polar, x=cat, y=val]
\end{VizFigure}

\chapter{\ApiText{Marks and paths}{Marken und Pfade}}

\ApiPkg{semio-viz-mark}
\ApiSig{\Cs{SemioVizMark}\{\meta{kind}\}[\meta{options}]\{\meta{x}, \meta{y}\}}
\ApiSig{\Cs{SemioVizPath}\{\meta{kind}\}[points=, curve=, closed=, \meta{options}]}
\ApiSig{\Cs{SemioVizText}\{\meta{kind}\}[\meta{options}]\{\meta{x}, \meta{y}\}\{\meta{content}\}}

\ApiText{%
  A mark is one primitive of taxonomy section 0: the smallest thing that carries ink. Every family
  and the plot engine draw through these three commands and nothing else, so a family's geometry is
  testable numerically without rasterising a page. Fifty-six mark kinds are registered.%
}{%
  Eine Marke ist ein Grundbaustein aus Taxonomieabschnitt 0: das kleinste Ding, das Farbe trägt. Jede
  Familie und die Plot-Maschine zeichnen über diese drei Befehle und über nichts sonst, die Geometrie
  einer Familie ist also numerisch prüfbar, ohne eine Seite zu rastern. Sechsundfünfzig Markenarten
  sind registriert.%
}

\SemioTableLong[text-size=8pt]{\ApiText{Mark kinds}{Markenarten}}{0.20,0.80}{\ApiText{Group}{Gruppe} & \ApiText{Kinds}{Arten}}{%
  \SemioTableRow{\ApiText{Points}{Punkte} & \Key{dot}, \Key{circle}, \Key{square}, \Key{rectangle}, \Key{triangle}, \Key{diamond}, \Key{cross}, \Key{plus}, \Key{star}, \Key{custom-glyph}, \Key{icon-mark}, \Key{image-mark}, \Key{text-mark}}
  \SemioTableRow{\ApiText{Lines}{Linien} & \Key{straight-line}, \Key{polyline}, \Key{step-line}, \Key{curved-line}, \Key{bezier-curve}, \Key{spline}, \Key{catmull-rom-spline}, \Key{basis-spline}, \Key{cardinal-spline}, \Key{monotone-spline}, \Key{closed-curve}}
  \SemioTableRow{\ApiText{Areas}{Flächen} & \Key{polygon}, \Key{filled-path}, \Key{ribbon}, \Key{band}, \Key{envelope}}
  \SemioTableRow{\ApiText{Arcs}{Bögen} & \Key{circular-arc}, \Key{elliptical-arc}, \Key{annular-arc}, \Key{sector}, \Key{wedge}}
  \SemioTableRow{\ApiText{Connectors}{Verbinder} & \Key{straight-connector}, \Key{orthogonal-connector}, \Key{curved-connector}, \Key{elbow-connector}, \Key{bundled-connector}, \Key{arrow}, \Key{bidirectional-arrow}}
  \SemioTableRow{\ApiText{Regions}{Regionen} & \Key{rectangular-region}, \Key{circular-region}, \Key{polygonal-region}, \Key{voronoi-region}, \Key{convex-hull}, \Key{concave-hull}}
  \SemioTableRow{\ApiText{Notation}{Notation} & \Key{label}, \Key{callout}, \Key{leader-line}, \Key{bracket}, \Key{brace}, \Key{highlight-region}, \Key{reference-line}, \Key{reference-band}, \Key{reference-point}}
}

\SemioTableLong[text-size=8pt]{\ApiText{Mark options}{Markenoptionen}}{0.22,0.14,0.16,0.48}{\ApiKey & \ApiType & \ApiDefault & \ApiMeaning}{%
  \SemioTableRow{\Key{size} & \Key{number} & \Key{6} & \ApiText{Mark area in square millimetres, with d3-symbol semantics.}{Markenfläche in Quadratmillimetern, mit der Semantik von d3-symbol.}}
  \SemioTableRow{\Key{fill} & \Key{colour} & \Key{semio-primary} & \ApiText{Fill colour, or \Key{none} for an unfilled mark.}{Füllfarbe oder \Key{none} für eine ungefüllte Marke.}}
  \SemioTableRow{\Key{stroke} & \Key{colour} & \Key{border-emphasized} & \ApiText{Stroke colour, or \Key{none}.}{Strichfarbe oder \Key{none}.}}
  \SemioTableRow{\Key{strokeWidth} & \Key{dimension} & \Key{stroke default} & \ApiText{Stroke width, taken from the design tokens.}{Strichstärke, aus den Gestaltungstoken.}}
  \SemioTableRow{\Key{opacity} & \Key{number} & \Key{1} & \ApiText{Draw opacity of the whole mark.}{Deckkraft der ganzen Marke.}}
  \SemioTableRow{\Key{rotate} & \Key{deg} & \Key{0} & \ApiText{Rotation about the mark's own point.}{Drehung um den eigenen Punkt der Marke.}}
  \SemioTableRow{\Key{pattern} & \Key{name} & \Key{none} & \ApiText{\Key{hatch}, \Key{crosshatch}, \Key{dots} or \Key{grid}: a grayscale-safe fill.}{\Key{hatch}, \Key{crosshatch}, \Key{dots} oder \Key{grid}: eine graustufensichere Füllung.}}
  \SemioTableRow{\Key{points} & \Key{list} & --- & \ApiText{\Key{x,y; x,y; …} in figure millimetres.}{\Key{x,y; x,y; …} in Figurenmillimetern.}}
  \SemioTableRow{\Key{points2} & \Key{list} & --- & \ApiText{The second boundary of a band, an envelope or a ribbon.}{Die zweite Begrenzung eines Bandes, einer Hüllkurve oder eines Zierbandes.}}
  \SemioTableRow{\Key{curve} & \Key{name} & \Key{linear} & \ApiText{Interpolator; see the curve table below.}{Interpolator; siehe die Kurventabelle weiter unten.}}
  \SemioTableRow{\Key{tension}, \Key{alpha}, \Key{beta} & \Key{number} & \Key{0}, \Key{0.5}, \Key{0.85} & \ApiText{Curve parameters of the cardinal, catmull-rom and bundle families.}{Kurvenparameter der Cardinal-, Catmull-Rom- und Bundle-Familien.}}
  \SemioTableRow{\Key{closed} & \Key{boolean} & \Key{false} & \ApiText{Closes the path after the last point.}{Schließt den Pfad nach dem letzten Punkt.}}
  \SemioTableRow{\Key{dash} & \Key{name} & \Key{none} & \ApiText{\Key{dashed}, \Key{dotted}, \Key{dash-dot} or an explicit dash pattern.}{\Key{dashed}, \Key{dotted}, \Key{dash-dot} oder ein explizites Strichmuster.}}
  \SemioTableRow{\Key{arrow} & \Key{name} & \Key{none} & \ApiText{\Key{end}, \Key{start} or \Key{both}.}{\Key{end}, \Key{start} oder \Key{both}.}}
  \SemioTableRow{\Key{arrowSize} & \Key{dimension} & \Key{1.6mm} & \ApiText{Arrowhead length.}{Länge der Pfeilspitze.}}
  \SemioTableRow{\Key{radius}, \Key{rx}, \Key{ry} & \Key{mm} & --- & \ApiText{Arc radii; \Key{rx} and \Key{ry} default to \Key{radius}.}{Bogenradien; \Key{rx} und \Key{ry} haben \Key{radius} als Vorgabe.}}
  \SemioTableRow{\Key{innerRadius}, \Key{outerRadius} & \Key{mm} & \Key{0}, \Key{radius} & \ApiText{Annulus radii.}{Radien eines Kreisrings.}}
  \SemioTableRow{\Key{startAngle}, \Key{endAngle} & \Key{rad} & \Key{0}, \Key{pi/2} & \ApiText{Clockwise from twelve o'clock.}{Im Uhrzeigersinn ab zwölf Uhr.}}
  \SemioTableRow{\Key{padAngle}, \Key{padRadius}, \Key{cornerRadius} & \Key{number} & \Key{0} & \ApiText{d3-arc padding and rounding.}{Abstand und Rundung nach d3-arc.}}
  \SemioTableRow{\Key{route} & \Key{name} & \Key{hv} & \ApiText{Order of an orthogonal connector: \Key{hv} or \Key{vh}.}{Reihenfolge eines rechtwinkligen Verbinders: \Key{hv} oder \Key{vh}.}}
  \SemioTableRow{\Key{curvature} & \Key{number} & \Key{0.4} & \ApiText{Bow of a curved connector, as a fraction of the span.}{Bauch eines gekrümmten Verbinders als Anteil der Spannweite.}}
  \SemioTableRow{\Key{via} & \Key{list} & --- & \ApiText{Waypoints of a bundled connector.}{Wegpunkte eines gebündelten Verbinders.}}
  \SemioTableRow{\Key{concavity} & \Key{number} & \Key{20} & \ApiText{Longest edge a concave hull will not dig into.}{Längste Kante, in die eine konkave Hülle nicht eindringt.}}
  \SemioTableRow{\Key{aspect} & \Key{number} & \Key{2} & \ApiText{Width-to-height ratio of the rectangle mark.}{Verhältnis von Breite zu Höhe der Rechteckmarke.}}
  \SemioTableRow{\Key{symbol} & \Key{name} & \Key{circle} & \ApiText{The d3 symbol drawn inside an icon mark.}{Das d3-Symbol, das in einer Symbolmarke gezeichnet wird.}}
  \SemioTableRow{\Key{file} & \Key{path} & --- & \ApiText{Image file of an image mark.}{Bilddatei einer Bildmarke.}}
  \SemioTableRow{\Key{extent} & \Key{lo,hi} & --- & \ApiText{Span of a reference band, a bracket or a brace.}{Spannweite eines Referenzbandes, einer eckigen oder geschweiften Klammer.}}
}

\SemioTableLong[text-size=8pt]{\ApiText{Curves}{Kurven}}{0.30,0.70}{\ApiValue & \ApiMeaning}{%
  \SemioTableRow{\Key{linear}, \Key{linear-closed} & \ApiText{Straight segments; the closed form joins the last point back to the first.}{Gerade Abschnitte; die geschlossene Form verbindet den letzten Punkt zurück zum ersten.}}
  \SemioTableRow{\Key{step}, \Key{step-before}, \Key{step-after} & \ApiText{Staircase with the riser in the middle, before or after the datum.}{Treppe mit der Stufe in der Mitte, vor oder nach dem Datum.}}
  \SemioTableRow{\Key{basis}, \Key{basis-open}, \Key{basis-closed} & \ApiText{Uniform cubic B-spline; the open form drops the end segments.}{Gleichmäßiger kubischer B-Spline; die offene Form lässt die Endabschnitte weg.}}
  \SemioTableRow{\Key{bundle} & \ApiText{Basis spline straightened towards the chord by \Key{beta}.}{Basis-Spline, um \Key{beta} zur Sehne hin begradigt.}}
  \SemioTableRow{\Key{cardinal}, \Key{cardinal-open}, \Key{cardinal-closed} & \ApiText{Cardinal spline with the given \Key{tension}.}{Cardinal-Spline mit der angegebenen \Key{tension}.}}
  \SemioTableRow{\Key{catmull-rom}, \Key{catmull-rom-open}, \Key{catmull-rom-closed} & \ApiText{Catmull-Rom spline parameterised by \Key{alpha}.}{Catmull-Rom-Spline, parametrisiert durch \Key{alpha}.}}
  \SemioTableRow{\Key{monotone-x}, \Key{monotone-y} & \ApiText{Cubic interpolation that preserves monotonicity in that channel.}{Kubische Interpolation, die die Monotonie in diesem Kanal erhält.}}
  \SemioTableRow{\Key{natural} & \ApiText{Natural cubic spline through every point.}{Natürlicher kubischer Spline durch jeden Punkt.}}
  \SemioTableRow{\Key{bezier} & \ApiText{Cubic Bézier through the explicit control points.}{Kubische Bézierkurve durch die expliziten Kontrollpunkte.}}
}

\begin{VizFigure}[title={\ApiText{Four curves over the same points}{Vier Kurven über denselben Punkten}}, width=80, height=44]
\SemioVizPath{polyline}[points={6,6; 22,32; 38,12; 54,30; 70,10}, curve=linear, dash=dotted]
\SemioVizPath{polyline}[points={6,6; 22,32; 38,12; 54,30; 70,10}, curve=monotone-x, stroke=semio-primary]
\SemioVizPath{polyline}[points={6,6; 22,32; 38,12; 54,30; 70,10}, curve=basis, stroke=semio-secondary]
\SemioVizPath{polyline}[points={6,6; 22,32; 38,12; 54,30; 70,10}, curve=step-after, stroke=semio-tertiary]
\end{VizFigure}

\section{\ApiText{Text marks}{Textmarken}}

\ApiText{%
  \Cs{SemioVizText} carries no ink of its own; its geometry is the anchor the text node is set on.
  Eight kinds fix the role and therefore the type style: \Key{cell}, \Key{module}, \Key{chip},
  \Key{label}, \Key{title}, \Key{value}, \Key{callout} and \Key{leader}.%
}{%
  \Cs{SemioVizText} trägt selbst keine Farbe; seine Geometrie ist der Ankerpunkt, auf den der
  Textknoten gesetzt wird. Acht Arten legen die Rolle und damit den Schriftstil fest: \Key{cell},
  \Key{module}, \Key{chip}, \Key{label}, \Key{title}, \Key{value}, \Key{callout} und \Key{leader}.%
}

\begin{VizFigure}[title={\ApiText{Every text kind}{Jede Textart}}, width=100, height=50]
\SemioVizText{cell}[width=25]{20,10}{Cell text}
\SemioVizText{module}[width=25]{65,10}{Module text}
\SemioVizText{chip}[width=25]{20,25}{Chip text}
\SemioVizText{label}[width=25]{50,25}{Label text}
\SemioVizText{title}[width=25]{5,40}{Title text}
\SemioVizText{value}[width=25]{80,40}{42}
\end{VizFigure}

\chapter{\ApiText{Shape generators}{Formgeneratoren}}

\ApiPkg{semio-viz-shape}
\ApiSig{\Cs{SemioVizLine}[points=, curve=, tension=, alpha=, beta=, defined=, name=, draw=]}
\ApiSig{\Cs{SemioVizArea}[points=, points2=, baseline=, curve=, defined=]}
\ApiSig{\Cs{SemioVizArc}[innerRadius=, outerRadius=, startAngle=, endAngle=, padAngle=, cornerRadius=, at=]}
\ApiSig{\Cs{SemioVizPie}[value=, sort=, startAngle=, endAngle=, padAngle=]}
\ApiSig{\Cs{SemioVizLink}[kind=, source=, target=]}
\ApiSig{\Cs{SemioVizSymbol}[type=, size=, at=]}
\ApiSig{\Cs{SemioVizRibbon}[radius=, startAngle=, endAngle=, targetRadius=, targetStartAngle=, targetEndAngle=, padAngle=, at=]}

\ApiText{%
  A generator computes a path from numbers. It draws that path by default and stores it under
  \Key{name}; with \Key{draw=false} it computes the numbers only, which is what the probe harness and
  every composed family use. Every option a mark understands is accepted here too, so a generator
  takes the full style vocabulary without duplicating it.%
}{%
  Ein Generator berechnet einen Pfad aus Zahlen. Er zeichnet diesen Pfad standardmäßig und legt ihn
  unter \Key{name} ab; mit \Key{draw=false} berechnet er nur die Zahlen, was die Prüfumgebung und
  jede zusammengesetzte Familie nutzen. Jede Option, die eine Marke versteht, wird auch hier
  angenommen, ein Generator übernimmt also das volle Stilvokabular, ohne es zu wiederholen.%
}

\SemioTableLong[text-size=8pt]{\ApiText{Generator options}{Generatoroptionen}}{0.20,0.14,0.16,0.50}{\ApiKey & \ApiType & \ApiDefault & \ApiMeaning}{%
  \SemioTableRow{\Key{name} & \Key{name} & --- & \ApiText{Stores the generated path under this name, for later reuse.}{Legt den erzeugten Pfad unter diesem Namen zur Wiederverwendung ab.}}
  \SemioTableRow{\Key{draw} & \Key{boolean} & \Key{true} & \ApiText{Renders the path; \Key{false} computes the numbers only.}{Zeichnet den Pfad; \Key{false} berechnet nur die Zahlen.}}
  \SemioTableRow{\Key{at} & \Key{x,y} & \Key{0,0} & \ApiText{Figure point the centred generators sit on.}{Figurenpunkt, auf dem die zentrierten Generatoren sitzen.}}
  \SemioTableRow{\Key{type} & \Key{name} & \Key{circle} & \ApiText{The d3 symbol type of \Cs{SemioVizSymbol}.}{Der d3-Symboltyp von \Cs{SemioVizSymbol}.}}
  \SemioTableRow{\Key{kind} & \Key{name} & \Key{horizontal} & \ApiText{Link flavour: \Key{horizontal}, \Key{vertical} or \Key{radial}.}{Verbindungsart: \Key{horizontal}, \Key{vertical} oder \Key{radial}.}}
  \SemioTableRow{\Key{source}, \Key{target} & \Key{x,y} & --- & \ApiText{The two endpoints of a link; \Key{angle,radius} when radial.}{Die beiden Endpunkte einer Verbindung; \Key{angle,radius} im radialen Fall.}}
  \SemioTableRow{\Key{value} & \Key{clist} & --- & \ApiText{The pie's values.}{Die Werte des Tortendiagramms.}}
  \SemioTableRow{\Key{sort} & \Key{name} & \Key{descending} & \ApiText{Pie layout order: \Key{descending}, \Key{ascending} or \Key{none}.}{Anordnung der Torte: \Key{descending}, \Key{ascending} oder \Key{none}.}}
  \SemioTableRow{\Key{baseline} & \Key{number} & --- & \ApiText{Constant lower boundary of an area when \Key{points2} is empty.}{Konstante untere Begrenzung einer Fläche, wenn \Key{points2} leer ist.}}
  \SemioTableRow{\Key{defined} & \Key{clist} & --- & \ApiText{Per-point 0/1 flags; a 0 breaks the line, exactly as d3's \Key{defined} does.}{0/1-Merker je Punkt; eine 0 unterbricht die Linie, genau wie d3s \Key{defined}.}}
}

\ApiText{%
  The results are readable: \Cs{SemioVizShapeResult} holds the path the last generator produced and
  \Cs{SemioVizShapePath}\{\meta{name}\} a named one, both as the d3 SVG token list.
  \Cs{SemioVizArcCentroid} holds the centroid of the last arc; \Cs{SemioVizPieStart}\{\meta{i}\},
  \Cs{SemioVizPieEnd}\{\meta{i}\} and \Cs{SemioVizPieCount} the angles of the last pie.%
}{%
  Die Ergebnisse sind lesbar: \Cs{SemioVizShapeResult} enthält den Pfad des letzten Generators und
  \Cs{SemioVizShapePath}\{\meta{Name}\} einen benannten, beide als d3-SVG-Tokenliste.
  \Cs{SemioVizArcCentroid} enthält den Schwerpunkt des letzten Bogens; \Cs{SemioVizPieStart}\{\meta{i}\},
  \Cs{SemioVizPieEnd}\{\meta{i}\} und \Cs{SemioVizPieCount} die Winkel der letzten Torte.%
}

\begin{VizFigure}[title={\ApiText{Arc, symbol and link}{Bogen, Symbol und Verbindung}}, width=80, height=44]
\SemioVizArc[at={20,22}, innerRadius=6, outerRadius=16, startAngle=0, endAngle=2.4, cornerRadius=1, fill=semio-primary]
\SemioVizSymbol[type=diamond, size=30, at={46,22}, fill=semio-secondary]
\SemioVizLink[kind=horizontal, source={58,10}, target={74,34}, stroke=semio-tertiary]
\end{VizFigure}

\chapter{\ApiText{Layout algorithms}{Layoutalgorithmen}}

\ApiPkg{semio-viz-layout}
\ApiSig{\Cs{SemioVizLayout}\{\meta{algorithm}\}\{\meta{in}\}\{\meta{out}\}[\meta{options}]}
\ApiSig{\Cs{SemioVizLayoutAlgorithms}}

\ApiText{%
  A layout turns a table, a hierarchy, a graph or a point set into positions. There is one flat
  namespace of algorithms: every domain kernel --- hierarchy, network, flow, geo, spatial and the
  transform kernel --- registers its own, so a caller never has to know which package implements
  one. \Cs{SemioVizLayoutAlgorithms} expands to the registered names, and an unknown name is an
  error that lists them. Thirty algorithms are registered:%
}{%
  Ein Layout macht aus einer Tabelle, einer Hierarchie, einem Graphen oder einer Punktmenge
  Positionen. Es gibt einen flachen Namensraum von Algorithmen: Jeder Fachkern --- Hierarchie,
  Netzwerk, Fluss, Geo, Raum und der Transformationskern --- registriert seine eigenen, ein Aufrufer
  muss also nie wissen, welches Paket einen umsetzt. \Cs{SemioVizLayoutAlgorithms} expandiert zu den
  registrierten Namen, und ein unbekannter Name ist ein Fehler, der sie auflistet. Dreißig
  Algorithmen sind registriert:%
}

\ApiSig{\SemioVizLayoutAlgorithms}

\SemioTableLong[text-size=8pt]{\ApiText{Layout algorithms by kernel}{Layoutalgorithmen nach Kern}}{0.20,0.80}{\ApiText{Kernel}{Kern} & \ApiText{Algorithms}{Algorithmen}}{%
  \SemioTableRow{\Key{hierarchy} & \Key{tree} \ApiText{(Reingold--Tilford/Buchheim)}{(Reingold--Tilford/Buchheim)}, \Key{cluster}, \Key{treemap}, \Key{partition}, \Key{pack}}
  \SemioTableRow{\Key{network} & \Key{force}, \Key{circular}, \Key{shell}, \Key{grid}, \Key{spectral}, \Key{layered} \ApiText{(Sugiyama)}{(Sugiyama)}, \Key{radial}, \Key{arc}, \Key{chord}, \Key{bundling}, \Key{adjacency}}
  \SemioTableRow{\Key{flow} & \Key{sankey}, \Key{alluvial}, \Key{parallel-sets}}
  \SemioTableRow{\Key{geo} & \Key{projection}}
  \SemioTableRow{\Key{spatial} & \Key{delaunay}, \Key{voronoi}, \Key{hull}, \Key{hexbin}, \Key{contour} \ApiText{(marching squares)}{(Marching Squares)}, \Key{density} \ApiText{(2D kernel density)}{(2D-Kerndichte)}, \Key{beeswarm}, \Key{jitter}}
  \SemioTableRow{\Key{transform} & \Key{stack}, \Key{bin}}
}

\section{\ApiText{Hierarchy layouts}{Hierarchielayouts}}

\SemioTableLong[text-size=8pt]{\ApiText{Hierarchy layout options}{Optionen der Hierarchielayouts}}{0.22,0.16,0.14,0.48}{\ApiKey & \ApiType & \ApiDefault & \ApiMeaning}{%
  \SemioTableRow{\Key{size} & \Key{\{w,h\}} & \Key{\{1,1\}} & \ApiText{Extent the layout is normalised into.}{Ausdehnung, auf die das Layout normiert wird.}}
  \SemioTableRow{\Key{nodeSize} & \Key{\{w,h\}} & --- & \ApiText{Fixed per-node spacing; disables \Key{size}.}{Feste Abstände je Knoten; schaltet \Key{size} ab.}}
  \SemioTableRow{\Key{separationSame} & \Key{number} & \Key{1} & \ApiText{Tree and cluster gap between siblings.}{Abstand zwischen Geschwistern bei Baum und Cluster.}}
  \SemioTableRow{\Key{separationOther} & \Key{number} & \Key{2} & \ApiText{Tree and cluster gap between cousins.}{Abstand zwischen Vettern bei Baum und Cluster.}}
  \SemioTableRow{\Key{separationDepth} & \Key{boolean} & \Key{false} & \ApiText{Divides the gap by the node depth, for radial trees.}{Teilt den Abstand durch die Knotentiefe, für radiale Bäume.}}
  \SemioTableRow{\Key{tile} & \Key{name} & \Key{squarify} & \ApiText{\Key{squarify}, \Key{resquarify}, \Key{slice}, \Key{dice}, \Key{slice-dice}, \Key{binary}, \Key{strip} or \Key{voronoi}.}{\Key{squarify}, \Key{resquarify}, \Key{slice}, \Key{dice}, \Key{slice-dice}, \Key{binary}, \Key{strip} oder \Key{voronoi}.}}
  \SemioTableRow{\Key{ratio} & \Key{number} & \Key{golden} & \ApiText{Target aspect ratio of \Key{squarify} and \Key{strip}.}{Angestrebtes Seitenverhältnis von \Key{squarify} und \Key{strip}.}}
  \SemioTableRow{\Key{padding} & \Key{number} & \Key{0} & \ApiText{Treemap inner and outer, partition and pack padding.}{Innerer und äußerer Abstand bei Treemap, Partition und Pack.}}
  \SemioTableRow{\Key{paddingInner} & \Key{number} & \Key{0} & \ApiText{Treemap gap between siblings.}{Treemap-Abstand zwischen Geschwistern.}}
  \SemioTableRow{\Key{paddingOuter} & \Key{number} & \Key{0} & \ApiText{Treemap gap inside a parent, on all four sides.}{Treemap-Abstand innerhalb eines Elternknotens, auf allen vier Seiten.}}
  \SemioTableRow{\Key{paddingTop} … \Key{paddingLeft} & \Key{number} & \Key{0} & \ApiText{The four sides individually.}{Die vier Seiten einzeln.}}
  \SemioTableRow{\Key{round} & \Key{boolean} & \Key{false} & \ApiText{Snaps rectangle corners to whole millimetres.}{Rastet Rechteckecken auf ganze Millimeter ein.}}
  \SemioTableRow{\Key{radius} & \Key{name|number} & \Key{auto} & \ApiText{Pack leaf radius: \Key{auto} is the square root of the value, \Key{value} the value itself, a number gives equal circles.}{Blattradius beim Pack: \Key{auto} ist die Wurzel des Werts, \Key{value} der Wert selbst, eine Zahl ergibt gleich große Kreise.}}
  \SemioTableRow{\Key{relaxation} & \Key{integer} & \Key{24} & \ApiText{Iterations of the voronoi tiling.}{Iterationen der Voronoi-Kachelung.}}
}

\begin{VizFigure}[title={\ApiText{Treemap and pack of one hierarchy}{Kacheldiagramm und Packung derselben Hierarchie}}, width=80, height=44]
\SemioVizChart{treemap}
\end{VizFigure}

\section{\ApiText{The force layout}{Das Kräftelayout}}

\ApiText{%
  \Key{force} reproduces \Key{d3.forceSimulation()} exactly, ticked by hand: every default is
  d3-force's own and the many-body force is evaluated once per unordered pair. \Key{theta} is
  deliberately absent --- there is no Barnes--Hut approximation, because approximating would give up
  the parity the test suite is built on. That exactness has a price; see the chapter on numerics.%
}{%
  \Key{force} bildet \Key{d3.forceSimulation()} exakt nach, von Hand getaktet: Jede Vorgabe ist die
  von d3-force, und die Vielkörperkraft wird einmal je ungeordnetem Paar ausgewertet. \Key{theta}
  fehlt bewusst --- es gibt keine Barnes-Hut-Näherung, denn eine Näherung gäbe die Übereinstimmung
  auf, auf der die Testsuite beruht. Diese Exaktheit hat einen Preis; siehe das Kapitel zur
  Numerik.%
}

\SemioTableLong[text-size=8pt]{\ApiText{Force layout options}{Optionen des Kräftelayouts}}{0.24,0.14,0.14,0.48}{\ApiKey & \ApiType & \ApiDefault & \ApiMeaning}{%
  \SemioTableRow{\Key{iterations} & \Key{integer} & \Key{300} & \ApiText{Ticks; the cost is proportional to iterations times (nodes squared plus edges).}{Takte; der Aufwand wächst mit Iterationen mal (Knoten zum Quadrat plus Kanten).}}
  \SemioTableRow{\Key{alpha} & \Key{number} & \Key{1} & \ApiText{Starting alpha.}{Anfangswert von Alpha.}}
  \SemioTableRow{\Key{alpha-min} & \Key{number} & \Key{0.001} & \ApiText{Only used to derive the automatic decay.}{Dient nur der Herleitung des automatischen Abfalls.}}
  \SemioTableRow{\Key{alpha-decay} & \Key{number} & \Key{-1} & \ApiText{\Key{-1} means d3's own value derived from \Key{alpha-min}.}{\Key{-1} bedeutet d3s eigenen, aus \Key{alpha-min} hergeleiteten Wert.}}
  \SemioTableRow{\Key{alpha-target} & \Key{number} & \Key{0} & \ApiText{Alpha decays towards this.}{Alpha fällt auf diesen Wert zu.}}
  \SemioTableRow{\Key{velocity-decay} & \Key{number} & \Key{0.4} & \ApiText{d3's public value; the multiplier is one minus it.}{d3s öffentlicher Wert; der Faktor ist eins minus dieser Wert.}}
  \SemioTableRow{\Key{forces} & \Key{clist} & \Key{link, many-body, center} & \ApiText{The forces, applied in the order they are listed.}{Die Kräfte, in der aufgeführten Reihenfolge angewandt.}}
  \SemioTableRow{\Key{link-distance} & \Key{number} & \Key{30} & \ApiText{Rest length of a link.}{Ruhelänge einer Kante.}}
  \SemioTableRow{\Key{link-strength} & \Key{number} & \Key{-1} & \ApiText{\Key{-1} means one over the smaller endpoint degree, as d3 does.}{\Key{-1} bedeutet eins durch den kleineren Endknotengrad, wie bei d3.}}
  \SemioTableRow{\Key{link-iterations} & \Key{integer} & \Key{1} & \ApiText{Passes of the link force per tick.}{Durchläufe der Kantenkraft je Takt.}}
  \SemioTableRow{\Key{many-body-strength} & \Key{number} & \Key{-30} & \ApiText{Charge of every node.}{Ladung jedes Knotens.}}
  \SemioTableRow{\Key{distance-min}, \Key{distance-max} & \Key{number} & \Key{1}, \Key{-1} & \ApiText{Many-body clamps; \Key{-1} means no upper cutoff.}{Begrenzungen der Vielkörperkraft; \Key{-1} bedeutet keine obere Schranke.}}
  \SemioTableRow{\Key{center-x}, \Key{center-y}, \Key{center-strength} & \Key{number} & \Key{0}, \Key{0}, \Key{1} & \ApiText{Where the centroid is pulled to, and how hard.}{Wohin der Schwerpunkt gezogen wird und wie stark.}}
  \SemioTableRow{\Key{collide-radius}, \Key{collide-strength}, \Key{collide-iterations} & \Key{number} & \Key{1}, \Key{0.7}, \Key{1} & \ApiText{One collision radius for every node.}{Ein Kollisionsradius für jeden Knoten.}}
  \SemioTableRow{\Key{x-target}, \Key{x-strength}, \Key{y-target}, \Key{y-strength} & \Key{number} & \Key{0}, \Key{0.1} & \ApiText{The per-axis springs.}{Die Federn je Achse.}}
  \SemioTableRow{\Key{radial-radius}, \Key{radial-x}, \Key{radial-y}, \Key{radial-strength} & \Key{number} & \Key{100}, \Key{0}, \Key{0}, \Key{0.1} & \ApiText{The circle the radial force pulls towards.}{Der Kreis, zu dem die radiale Kraft zieht.}}
  \SemioTableRow{\Key{seed} & \Key{number} & \Key{1} & \ApiText{Seed of the generator behind d3's jiggle; the layout is deterministic.}{Startwert des Generators hinter d3s jiggle; das Layout ist deterministisch.}}
}

\section{\ApiText{Flow layouts}{Flusslayouts}}

\SemioTableLong[text-size=8pt]{\ApiText{Sankey, alluvial and parallel-sets options}{Optionen von Sankey, Alluvial und Parallelmengen}}{0.20,0.16,0.16,0.48}{\ApiKey & \ApiType & \ApiDefault & \ApiMeaning}{%
  \SemioTableRow{\Key{node-width} & \Key{number} & \Key{24} & \ApiText{Width of a node rectangle, in layout units.}{Breite eines Knotenrechtecks, in Layouteinheiten.}}
  \SemioTableRow{\Key{node-padding} & \Key{number} & \Key{8} & \ApiText{Vertical gap between the nodes of one column.}{Senkrechter Abstand zwischen den Knoten einer Spalte.}}
  \SemioTableRow{\Key{node-align} & \Key{name} & \Key{justify} & \ApiText{Column assignment: \Key{left}, \Key{right}, \Key{center} or \Key{justify}.}{Spaltenzuweisung: \Key{left}, \Key{right}, \Key{center} oder \Key{justify}.}}
  \SemioTableRow{\Key{node-sort} & \Key{name} & \Key{auto} & \ApiText{\Key{auto} is d3's ordering of each column by its lower edge.}{\Key{auto} ist d3s Sortierung jeder Spalte nach ihrer Unterkante.}}
  \SemioTableRow{\Key{link-sort} & \Key{name} & \Key{auto} & \ApiText{\Key{auto} is d3's link ordering by partner breadth.}{\Key{auto} ist d3s Kantensortierung nach der Breite des Partners.}}
  \SemioTableRow{\Key{iterations} & \Key{integer} & \Key{6} & \ApiText{Relaxation passes.}{Entspannungsdurchläufe.}}
  \SemioTableRow{\Key{x0}, \Key{y0}, \Key{x1}, \Key{y1} & \Key{number} & \Key{0}, \Key{0}, \Key{1}, \Key{1} & \ApiText{The extent the layout writes into.}{Die Ausdehnung, in die das Layout schreibt.}}
}

\ApiText{%
  \Key{alluvial} presets \Key{node-align=left} and \Key{node-sort=none}; \Key{parallel-sets}
  additionally presets \Key{node-padding=0}. Otherwise the three share one implementation and one
  vocabulary.%
}{%
  \Key{alluvial} setzt \Key{node-align=left} und \Key{node-sort=none} vor; \Key{parallel-sets} setzt
  zusätzlich \Key{node-padding=0} vor. Ansonsten teilen sich die drei eine Umsetzung und ein
  Vokabular.%
}

\section{\ApiText{Spatial layouts}{Raumlayouts}}

\SemioTableLong[text-size=8pt]{\ApiText{Spatial layout options}{Optionen der Raumlayouts}}{0.20,0.16,0.20,0.44}{\ApiKey & \ApiType & \ApiDefault & \ApiMeaning}{%
  \SemioTableRow{\Key{x0}, \Key{y0}, \Key{x1}, \Key{y1} & \Key{number} & \Key{0}, \Key{0}, \Key{100}, \Key{100} & \ApiText{Clip rectangle of the voronoi cells.}{Beschneidungsrechteck der Voronoi-Zellen.}}
  \SemioTableRow{\Key{radius} & \Key{number} & \Key{4} & \ApiText{Hexbin radius and beeswarm collision radius.}{Hexbin-Radius und Kollisionsradius des Bienenschwarms.}}
  \SemioTableRow{\Key{amount} & \Key{number} & \Key{2} & \ApiText{Jitter amplitude.}{Streuungsweite des Jitters.}}
  \SemioTableRow{\Key{threshold} & \Key{number} & \Key{0.5} & \ApiText{Contour isoline level.}{Höhe der Isolinie.}}
  \SemioTableRow{\Key{bandwidth} & \Key{number} & \Key{8} & \ApiText{Density kernel bandwidth.}{Bandbreite des Dichtekerns.}}
  \SemioTableRow{\Key{cell} & \Key{number} & \Key{4} & \ApiText{Density grid step.}{Schrittweite des Dichtegitters.}}
  \SemioTableRow{\Key{width}, \Key{height} & \Key{integer} & \Key{20} & \ApiText{Density grid size in cells.}{Größe des Dichtegitters in Zellen.}}
}

\ApiText{%
  The spatial kernel is also reachable directly, which is what a family that needs the numbers rather
  than a table does: \Cs{SemioVizPointSet}, \Cs{SemioVizValueGrid}, \Cs{SemioVizDelaunay},
  \Cs{SemioVizVoronoi}, \Cs{SemioVizHull}, \Cs{SemioVizHexbin}, \Cs{SemioVizBeeswarm},
  \Cs{SemioVizJitter}, \Cs{SemioVizContour} and \Cs{SemioVizDensity}.%
}{%
  Der Raumkern ist auch direkt erreichbar, was eine Familie nutzt, die die Zahlen statt einer Tabelle
  braucht: \Cs{SemioVizPointSet}, \Cs{SemioVizValueGrid}, \Cs{SemioVizDelaunay},
  \Cs{SemioVizVoronoi}, \Cs{SemioVizHull}, \Cs{SemioVizHexbin}, \Cs{SemioVizBeeswarm},
  \Cs{SemioVizJitter}, \Cs{SemioVizContour} und \Cs{SemioVizDensity}.%
}

\ApiSig{\Cs{SemioVizQuadtree}\{\meta{table}\}[x=, y=] \ \Cs{SemioVizQuadtreeFind}\{\meta{x}\}\{\meta{y}\}[\meta{radius}]}

\ApiText{%
  The spatial index the collision layouts stand on. \Cs{SemioVizQuadtree} builds d3-quadtree's tree
  over two columns of a table: the covering square is not fitted to the data but grown by doubling
  towards any point that falls outside it, which is what makes two implementations agree on the
  extent and not merely on the answers. \Cs{SemioVizQuadtreeFind} typesets the one-based row index
  of the nearest point within the search radius, or \Key{0} when the tree holds none; it walks the
  tree nearest quadrant first and prunes every quadrant that cannot hold anything closer than the
  best candidate so far.%
}{%
  Der Raumindex, auf dem die Kollisionslayouts stehen. \Cs{SemioVizQuadtree} baut den Baum von
  d3-quadtree über zwei Spalten einer Tabelle: das überdeckende Quadrat wird nicht an die Daten
  angepasst, sondern durch Verdopplung auf jeden Punkt außerhalb erweitert --- erst das lässt zwei
  Implementierungen in der Ausdehnung und nicht nur in den Antworten übereinstimmen.
  \Cs{SemioVizQuadtreeFind} setzt den einsbasierten Zeilenindex des nächsten Punktes innerhalb des
  Suchradius, oder \Key{0}, wenn der Baum keinen enthält; die Suche läuft quadrantenweise vom
  nächsten aus und schneidet jeden Quadranten ab, der nichts Näheres als den bisher besten Kandidaten
  enthalten kann.%
}

\section{\ApiText{The geographic kernel}{Der geographische Kern}}

\ApiPkg{semio-viz-geo}
\ApiSig{\Cs{SemioVizProjection}\{\meta{name}\}[kind=, scale=, translate=, center=, rotate=, parallels=, clipAngle=, clipExtent=, precision=]}
\ApiSig{\Cs{SemioVizGeoCollection}\{\meta{name}\} \ \Cs{SemioVizGeoPolygon} \ \Cs{SemioVizGeoLine} \ \Cs{SemioVizGeoPointPart} \ \Cs{SemioVizGeoFromTable}}
\ApiSig{\Cs{SemioVizGeoPath}\{\meta{projection}\}\{\meta{geometry}\}[\meta{style}] \ \Cs{SemioVizGraticule}\{\meta{projection}\}[\meta{keys}]}
\ApiSig{\Cs{SemioVizGeoProject}\{\meta{name}\}\{\meta{lon}\}\{\meta{lat}\} \ \Cs{SemioVizGeoInvert} \ \Cs{SemioVizGeoX} \ \Cs{SemioVizGeoY} \ \Cs{SemioVizGeoLon} \ \Cs{SemioVizGeoLat}}
\ApiSig{\Cs{SemioVizGeoFit}, \Cs{SemioVizGeoFitSize}, \Cs{SemioVizGeoFitWidth}, \Cs{SemioVizGeoArea}, \Cs{SemioVizGeoDistance}, \Cs{SemioVizGeoCentroid}, \Cs{SemioVizGeoInterpolate}, \Cs{SemioVizGeoCircle}}

\ApiText{%
  A geometry is a named collection of rings, lines and points in degrees; a projection is a named
  mapping from degrees to millimetres. Fourteen projection kinds are registered, and every default
  equals d3's default for that kind, so an unconfigured projection reproduces d3 exactly. The
  spherical measures --- area, great-circle distance, centroid, interpolation and the small circle
  --- are the d3-geo ones and leave their result in the expandable accessors.%
}{%
  Eine Geometrie ist eine benannte Sammlung von Ringen, Linien und Punkten in Grad; eine Projektion
  ist eine benannte Abbildung von Grad auf Millimeter. Vierzehn Projektionsarten sind registriert,
  und jede Vorgabe entspricht d3s Vorgabe für diese Art, eine unkonfigurierte Projektion bildet d3
  also exakt nach. Die sphärischen Maße --- Fläche, Großkreisentfernung, Schwerpunkt, Interpolation
  und der Kleinkreis --- sind die von d3-geo und hinterlassen ihr Ergebnis in den expandierbaren
  Zugriffsbefehlen.%
}

\SemioTableLong[text-size=8pt]{\ApiText{Projection options}{Projektionsoptionen}}{0.18,0.16,0.16,0.50}{\ApiKey & \ApiType & \ApiDefault & \ApiMeaning}{%
  \SemioTableRow{\Key{kind} & \Key{choice} & \Key{equirectangular} & \ApiText{\Key{mercator}, \Key{transverse-mercator}, \Key{equal-earth}, \Key{natural-earth1}, \Key{orthographic}, \Key{stereographic}, \Key{gnomonic}, \Key{azimuthal-equal-area}, \Key{azimuthal-equidistant}, \Key{conic-conformal}, \Key{conic-equal-area}, \Key{conic-equidistant} and \Key{albers}.}{\Key{mercator}, \Key{transverse-mercator}, \Key{equal-earth}, \Key{natural-earth1}, \Key{orthographic}, \Key{stereographic}, \Key{gnomonic}, \Key{azimuthal-equal-area}, \Key{azimuthal-equidistant}, \Key{conic-conformal}, \Key{conic-equal-area}, \Key{conic-equidistant} und \Key{albers}.}}
  \SemioTableRow{\Key{scale} & \Key{number} & \ApiText{d3's own}{d3s eigener} & \ApiText{The projection scale factor.}{Der Maßstabsfaktor der Projektion.}}
  \SemioTableRow{\Key{translate} & \Key{x,y} & \Key{480,250} & \ApiText{Plane translation in millimetres --- d3's default.}{Verschiebung in der Ebene in Millimetern --- d3s Vorgabe.}}
  \SemioTableRow{\Key{center} & \Key{lon,lat} & \ApiText{d3's own}{d3s eigener} & \ApiText{The centre of the projection, in degrees.}{Der Mittelpunkt der Projektion, in Grad.}}
  \SemioTableRow{\Key{rotate} & \Key{l,p,g} & \ApiText{d3's own}{d3s eigene} & \ApiText{The three rotation angles; \Key{transverse-mercator} adds d3's own gamma of ninety degrees.}{Die drei Drehwinkel; \Key{transverse-mercator} ergänzt d3s eigenes Gamma von neunzig Grad.}}
  \SemioTableRow{\Key{parallels} & \Key{p0,p1} & \ApiText{d3's own}{d3s eigene} & \ApiText{The standard parallels of the conic kinds.}{Die Standardparallelen der Kegelprojektionen.}}
  \SemioTableRow{\Key{clipAngle} & \Key{number} & --- & \ApiText{Great-circle clip radius in degrees; empty disables clipping.}{Großkreis-Beschneidungsradius in Grad; leer schaltet das Beschneiden ab.}}
  \SemioTableRow{\Key{clipExtent} & \Key{x0,y0,x1,y1} & --- & \ApiText{Plane clip rectangle; empty disables clipping.}{Beschneidungsrechteck in der Ebene; leer schaltet das Beschneiden ab.}}
  \SemioTableRow{\Key{precision} & \Key{number} & \Key{sqrt(0.5)} & \ApiText{Adaptive resampling tolerance in plane units; \Key{0} disables resampling. This is the knob that decides how expensive a map is.}{Toleranz des adaptiven Nachtastens in Ebeneneinheiten; \Key{0} schaltet das Nachtasten ab. Dieser Regler bestimmt, wie teuer eine Karte wird.}}
}

\begin{VizFigure}[title={\ApiText{A projected geometry with a graticule}{Eine projizierte Geometrie mit Gradnetz}}, width=90, height=50]
\SemioVizProjection{apimap}[kind=equal-earth]
\SemioVizGeoFromTable{apigeo}{demo-geo}[ring=ring, lon=lon, lat=lat]
\SemioVizGeoFitSize{apimap}{apigeo}{80}{40}
\SemioVizGraticule{apimap}[step={30,30}, precision=30]
\SemioVizGeoPath{apimap}{apigeo}
\end{VizFigure}

\section{\ApiText{Direct kernel entry points}{Direkte Kerneinstiege}}

\ApiText{%
  \Cs{SemioVizLayout} is the grammar element; a family that needs the intermediate numbers rather
  than a positions table calls its kernel directly. The hierarchy kernel offers
  \Cs{SemioVizHierarchyBuild}\{\meta{name}\}[\meta{keys}] to stratify a table (by \Key{id} and
  \Key{parent}, or by a delimited \Key{path} column),
  \Cs{SemioVizHierarchyLayout}\{\meta{algorithm}\}\{\meta{in}\}\{\meta{out}\}[\meta{keys}],
  the two iterators \Cs{SemioVizHierarchyMapNodes} and \Cs{SemioVizHierarchyMapLinks}, and the
  expandable field accessor \Cs{SemioVizHierarchyField}\{\meta{name}\}\{\meta{field}\}\{\meta{index}\}.
  The network kernel offers \Cs{SemioVizGraphMetrics} (degree, in and out degree, component and
  layer per node), \Cs{SemioVizGraphOrder} (an ordered traversal table) and the deterministic pair
  \Cs{SemioVizGraphSeed} / \Cs{SemioVizGraphRandom}, whose documented default seed is one.%
}{%
  \Cs{SemioVizLayout} ist das Grammatikelement; eine Familie, die die Zwischenzahlen statt einer
  Positionstabelle braucht, ruft ihren Kern direkt auf. Der Hierarchiekern bietet
  \Cs{SemioVizHierarchyBuild}\{\meta{Name}\}[\meta{Schlüssel}], um eine Tabelle zu schichten (über
  \Key{id} und \Key{parent} oder über eine getrennte \Key{path}-Spalte),
  \Cs{SemioVizHierarchyLayout}\{\meta{Algorithmus}\}\{\meta{ein}\}\{\meta{aus}\}[\meta{Schlüssel}],
  die beiden Iteratoren \Cs{SemioVizHierarchyMapNodes} und \Cs{SemioVizHierarchyMapLinks} sowie den
  expandierbaren Feldzugriff
  \Cs{SemioVizHierarchyField}\{\meta{Name}\}\{\meta{Feld}\}\{\meta{Index}\}. Der Netzwerkkern bietet
  \Cs{SemioVizGraphMetrics} (Grad, Ein- und Ausgangsgrad, Komponente und Ebene je Knoten),
  \Cs{SemioVizGraphOrder} (eine geordnete Durchlauftabelle) und das deterministische Paar
  \Cs{SemioVizGraphSeed} / \Cs{SemioVizGraphRandom}, dessen dokumentierter Startwert eins ist.%
}

\chapter{\ApiText{Guides}{Hilfselemente}}

\ApiPkg{semio-viz-guide}
\ApiSig{\Cs{SemioVizAxis}[scale=, orient=, title=, ticks=, tickFormat=, grid=, …]}
\ApiSig{\Cs{SemioVizGrid}[x=, y=, minor=, ticks=, style=]}
\ApiSig{\Cs{SemioVizLegend}[scale=, kind=, title=, orient=, columns=, at=]}

\ApiText{%
  A guide describes a scale. It takes a scale by name and derives its geometry from that scale
  alone, so an axis can never disagree with the marks: both read the same mapping. An axis is
  cartesian, polar or drawn along an arbitrary segment.%
}{%
  Ein Hilfselement beschreibt eine Skala. Es nimmt eine Skala über ihren Namen und leitet seine
  Geometrie allein aus dieser Skala ab; eine Achse kann also nie im Widerspruch zu den Marken stehen,
  denn beide lesen dieselbe Abbildung. Eine Achse ist kartesisch, polar oder entlang einer beliebigen
  Strecke gezeichnet.%
}

\SemioTableLong[text-size=8pt]{\ApiText{Axis options}{Achsenoptionen}}{0.20,0.14,0.16,0.50}{\ApiKey & \ApiType & \ApiDefault & \ApiMeaning}{%
  \SemioTableRow{\Key{scale} & \Key{name} & \Key{x} & \ApiText{The scale whose domain the axis describes.}{Die Skala, deren Bereich die Achse beschreibt.}}
  \SemioTableRow{\Key{orient} & \Key{choice} & \Key{bottom} & \ApiText{\Key{bottom}, \Key{top}, \Key{left}, \Key{right}, \Key{segment}, \Key{angular} or \Key{radial}.}{\Key{bottom}, \Key{top}, \Key{left}, \Key{right}, \Key{segment}, \Key{angular} oder \Key{radial}.}}
  \SemioTableRow{\Key{at} & \Key{mm} & --- & \ApiText{Cross-axis position of the axis line; empty means the frame edge.}{Position der Achsenlinie quer zur Achse; leer bedeutet die Rahmenkante.}}
  \SemioTableRow{\Key{offset} & \Key{mm} & \Key{0} & \ApiText{Outward shift of the axis line, for dual axes.}{Verschiebung der Achsenlinie nach außen, für Doppelachsen.}}
  \SemioTableRow{\Key{from}, \Key{to} & \Key{x,y} & --- & \ApiText{Segment endpoints when \Key{orient=segment}.}{Endpunkte der Strecke bei \Key{orient=segment}.}}
  \SemioTableRow{\Key{center} & \Key{x,y} & \ApiText{frame centre}{Rahmenmitte} & \ApiText{Polar centre for \Key{angular} and \Key{radial}.}{Polarmittelpunkt für \Key{angular} und \Key{radial}.}}
  \SemioTableRow{\Key{innerRadius}, \Key{outerRadius} & \Key{mm} & \Key{0}, \ApiText{fitted}{eingepasst} & \ApiText{Polar radii of the axis.}{Polarradien der Achse.}}
  \SemioTableRow{\Key{startAngle}, \Key{endAngle} & \Key{deg} & \Key{0}, \Key{360} & \ApiText{Angular sweep.}{Winkelbereich.}}
  \SemioTableRow{\Key{angle} & \Key{deg} & \Key{90} & \ApiText{Direction of a radial axis.}{Richtung einer radialen Achse.}}
  \SemioTableRow{\Key{title}, \Key{titleAnchor}, \Key{titleGap} & \Key{text}, \Key{choice}, \Key{mm} & ---, \Key{middle}, \Key{5} & \ApiText{Axis title, its anchor and its distance from the tick labels.}{Achsentitel, sein Anker und sein Abstand zur Teilstrichbeschriftung.}}
  \SemioTableRow{\Key{ticks} & \Key{integer} & \Key{5} & \ApiText{Target tick count for continuous scales.}{Angestrebte Teilstrichzahl bei stetigen Skalen.}}
  \SemioTableRow{\Key{tickValues} & \Key{clist} & --- & \ApiText{Explicit tick values; overrides \Key{ticks}.}{Explizite Teilstrichwerte; übergeht \Key{ticks}.}}
  \SemioTableRow{\Key{tickFormat} & \Key{spec} & --- & \ApiText{Format specifier for the tick labels.}{Formatangabe für die Teilstrichbeschriftung.}}
  \SemioTableRow{\Key{tickSize}, \Key{tickSizeOuter}, \Key{tickPadding} & \Key{mm} & \Key{1.4}, \Key{1.4}, \Key{0.8} & \ApiText{Inner tick length, end caps of the domain line, gap to the label.}{Innere Teilstrichlänge, Endkappen der Achsenlinie, Abstand zur Beschriftung.}}
  \SemioTableRow{\Key{minor}, \Key{minorTicks}, \Key{minorSize} & \Key{boolean}, \Key{integer}, \Key{mm} & \Key{false}, \Key{4}, \Key{0.7} & \ApiText{Minor ticks and their subdivision.}{Nebenteilstriche und ihre Unterteilung.}}
  \SemioTableRow{\Key{grid}, \Key{gridLength} & \Key{boolean}, \Key{mm} & \Key{false}, --- & \ApiText{A grid line at every major tick; empty length spans the plot area.}{Eine Gitterlinie an jedem Hauptteilstrich; leere Länge überspannt die Zeichenfläche.}}
  \SemioTableRow{\Key{domainLine} & \Key{boolean} & \Key{true} & \ApiText{Draws the axis line itself.}{Zeichnet die Achsenlinie selbst.}}
  \SemioTableRow{\Key{labels}, \Key{labelRotate}, \Key{labelAlign} & \Key{boolean}, \Key{deg}, \Key{choice} & \Key{true}, \Key{0}, \Key{auto} & \ApiText{Tick labels, their rotation and their alignment.}{Teilstrichbeschriftung, ihre Drehung und ihre Ausrichtung.}}
  \SemioTableRow{\Key{broken}, \Key{breakGap} & \Key{lo,hi}, \Key{mm} & ---, \Key{2} & \ApiText{A data interval removed from the axis, and the width it collapses to.}{Ein aus der Achse entferntes Datenintervall und die Breite, auf die es zusammenfällt.}}
}

\SemioTableLong[text-size=8pt]{\ApiText{Legend options}{Legendenoptionen}}{0.20,0.14,0.16,0.50}{\ApiKey & \ApiType & \ApiDefault & \ApiMeaning}{%
  \SemioTableRow{\Key{scale} & \Key{name} & --- & \ApiText{The scale the entries come from.}{Die Skala, aus der die Einträge stammen.}}
  \SemioTableRow{\Key{kind} & \Key{choice} & \Key{swatch} & \ApiText{\Key{swatch}, \Key{line}, \Key{symbol}, \Key{gradient}, \Key{size} or \Key{pattern}.}{\Key{swatch}, \Key{line}, \Key{symbol}, \Key{gradient}, \Key{size} oder \Key{pattern}.}}
  \SemioTableRow{\Key{title} & \Key{text} & --- & \ApiText{Legend title.}{Titel der Legende.}}
  \SemioTableRow{\Key{orient} & \Key{choice} & \Key{vertical} & \ApiText{\Key{horizontal} or \Key{vertical}.}{\Key{horizontal} oder \Key{vertical}.}}
  \SemioTableRow{\Key{columns} & \Key{integer} & \Key{1} & \ApiText{Entries per row or per column block.}{Einträge je Zeile bzw. Spaltenblock.}}
  \SemioTableRow{\Key{at}, \Key{anchor} & \Key{x,y}, \Key{anchor} & ---, \Key{north west} & \ApiText{Placement; empty \Key{at} is the frame's top-right corner.}{Platzierung; leeres \Key{at} ist die obere rechte Rahmenecke.}}
  \SemioTableRow{\Key{itemGap}, \Key{swatchSize}, \Key{labelGap} & \Key{mm} & \Key{1.2}, \Key{2.6}, \Key{1.2} & \ApiText{Gap between entries, swatch edge length, gap to the label.}{Abstand zwischen Einträgen, Kantenlänge des Farbfelds, Abstand zur Beschriftung.}}
  \SemioTableRow{\Key{format} & \Key{spec} & --- & \ApiText{Format specifier for numeric entries.}{Formatangabe für numerische Einträge.}}
  \SemioTableRow{\Key{ticks}, \Key{length} & \Key{integer}, \Key{mm} & \Key{5}, \Key{24} & \ApiText{Entry count and bar length of gradient and size legends.}{Eintragszahl und Balkenlänge von Verlaufs- und Größenlegenden.}}
  \SemioTableRow{\Key{symbol} & \Key{choice} & \Key{circle} & \ApiText{\Key{circle}, \Key{square}, \Key{triangle} or \Key{diamond}.}{\Key{circle}, \Key{square}, \Key{triangle} oder \Key{diamond}.}}
}

\ApiText{%
  \Cs{SemioVizGrid} is the short form: it draws grid lines at the ticks of the named \Key{x} and
  \Key{y} scales and takes \Key{minor}, \Key{ticks} and \Key{style}.%
}{%
  \Cs{SemioVizGrid} ist die Kurzform: Er zeichnet Gitterlinien an den Teilstrichen der benannten
  Skalen \Key{x} und \Key{y} und nimmt \Key{minor}, \Key{ticks} und \Key{style}.%
}

\begin{VizFigure}[title={\ApiText{Axis, grid and legend}{Achse, Gitter und Legende}}, width=90, height=48]
\SemioVizScale{apig}{linear}{0, 10}{8, 70}[nice, ticks=5]
\SemioVizScale{apigy}{linear}{0, 10}{6, 38}[nice, ticks=4]
\SemioVizGrid[x=apig, y=apigy, minor=true]
\SemioVizAxis[scale=apig, orient=bottom, title={x}]
\SemioVizAxis[scale=apigy, orient=left, title={y}]
\SemioVizLegend[kind=gradient, scale=apig, title={value}]
\end{VizFigure}

\chapter{\ApiText{Annotation and labels}{Anmerkungen und Beschriftungen}}

\ApiPkg{semio-viz-annotation, semio-viz-label}
\ApiSig{\Cs{SemioVizAnnotate}\{\meta{kind}\}[\meta{options}]}
\ApiSig{\Cs{SemioVizLabel}[\meta{options}] \ \Cs{SemioVizLabelSeries}[\meta{options}] \ \Cs{SemioVizLabelsPlace}[\meta{options}]}

\ApiText{%
  An annotation is anchored either in plot millimetres (\Key{at}) or in data units (\Key{x},
  \Key{y}), the latter mapped through the named scales, so a note stays on its datum when the scale
  changes. Eleven kinds exist: \Key{label}, \Key{callout}, \Key{leader}, \Key{bracket}, \Key{brace},
  \Key{highlight}, \Key{reference-line}, \Key{reference-band}, \Key{reference-point},
  \Key{event-marker} and \Key{threshold}.%
}{%
  Eine Anmerkung wird entweder in Zeichnungsmillimetern (\Key{at}) oder in Dateneinheiten (\Key{x},
  \Key{y}) verankert, letztere über die benannten Skalen abgebildet; eine Notiz bleibt so bei ihrem
  Datum, wenn sich die Skala ändert. Es gibt elf Arten: \Key{label}, \Key{callout}, \Key{leader},
  \Key{bracket}, \Key{brace}, \Key{highlight}, \Key{reference-line}, \Key{reference-band},
  \Key{reference-point}, \Key{event-marker} und \Key{threshold}.%
}

\SemioTableLong[text-size=8pt]{\ApiText{Annotation options}{Anmerkungsoptionen}}{0.18,0.14,0.16,0.52}{\ApiKey & \ApiType & \ApiDefault & \ApiMeaning}{%
  \SemioTableRow{\Key{at} & \Key{x,y} & --- & \ApiText{Anchor in plot millimetres; wins over the data coordinate.}{Anker in Zeichnungsmillimetern; hat Vorrang vor der Datenkoordinate.}}
  \SemioTableRow{\Key{to}, \Key{from} & \Key{x,y} & ---, \Key{at} & \ApiText{Second and first point: callout text position, leader end, bracket ends.}{Zweiter und erster Punkt: Textposition der Sprechblase, Ende der Führungslinie, Enden der Klammer.}}
  \SemioTableRow{\Key{x}, \Key{y}, \Key{x2}, \Key{y2} & \Key{data} & --- & \ApiText{Anchor and second corner in data units.}{Anker und zweite Ecke in Dateneinheiten.}}
  \SemioTableRow{\Key{xScale}, \Key{yScale} & \Key{name} & \Key{x}, \Key{y} & \ApiText{The scales the data coordinates are mapped through.}{Die Skalen, über die die Datenkoordinaten abgebildet werden.}}
  \SemioTableRow{\Key{text}, \Key{width} & \Key{text}, \Key{mm} & --- & \ApiText{Annotation text and its wrapping width.}{Anmerkungstext und seine Umbruchbreite.}}
  \SemioTableRow{\Key{anchor}, \Key{offset}, \Key{rotate} & \Key{anchor}, \Key{dx,dy}, \Key{deg} & \Key{south west}, \Key{0,0}, \Key{0} & \ApiText{Placement of the text node.}{Platzierung des Textknotens.}}
  \SemioTableRow{\Key{shape}, \Key{points} & \Key{choice}, \Key{list} & \Key{rect}, --- & \ApiText{\Key{rect}, \Key{ellipse} or \Key{polygon} for a highlight, with its outline.}{\Key{rect}, \Key{ellipse} oder \Key{polygon} für eine Hervorhebung, mit ihrem Umriss.}}
  \SemioTableRow{\Key{depth}, \Key{size} & \Key{mm} & \Key{1.6}, \Key{1} & \ApiText{Bracket arm length or brace amplitude, and marker radius.}{Armlänge der Klammer bzw. Auslenkung der geschweiften Klammer, und Markerradius.}}
  \SemioTableRow{\Key{color} & \Key{colour} & \ApiText{theme role}{Themenrolle} & \ApiText{Empty takes the annotation role from the theme.}{Leer nimmt die Anmerkungsrolle aus dem Thema.}}
  \SemioTableRow{\Key{fillOpacity} & \Key{number} & \Key{0.12} & \ApiText{Fill opacity of highlights and bands.}{Deckkraft der Füllung von Hervorhebungen und Bändern.}}
  \SemioTableRow{\Key{dash}, \Key{arrow}, \Key{span} & \Key{boolean} & \Key{false}, \Key{true}, \Key{true} & \ApiText{Dashed stroke, arrow tip, and whether a reference line spans the whole plot area.}{Gestrichelte Linie, Pfeilspitze, und ob eine Referenzlinie die ganze Zeichenfläche überspannt.}}
}

\ApiText{%
  Labels are different: they are queued, not drawn. \Cs{SemioVizLabel} and
  \Cs{SemioVizLabelSeries} add candidates, and \Cs{SemioVizLabelsPlace} then solves the whole set at
  once, displacing on collision, drawing a leader where a label had to move and hiding a label that
  found no room. That is why direct labelling is safe on a dense chart.%
}{%
  Beschriftungen sind anders: Sie werden eingereiht, nicht gezeichnet. \Cs{SemioVizLabel} und
  \Cs{SemioVizLabelSeries} fügen Kandidaten hinzu, und \Cs{SemioVizLabelsPlace} löst dann die ganze
  Menge auf einmal auf, verschiebt bei Überschneidung, zeichnet eine Führungslinie, wo eine
  Beschriftung ausweichen musste, und blendet eine aus, die keinen Platz fand. Deshalb ist direkte
  Beschriftung auch auf einem dichten Diagramm sicher.%
}

\SemioTableLong[text-size=8pt]{\ApiText{Label options}{Beschriftungsoptionen}}{0.18,0.14,0.16,0.52}{\ApiKey & \ApiType & \ApiDefault & \ApiMeaning}{%
  \SemioTableRow{\Key{at}, \Key{x}, \Key{y} & \Key{x,y}, \Key{data} & --- & \ApiText{Anchor in millimetres or in data units.}{Anker in Millimetern oder in Dateneinheiten.}}
  \SemioTableRow{\Key{xScale}, \Key{yScale} & \Key{name} & \Key{x}, \Key{y} & \ApiText{The scales the data coordinates are mapped through.}{Die Skalen, über die die Datenkoordinaten abgebildet werden.}}
  \SemioTableRow{\Key{text} & \Key{text} & --- & \ApiText{Label text; empty on a value label means the formatted value.}{Beschriftungstext; leer bei einer Wertbeschriftung bedeutet den formatierten Wert.}}
  \SemioTableRow{\Key{value}, \Key{labelFormat} & \Key{number}, \Key{spec} & --- & \ApiText{The value and the specifier it is rendered through.}{Der Wert und die Formatangabe, mit der er gesetzt wird.}}
  \SemioTableRow{\Key{priority} & \Key{integer} & \Key{0} & \ApiText{Higher wins a collision.}{Höher gewinnt eine Überschneidung.}}
  \SemioTableRow{\Key{anchor} & \Key{choice} & \Key{auto} & \ApiText{\Key{auto} places the label away from the plot centre.}{\Key{auto} setzt die Beschriftung von der Zeichnungsmitte weg.}}
  \SemioTableRow{\Key{offset} & \Key{dx,dy} & \Key{0,1} & \ApiText{Shift from the anchor before collision solving.}{Verschiebung vom Anker vor der Kollisionsauflösung.}}
  \SemioTableRow{\Key{leader} & \Key{boolean} & \Key{true} & \ApiText{Draws a leader line when the label had to be displaced.}{Zeichnet eine Führungslinie, wenn die Beschriftung ausweichen musste.}}
  \SemioTableRow{\Key{padding}, \Key{step}, \Key{tries} & \Key{mm}, \Key{mm}, \Key{integer} & \Key{0.4}, \Key{1}, \Key{12} & \ApiText{Enforced gap, displacement increment, and candidates tried before hiding.}{Erzwungener Abstand, Verschiebungsschritt und Anzahl Versuche vor dem Ausblenden.}}
  \SemioTableRow{\Key{color} & \Key{colour} & \ApiText{theme role}{Themenrolle} & \ApiText{Text colour.}{Textfarbe.}}
}

\begin{VizFigure}[title={\ApiText{Annotated bars}{Beschriftete Balken}}, width=90, height=48]
\SemioVizPlot[data=demo, mark=bar, x=cat, y=val, guide={axis=x, axis=y}]
\SemioVizAnnotate{reference-line}[y=6, yScale=plot-y, dash=true, text={mean}]
\SemioVizAnnotate{callout}[at={30,30}, to={50,38}, text={\ApiText{A callout}{Eine Sprechblase}}]
\end{VizFigure}

\chapter{\ApiText{Facet}{Facettierung}}

\ApiPkg{semio-viz-facet}
\ApiSig{\Cs{begin}\{VizFacet\}[wrap=, row=, column=, columns=, sharedX=, sharedY=, gap=]}
\ApiSig{\Cs{SemioVizFacet}[\meta{keys}]\{\meta{body}\}}
\ApiSig{\Cs{SemioVizFacetValue} \ \Cs{SemioVizFacetIndex} \ \Cs{SemioVizFacetAxes}[\meta{keys}]}

\ApiText{%
  A facet repeats its body once per group on a computed panel grid inside the current figure. The
  body reads the current group through \Cs{SemioVizFacetValue}, so the same recipe draws every
  panel. Shared axes are drawn only where they are needed, which is what makes small multiples
  readable.%
}{%
  Eine Facettierung wiederholt ihren Rumpf einmal je Gruppe auf einem berechneten Kachelraster
  innerhalb der aktuellen Abbildung. Der Rumpf liest die aktuelle Gruppe über
  \Cs{SemioVizFacetValue}, dasselbe Rezept zeichnet also jede Kachel. Geteilte Achsen werden nur dort
  gezeichnet, wo sie gebraucht werden, und genau das macht Kleinserien lesbar.%
}

\SemioTableLong[text-size=8pt]{\ApiText{Facet options}{Facettierungsoptionen}}{0.18,0.14,0.16,0.52}{\ApiKey & \ApiType & \ApiDefault & \ApiMeaning}{%
  \SemioTableRow{\Key{wrap} & \Key{clist} & --- & \ApiText{Group values laid out left to right and wrapped into \Key{columns}.}{Gruppenwerte, von links nach rechts gelegt und auf \Key{columns} umgebrochen.}}
  \SemioTableRow{\Key{row}, \Key{column} & \Key{clist} & --- & \ApiText{Group values, one panel row or one panel column each.}{Gruppenwerte, je eine Kachelzeile bzw. Kachelspalte.}}
  \SemioTableRow{\Key{columns} & \Key{integer} & \Key{3} & \ApiText{Panels per row for \Key{wrap}.}{Kacheln je Zeile bei \Key{wrap}.}}
  \SemioTableRow{\Key{sharedX}, \Key{sharedY} & \Key{boolean} & \Key{true} & \ApiText{Draws the x axis only on the bottom row, the y axis only on the leftmost column.}{Zeichnet die x-Achse nur in der untersten Zeile, die y-Achse nur in der linken Spalte.}}
  \SemioTableRow{\Key{gap} & \Key{mm} & \Key{4} & \ApiText{Gap between panels.}{Abstand zwischen Kacheln.}}
  \SemioTableRow{\Key{header}, \Key{headerGap} & \Key{boolean}, \Key{mm} & \Key{true}, \Key{3.4} & \ApiText{The group value above each panel, and the height reserved for it.}{Der Gruppenwert über jeder Kachel und die dafür vorgesehene Höhe.}}
  \SemioTableRow{\Key{width}, \Key{height} & \Key{mm} & \ApiText{the figure}{die Abbildung} & \ApiText{Total size of the facet block.}{Gesamtgröße des Facettenblocks.}}
  \SemioTableRow{\Key{xScale}, \Key{yScale} & \Key{name} & \Key{x}, \Key{y} & \ApiText{The scales re-ranged to each panel's extent.}{Die Skalen, die auf die Ausdehnung jeder Kachel umgespannt werden.}}
  \SemioTableRow{\Key{pad} & \Key{mm} & \Key{6} & \ApiText{Plot inset inside each panel.}{Zeichnungseinzug innerhalb jeder Kachel.}}
}

\chapter{\ApiText{Composition}{Komposition}}

\ApiPkg{semio-viz, semio-viz-composition}

\ApiText{%
  Composition is taxonomy section 77: how several drawings share one page. The rule throughout is
  that a block declares its own size and the container places it, so no recipe ever contains a
  literal coordinate for what sits beside what.%
}{%
  Komposition ist Taxonomieabschnitt 77: wie sich mehrere Zeichnungen eine Seite teilen. Durchgehend
  gilt: Ein Block erklärt seine eigene Größe und der Behälter platziert ihn, kein Rezept enthält also
  je eine wörtliche Koordinate dafür, was neben was liegt.%
}

\SemioTableLong[text-size=8pt]{\ApiText{Composition blocks}{Kompositionsblöcke}}{0.28,0.72}{\ApiText{Block}{Block} & \ApiMeaning}{%
  \SemioTableRow{\Key{VizFigure} & \ApiText{The frame: a figure window with a millimetre canvas.}{Der Rahmen: ein Abbildungsfenster mit Millimeterleinwand.}}
  \SemioTableRow{\Key{VizSection} & \ApiText{A titled band, auto-stacked below the previous one.}{Ein betitelter Streifen, automatisch unter den vorigen gestapelt.}}
  \SemioTableRow{\Key{VizColumn} & \ApiText{A column of its own declared width inside a section, with an optional note.}{Eine Spalte eigener erklärter Breite innerhalb eines Streifens, mit optionaler Notiz.}}
  \SemioTableRow{\Key{VizLayer} & \ApiText{One drawing layer: same frame, same scales, its own style and opacity.}{Eine Zeichenebene: gleicher Rahmen, gleiche Skalen, eigener Stil und eigene Deckkraft.}}
  \SemioTableRow{\Key{VizOverlay}, \Cs{SemioVizOverlay} & \ApiText{Superimposed layers sharing one frame.}{Übereinanderliegende Ebenen, die sich einen Rahmen teilen.}}
  \SemioTableRow{\Key{VizConcat}, \Cs{SemioVizConcatItem} & \ApiText{Items placed side by side or stacked; keys \Key{direction} (default \Key{horizontal}) and \Key{gap} (default 4).}{Elemente nebeneinander oder gestapelt; Schlüssel \Key{direction} (Vorgabe \Key{horizontal}) und \Key{gap} (Vorgabe 4).}}
  \SemioTableRow{\Cs{SemioVizInset} & \ApiText{A sub-frame that optionally marks and links the host region it magnifies; keys \Key{at}, \Key{zoom}, \Key{width}, \Key{height}, \Key{pad}, \Key{frame}, \Key{link}.}{Ein Teilrahmen, der die vergrößerte Wirtsregion auf Wunsch markiert und verbindet; Schlüssel \Key{at}, \Key{zoom}, \Key{width}, \Key{height}, \Key{pad}, \Key{frame}, \Key{link}.}}
  \SemioTableRow{\Key{VizDashboard}, \Cs{SemioVizCell} & \ApiText{A fixed cell grid; keys \Key{columns} (3), \Key{rows} (2), \Key{gap} (3), and per cell \Key{colspan}, \Key{rowspan}.}{Ein festes Zellenraster; Schlüssel \Key{columns} (3), \Key{rows} (2), \Key{gap} (3) und je Zelle \Key{colspan}, \Key{rowspan}.}}
  \SemioTableRow{\Cs{SemioVizCrossLine} & \ApiText{A reference line across the whole composition, in composition millimetres.}{Eine Referenzlinie quer über die ganze Komposition, in Kompositionsmillimetern.}}
  \SemioTableRow{\Cs{SemioVizCellGrid}, \Cs{SemioVizSpanRows}, \Cs{SemioVizStageChain} & \ApiText{Table-driven blocks: a grid of labelled cells, one labelled span row per entry, and a source chip followed by a chain of stage boxes.}{Tabellengetriebene Blöcke: ein Raster beschrifteter Zellen, je Eintrag eine beschriftete Spannzeile, und ein Quellenchip gefolgt von einer Kette von Stufenkästen.}}
  \SemioTableRow{\Cs{SemioVizNoteBelow} & \ApiText{Caption text placed below a captured anchor, so a note follows the drawing's real extent.}{Bildunterschrift unter einem festgehaltenen Anker, eine Notiz folgt so der tatsächlichen Ausdehnung der Zeichnung.}}
}

\begin{VizFigure}[title={\ApiText{Sections, columns and a note}{Streifen, Spalten und eine Notiz}}, width=100]
\begin{VizSection}[title={\ApiText{A \ One section}{A \ Ein Streifen}}]
\SemioVizChart{icon-array}[columns=5]
\end{VizSection}
\begin{VizSection}[title={\ApiText{B \ Two columns}{B \ Zwei Spalten}}]
\begin{VizColumn}[width=48, title={\ApiText{Left}{Links}}]
\SemioVizText{label}[width=40]{0,0}{\ApiText{A text mark inside a column.}{Eine Textmarke in einer Spalte.}}
\end{VizColumn}
\begin{VizColumn}[width=48, title={\ApiText{Right}{Rechts}}, note={\ApiText{Caption text before the chart.}{Bildunterschrift vor dem Diagramm.}}]
\SemioVizChart{horizontal-bar-chart}[axes=false]
\coordinate (api-bar-end) at (0, 0);
\SemioVizNoteBelow{api-bar-end}{48}{\ApiText{A note placed from the chart's own drawn extent.}{Eine Notiz, aus der tatsächlichen Ausdehnung des Diagramms heraus platziert.}}
\end{VizColumn}
\end{VizSection}
\end{VizFigure}

\chapter{\ApiText{Theme}{Thema}}

\ApiPkg{semio-viz-theme}
\ApiSig{\Cs{SemioVizTheme}\{\meta{name}\}[\meta{keys}] \ \Cs{SemioVizThemeSet}[\meta{keys}]}

\ApiText{%
  Every renderer of the library reads its colours, strokes and fonts here and nowhere else, so a
  document changes its whole visual language through these keys and no figure has to be edited. Three
  named presets exist: \Key{default}, \Key{grayscale} and \Key{pattern}; \Cs{SemioVizThemeSet}
  overrides individual keys without changing the preset.%
}{%
  Jeder Renderer der Bibliothek liest seine Farben, Striche und Schriften hier und sonst nirgends;
  ein Dokument ändert seine ganze Bildsprache also über diese Schlüssel, und keine Abbildung muss
  angefasst werden. Es gibt drei benannte Voreinstellungen: \Key{default}, \Key{grayscale} und
  \Key{pattern}; \Cs{SemioVizThemeSet} überschreibt einzelne Schlüssel, ohne die Voreinstellung zu
  wechseln.%
}

\SemioTableLong[text-size=8pt]{\ApiText{Theme options}{Themenoptionen}}{0.18,0.14,0.16,0.52}{\ApiKey & \ApiType & \ApiDefault & \ApiMeaning}{%
  \SemioTableRow{\Key{appearance} & \Key{light|dark} & \ApiText{document theme}{Dokumentthema} & \ApiText{An explicit key wins; otherwise the document theme decides.}{Ein ausdrücklicher Schlüssel gewinnt; sonst entscheidet das Dokumentthema.}}
  \SemioTableRow{\Key{palette} & \Key{name} & \Key{presence} & \ApiText{Categorical palette: \Key{presence} or \Key{brand}.}{Kategoriale Palette: \Key{presence} oder \Key{brand}.}}
  \SemioTableRow{\Key{scheme} & \Key{name} & \Key{primary} & \ApiText{The sequential or diverging ramp.}{Die sequentielle oder divergierende Rampe.}}
  \SemioTableRow{\Key{interpolator} & \Key{name} & \Key{oklab} & \ApiText{Colour space the ramp is sampled in.}{Farbraum, in dem die Rampe abgetastet wird.}}
  \SemioTableRow{\Key{grayscale} & \Key{boolean} & \Key{false} & \ApiText{Replaces the hues by the gray ramp.}{Ersetzt die Farbtöne durch die Graurampe.}}
  \SemioTableRow{\Key{patterns} & \Key{boolean} & \Key{false} & \ApiText{Pairs every categorical slot with a hatch, so a figure survives being photocopied.}{Verbindet jeden kategorialen Platz mit einer Schraffur, eine Abbildung übersteht so das Kopieren.}}
  \SemioTableRow{\Key{stroke} & \Key{name} & \Key{default} & \ApiText{\Key{hairline}, \Key{default} or \Key{focus}, from the design tokens.}{\Key{hairline}, \Key{default} oder \Key{focus}, aus den Gestaltungstoken.}}
  \SemioTableRow{\Key{font} & \Key{name} & \Key{sans} & \ApiText{\Key{sans}, \Key{serif} or \Key{mono}, from \Key{semio-fonts}.}{\Key{sans}, \Key{serif} oder \Key{mono}, aus \Key{semio-fonts}.}}
}

\begin{VizFigure}[title={\ApiText{The same chart, grayscale and patterned}{Dasselbe Diagramm, Graustufen und Schraffur}}, width=90, height=44]
\SemioVizThemeSet[grayscale=true, patterns=true]
\SemioVizPlot[data=demo, mark=bar, x=cat, y=val, guide={axis=x, axis=y}]
\SemioVizThemeSet[grayscale=false, patterns=false]
\end{VizFigure}

\chapter{\ApiText{The plot grammar}{Die Plot-Grammatik}}

\ApiPkg{semio-viz-plot}
\ApiSig{\Cs{SemioVizPlot}[data=, mark=, coordinate=, \meta{encodings}, transform=, layout=, facet=, guide=, annotation=, theme=]}
\ApiSig{\Cs{begin}\{VizPlot\}[\meta{shared defaults}] \ \Cs{SemioVizLayer}[\meta{keys}] … \ \Cs{end}\{VizPlot\}}

\ApiText{%
  \Cs{SemioVizPlot} is the whole grammar in one call: it reads a table, runs the declared transform
  and layout on it, installs a coordinate system, binds the encoding channels to columns and scales,
  draws a mark per row and then the guides and annotations. The environment form is the layered
  version: its keys are the shared defaults and each \Cs{SemioVizLayer} overrides what it needs. The
  two positional scales are registered as \Key{plot-x} and \Key{plot-y}, so a hand-written
  \Cs{SemioVizAxis}\Key{[scale=plot-x]} matches the marks exactly.%
}{%
  \Cs{SemioVizPlot} ist die ganze Grammatik in einem Aufruf: Er liest eine Tabelle, führt die
  erklärte Transformation und das Layout darauf aus, installiert ein Koordinatensystem, bindet die
  Kodierungskanäle an Spalten und Skalen, zeichnet je Zeile eine Marke und danach die Hilfselemente
  und Anmerkungen. Die Umgebungsform ist die geschichtete Fassung: Ihre Schlüssel sind die geteilten
  Vorgaben, und jede \Cs{SemioVizLayer} überschreibt, was sie braucht. Die beiden Positionsskalen
  sind als \Key{plot-x} und \Key{plot-y} registriert, eine von Hand geschriebene
  \Cs{SemioVizAxis}\Key{[scale=plot-x]} passt also genau zu den Marken.%
}

\SemioTableLong[text-size=8pt]{\ApiText{Plot keys}{Plot-Schlüssel}}{0.18,0.14,0.16,0.52}{\ApiKey & \ApiType & \ApiDefault & \ApiMeaning}{%
  \SemioTableRow{\Key{data} & \Key{name} & \Key{demo} & \ApiText{The table to read.}{Die zu lesende Tabelle.}}
  \SemioTableRow{\Key{mark} & \Key{name} & \Key{point} & \ApiText{\Key{bar}, \Key{line}, \Key{area} or any section 0 mark kind.}{\Key{bar}, \Key{line}, \Key{area} oder jede Markenart aus Abschnitt 0.}}
  \SemioTableRow{\Key{coordinate} & \Key{name} & \Key{cartesian} & \ApiText{Any coordinate kind.}{Jede Koordinatenart.}}
  \SemioTableRow{\Key{transform} & \Key{list} & --- & \ApiText{A \Cs{SemioVizTransform} option list, run on the plot's own table.}{Eine Optionsliste für \Cs{SemioVizTransform}, auf der eigenen Tabelle des Plots ausgeführt.}}
  \SemioTableRow{\Key{layout} & \Key{name} & --- & \ApiText{A registered layout algorithm, run after the transform.}{Ein registrierter Layoutalgorithmus, nach der Transformation ausgeführt.}}
  \SemioTableRow{\Key{facet} & \Key{list} & --- & \ApiText{A facet option list; the plot becomes small multiples.}{Eine Facettierungsliste; der Plot wird zur Kleinserie.}}
  \SemioTableRow{\Key{guide} & \Key{list} & --- & \ApiText{\Key{axis=x}, \Key{axis=y}, \Key{grid}, \Key{legend}.}{\Key{axis=x}, \Key{axis=y}, \Key{grid}, \Key{legend}.}}
  \SemioTableRow{\Key{annotation} & \Key{list} & --- & \ApiText{Annotations drawn after the marks.}{Anmerkungen, nach den Marken gezeichnet.}}
  \SemioTableRow{\Key{theme} & \Key{list} & --- & \ApiText{Theme keys for this plot only.}{Themenschlüssel nur für diesen Plot.}}
  \SemioTableRow{\Key{curve} & \Key{name} & \Key{linear} & \ApiText{Interpolator of a line or area mark.}{Interpolator einer Linien- oder Flächenmarke.}}
}

\ApiText{%
  Everything else in the option list is an encoding channel. Sixteen channels exist, each taking a
  column name or a braced pair \Key{\{column, scale=<name>\}} that binds the channel to a named
  scale instead of the automatic one:%
}{%
  Alles Übrige in der Optionsliste ist ein Kodierungskanal. Es gibt sechzehn Kanäle, jeder nimmt
  einen Spaltennamen oder ein geklammertes Paar \Key{\{Spalte, scale=<Name>\}}, das den Kanal an eine
  benannte Skala statt an die automatische bindet:%
}

\SemioTableLong[text-size=8pt]{\ApiText{Encoding channels}{Kodierungskanäle}}{0.18,0.82}{\ApiText{Channel}{Kanal} & \ApiMeaning}{%
  \SemioTableRow{\Key{x}, \Key{y} & \ApiText{The two positional channels of the active coordinate system.}{Die beiden Positionskanäle des aktiven Koordinatensystems.}}
  \SemioTableRow{\Key{x2}, \Key{y2} & \ApiText{The second position: bar extents, ranges, bands, rules.}{Die zweite Position: Balkenausdehnungen, Spannen, Bänder, Linien.}}
  \SemioTableRow{\Key{angle}, \Key{radius} & \ApiText{The polar pair, when the coordinate system is polar.}{Das Polarpaar, wenn das Koordinatensystem polar ist.}}
  \SemioTableRow{\Key{size} & \ApiText{Mark area; a square-root scale is the correct default.}{Markenfläche; eine Wurzelskala ist die richtige Vorgabe.}}
  \SemioTableRow{\Key{shape} & \ApiText{Symbol type per category.}{Symboltyp je Kategorie.}}
  \SemioTableRow{\Key{fill}, \Key{stroke} & \ApiText{Colour channels, through a categorical palette or a ramp.}{Farbkanäle, über eine kategoriale Palette oder eine Rampe.}}
  \SemioTableRow{\Key{opacity} & \ApiText{Draw opacity per row.}{Deckkraft je Zeile.}}
  \SemioTableRow{\Key{text} & \ApiText{The column a text mark or a value label prints.}{Die Spalte, die eine Textmarke oder eine Wertbeschriftung ausgibt.}}
  \SemioTableRow{\Key{dash} & \ApiText{Dash pattern per category --- the grayscale-safe partner of \Key{stroke}.}{Strichmuster je Kategorie --- der graustufensichere Partner von \Key{stroke}.}}
  \SemioTableRow{\Key{width} & \ApiText{Mark width where it is not derived from a band scale.}{Markenbreite, wo sie nicht aus einer Bandskala folgt.}}
  \SemioTableRow{\Key{order} & \ApiText{Drawing and stacking order.}{Zeichen- und Stapelreihenfolge.}}
  \SemioTableRow{\Key{detail} & \ApiText{Splits the data into series without giving them a visual channel.}{Teilt die Daten in Reihen, ohne ihnen einen Sichtkanal zu geben.}}
}

\begin{VizFigure}[title={\ApiText{One plot, three layers}{Ein Plot, drei Ebenen}}, width=90, height=48]
\begin{VizPlot}[data=demo, coordinate=cartesian, x=cat, y=val, guide={axis=x, axis=y}]
\SemioVizLayer[mark=bar]
\SemioVizLayer[mark=line, curve=monotone-x]
\SemioVizLayer[mark=point, size=val]
\end{VizPlot}
\end{VizFigure}

\chapter{\ApiText{Chart kinds and families}{Diagrammarten und Familien}}

\ApiPkg{semio-viz-catalog, semio-viz-catalog-labels, semio-viz-family}
\ApiSig{\Cs{SemioVizChart}\{\meta{slug}\}[\meta{overrides}]}
\ApiSig{\Cs{SemioVizChartKind}\{\meta{slug}\}\{\meta{family}\}\{\meta{options}\}}
\ApiSig{\Cs{SemioVizKindCovers}\{\meta{slug}\}\{\meta{leaf ids}\} \ \Cs{SemioVizKindTitle}\{\meta{slug}\} \ \Cs{SemioVizKindLabel}\{\meta{slug}\}\{en\}\{de\}}
\ApiSig{\Cs{SemioVizFamily}\{\meta{name}\}\{\meta{code}\} \ \Cs{SemioVizRunFamily}\{\meta{name}\}[\meta{options}] \ \Cs{SemioVizDemo}\{\meta{name}\}}
\ApiSig{\Cs{SemioVizLocalized}\{en\}\{de\}}

\ApiText{%
  A chart kind is a preset, not code. The handcrafted catalogue
  \Key{\E{🖼}assets/\E{🔣}viz-catalog.json} names, for every kind, its family, its option list, its demo
  table, its title in both languages and the taxonomy leaves it covers; the generator turns that into
  \Key{semio-viz-catalog.sty} (one \Cs{SemioVizChartKind} and one \Cs{SemioVizKindCovers} per kind)
  and \Key{semio-viz-catalog-labels.sty} (one \Cs{SemioVizKindLabel} per kind). Neither file is ever
  edited by hand.%
}{%
  Eine Diagrammart ist eine Voreinstellung, kein Code. Der handgefertigte Katalog
  \Key{\E{🖼}assets/\E{🔣}viz-catalog.json} nennt zu jeder Art ihre Familie, ihre Optionsliste, ihre
  Demotabelle, ihren Titel in beiden Sprachen und die Taxonomieblätter, die sie abdeckt; der Generator
  macht daraus \Key{semio-viz-catalog.sty} (je Art ein \Cs{SemioVizChartKind} und ein
  \Cs{SemioVizKindCovers}) und \Key{semio-viz-catalog-labels.sty} (je Art ein
  \Cs{SemioVizKindLabel}). Keine der beiden Dateien wird je von Hand bearbeitet.%
}

\ApiText{%
  \Cs{SemioVizChart}\{\meta{slug}\} draws the preset; the optional list overrides individual family
  options for this one call, which is how a gallery figure differs from a document figure without a
  second catalogue entry. \Cs{SemioVizRunFamily} bypasses the catalogue and calls a family directly,
  and \Cs{SemioVizDemo} draws a mark primitive or a chart kind on its own demo table.
  \Cs{SemioVizKindTitle} prints a kind's title in the document language, and
  \Cs{SemioVizLocalized}\{en\}\{de\} is the same mechanism for any generated heading.%
}{%
  \Cs{SemioVizChart}\{\meta{Slug}\} zeichnet die Voreinstellung; die optionale Liste überschreibt
  einzelne Familienoptionen für diesen einen Aufruf --- so unterscheidet sich eine Galerieabbildung
  von einer Dokumentabbildung ohne zweiten Katalogeintrag. \Cs{SemioVizRunFamily} umgeht den Katalog
  und ruft eine Familie direkt auf, und \Cs{SemioVizDemo} zeichnet einen Markengrundbaustein oder
  eine Diagrammart auf ihrer eigenen Demotabelle. \Cs{SemioVizKindTitle} gibt den Titel einer Art in
  der Dokumentsprache aus, und \Cs{SemioVizLocalized}\{en\}\{de\} ist dasselbe für jede erzeugte
  Überschrift.%
}

\begin{VizFigure}[title={\SemioVizKindTitle{dot-plot}}, width=90, height=44]
\SemioVizChart{dot-plot}
\end{VizFigure}

\begin{VizFigure}[title={\ApiText{A mark primitive on its own demo table}{Ein Markengrundbaustein auf seiner eigenen Demotabelle}}, width=90, height=44]
\SemioVizDemo{dot}
\end{VizFigure}

%region 🔖️GeneratedNamespaces

\chapter{\ApiText{The namespaces}{Die Namensräume}}

\ApiText{%
  Taxonomy section 76 divides the library into package namespaces, and every chart kind of
  the catalogue belongs to exactly one family of exactly one namespace. A family is a
  renderer registered with \Cs{SemioVizFamily}; its options are l3keys in
  \Key{semio / viz / family / <name>} and are listed below with their type, their default
  and their meaning. Two options are universal and therefore not repeated in every table:
  \Key{variant} is the catalogue slug the family is rendering, and \Key{data} names the
  table it reads.%
}{%
  Taxonomieabschnitt 76 teilt die Bibliothek in Paketnamensräume, und jede Diagrammart des
  Katalogs gehört zu genau einer Familie genau eines Namensraums. Eine Familie ist ein mit
  \Cs{SemioVizFamily} registrierter Renderer; ihre Optionen sind l3keys in
  \Key{semio / viz / family / <Name>} und stehen unten mit Typ, Vorgabe und Bedeutung. Zwei
  Optionen sind allgemein und werden deshalb nicht in jeder Tabelle wiederholt:
  \Key{variant} ist der Katalog-Slug, den die Familie gerade zeichnet, und \Key{data}
  benennt die Tabelle, die sie liest.%
}

\section{\ApiText{Kernel}{Kern}}

\ApiText{The grammar itself: the families that expose a kernel capability — a scale, an axis, a mark, a shape, a layout algorithm — as a chart kind of its own, so the catalogue covers taxonomy sections 74 to 79 the same way it covers a bar chart.}{Die Grammatik selbst: die Familien, die eine Kernfähigkeit — eine Skala, eine Achse, eine Marke, eine Form, einen Layoutalgorithmus — als eigene Diagrammart anbieten, damit der Katalog die Taxonomieabschnitte 74 bis 79 genauso abdeckt wie ein Balkendiagramm.}

\subsection{\Key{axis}}

\SemioTableLong[text-size=8pt]{\Key{axis}}{0.20,0.13,0.15,0.52}{\ApiKey & \ApiType & \ApiDefault & \ApiMeaning}{%
  \SemioTableRow{\Key{x} & \Key{string} & --- & \ApiText{Column mapped onto the horizontal axis.}{Spalte, die auf die waagerechte Achse abgebildet wird.}}
  \SemioTableRow{\Key{y} & \Key{string} & --- & \ApiText{Column mapped onto the vertical axis.}{Spalte, die auf die senkrechte Achse abgebildet wird.}}
  \SemioTableRow{\Key{y2} & \Key{string} & --- & \ApiText{Column holding the upper bound of a span.}{Spalte mit der oberen Grenze einer Spanne.}}
  \SemioTableRow{\Key{ticks} & \Key{integer} & --- & \ApiText{Target tick count for continuous scales, default 5.}{Angestrebte Teilstrichzahl stetiger Skalen.}}
  \SemioTableRow{\Key{grid} & \Key{boolean} & --- & \ApiText{Draw a grid line at every major tick, default false.}{An jedem Hauptteilstrich eine Gitterlinie zeichnen.}}
  \SemioTableRow{\Key{minor} & \Key{boolean} & --- & \ApiText{Draw minor ticks, default false.}{Zwischenteilstriche zeichnen.}}
  \SemioTableRow{\Key{title} & \Key{string} & --- & \ApiText{Axis title.}{Titel der Achse.}}
  \SemioTableRow{\Key{broken} & \Key{string} & --- & \ApiText{Data interval removed from the axis, default empty.}{Datenintervall, das aus der Achse herausgenommen wird.}}
}

\subsection{\Key{composition}}

\SemioTableLong[text-size=8pt]{\Key{composition}}{0.20,0.13,0.15,0.52}{\ApiKey & \ApiType & \ApiDefault & \ApiMeaning}{%
  \SemioTableRow{\Key{x} & \Key{string} & --- & \ApiText{Column mapped onto the horizontal axis.}{Spalte, die auf die waagerechte Achse abgebildet wird.}}
  \SemioTableRow{\Key{y} & \Key{string} & --- & \ApiText{Column mapped onto the vertical axis.}{Spalte, die auf die senkrechte Achse abgebildet wird.}}
  \SemioTableRow{\Key{y2} & \Key{string} & --- & \ApiText{Column holding the upper bound of a span.}{Spalte mit der oberen Grenze einer Spanne.}}
  \SemioTableRow{\Key{groups} & \Key{string} & --- & \ApiText{Group boxes drawn behind the members.}{Gruppenrahmen hinter den Mitgliedern.}}
  \SemioTableRow{\Key{columns} & \Key{integer} & --- & \ApiText{Number of columns of the layout grid.}{Anzahl der Spalten des Anordnungsgitters.}}
  \SemioTableRow{\Key{gap} & \Key{number} & --- & \ApiText{Millimetre gap between adjacent parts.}{Millimeterabstand zwischen benachbarten Teilen.}}
  \SemioTableRow{\Key{ticks} & \Key{integer} & --- & \ApiText{Number of ticks drawn per axis.}{Anzahl der Teilstriche je Achse.}}
}

\subsection{\Key{encoding}}

\ApiText{This family takes no options beyond \Key{variant} and \Key{data}.}{Diese Familie nimmt außer \Key{variant} und \Key{data} keine Optionen.}

\subsection{\Key{figure}}

\ApiText{This family takes no options beyond \Key{variant} and \Key{data}.}{Diese Familie nimmt außer \Key{variant} und \Key{data} keine Optionen.}

\subsection{\Key{grammar}}

\ApiText{This family takes no options beyond \Key{variant} and \Key{data}.}{Diese Familie nimmt außer \Key{variant} und \Key{data} keine Optionen.}

\subsection{\Key{layout-algorithm}}

\ApiText{This family takes no options beyond \Key{variant} and \Key{data}.}{Diese Familie nimmt außer \Key{variant} und \Key{data} keine Optionen.}

\subsection{\Key{mark-annotation}}

\ApiText{This family takes no options beyond \Key{variant} and \Key{data}.}{Diese Familie nimmt außer \Key{variant} und \Key{data} keine Optionen.}

\subsection{\Key{mark-arc}}

\ApiText{This family takes no options beyond \Key{variant} and \Key{data}.}{Diese Familie nimmt außer \Key{variant} und \Key{data} keine Optionen.}

\subsection{\Key{mark-area}}

\ApiText{This family takes no options beyond \Key{variant} and \Key{data}.}{Diese Familie nimmt außer \Key{variant} und \Key{data} keine Optionen.}

\subsection{\Key{mark-connector}}

\ApiText{This family takes no options beyond \Key{variant} and \Key{data}.}{Diese Familie nimmt außer \Key{variant} und \Key{data} keine Optionen.}

\subsection{\Key{mark-line}}

\ApiText{This family takes no options beyond \Key{variant} and \Key{data}.}{Diese Familie nimmt außer \Key{variant} und \Key{data} keine Optionen.}

\subsection{\Key{mark-point}}

\ApiText{This family takes no options beyond \Key{variant} and \Key{data}.}{Diese Familie nimmt außer \Key{variant} und \Key{data} keine Optionen.}

\subsection{\Key{mark-region}}

\ApiText{This family takes no options beyond \Key{variant} and \Key{data}.}{Diese Familie nimmt außer \Key{variant} und \Key{data} keine Optionen.}

\subsection{\Key{namespace}}

\ApiText{This family takes no options beyond \Key{variant} and \Key{data}.}{Diese Familie nimmt außer \Key{variant} und \Key{data} keine Optionen.}

\subsection{\Key{scale}}

\SemioTableLong[text-size=8pt]{\Key{scale}}{0.20,0.13,0.15,0.52}{\ApiKey & \ApiType & \ApiDefault & \ApiMeaning}{%
  \SemioTableRow{\Key{domain} & \Key{string} & --- & \ApiText{Input domain 'min,max' of the scale or function.}{Eingabebereich 'min,max' der Skala bzw. Funktion.}}
  \SemioTableRow{\Key{nice} & \Key{string} & --- & \ApiText{Extend the domain to round tick values.}{Den Wertebereich auf runde Teilstrichwerte erweitern.}}
  \SemioTableRow{\Key{range} & \Key{string} & --- & \ApiText{Output range 'min,max' of the scale or sweep.}{Ausgabebereich 'min,max' der Skala bzw. des Durchlaufs.}}
}

\subsection{\Key{shape}}

\ApiText{This family takes no options beyond \Key{variant} and \Key{data}.}{Diese Familie nimmt außer \Key{variant} und \Key{data} keine Optionen.}

\subsection{\Key{transform-data}}

\ApiText{This family takes no options beyond \Key{variant} and \Key{data}.}{Diese Familie nimmt außer \Key{variant} und \Key{data} keine Optionen.}

\subsection{\Key{transform-statistical}}

\ApiText{This family takes no options beyond \Key{variant} and \Key{data}.}{Diese Familie nimmt außer \Key{variant} und \Key{data} keine Optionen.}

\section{\ApiText{Charts}{Diagramme}}

\ApiText{The cartesian and polar chart namespaces: bar, line, area, scatter, financial, timeline, dashboard, funnel, distribution, statistical, polar, matrix and table.}{Die kartesischen und polaren Diagrammnamensräume: Balken, Linie, Fläche, Streuung, Finanzen, Zeitstrahl, Dashboard, Trichter, Verteilung, Statistik, Polar, Matrix und Tabelle.}

\subsection{\Key{analytical}}

\ApiText{This family takes no options beyond \Key{variant} and \Key{data}.}{Diese Familie nimmt außer \Key{variant} und \Key{data} keine Optionen.}

\subsection{\Key{annotated}}

\SemioTableLong[text-size=8pt]{\Key{annotated}}{0.20,0.13,0.15,0.52}{\ApiKey & \ApiType & \ApiDefault & \ApiMeaning}{%
  \SemioTableRow{\Key{mode} & \Key{string} & \Key{callout} & \ApiText{How several series share one band.}{Wie sich mehrere Reihen ein Band teilen.}}
  \SemioTableRow{\Key{at} & \Key{number} & \Key{3} & \ApiText{The x position the annotation attaches to.}{Die x-Position, an der die Anmerkung ansetzt.}}
  \SemioTableRow{\Key{value} & \Key{number} & \Key{6} & \ApiText{The threshold / reference value.}{Der Schwellen- bzw. Referenzwert.}}
  \SemioTableRow{\Key{note} & \Key{string} & --- & \ApiText{The annotation text; empty picks the localized default.}{Der Text der Anmerkung; leer wählt die lokalisierte Vorgabe.}}
  \SemioTableRow{\Key{x} & \Key{string} & \Key{cat} & \ApiText{Independent / category column.}{Spalte der unabhängigen Größe bzw. der Kategorie.}}
  \SemioTableRow{\Key{y} & \Key{string} & \Key{val} & \ApiText{Value column.}{Spalte mit dem Wert.}}
  \SemioTableRow{\Key{y2} & \Key{string} & --- & \ApiText{Upper value column of span modes.}{Spalte des oberen Werts der Spannenmodi.}}
  \SemioTableRow{\Key{group} & \Key{string} & --- & \ApiText{Series / stack key column.}{Spalte des Reihen- bzw. Stapelschlüssels.}}
  \SemioTableRow{\Key{series} & \Key{string} & --- & \ApiText{Column that splits the rows into series.}{Spalte, die die Zeilen in Reihen aufteilt.}}
  \SemioTableRow{\Key{label} & \Key{string} & --- & \ApiText{Text column for direct labels.}{Textspalte für Direktbeschriftungen.}}
  \SemioTableRow{\Key{padLeft} & \Key{number} & --- & \ApiText{Plot inset inside the canvas.}{Einrückung der Zeichenfläche innerhalb der Leinwand.}}
  \SemioTableRow{\Key{padRight} & \Key{number} & --- & \ApiText{Plot inset inside the canvas.}{Einrückung der Zeichenfläche innerhalb der Leinwand.}}
  \SemioTableRow{\Key{padTop} & \Key{number} & --- & \ApiText{Plot inset inside the canvas.}{Einrückung der Zeichenfläche innerhalb der Leinwand.}}
  \SemioTableRow{\Key{padBottom} & \Key{number} & --- & \ApiText{Plot inset inside the canvas.}{Einrückung der Zeichenfläche innerhalb der Leinwand.}}
  \SemioTableRow{\Key{yMin} & \Key{string} & \Key{from the data} & \ApiText{Value-axis domain override.}{Vorgabe des Wertebereichs der Werteachse.}}
  \SemioTableRow{\Key{yMax} & \Key{string} & \Key{from the data} & \ApiText{Value-axis domain override.}{Vorgabe des Wertebereichs der Werteachse.}}
  \SemioTableRow{\Key{axes} & \Key{boolean} & \Key{True} & \ApiText{Axis lines, tick labels, legend.}{Achsenlinien, Teilstrichbeschriftungen und Legende.}}
  \SemioTableRow{\Key{grid} & \Key{string} & \Key{none} & \ApiText{Grid lines.}{Gitterlinien hinter den Marken.}}
  \SemioTableRow{\Key{labels} & \Key{string} & \Key{none} & \ApiText{Per-mark labels.}{Beschriftungen an den einzelnen Marken.}}
  \SemioTableRow{\Key{sort} & \Key{string} & \Key{none} & \ApiText{Category order by total value.}{Reihenfolge der Kategorien nach Gesamtwert.}}
  \SemioTableRow{\Key{curve} & \Key{string} & --- & \ApiText{Interpolation between points.}{Interpolation zwischen den Punkten.}}
}

\subsection{\Key{area}}

\SemioTableLong[text-size=8pt]{\Key{area}}{0.20,0.13,0.15,0.52}{\ApiKey & \ApiType & \ApiDefault & \ApiMeaning}{%
  \SemioTableRow{\Key{stack} & \Key{string} & \Key{none} & \ApiText{How series share the value axis.}{Wie sich die Reihen die Werteachse teilen.}}
  \SemioTableRow{\Key{offset} & \Key{string} & \Key{zero} & \ApiText{D3-stack offset the stack uses.}{D3-stack-Versatz, den der Stapel verwendet.}}
  \SemioTableRow{\Key{baseline} & \Key{number} & \Key{0} & \ApiText{Value the first layer sits on.}{Wert, auf dem die erste Schicht aufsitzt.}}
  \SemioTableRow{\Key{horizon} & \Key{integer} & \Key{0} & \ApiText{Bands of a horizon chart, 0 = off.}{Bänder eines Horizontdiagramms, 0 = aus.}}
  \SemioTableRow{\Key{difference} & \Key{boolean} & \Key{False} & \ApiText{Two series with a signed band.}{Zwei Reihen mit einem vorzeichenbehafteten Band.}}
  \SemioTableRow{\Key{range} & \Key{boolean} & \Key{False} & \ApiText{Lo/hi band instead of a value area.}{Lo/hi-Band statt einer Wertfläche.}}
  \SemioTableRow{\Key{lo} & \Key{string} & \Key{lo / hi} & \ApiText{The range/confidence columns.}{Die Spalten der Spanne bzw. des Konfidenzbereichs.}}
  \SemioTableRow{\Key{hi} & \Key{string} & \Key{lo / hi} & \ApiText{The range/confidence columns.}{Die Spalten der Spanne bzw. des Konfidenzbereichs.}}
  \SemioTableRow{\Key{x} & \Key{string} & \Key{cat} & \ApiText{Independent / category column.}{Spalte der unabhängigen Größe bzw. der Kategorie.}}
  \SemioTableRow{\Key{y} & \Key{string} & \Key{val} & \ApiText{Value column.}{Spalte mit dem Wert.}}
  \SemioTableRow{\Key{y2} & \Key{string} & --- & \ApiText{Upper value column of span modes.}{Spalte des oberen Werts der Spannenmodi.}}
  \SemioTableRow{\Key{group} & \Key{string} & --- & \ApiText{Series / stack key column.}{Spalte des Reihen- bzw. Stapelschlüssels.}}
  \SemioTableRow{\Key{series} & \Key{string} & --- & \ApiText{Column that splits the rows into series.}{Spalte, die die Zeilen in Reihen aufteilt.}}
  \SemioTableRow{\Key{label} & \Key{string} & --- & \ApiText{Text column for direct labels.}{Textspalte für Direktbeschriftungen.}}
  \SemioTableRow{\Key{padLeft} & \Key{number} & --- & \ApiText{Plot inset inside the canvas.}{Einrückung der Zeichenfläche innerhalb der Leinwand.}}
  \SemioTableRow{\Key{padRight} & \Key{number} & --- & \ApiText{Plot inset inside the canvas.}{Einrückung der Zeichenfläche innerhalb der Leinwand.}}
  \SemioTableRow{\Key{padTop} & \Key{number} & --- & \ApiText{Plot inset inside the canvas.}{Einrückung der Zeichenfläche innerhalb der Leinwand.}}
  \SemioTableRow{\Key{padBottom} & \Key{number} & --- & \ApiText{Plot inset inside the canvas.}{Einrückung der Zeichenfläche innerhalb der Leinwand.}}
  \SemioTableRow{\Key{yMin} & \Key{string} & \Key{from the data} & \ApiText{Value-axis domain override.}{Vorgabe des Wertebereichs der Werteachse.}}
  \SemioTableRow{\Key{yMax} & \Key{string} & \Key{from the data} & \ApiText{Value-axis domain override.}{Vorgabe des Wertebereichs der Werteachse.}}
  \SemioTableRow{\Key{axes} & \Key{boolean} & \Key{True} & \ApiText{Axis lines, tick labels, legend.}{Achsenlinien, Teilstrichbeschriftungen und Legende.}}
  \SemioTableRow{\Key{grid} & \Key{string} & \Key{none} & \ApiText{Grid lines.}{Gitterlinien hinter den Marken.}}
  \SemioTableRow{\Key{labels} & \Key{string} & \Key{none} & \ApiText{Per-mark labels.}{Beschriftungen an den einzelnen Marken.}}
  \SemioTableRow{\Key{sort} & \Key{string} & \Key{none} & \ApiText{Category order by total value.}{Reihenfolge der Kategorien nach Gesamtwert.}}
  \SemioTableRow{\Key{curve} & \Key{string} & --- & \ApiText{Interpolation of the boundary.}{Interpolation der Begrenzung.}}
}

\subsection{\Key{axisplot}}

\SemioTableLong[text-size=8pt]{\Key{axisplot}}{0.20,0.13,0.15,0.52}{\ApiKey & \ApiType & \ApiDefault & \ApiMeaning}{%
  \SemioTableRow{\Key{scale} & \Key{string} & \Key{linear} & \ApiText{Named scale the values pass through.}{Benannte Skala, durch die die Werte laufen.}}
  \SemioTableRow{\Key{mode} & \Key{string} & \Key{single} & \ApiText{How several series share one band.}{Wie sich mehrere Reihen ein Band teilen.}}
  \SemioTableRow{\Key{breakAt} & \Key{number} & \Key{5} & \ApiText{Where a broken/discontinuous axis is interrupted.}{Wo eine unterbrochene Achse aussetzt.}}
  \SemioTableRow{\Key{x} & \Key{string} & \Key{cat} & \ApiText{Independent / category column.}{Spalte der unabhängigen Größe bzw. der Kategorie.}}
  \SemioTableRow{\Key{y} & \Key{string} & \Key{val} & \ApiText{Value column.}{Spalte mit dem Wert.}}
  \SemioTableRow{\Key{y2} & \Key{string} & --- & \ApiText{Upper value column of span modes.}{Spalte des oberen Werts der Spannenmodi.}}
  \SemioTableRow{\Key{group} & \Key{string} & --- & \ApiText{Series / stack key column.}{Spalte des Reihen- bzw. Stapelschlüssels.}}
  \SemioTableRow{\Key{series} & \Key{string} & --- & \ApiText{Column that splits the rows into series.}{Spalte, die die Zeilen in Reihen aufteilt.}}
  \SemioTableRow{\Key{label} & \Key{string} & --- & \ApiText{Text column for direct labels.}{Textspalte für Direktbeschriftungen.}}
  \SemioTableRow{\Key{padLeft} & \Key{number} & --- & \ApiText{Plot inset inside the canvas.}{Einrückung der Zeichenfläche innerhalb der Leinwand.}}
  \SemioTableRow{\Key{padRight} & \Key{number} & --- & \ApiText{Plot inset inside the canvas.}{Einrückung der Zeichenfläche innerhalb der Leinwand.}}
  \SemioTableRow{\Key{padTop} & \Key{number} & --- & \ApiText{Plot inset inside the canvas.}{Einrückung der Zeichenfläche innerhalb der Leinwand.}}
  \SemioTableRow{\Key{padBottom} & \Key{number} & --- & \ApiText{Plot inset inside the canvas.}{Einrückung der Zeichenfläche innerhalb der Leinwand.}}
  \SemioTableRow{\Key{yMin} & \Key{string} & \Key{from the data} & \ApiText{Value-axis domain override.}{Vorgabe des Wertebereichs der Werteachse.}}
  \SemioTableRow{\Key{yMax} & \Key{string} & \Key{from the data} & \ApiText{Value-axis domain override.}{Vorgabe des Wertebereichs der Werteachse.}}
  \SemioTableRow{\Key{axes} & \Key{boolean} & \Key{True} & \ApiText{Axis lines, tick labels, legend.}{Achsenlinien, Teilstrichbeschriftungen und Legende.}}
  \SemioTableRow{\Key{grid} & \Key{string} & \Key{none} & \ApiText{Grid lines.}{Gitterlinien hinter den Marken.}}
  \SemioTableRow{\Key{labels} & \Key{string} & \Key{none} & \ApiText{Per-mark labels.}{Beschriftungen an den einzelnen Marken.}}
  \SemioTableRow{\Key{sort} & \Key{string} & \Key{none} & \ApiText{Category order by total value.}{Reihenfolge der Kategorien nach Gesamtwert.}}
  \SemioTableRow{\Key{curve} & \Key{string} & --- & \ApiText{Interpolation between points.}{Interpolation zwischen den Punkten.}}
}

\subsection{\Key{bar}}

\SemioTableLong[text-size=8pt]{\Key{bar}}{0.20,0.13,0.15,0.52}{\ApiKey & \ApiType & \ApiDefault & \ApiMeaning}{%
  \SemioTableRow{\Key{orient} & \Key{string} & \Key{vertical} & \ApiText{Bar direction.}{Richtung der Balken.}}
  \SemioTableRow{\Key{mode} & \Key{string} & \Key{grouped} & \ApiText{How several series share one band.}{Wie sich mehrere Reihen ein Band teilen.}}
  \SemioTableRow{\Key{padding} & \Key{number} & \Key{0.2} & \ApiText{Fraction of the band left empty.}{Anteil des Bands, der leer bleibt.}}
  \SemioTableRow{\Key{barWidth} & \Key{number} & --- & \ApiText{Explicit bar thickness override.}{Ausdrückliche Vorgabe der Balkendicke.}}
  \SemioTableRow{\Key{cornerRadius} & \Key{number} & \Key{0} & \ApiText{Rounded bar corners.}{Abgerundete Balkenecken.}}
  \SemioTableRow{\Key{lollipop} & \Key{boolean} & \Key{False} & \ApiText{Stem plus head instead of a filled bar.}{Stiel mit Kopf statt eines gefüllten Balkens.}}
  \SemioTableRow{\Key{mirror} & \Key{boolean} & \Key{False} & \ApiText{Value also mirrored below the baseline.}{Der Wert wird auch unter der Grundlinie gespiegelt.}}
  \SemioTableRow{\Key{paired} & \Key{boolean} & \Key{False} & \ApiText{Series drawn back to back.}{Reihen, Rücken an Rücken gezeichnet.}}
  \SemioTableRow{\Key{baseline} & \Key{number} & \Key{0} & \ApiText{Value the bars grow from.}{Wert, von dem aus die Balken wachsen.}}
  \SemioTableRow{\Key{x} & \Key{string} & \Key{cat} & \ApiText{Independent / category column.}{Spalte der unabhängigen Größe bzw. der Kategorie.}}
  \SemioTableRow{\Key{y} & \Key{string} & \Key{val} & \ApiText{Value column.}{Spalte mit dem Wert.}}
  \SemioTableRow{\Key{y2} & \Key{string} & --- & \ApiText{Upper value column of span modes.}{Spalte des oberen Werts der Spannenmodi.}}
  \SemioTableRow{\Key{group} & \Key{string} & --- & \ApiText{Series / stack key column.}{Spalte des Reihen- bzw. Stapelschlüssels.}}
  \SemioTableRow{\Key{series} & \Key{string} & --- & \ApiText{Column that splits the rows into series.}{Spalte, die die Zeilen in Reihen aufteilt.}}
  \SemioTableRow{\Key{label} & \Key{string} & --- & \ApiText{Text column for direct labels.}{Textspalte für Direktbeschriftungen.}}
  \SemioTableRow{\Key{padLeft} & \Key{number} & --- & \ApiText{Plot inset inside the canvas.}{Einrückung der Zeichenfläche innerhalb der Leinwand.}}
  \SemioTableRow{\Key{padRight} & \Key{number} & --- & \ApiText{Plot inset inside the canvas.}{Einrückung der Zeichenfläche innerhalb der Leinwand.}}
  \SemioTableRow{\Key{padTop} & \Key{number} & --- & \ApiText{Plot inset inside the canvas.}{Einrückung der Zeichenfläche innerhalb der Leinwand.}}
  \SemioTableRow{\Key{padBottom} & \Key{number} & --- & \ApiText{Plot inset inside the canvas.}{Einrückung der Zeichenfläche innerhalb der Leinwand.}}
  \SemioTableRow{\Key{yMin} & \Key{string} & \Key{from the data} & \ApiText{Value-axis domain override.}{Vorgabe des Wertebereichs der Werteachse.}}
  \SemioTableRow{\Key{yMax} & \Key{string} & \Key{from the data} & \ApiText{Value-axis domain override.}{Vorgabe des Wertebereichs der Werteachse.}}
  \SemioTableRow{\Key{axes} & \Key{boolean} & \Key{True} & \ApiText{Axis lines, tick labels, legend.}{Achsenlinien, Teilstrichbeschriftungen und Legende.}}
  \SemioTableRow{\Key{grid} & \Key{string} & \Key{none} & \ApiText{Grid lines.}{Gitterlinien hinter den Marken.}}
  \SemioTableRow{\Key{labels} & \Key{string} & \Key{none} & \ApiText{Per-mark labels.}{Beschriftungen an den einzelnen Marken.}}
  \SemioTableRow{\Key{sort} & \Key{string} & \Key{none} & \ApiText{Category order by total value.}{Reihenfolge der Kategorien nach Gesamtwert.}}
  \SemioTableRow{\Key{curve} & \Key{string} & --- & \ApiText{Interpolation used between consecutive points.}{Interpolation zwischen aufeinanderfolgenden Punkten.}}
}

\subsection{\Key{box}}

\SemioTableLong[text-size=8pt]{\Key{box}}{0.20,0.13,0.15,0.52}{\ApiKey & \ApiType & \ApiDefault & \ApiMeaning}{%
  \SemioTableRow{\Key{orient} & \Key{string} & --- & \ApiText{Orientation of the marks.}{Ausrichtung der Marken.}}
  \SemioTableRow{\Key{notched} & \Key{boolean} & --- & \ApiText{Draw the median notch of the box.}{Die Mediankerbe des Kastens zeichnen.}}
  \SemioTableRow{\Key{varwidth} & \Key{boolean} & --- & \ApiText{Scale the box width by the group size.}{Die Kastenbreite mit der Gruppengröße skalieren.}}
  \SemioTableRow{\Key{letter} & \Key{integer} & --- & \ApiText{Depth of the letter-value summary.}{Tiefe der Letter-Value-Zusammenfassung.}}
  \SemioTableRow{\Key{grouped} & \Key{boolean} & --- & \ApiText{Place one box per group side by side.}{Je Gruppe einen Kasten nebeneinander stellen.}}
  \SemioTableRow{\Key{whisker} & \Key{number} & --- & \ApiText{Whisker length as a multiple of the interquartile range.}{Länge der Antenne als Vielfaches des Interquartilsabstands.}}
}

\subsection{\Key{composite}}

\SemioTableLong[text-size=8pt]{\Key{composite}}{0.20,0.13,0.15,0.52}{\ApiKey & \ApiType & \ApiDefault & \ApiMeaning}{%
  \SemioTableRow{\Key{layout} & \Key{string} & --- & \ApiText{Layout algorithm that places the nodes.}{Layoutalgorithmus, der die Knoten platziert.}}
  \SemioTableRow{\Key{main} & \Key{string} & --- & \ApiText{Name of the chart drawn in the main panel.}{Name des Diagramms im Hauptfeld.}}
  \SemioTableRow{\Key{mainopts} & \Key{string} & --- & \ApiText{Option list passed to the main panel's chart.}{Optionsliste, die an das Diagramm des Hauptfelds geht.}}
  \SemioTableRow{\Key{side} & \Key{string} & --- & \ApiText{Which side the secondary element sits on.}{Auf welcher Seite das Nebenelement sitzt.}}
  \SemioTableRow{\Key{sideopts} & \Key{string} & --- & \ApiText{Option list passed to the side panel's chart.}{Optionsliste, die an das Diagramm des Seitenfelds geht.}}
  \SemioTableRow{\Key{ratio} & \Key{number} & --- & \ApiText{Fraction of the frame the main panel takes.}{Anteil des Rahmens, den das Hauptfeld einnimmt.}}
  \SemioTableRow{\Key{gap} & \Key{number} & --- & \ApiText{Millimetre gap between adjacent parts.}{Millimeterabstand zwischen benachbarten Teilen.}}
}

\subsection{\Key{composition-bar}}

\ApiText{This family takes no options beyond \Key{variant} and \Key{data}.}{Diese Familie nimmt außer \Key{variant} und \Key{data} keine Optionen.}

\subsection{\Key{control}}

\SemioTableLong[text-size=8pt]{\Key{control}}{0.20,0.13,0.15,0.52}{\ApiKey & \ApiType & \ApiDefault & \ApiMeaning}{%
  \SemioTableRow{\Key{chart} & \Key{string} & --- & \ApiText{Control-chart rule set the limits follow.}{Regelwerk der Qualitätsregelkarte, dem die Grenzen folgen.}}
  \SemioTableRow{\Key{timecol} & \Key{string} & --- & \ApiText{Column holding the observation time.}{Spalte mit dem Beobachtungszeitpunkt.}}
  \SemioTableRow{\Key{valuecol} & \Key{string} & --- & \ApiText{Column mapped onto the value.}{Spalte, die auf den Wert abgebildet wird.}}
  \SemioTableRow{\Key{sizecol} & \Key{string} & --- & \ApiText{Column mapped onto the mark size.}{Spalte, die auf die Markengröße abgebildet wird.}}
  \SemioTableRow{\Key{sigma} & \Key{number} & --- & \ApiText{Width of the control limits in standard deviations.}{Breite der Eingriffsgrenzen in Standardabweichungen.}}
  \SemioTableRow{\Key{lambda} & \Key{number} & --- & \ApiText{Smoothing weight of the EWMA control chart.}{Glättungsgewicht der EWMA-Regelkarte.}}
  \SemioTableRow{\Key{target} & \Key{number} & --- & \ApiText{Column holding the target node of an edge, or the target value.}{Spalte mit dem Zielknoten einer Kante oder mit dem Zielwert.}}
}

\subsection{\Key{coordplot}}

\SemioTableLong[text-size=8pt]{\Key{coordplot}}{0.20,0.13,0.15,0.52}{\ApiKey & \ApiType & \ApiDefault & \ApiMeaning}{%
  \SemioTableRow{\Key{coordinate} & \Key{string} & \Key{cartesian} & \ApiText{Coordinate system the marks are placed in.}{Koordinatensystem, in dem die Marken platziert werden.}}
  \SemioTableRow{\Key{axesCount} & \Key{integer} & \Key{4} & \ApiText{Axes a parallel-coordinate plot draws.}{Achsen, die ein Parallelkoordinaten-Diagramm zeichnet.}}
  \SemioTableRow{\Key{x} & \Key{string} & \Key{cat} & \ApiText{Independent / category column.}{Spalte der unabhängigen Größe bzw. der Kategorie.}}
  \SemioTableRow{\Key{y} & \Key{string} & \Key{val} & \ApiText{Value column.}{Spalte mit dem Wert.}}
  \SemioTableRow{\Key{y2} & \Key{string} & --- & \ApiText{Upper value column of span modes.}{Spalte des oberen Werts der Spannenmodi.}}
  \SemioTableRow{\Key{group} & \Key{string} & --- & \ApiText{Series / stack key column.}{Spalte des Reihen- bzw. Stapelschlüssels.}}
  \SemioTableRow{\Key{series} & \Key{string} & --- & \ApiText{Column that splits the rows into series.}{Spalte, die die Zeilen in Reihen aufteilt.}}
  \SemioTableRow{\Key{label} & \Key{string} & --- & \ApiText{Text column for direct labels.}{Textspalte für Direktbeschriftungen.}}
  \SemioTableRow{\Key{padLeft} & \Key{number} & --- & \ApiText{Plot inset inside the canvas.}{Einrückung der Zeichenfläche innerhalb der Leinwand.}}
  \SemioTableRow{\Key{padRight} & \Key{number} & --- & \ApiText{Plot inset inside the canvas.}{Einrückung der Zeichenfläche innerhalb der Leinwand.}}
  \SemioTableRow{\Key{padTop} & \Key{number} & --- & \ApiText{Plot inset inside the canvas.}{Einrückung der Zeichenfläche innerhalb der Leinwand.}}
  \SemioTableRow{\Key{padBottom} & \Key{number} & --- & \ApiText{Plot inset inside the canvas.}{Einrückung der Zeichenfläche innerhalb der Leinwand.}}
  \SemioTableRow{\Key{yMin} & \Key{string} & \Key{from the data} & \ApiText{Value-axis domain override.}{Vorgabe des Wertebereichs der Werteachse.}}
  \SemioTableRow{\Key{yMax} & \Key{string} & \Key{from the data} & \ApiText{Value-axis domain override.}{Vorgabe des Wertebereichs der Werteachse.}}
  \SemioTableRow{\Key{axes} & \Key{boolean} & \Key{True} & \ApiText{Axis lines, tick labels, legend.}{Achsenlinien, Teilstrichbeschriftungen und Legende.}}
  \SemioTableRow{\Key{grid} & \Key{string} & --- & \ApiText{Grid lines.}{Gitterlinien hinter den Marken.}}
  \SemioTableRow{\Key{labels} & \Key{string} & \Key{none} & \ApiText{Per-mark labels.}{Beschriftungen an den einzelnen Marken.}}
  \SemioTableRow{\Key{sort} & \Key{string} & \Key{none} & \ApiText{Category order by total value.}{Reihenfolge der Kategorien nach Gesamtwert.}}
  \SemioTableRow{\Key{curve} & \Key{string} & --- & \ApiText{Interpolation used between consecutive points.}{Interpolation zwischen aufeinanderfolgenden Punkten.}}
}

\subsection{\Key{correlation}}

\SemioTableLong[text-size=8pt]{\Key{correlation}}{0.20,0.13,0.15,0.52}{\ApiKey & \ApiType & \ApiDefault & \ApiMeaning}{%
  \SemioTableRow{\Key{columns} & \Key{string} & --- & \ApiText{Number of columns of the layout grid.}{Anzahl der Spalten des Anordnungsgitters.}}
  \SemioTableRow{\Key{cell} & \Key{string} & --- & \ApiText{Cell size in millimetres, or the column that carries it.}{Zellgröße in Millimetern oder die Spalte, die sie trägt.}}
  \SemioTableRow{\Key{triangle} & \Key{string} & --- & \ApiText{Which triangle of the matrix is drawn.}{Welches Dreieck der Matrix gezeichnet wird.}}
  \SemioTableRow{\Key{values} & \Key{boolean} & --- & \ApiText{Print the coefficients inside the cells.}{Die Koeffizienten in die Zellen schreiben.}}
  \SemioTableRow{\Key{labels} & \Key{boolean} & --- & \ApiText{Draw the captions.}{Die Beschriftungen zeichnen.}}
}

\subsection{\Key{coxcomb}}

\SemioTableLong[text-size=8pt]{\Key{coxcomb}}{0.20,0.13,0.15,0.52}{\ApiKey & \ApiType & \ApiDefault & \ApiMeaning}{%
  \SemioTableRow{\Key{area} & \Key{boolean} & --- & \ApiText{Encode the value in the sector area instead of its radius.}{Den Wert über die Sektorfläche statt über den Radius codieren.}}
  \SemioTableRow{\Key{angle} & \Key{string} & --- & \ApiText{Angular position, in degrees or as the column that drives it.}{Winkelposition, in Grad oder als die Spalte, die sie steuert.}}
  \SemioTableRow{\Key{radiusv} & \Key{string} & --- & \ApiText{Column mapped onto the radius.}{Spalte, die auf den Radius abgebildet wird.}}
  \SemioTableRow{\Key{inner} & \Key{number} & --- & \ApiText{Inner radius in millimetres; larger values open a hole.}{Innenradius in Millimetern; größere Werte öffnen ein Loch.}}
  \SemioTableRow{\Key{stacked} & \Key{boolean} & --- & \ApiText{Stack the series inside one sector.}{Die Reihen in einem Sektor stapeln.}}
  \SemioTableRow{\Key{second} & \Key{string} & --- & \ApiText{Column of the second series drawn on the same rose.}{Spalte der zweiten Reihe auf derselben Rose.}}
}

\subsection{\Key{curve}}

\SemioTableLong[text-size=8pt]{\Key{curve}}{0.20,0.13,0.15,0.52}{\ApiKey & \ApiType & \ApiDefault & \ApiMeaning}{%
  \SemioTableRow{\Key{mode} & \Key{string} & \Key{supply-demand} & \ApiText{How several series share one band.}{Wie sich mehrere Reihen ein Band teilen.}}
  \SemioTableRow{\Key{samples} & \Key{integer} & \Key{24} & \ApiText{Points the curve is sampled at.}{Punkte, an denen die Kurve abgetastet wird.}}
  \SemioTableRow{\Key{shift} & \Key{number} & \Key{0} & \ApiText{Shifts the second curve of a two-curve mode.}{Verschiebt die zweite Kurve eines Zweikurvenmodus.}}
  \SemioTableRow{\Key{x} & \Key{string} & \Key{cat} & \ApiText{Independent / category column.}{Spalte der unabhängigen Größe bzw. der Kategorie.}}
  \SemioTableRow{\Key{y} & \Key{string} & \Key{val} & \ApiText{Value column.}{Spalte mit dem Wert.}}
  \SemioTableRow{\Key{y2} & \Key{string} & --- & \ApiText{Upper value column of span modes.}{Spalte des oberen Werts der Spannenmodi.}}
  \SemioTableRow{\Key{group} & \Key{string} & --- & \ApiText{Series / stack key column.}{Spalte des Reihen- bzw. Stapelschlüssels.}}
  \SemioTableRow{\Key{series} & \Key{string} & --- & \ApiText{Column that splits the rows into series.}{Spalte, die die Zeilen in Reihen aufteilt.}}
  \SemioTableRow{\Key{label} & \Key{string} & --- & \ApiText{Text column for direct labels.}{Textspalte für Direktbeschriftungen.}}
  \SemioTableRow{\Key{padLeft} & \Key{number} & --- & \ApiText{Plot inset inside the canvas.}{Einrückung der Zeichenfläche innerhalb der Leinwand.}}
  \SemioTableRow{\Key{padRight} & \Key{number} & --- & \ApiText{Plot inset inside the canvas.}{Einrückung der Zeichenfläche innerhalb der Leinwand.}}
  \SemioTableRow{\Key{padTop} & \Key{number} & --- & \ApiText{Plot inset inside the canvas.}{Einrückung der Zeichenfläche innerhalb der Leinwand.}}
  \SemioTableRow{\Key{padBottom} & \Key{number} & --- & \ApiText{Plot inset inside the canvas.}{Einrückung der Zeichenfläche innerhalb der Leinwand.}}
  \SemioTableRow{\Key{yMin} & \Key{string} & \Key{from the data} & \ApiText{Value-axis domain override.}{Vorgabe des Wertebereichs der Werteachse.}}
  \SemioTableRow{\Key{yMax} & \Key{string} & \Key{from the data} & \ApiText{Value-axis domain override.}{Vorgabe des Wertebereichs der Werteachse.}}
  \SemioTableRow{\Key{axes} & \Key{boolean} & \Key{True} & \ApiText{Axis lines, tick labels, legend.}{Achsenlinien, Teilstrichbeschriftungen und Legende.}}
  \SemioTableRow{\Key{grid} & \Key{string} & \Key{none} & \ApiText{Grid lines.}{Gitterlinien hinter den Marken.}}
  \SemioTableRow{\Key{labels} & \Key{string} & \Key{none} & \ApiText{Per-mark labels.}{Beschriftungen an den einzelnen Marken.}}
  \SemioTableRow{\Key{sort} & \Key{string} & \Key{none} & \ApiText{Category order by total value.}{Reihenfolge der Kategorien nach Gesamtwert.}}
  \SemioTableRow{\Key{curve} & \Key{string} & --- & \ApiText{Interpolation used between consecutive points.}{Interpolation zwischen aufeinanderfolgenden Punkten.}}
}

\subsection{\Key{density}}

\SemioTableLong[text-size=8pt]{\Key{density}}{0.20,0.13,0.15,0.52}{\ApiKey & \ApiType & \ApiDefault & \ApiMeaning}{%
  \SemioTableRow{\Key{bandwidth} & \Key{number} & --- & \ApiText{Kernel bandwidth of the density estimate; '0' uses Silverman's rule.}{Bandbreite des Kerns der Dichteschätzung; '0' nutzt Silvermans Regel.}}
  \SemioTableRow{\Key{kernel} & \Key{string} & --- & \ApiText{Kernel function of the density estimate.}{Kernfunktion der Dichteschätzung.}}
  \SemioTableRow{\Key{samples} & \Key{integer} & --- & \ApiText{Number of samples the curve or field is evaluated at.}{Anzahl der Stützstellen, an denen Kurve oder Feld ausgewertet werden.}}
  \SemioTableRow{\Key{filled} & \Key{boolean} & --- & \ApiText{Fill the bands between contour levels.}{Die Bänder zwischen den Höhenlinien füllen.}}
  \SemioTableRow{\Key{multiple} & \Key{boolean} & --- & \ApiText{Draw one density per group.}{Je Gruppe eine Dichte zeichnen.}}
  \SemioTableRow{\Key{ridgeline} & \Key{boolean} & --- & \ApiText{Stack the group densities as a ridgeline.}{Die Gruppendichten als Gebirgszug stapeln.}}
  \SemioTableRow{\Key{offset} & \Key{number} & --- & \ApiText{Baseline offset between stacked densities.}{Grundlinienversatz zwischen gestapelten Dichten.}}
  \SemioTableRow{\Key{cumulative} & \Key{boolean} & --- & \ApiText{Accumulate the values instead of showing them per bin.}{Die Werte kumulieren statt sie je Klasse zu zeigen.}}
}

\subsection{\Key{dot}}

\SemioTableLong[text-size=8pt]{\Key{dot}}{0.20,0.13,0.15,0.52}{\ApiKey & \ApiType & \ApiDefault & \ApiMeaning}{%
  \SemioTableRow{\Key{mode} & \Key{string} & \Key{dot} & \ApiText{How several series share one band.}{Wie sich mehrere Reihen ein Band teilen.}}
  \SemioTableRow{\Key{orient} & \Key{string} & \Key{horizontal} & \ApiText{Bar direction.}{Richtung der Balken.}}
  \SemioTableRow{\Key{connect} & \Key{string} & \Key{from the mode} & \ApiText{Rule between the dots of one category.}{Linie zwischen den Punkten einer Kategorie.}}
  \SemioTableRow{\Key{size} & \Key{number} & \Key{1.2} & \ApiText{Dot radius.}{Radius eines Punktes.}}
  \SemioTableRow{\Key{x} & \Key{string} & \Key{cat} & \ApiText{Independent / category column.}{Spalte der unabhängigen Größe bzw. der Kategorie.}}
  \SemioTableRow{\Key{y} & \Key{string} & \Key{val} & \ApiText{Value column.}{Spalte mit dem Wert.}}
  \SemioTableRow{\Key{y2} & \Key{string} & --- & \ApiText{Upper value column of span modes.}{Spalte des oberen Werts der Spannenmodi.}}
  \SemioTableRow{\Key{group} & \Key{string} & --- & \ApiText{Series / stack key column.}{Spalte des Reihen- bzw. Stapelschlüssels.}}
  \SemioTableRow{\Key{series} & \Key{string} & --- & \ApiText{Column that splits the rows into series.}{Spalte, die die Zeilen in Reihen aufteilt.}}
  \SemioTableRow{\Key{label} & \Key{string} & --- & \ApiText{Text column for direct labels.}{Textspalte für Direktbeschriftungen.}}
  \SemioTableRow{\Key{padLeft} & \Key{number} & --- & \ApiText{Plot inset inside the canvas.}{Einrückung der Zeichenfläche innerhalb der Leinwand.}}
  \SemioTableRow{\Key{padRight} & \Key{number} & --- & \ApiText{Plot inset inside the canvas.}{Einrückung der Zeichenfläche innerhalb der Leinwand.}}
  \SemioTableRow{\Key{padTop} & \Key{number} & --- & \ApiText{Plot inset inside the canvas.}{Einrückung der Zeichenfläche innerhalb der Leinwand.}}
  \SemioTableRow{\Key{padBottom} & \Key{number} & --- & \ApiText{Plot inset inside the canvas.}{Einrückung der Zeichenfläche innerhalb der Leinwand.}}
  \SemioTableRow{\Key{yMin} & \Key{string} & \Key{from the data} & \ApiText{Value-axis domain override.}{Vorgabe des Wertebereichs der Werteachse.}}
  \SemioTableRow{\Key{yMax} & \Key{string} & \Key{from the data} & \ApiText{Value-axis domain override.}{Vorgabe des Wertebereichs der Werteachse.}}
  \SemioTableRow{\Key{axes} & \Key{boolean} & \Key{True} & \ApiText{Axis lines, tick labels, legend.}{Achsenlinien, Teilstrichbeschriftungen und Legende.}}
  \SemioTableRow{\Key{grid} & \Key{string} & \Key{none} & \ApiText{Grid lines.}{Gitterlinien hinter den Marken.}}
  \SemioTableRow{\Key{labels} & \Key{string} & \Key{none} & \ApiText{Per-mark labels.}{Beschriftungen an den einzelnen Marken.}}
  \SemioTableRow{\Key{sort} & \Key{string} & \Key{none} & \ApiText{Category order by total value.}{Reihenfolge der Kategorien nach Gesamtwert.}}
  \SemioTableRow{\Key{curve} & \Key{string} & --- & \ApiText{Interpolation used between consecutive points.}{Interpolation zwischen aufeinanderfolgenden Punkten.}}
}

\subsection{\Key{ecdf}}

\SemioTableLong[text-size=8pt]{\Key{ecdf}}{0.20,0.13,0.15,0.52}{\ApiKey & \ApiType & \ApiDefault & \ApiMeaning}{%
  \SemioTableRow{\Key{complementary} & \Key{boolean} & --- & \ApiText{Draw the complementary distribution 1−F(x).}{Die komplementäre Verteilung 1−F(x) zeichnen.}}
  \SemioTableRow{\Key{step} & \Key{boolean} & --- & \ApiText{Spacing between consecutive steps.}{Abstand zwischen aufeinanderfolgenden Schritten.}}
  \SemioTableRow{\Key{points} & \Key{boolean} & --- & \ApiText{Draw the individual data points.}{Die einzelnen Datenpunkte zeichnen.}}
}

\subsection{\Key{election}}

\ApiText{This family takes no options beyond \Key{variant} and \Key{data}.}{Diese Familie nimmt außer \Key{variant} und \Key{data} keine Optionen.}

\subsection{\Key{embedding}}

\SemioTableLong[text-size=8pt]{\Key{embedding}}{0.20,0.13,0.15,0.52}{\ApiKey & \ApiType & \ApiDefault & \ApiMeaning}{%
  \SemioTableRow{\Key{ex} & \Key{string} & --- & \ApiText{Column holding the first embedding coordinate.}{Spalte mit der ersten Einbettungskoordinate.}}
  \SemioTableRow{\Key{ey} & \Key{string} & --- & \ApiText{Column holding the second embedding coordinate.}{Spalte mit der zweiten Einbettungskoordinate.}}
  \SemioTableRow{\Key{cluster} & \Key{string} & --- & \ApiText{Column holding the cluster assignment of a point.}{Spalte mit der Clusterzuordnung eines Punktes.}}
  \SemioTableRow{\Key{loading} & \Key{string} & --- & \ApiText{Column holding the loading of a variable.}{Spalte mit der Ladung einer Variablen.}}
  \SemioTableRow{\Key{arrows} & \Key{boolean} & --- & \ApiText{Draw the projection arrows of the loading vectors.}{Die Projektionspfeile der Ladungsvektoren zeichnen.}}
  \SemioTableRow{\Key{hulls} & \Key{boolean} & --- & \ApiText{Draw the convex hull of every cluster.}{Die konvexe Hülle jedes Clusters zeichnen.}}
  \SemioTableRow{\Key{radius} & \Key{number} & --- & \ApiText{Radius in millimetres, or the column that drives it.}{Radius in Millimetern oder die Spalte, die ihn steuert.}}
}

\subsection{\Key{evaluation}}

\SemioTableLong[text-size=8pt]{\Key{evaluation}}{0.20,0.13,0.15,0.52}{\ApiKey & \ApiType & \ApiDefault & \ApiMeaning}{%
  \SemioTableRow{\Key{curve} & \Key{string} & --- & \ApiText{Interpolation used between consecutive points.}{Interpolation zwischen aufeinanderfolgenden Punkten.}}
  \SemioTableRow{\Key{score} & \Key{string} & --- & \ApiText{Column holding the model score.}{Spalte mit dem Modellwert.}}
  \SemioTableRow{\Key{truth} & \Key{string} & --- & \ApiText{Column holding the true class.}{Spalte mit der wahren Klasse.}}
  \SemioTableRow{\Key{diagonal} & \Key{boolean} & --- & \ApiText{What the diagonal cells show.}{Was die Diagonalzellen zeigen.}}
  \SemioTableRow{\Key{points} & \Key{boolean} & --- & \ApiText{Draw the individual data points.}{Die einzelnen Datenpunkte zeichnen.}}
  \SemioTableRow{\Key{bars} & \Key{integer} & --- & \ApiText{Number of bars in the evaluation panel.}{Anzahl der Balken im Bewertungsfeld.}}
  \SemioTableRow{\Key{height} & \Key{string} & --- & \ApiText{Frame height in millimetres.}{Rahmenhöhe in Millimetern.}}
}

\subsection{\Key{facet}}

\SemioTableLong[text-size=8pt]{\Key{facet}}{0.20,0.13,0.15,0.52}{\ApiKey & \ApiType & \ApiDefault & \ApiMeaning}{%
  \SemioTableRow{\Key{mode} & \Key{string} & --- & \ApiText{Which member of the family the preset renders.}{Welches Mitglied der Familie die Voreinstellung zeichnet.}}
  \SemioTableRow{\Key{panel} & \Key{string} & --- & \ApiText{Chart drawn inside every facet panel.}{Diagramm, das in jedem Facettenfeld gezeichnet wird.}}
  \SemioTableRow{\Key{opts} & \Key{string} & --- & \ApiText{Option list forwarded to every facet.}{Optionsliste, die an jede Facette weitergereicht wird.}}
  \SemioTableRow{\Key{by} & \Key{string} & --- & \ApiText{Column the facets are split by.}{Spalte, nach der die Facetten aufgeteilt werden.}}
  \SemioTableRow{\Key{columns} & \Key{integer} & --- & \ApiText{Panels per row for wrap=, default 3.}{Felder je Zeile bei wrap=.}}
  \SemioTableRow{\Key{gap} & \Key{number} & --- & \ApiText{Gap between panels, default 4.}{Abstand zwischen den Feldern.}}
  \SemioTableRow{\Key{shared} & \Key{boolean} & --- & \ApiText{Share the scales across all facets.}{Die Skalen über alle Facetten teilen.}}
}

\subsection{\Key{financial}}

\SemioTableLong[text-size=8pt]{\Key{financial}}{0.20,0.13,0.15,0.52}{\ApiKey & \ApiType & \ApiDefault & \ApiMeaning}{%
  \SemioTableRow{\Key{mode} & \Key{string} & \Key{candlestick} & \ApiText{How several series share one band.}{Wie sich mehrere Reihen ein Band teilen.}}
  \SemioTableRow{\Key{overlay} & \Key{string} & \Key{none} & \ApiText{Price overlay drawn on the same panel.}{Kursüberlagerung im selben Feld.}}
  \SemioTableRow{\Key{indicator} & \Key{string} & \Key{none} & \ApiText{Technical indicator drawn on the price panel.}{Technischer Indikator im Kursfeld.}}
  \SemioTableRow{\Key{window} & \Key{integer} & \Key{3} & \ApiText{Window of the moving average and the bands.}{Fenster des gleitenden Mittels und der Bänder.}}
  \SemioTableRow{\Key{boxSize} & \Key{number} & \Key{1} & \ApiText{Brick / box height of renko and point-and-figure.}{Ziegel- bzw. Kastenhöhe von Renko und Point-and-Figure.}}
  \SemioTableRow{\Key{open} & \Key{string} & --- & \ApiText{Column holding the opening price.}{Spalte mit dem Eröffnungskurs.}}
  \SemioTableRow{\Key{high} & \Key{string} & --- & \ApiText{Column holding the session high.}{Spalte mit dem Tageshoch.}}
  \SemioTableRow{\Key{low} & \Key{string} & --- & \ApiText{Column holding the session low.}{Spalte mit dem Tagestief.}}
  \SemioTableRow{\Key{close} & \Key{string} & --- & \ApiText{Column holding the closing price.}{Spalte mit dem Schlusskurs.}}
  \SemioTableRow{\Key{volume} & \Key{string} & --- & \ApiText{Column holding the traded volume.}{Spalte mit dem gehandelten Volumen.}}
  \SemioTableRow{\Key{x} & \Key{string} & \Key{cat} & \ApiText{Independent / category column.}{Spalte der unabhängigen Größe bzw. der Kategorie.}}
  \SemioTableRow{\Key{y} & \Key{string} & \Key{val} & \ApiText{Value column.}{Spalte mit dem Wert.}}
  \SemioTableRow{\Key{y2} & \Key{string} & --- & \ApiText{Upper value column of span modes.}{Spalte des oberen Werts der Spannenmodi.}}
  \SemioTableRow{\Key{group} & \Key{string} & --- & \ApiText{Series / stack key column.}{Spalte des Reihen- bzw. Stapelschlüssels.}}
  \SemioTableRow{\Key{series} & \Key{string} & --- & \ApiText{Column that splits the rows into series.}{Spalte, die die Zeilen in Reihen aufteilt.}}
  \SemioTableRow{\Key{label} & \Key{string} & --- & \ApiText{Text column for direct labels.}{Textspalte für Direktbeschriftungen.}}
  \SemioTableRow{\Key{padLeft} & \Key{number} & --- & \ApiText{Plot inset inside the canvas.}{Einrückung der Zeichenfläche innerhalb der Leinwand.}}
  \SemioTableRow{\Key{padRight} & \Key{number} & --- & \ApiText{Plot inset inside the canvas.}{Einrückung der Zeichenfläche innerhalb der Leinwand.}}
  \SemioTableRow{\Key{padTop} & \Key{number} & --- & \ApiText{Plot inset inside the canvas.}{Einrückung der Zeichenfläche innerhalb der Leinwand.}}
  \SemioTableRow{\Key{padBottom} & \Key{number} & --- & \ApiText{Plot inset inside the canvas.}{Einrückung der Zeichenfläche innerhalb der Leinwand.}}
  \SemioTableRow{\Key{yMin} & \Key{string} & \Key{from the data} & \ApiText{Value-axis domain override.}{Vorgabe des Wertebereichs der Werteachse.}}
  \SemioTableRow{\Key{yMax} & \Key{string} & \Key{from the data} & \ApiText{Value-axis domain override.}{Vorgabe des Wertebereichs der Werteachse.}}
  \SemioTableRow{\Key{axes} & \Key{boolean} & \Key{True} & \ApiText{Axis lines, tick labels, legend.}{Achsenlinien, Teilstrichbeschriftungen und Legende.}}
  \SemioTableRow{\Key{grid} & \Key{string} & \Key{none} & \ApiText{Grid lines.}{Gitterlinien hinter den Marken.}}
  \SemioTableRow{\Key{labels} & \Key{string} & \Key{none} & \ApiText{Per-mark labels.}{Beschriftungen an den einzelnen Marken.}}
  \SemioTableRow{\Key{sort} & \Key{string} & \Key{none} & \ApiText{Category order by total value.}{Reihenfolge der Kategorien nach Gesamtwert.}}
  \SemioTableRow{\Key{curve} & \Key{string} & --- & \ApiText{Interpolation used between consecutive points.}{Interpolation zwischen aufeinanderfolgenden Punkten.}}
}

\subsection{\Key{funnel}}

\SemioTableLong[text-size=8pt]{\Key{funnel}}{0.20,0.13,0.15,0.52}{\ApiKey & \ApiType & \ApiDefault & \ApiMeaning}{%
  \SemioTableRow{\Key{mode} & \Key{string} & \Key{funnel} & \ApiText{How several series share one band.}{Wie sich mehrere Reihen ein Band teilen.}}
  \SemioTableRow{\Key{stage} & \Key{string} & \Key{stage} & \ApiText{The stage label column.}{Die Spalte mit der Stufenbeschriftung.}}
  \SemioTableRow{\Key{value} & \Key{string} & \Key{value} & \ApiText{The stage size column.}{Die Spalte mit der Stufengröße.}}
  \SemioTableRow{\Key{orient} & \Key{string} & \Key{vertical} & \ApiText{Bar direction.}{Richtung der Balken.}}
  \SemioTableRow{\Key{gap} & \Key{number} & \Key{1} & \ApiText{Space between two stages.}{Abstand zwischen zwei Stufen.}}
  \SemioTableRow{\Key{dropoff} & \Key{boolean} & \Key{False} & \ApiText{Render the lost share of every stage as well.}{Auch den verlorenen Anteil jeder Stufe zeichnen.}}
  \SemioTableRow{\Key{x} & \Key{string} & \Key{cat} & \ApiText{Independent / category column.}{Spalte der unabhängigen Größe bzw. der Kategorie.}}
  \SemioTableRow{\Key{y} & \Key{string} & \Key{val} & \ApiText{Value column.}{Spalte mit dem Wert.}}
  \SemioTableRow{\Key{y2} & \Key{string} & --- & \ApiText{Upper value column of span modes.}{Spalte des oberen Werts der Spannenmodi.}}
  \SemioTableRow{\Key{group} & \Key{string} & --- & \ApiText{Series / stack key column.}{Spalte des Reihen- bzw. Stapelschlüssels.}}
  \SemioTableRow{\Key{series} & \Key{string} & --- & \ApiText{Column that splits the rows into series.}{Spalte, die die Zeilen in Reihen aufteilt.}}
  \SemioTableRow{\Key{label} & \Key{string} & --- & \ApiText{Text column for direct labels.}{Textspalte für Direktbeschriftungen.}}
  \SemioTableRow{\Key{padLeft} & \Key{number} & --- & \ApiText{Plot inset inside the canvas.}{Einrückung der Zeichenfläche innerhalb der Leinwand.}}
  \SemioTableRow{\Key{padRight} & \Key{number} & --- & \ApiText{Plot inset inside the canvas.}{Einrückung der Zeichenfläche innerhalb der Leinwand.}}
  \SemioTableRow{\Key{padTop} & \Key{number} & --- & \ApiText{Plot inset inside the canvas.}{Einrückung der Zeichenfläche innerhalb der Leinwand.}}
  \SemioTableRow{\Key{padBottom} & \Key{number} & --- & \ApiText{Plot inset inside the canvas.}{Einrückung der Zeichenfläche innerhalb der Leinwand.}}
  \SemioTableRow{\Key{yMin} & \Key{string} & \Key{from the data} & \ApiText{Value-axis domain override.}{Vorgabe des Wertebereichs der Werteachse.}}
  \SemioTableRow{\Key{yMax} & \Key{string} & \Key{from the data} & \ApiText{Value-axis domain override.}{Vorgabe des Wertebereichs der Werteachse.}}
  \SemioTableRow{\Key{axes} & \Key{boolean} & \Key{True} & \ApiText{Axis lines, tick labels, legend.}{Achsenlinien, Teilstrichbeschriftungen und Legende.}}
  \SemioTableRow{\Key{grid} & \Key{string} & \Key{none} & \ApiText{Grid lines.}{Gitterlinien hinter den Marken.}}
  \SemioTableRow{\Key{labels} & \Key{string} & --- & \ApiText{Per-mark labels.}{Beschriftungen an den einzelnen Marken.}}
  \SemioTableRow{\Key{sort} & \Key{string} & \Key{none} & \ApiText{Category order by total value.}{Reihenfolge der Kategorien nach Gesamtwert.}}
  \SemioTableRow{\Key{curve} & \Key{string} & --- & \ApiText{Interpolation used between consecutive points.}{Interpolation zwischen aufeinanderfolgenden Punkten.}}
}

\subsection{\Key{gauge}}

\SemioTableLong[text-size=8pt]{\Key{gauge}}{0.20,0.13,0.15,0.52}{\ApiKey & \ApiType & \ApiDefault & \ApiMeaning}{%
  \SemioTableRow{\Key{mode} & \Key{string} & \Key{radial} & \ApiText{How several series share one band.}{Wie sich mehrere Reihen ein Band teilen.}}
  \SemioTableRow{\Key{value} & \Key{string} & \Key{value} & \ApiText{The measured column.}{Die Spalte mit dem Messwert.}}
  \SemioTableRow{\Key{target} & \Key{string} & \Key{target} & \ApiText{The target marker.}{Die Zielmarke.}}
  \SemioTableRow{\Key{min} & \Key{string} & \Key{min / max} & \ApiText{The gauge domain columns.}{Die Spalten des Anzeigebereichs.}}
  \SemioTableRow{\Key{max} & \Key{string} & \Key{min / max} & \ApiText{The gauge domain columns.}{Die Spalten des Anzeigebereichs.}}
  \SemioTableRow{\Key{startAngle} & \Key{number} & \Key{180} & \ApiText{Where the arc begins.}{Wo der Bogen beginnt.}}
  \SemioTableRow{\Key{endAngle} & \Key{number} & \Key{0} & \ApiText{Where the arc ends.}{Wo der Bogen endet.}}
  \SemioTableRow{\Key{thickness} & \Key{number} & \Key{3} & \ApiText{Ring or bar thickness.}{Dicke des Rings bzw. Balkens.}}
  \SemioTableRow{\Key{ticks} & \Key{integer} & \Key{5} & \ApiText{Scale ticks on the arc.}{Skalenteilstriche auf dem Bogen.}}
  \SemioTableRow{\Key{x} & \Key{string} & \Key{cat} & \ApiText{Independent / category column.}{Spalte der unabhängigen Größe bzw. der Kategorie.}}
  \SemioTableRow{\Key{y} & \Key{string} & \Key{val} & \ApiText{Value column.}{Spalte mit dem Wert.}}
  \SemioTableRow{\Key{y2} & \Key{string} & --- & \ApiText{Upper value column of span modes.}{Spalte des oberen Werts der Spannenmodi.}}
  \SemioTableRow{\Key{group} & \Key{string} & --- & \ApiText{Series / stack key column.}{Spalte des Reihen- bzw. Stapelschlüssels.}}
  \SemioTableRow{\Key{series} & \Key{string} & --- & \ApiText{Column that splits the rows into series.}{Spalte, die die Zeilen in Reihen aufteilt.}}
  \SemioTableRow{\Key{label} & \Key{string} & --- & \ApiText{Text column for direct labels.}{Textspalte für Direktbeschriftungen.}}
  \SemioTableRow{\Key{padLeft} & \Key{number} & --- & \ApiText{Plot inset inside the canvas.}{Einrückung der Zeichenfläche innerhalb der Leinwand.}}
  \SemioTableRow{\Key{padRight} & \Key{number} & --- & \ApiText{Plot inset inside the canvas.}{Einrückung der Zeichenfläche innerhalb der Leinwand.}}
  \SemioTableRow{\Key{padTop} & \Key{number} & --- & \ApiText{Plot inset inside the canvas.}{Einrückung der Zeichenfläche innerhalb der Leinwand.}}
  \SemioTableRow{\Key{padBottom} & \Key{number} & --- & \ApiText{Plot inset inside the canvas.}{Einrückung der Zeichenfläche innerhalb der Leinwand.}}
  \SemioTableRow{\Key{yMin} & \Key{string} & \Key{from the data} & \ApiText{Value-axis domain override.}{Vorgabe des Wertebereichs der Werteachse.}}
  \SemioTableRow{\Key{yMax} & \Key{string} & \Key{from the data} & \ApiText{Value-axis domain override.}{Vorgabe des Wertebereichs der Werteachse.}}
  \SemioTableRow{\Key{axes} & \Key{boolean} & \Key{True} & \ApiText{Axis lines, tick labels, legend.}{Achsenlinien, Teilstrichbeschriftungen und Legende.}}
  \SemioTableRow{\Key{grid} & \Key{string} & \Key{none} & \ApiText{Grid lines.}{Gitterlinien hinter den Marken.}}
  \SemioTableRow{\Key{labels} & \Key{string} & \Key{none} & \ApiText{Per-mark labels.}{Beschriftungen an den einzelnen Marken.}}
  \SemioTableRow{\Key{sort} & \Key{string} & \Key{none} & \ApiText{Category order by total value.}{Reihenfolge der Kategorien nach Gesamtwert.}}
  \SemioTableRow{\Key{curve} & \Key{string} & --- & \ApiText{Interpolation used between consecutive points.}{Interpolation zwischen aufeinanderfolgenden Punkten.}}
}

\subsection{\Key{glyph}}

\SemioTableLong[text-size=8pt]{\Key{glyph}}{0.20,0.13,0.15,0.52}{\ApiKey & \ApiType & \ApiDefault & \ApiMeaning}{%
  \SemioTableRow{\Key{kind} & \Key{string} & --- & \ApiText{Which member of the family is drawn.}{Welches Mitglied der Familie gezeichnet wird.}}
  \SemioTableRow{\Key{axes} & \Key{string} & --- & \ApiText{Columns mapped onto the glyph or parallel axes.}{Spalten, die auf die Glyphen- bzw. Parallelachsen abgebildet werden.}}
  \SemioTableRow{\Key{label} & \Key{string} & --- & \ApiText{Column holding the visible caption.}{Spalte mit der sichtbaren Beschriftung.}}
  \SemioTableRow{\Key{columns} & \Key{integer} & --- & \ApiText{Number of columns of the layout grid.}{Anzahl der Spalten des Anordnungsgitters.}}
  \SemioTableRow{\Key{samples} & \Key{integer} & --- & \ApiText{Number of samples the curve or field is evaluated at.}{Anzahl der Stützstellen, an denen Kurve oder Feld ausgewertet werden.}}
}

\subsection{\Key{heatmap}}

\SemioTableLong[text-size=8pt]{\Key{heatmap}}{0.20,0.13,0.15,0.52}{\ApiKey & \ApiType & \ApiDefault & \ApiMeaning}{%
  \SemioTableRow{\Key{rowcol} & \Key{string} & --- & \ApiText{Column of the source table that names the heatmap rows.}{Spalte der Quelltabelle, die die Heatmap-Zeilen benennt.}}
  \SemioTableRow{\Key{colcol} & \Key{string} & --- & \ApiText{Column of the source table that names the heatmap columns.}{Spalte der Quelltabelle, die die Heatmap-Spalten benennt.}}
  \SemioTableRow{\Key{valuecol} & \Key{string} & --- & \ApiText{Column mapped onto the value.}{Spalte, die auf den Wert abgebildet wird.}}
  \SemioTableRow{\Key{cell} & \Key{string} & --- & \ApiText{Cell size in millimetres, or the column that carries it.}{Zellgröße in Millimetern oder die Spalte, die sie trägt.}}
  \SemioTableRow{\Key{scheme} & \Key{string} & --- & \ApiText{Named sequential or diverging colour scheme.}{Benanntes sequenzielles oder divergierendes Farbschema.}}
  \SemioTableRow{\Key{annotate} & \Key{boolean} & --- & \ApiText{Print the cell values on top of the shaded cells.}{Die Zellwerte auf die eingefärbten Zellen schreiben.}}
  \SemioTableRow{\Key{dendrogram} & \Key{string} & --- & \ApiText{Which margins carry a clustering dendrogram.}{Welche Ränder ein Cluster-Dendrogramm tragen.}}
  \SemioTableRow{\Key{marginals} & \Key{boolean} & --- & \ApiText{Draw the marginal summaries next to the matrix.}{Die Randzusammenfassungen neben der Matrix zeichnen.}}
  \SemioTableRow{\Key{labels} & \Key{boolean} & --- & \ApiText{Draw the captions.}{Die Beschriftungen zeichnen.}}
  \SemioTableRow{\Key{padding} & \Key{number} & --- & \ApiText{Padding between adjacent elements.}{Abstand zwischen benachbarten Elementen.}}
  \SemioTableRow{\Key{sparse} & \Key{boolean} & --- & \ApiText{Leave the cells without a value empty.}{Zellen ohne Wert leer lassen.}}
}

\subsection{\Key{hexbin}}

\SemioTableLong[text-size=8pt]{\Key{hexbin}}{0.20,0.13,0.15,0.52}{\ApiKey & \ApiType & \ApiDefault & \ApiMeaning}{%
  \SemioTableRow{\Key{xcol} & \Key{string} & --- & \ApiText{Column mapped onto the horizontal position.}{Spalte, die auf die waagerechte Position abgebildet wird.}}
  \SemioTableRow{\Key{ycol} & \Key{string} & --- & \ApiText{Column mapped onto the vertical position.}{Spalte, die auf die senkrechte Position abgebildet wird.}}
  \SemioTableRow{\Key{radius} & \Key{number} & --- & \ApiText{Radius in millimetres, or the column that drives it.}{Radius in Millimetern oder die Spalte, die ihn steuert.}}
  \SemioTableRow{\Key{mark} & \Key{string} & --- & \ApiText{Mark kind the data points are drawn with.}{Markenart, mit der die Datenpunkte gezeichnet werden.}}
  \SemioTableRow{\Key{density} & \Key{boolean} & --- & \ApiText{Normalise the counts to a density.}{Die Häufigkeiten zu einer Dichte normieren.}}
  \SemioTableRow{\Key{levels} & \Key{integer} & --- & \ApiText{Number of contour levels.}{Anzahl der Höhenlinienniveaus.}}
}

\subsection{\Key{histogram}}

\SemioTableLong[text-size=8pt]{\Key{histogram}}{0.20,0.13,0.15,0.52}{\ApiKey & \ApiType & \ApiDefault & \ApiMeaning}{%
  \SemioTableRow{\Key{bins} & \Key{integer} & --- & \ApiText{Number of histogram bins.}{Anzahl der Histogrammklassen.}}
  \SemioTableRow{\Key{thresholds} & \Key{string} & --- & \ApiText{Explicit bin thresholds.}{Explizite Klassengrenzen.}}
  \SemioTableRow{\Key{mode} & \Key{string} & --- & \ApiText{Which member of the family the preset renders.}{Welches Mitglied der Familie die Voreinstellung zeichnet.}}
  \SemioTableRow{\Key{relative} & \Key{boolean} & --- & \ApiText{Show relative frequencies instead of counts.}{Relative Häufigkeiten statt Absolutzahlen zeigen.}}
  \SemioTableRow{\Key{density} & \Key{boolean} & --- & \ApiText{Normalise the counts to a density.}{Die Häufigkeiten zu einer Dichte normieren.}}
  \SemioTableRow{\Key{cumulative} & \Key{boolean} & --- & \ApiText{Accumulate the values instead of showing them per bin.}{Die Werte kumulieren statt sie je Klasse zu zeigen.}}
  \SemioTableRow{\Key{variablewidth} & \Key{boolean} & --- & \ApiText{Allow bins of different widths.}{Klassen unterschiedlicher Breite zulassen.}}
  \SemioTableRow{\Key{coordinate} & \Key{string} & --- & \ApiText{Coordinate system the marks are placed in.}{Koordinatensystem, in dem die Marken platziert werden.}}
  \SemioTableRow{\Key{gap} & \Key{number} & --- & \ApiText{Millimetre gap between adjacent parts.}{Millimeterabstand zwischen benachbarten Teilen.}}
}

\subsection{\Key{kpi}}

\SemioTableLong[text-size=8pt]{\Key{kpi}}{0.20,0.13,0.15,0.52}{\ApiKey & \ApiType & \ApiDefault & \ApiMeaning}{%
  \SemioTableRow{\Key{mode} & \Key{string} & \Key{card} & \ApiText{How several series share one band.}{Wie sich mehrere Reihen ein Band teilen.}}
  \SemioTableRow{\Key{value} & \Key{string} & \Key{value} & \ApiText{The measured column.}{Die Spalte mit dem Messwert.}}
  \SemioTableRow{\Key{target} & \Key{string} & \Key{target} & \ApiText{The comparison column.}{Die Spalte mit dem Vergleichswert.}}
  \SemioTableRow{\Key{delta} & \Key{string} & \Key{delta} & \ApiText{The change column.}{Die Spalte mit der Veränderung.}}
  \SemioTableRow{\Key{unit} & \Key{string} & --- & \ApiText{Suffix printed after the number.}{Zusatz, der hinter die Zahl gesetzt wird.}}
  \SemioTableRow{\Key{columns} & \Key{integer} & \Key{3} & \ApiText{Cards per row.}{Karten je Zeile.}}
  \SemioTableRow{\Key{x} & \Key{string} & \Key{cat} & \ApiText{Independent / category column.}{Spalte der unabhängigen Größe bzw. der Kategorie.}}
  \SemioTableRow{\Key{y} & \Key{string} & \Key{val} & \ApiText{Value column.}{Spalte mit dem Wert.}}
  \SemioTableRow{\Key{y2} & \Key{string} & --- & \ApiText{Upper value column of span modes.}{Spalte des oberen Werts der Spannenmodi.}}
  \SemioTableRow{\Key{group} & \Key{string} & --- & \ApiText{Series / stack key column.}{Spalte des Reihen- bzw. Stapelschlüssels.}}
  \SemioTableRow{\Key{series} & \Key{string} & --- & \ApiText{Column that splits the rows into series.}{Spalte, die die Zeilen in Reihen aufteilt.}}
  \SemioTableRow{\Key{label} & \Key{string} & \Key{label} & \ApiText{The caption column.}{Die Spalte mit der Bildunterschrift.}}
  \SemioTableRow{\Key{padLeft} & \Key{number} & --- & \ApiText{Plot inset inside the canvas.}{Einrückung der Zeichenfläche innerhalb der Leinwand.}}
  \SemioTableRow{\Key{padRight} & \Key{number} & --- & \ApiText{Plot inset inside the canvas.}{Einrückung der Zeichenfläche innerhalb der Leinwand.}}
  \SemioTableRow{\Key{padTop} & \Key{number} & --- & \ApiText{Plot inset inside the canvas.}{Einrückung der Zeichenfläche innerhalb der Leinwand.}}
  \SemioTableRow{\Key{padBottom} & \Key{number} & --- & \ApiText{Plot inset inside the canvas.}{Einrückung der Zeichenfläche innerhalb der Leinwand.}}
  \SemioTableRow{\Key{yMin} & \Key{string} & \Key{from the data} & \ApiText{Value-axis domain override.}{Vorgabe des Wertebereichs der Werteachse.}}
  \SemioTableRow{\Key{yMax} & \Key{string} & \Key{from the data} & \ApiText{Value-axis domain override.}{Vorgabe des Wertebereichs der Werteachse.}}
  \SemioTableRow{\Key{axes} & \Key{boolean} & \Key{True} & \ApiText{Axis lines, tick labels, legend.}{Achsenlinien, Teilstrichbeschriftungen und Legende.}}
  \SemioTableRow{\Key{grid} & \Key{string} & \Key{none} & \ApiText{Grid lines.}{Gitterlinien hinter den Marken.}}
  \SemioTableRow{\Key{labels} & \Key{string} & \Key{none} & \ApiText{Per-mark labels.}{Beschriftungen an den einzelnen Marken.}}
  \SemioTableRow{\Key{sort} & \Key{string} & \Key{none} & \ApiText{Category order by total value.}{Reihenfolge der Kategorien nach Gesamtwert.}}
  \SemioTableRow{\Key{curve} & \Key{string} & --- & \ApiText{Interpolation used between consecutive points.}{Interpolation zwischen aufeinanderfolgenden Punkten.}}
}

\subsection{\Key{line}}

\SemioTableLong[text-size=8pt]{\Key{line}}{0.20,0.13,0.15,0.52}{\ApiKey & \ApiType & \ApiDefault & \ApiMeaning}{%
  \SemioTableRow{\Key{mode} & \Key{string} & \Key{plain} & \ApiText{Value transform before mapping.}{Werttransformation vor der Abbildung.}}
  \SemioTableRow{\Key{markers} & \Key{boolean} & \Key{False} & \ApiText{Dots on every data point.}{Punkte an jedem Datenpunkt.}}
  \SemioTableRow{\Key{highlight} & \Key{string} & --- & \ApiText{One series in the accent colour.}{Eine Reihe in der Akzentfarbe.}}
  \SemioTableRow{\Key{confidence} & \Key{boolean} & \Key{False} & \ApiText{Draw the lo/hi band behind.}{Das lo/hi-Band dahinter zeichnen.}}
  \SemioTableRow{\Key{lo} & \Key{string} & \Key{lo / hi} & \ApiText{The confidence band columns.}{Die Spalten des Konfidenzbands.}}
  \SemioTableRow{\Key{hi} & \Key{string} & \Key{lo / hi} & \ApiText{The confidence band columns.}{Die Spalten des Konfidenzbands.}}
  \SemioTableRow{\Key{x} & \Key{string} & \Key{cat} & \ApiText{Independent / category column.}{Spalte der unabhängigen Größe bzw. der Kategorie.}}
  \SemioTableRow{\Key{y} & \Key{string} & \Key{val} & \ApiText{Value column.}{Spalte mit dem Wert.}}
  \SemioTableRow{\Key{y2} & \Key{string} & --- & \ApiText{Upper value column of span modes.}{Spalte des oberen Werts der Spannenmodi.}}
  \SemioTableRow{\Key{group} & \Key{string} & --- & \ApiText{Series / stack key column.}{Spalte des Reihen- bzw. Stapelschlüssels.}}
  \SemioTableRow{\Key{series} & \Key{string} & --- & \ApiText{Column that splits the rows into series.}{Spalte, die die Zeilen in Reihen aufteilt.}}
  \SemioTableRow{\Key{label} & \Key{string} & --- & \ApiText{Text column for direct labels.}{Textspalte für Direktbeschriftungen.}}
  \SemioTableRow{\Key{padLeft} & \Key{number} & --- & \ApiText{Plot inset inside the canvas.}{Einrückung der Zeichenfläche innerhalb der Leinwand.}}
  \SemioTableRow{\Key{padRight} & \Key{number} & --- & \ApiText{Plot inset inside the canvas.}{Einrückung der Zeichenfläche innerhalb der Leinwand.}}
  \SemioTableRow{\Key{padTop} & \Key{number} & --- & \ApiText{Plot inset inside the canvas.}{Einrückung der Zeichenfläche innerhalb der Leinwand.}}
  \SemioTableRow{\Key{padBottom} & \Key{number} & --- & \ApiText{Plot inset inside the canvas.}{Einrückung der Zeichenfläche innerhalb der Leinwand.}}
  \SemioTableRow{\Key{yMin} & \Key{string} & \Key{from the data} & \ApiText{Value-axis domain override.}{Vorgabe des Wertebereichs der Werteachse.}}
  \SemioTableRow{\Key{yMax} & \Key{string} & \Key{from the data} & \ApiText{Value-axis domain override.}{Vorgabe des Wertebereichs der Werteachse.}}
  \SemioTableRow{\Key{axes} & \Key{boolean} & \Key{True} & \ApiText{Axis lines, tick labels, legend.}{Achsenlinien, Teilstrichbeschriftungen und Legende.}}
  \SemioTableRow{\Key{grid} & \Key{string} & \Key{none} & \ApiText{Grid lines.}{Gitterlinien hinter den Marken.}}
  \SemioTableRow{\Key{labels} & \Key{string} & \Key{none} & \ApiText{Per-mark labels.}{Beschriftungen an den einzelnen Marken.}}
  \SemioTableRow{\Key{sort} & \Key{string} & \Key{none} & \ApiText{Category order by total value.}{Reihenfolge der Kategorien nach Gesamtwert.}}
  \SemioTableRow{\Key{curve} & \Key{string} & \Key{linear} & \ApiText{Interpolation between points.}{Interpolation zwischen den Punkten.}}
}

\subsection{\Key{marimekko}}

\SemioTableLong[text-size=8pt]{\Key{marimekko}}{0.20,0.13,0.15,0.52}{\ApiKey & \ApiType & \ApiDefault & \ApiMeaning}{%
  \SemioTableRow{\Key{mode} & \Key{string} & --- & \ApiText{Which member of the family the preset renders.}{Welches Mitglied der Familie die Voreinstellung zeichnet.}}
  \SemioTableRow{\Key{column} & \Key{string} & --- & \ApiText{Column that splits the data into Marimekko columns.}{Spalte, die die Daten in Marimekko-Spalten aufteilt.}}
  \SemioTableRow{\Key{part} & \Key{string} & --- & \ApiText{Column that splits a column into its parts.}{Spalte, die eine Spalte in ihre Teile zerlegt.}}
  \SemioTableRow{\Key{amount} & \Key{string} & --- & \ApiText{Column holding the quantity each cell or segment encodes.}{Spalte mit der Menge, die jede Zelle oder jedes Segment codiert.}}
  \SemioTableRow{\Key{gap} & \Key{number} & --- & \ApiText{Millimetre gap between adjacent parts.}{Millimeterabstand zwischen benachbarten Teilen.}}
}

\subsection{\Key{monitoring}}

\ApiText{This family takes no options beyond \Key{variant} and \Key{data}.}{Diese Familie nimmt außer \Key{variant} und \Key{data} keine Optionen.}

\subsection{\Key{narrative}}

\SemioTableLong[text-size=8pt]{\Key{narrative}}{0.20,0.13,0.15,0.52}{\ApiKey & \ApiType & \ApiDefault & \ApiMeaning}{%
  \SemioTableRow{\Key{mode} & \Key{string} & \Key{storyline} & \ApiText{How several series share one band.}{Wie sich mehrere Reihen ein Band teilen.}}
  \SemioTableRow{\Key{panels} & \Key{integer} & \Key{4} & \ApiText{Panels of the panel/reveal/sequence modes.}{Felder der Modi panel, reveal und sequence.}}
  \SemioTableRow{\Key{x} & \Key{string} & \Key{cat} & \ApiText{Independent / category column.}{Spalte der unabhängigen Größe bzw. der Kategorie.}}
  \SemioTableRow{\Key{y} & \Key{string} & \Key{val} & \ApiText{Value column.}{Spalte mit dem Wert.}}
  \SemioTableRow{\Key{y2} & \Key{string} & --- & \ApiText{Upper value column of span modes.}{Spalte des oberen Werts der Spannenmodi.}}
  \SemioTableRow{\Key{group} & \Key{string} & --- & \ApiText{Series / stack key column.}{Spalte des Reihen- bzw. Stapelschlüssels.}}
  \SemioTableRow{\Key{series} & \Key{string} & --- & \ApiText{Column that splits the rows into series.}{Spalte, die die Zeilen in Reihen aufteilt.}}
  \SemioTableRow{\Key{label} & \Key{string} & --- & \ApiText{Text column for direct labels.}{Textspalte für Direktbeschriftungen.}}
  \SemioTableRow{\Key{padLeft} & \Key{number} & --- & \ApiText{Plot inset inside the canvas.}{Einrückung der Zeichenfläche innerhalb der Leinwand.}}
  \SemioTableRow{\Key{padRight} & \Key{number} & --- & \ApiText{Plot inset inside the canvas.}{Einrückung der Zeichenfläche innerhalb der Leinwand.}}
  \SemioTableRow{\Key{padTop} & \Key{number} & --- & \ApiText{Plot inset inside the canvas.}{Einrückung der Zeichenfläche innerhalb der Leinwand.}}
  \SemioTableRow{\Key{padBottom} & \Key{number} & --- & \ApiText{Plot inset inside the canvas.}{Einrückung der Zeichenfläche innerhalb der Leinwand.}}
  \SemioTableRow{\Key{yMin} & \Key{string} & \Key{from the data} & \ApiText{Value-axis domain override.}{Vorgabe des Wertebereichs der Werteachse.}}
  \SemioTableRow{\Key{yMax} & \Key{string} & \Key{from the data} & \ApiText{Value-axis domain override.}{Vorgabe des Wertebereichs der Werteachse.}}
  \SemioTableRow{\Key{axes} & \Key{boolean} & \Key{True} & \ApiText{Axis lines, tick labels, legend.}{Achsenlinien, Teilstrichbeschriftungen und Legende.}}
  \SemioTableRow{\Key{grid} & \Key{string} & \Key{none} & \ApiText{Grid lines.}{Gitterlinien hinter den Marken.}}
  \SemioTableRow{\Key{labels} & \Key{string} & \Key{none} & \ApiText{Per-mark labels.}{Beschriftungen an den einzelnen Marken.}}
  \SemioTableRow{\Key{sort} & \Key{string} & \Key{none} & \ApiText{Category order by total value.}{Reihenfolge der Kategorien nach Gesamtwert.}}
  \SemioTableRow{\Key{curve} & \Key{string} & --- & \ApiText{Interpolation used between consecutive points.}{Interpolation zwischen aufeinanderfolgenden Punkten.}}
}

\subsection{\Key{niche}}

\ApiText{This family takes no options beyond \Key{variant} and \Key{data}.}{Diese Familie nimmt außer \Key{variant} und \Key{data} keine Optionen.}

\subsection{\Key{parallel}}

\SemioTableLong[text-size=8pt]{\Key{parallel}}{0.20,0.13,0.15,0.52}{\ApiKey & \ApiType & \ApiDefault & \ApiMeaning}{%
  \SemioTableRow{\Key{axes} & \Key{string} & --- & \ApiText{Columns mapped onto the glyph or parallel axes.}{Spalten, die auf die Glyphen- bzw. Parallelachsen abgebildet werden.}}
  \SemioTableRow{\Key{label} & \Key{string} & --- & \ApiText{Column holding the visible caption.}{Spalte mit der sichtbaren Beschriftung.}}
  \SemioTableRow{\Key{curve} & \Key{boolean} & --- & \ApiText{Interpolation used between consecutive points.}{Interpolation zwischen aufeinanderfolgenden Punkten.}}
  \SemioTableRow{\Key{ticks} & \Key{integer} & --- & \ApiText{Number of ticks drawn per axis.}{Anzahl der Teilstriche je Achse.}}
}

\subsection{\Key{parliament}}

\SemioTableLong[text-size=8pt]{\Key{parliament}}{0.20,0.13,0.15,0.52}{\ApiKey & \ApiType & \ApiDefault & \ApiMeaning}{%
  \SemioTableRow{\Key{rows} & \Key{integer} & --- & \ApiText{Number of rows of the layout grid.}{Anzahl der Zeilen des Anordnungsgitters.}}
  \SemioTableRow{\Key{spread} & \Key{number} & --- & \ApiText{Deterministic spread of the placed marks.}{Deterministische Streuung der platzierten Marken.}}
  \SemioTableRow{\Key{inner} & \Key{number} & --- & \ApiText{Inner radius in millimetres; larger values open a hole.}{Innenradius in Millimetern; größere Werte öffnen ein Loch.}}
  \SemioTableRow{\Key{seat} & \Key{number} & --- & \ApiText{Millimetre radius of one seat.}{Millimeterradius eines Sitzes.}}
  \SemioTableRow{\Key{party} & \Key{string} & --- & \ApiText{Column holding the party of a seat block.}{Spalte mit der Partei eines Sitzblocks.}}
  \SemioTableRow{\Key{seats} & \Key{string} & --- & \ApiText{Column holding the seat count of a block.}{Spalte mit der Sitzzahl eines Blocks.}}
}

\subsection{\Key{performance}}

\ApiText{This family takes no options beyond \Key{variant} and \Key{data}.}{Diese Familie nimmt außer \Key{variant} und \Key{data} keine Optionen.}

\subsection{\Key{pie}}

\SemioTableLong[text-size=8pt]{\Key{pie}}{0.20,0.13,0.15,0.52}{\ApiKey & \ApiType & \ApiDefault & \ApiMeaning}{%
  \SemioTableRow{\Key{inner} & \Key{number} & --- & \ApiText{Inner radius in millimetres; larger values open a hole.}{Innenradius in Millimetern; größere Werte öffnen ein Loch.}}
  \SemioTableRow{\Key{outer} & \Key{number} & --- & \ApiText{Outer radius in millimetres.}{Außenradius in Millimetern.}}
  \SemioTableRow{\Key{start} & \Key{number} & --- & \ApiText{Start angle of the arc in degrees.}{Startwinkel des Bogens in Grad.}}
  \SemioTableRow{\Key{end} & \Key{number} & --- & \ApiText{End angle of the arc in degrees.}{Endwinkel des Bogens in Grad.}}
  \SemioTableRow{\Key{padangle} & \Key{number} & --- & \ApiText{Degrees of padding between adjacent slices.}{Gradabstand zwischen benachbarten Segmenten.}}
  \SemioTableRow{\Key{sort} & \Key{string} & --- & \ApiText{Order the entries are laid out in.}{Reihenfolge, in der die Einträge angeordnet werden.}}
  \SemioTableRow{\Key{explode} & \Key{number} & --- & \ApiText{Radial offset of the exploded slices in millimetres.}{Radialer Versatz der herausgezogenen Segmente in Millimetern.}}
  \SemioTableRow{\Key{rings} & \Key{integer} & --- & \ApiText{Number of concentric rings drawn.}{Anzahl der gezeichneten konzentrischen Ringe.}}
  \SemioTableRow{\Key{gauge} & \Key{boolean} & --- & \ApiText{Render the ring as a gauge with an open sector.}{Den Ring als Messuhr mit offenem Sektor zeichnen.}}
  \SemioTableRow{\Key{labels} & \Key{boolean} & --- & \ApiText{Draw the captions.}{Die Beschriftungen zeichnen.}}
}

\subsection{\Key{polar}}

\SemioTableLong[text-size=8pt]{\Key{polar}}{0.20,0.13,0.15,0.52}{\ApiKey & \ApiType & \ApiDefault & \ApiMeaning}{%
  \SemioTableRow{\Key{mark} & \Key{string} & --- & \ApiText{Mark kind the data points are drawn with.}{Markenart, mit der die Datenpunkte gezeichnet werden.}}
  \SemioTableRow{\Key{angle} & \Key{string} & --- & \ApiText{Angular position, in degrees or as the column that drives it.}{Winkelposition, in Grad oder als die Spalte, die sie steuert.}}
  \SemioTableRow{\Key{radiusv} & \Key{string} & --- & \ApiText{Column mapped onto the radius.}{Spalte, die auf den Radius abgebildet wird.}}
  \SemioTableRow{\Key{inner} & \Key{number} & --- & \ApiText{Inner radius in millimetres; larger values open a hole.}{Innenradius in Millimetern; größere Werte öffnen ein Loch.}}
  \SemioTableRow{\Key{start} & \Key{number} & --- & \ApiText{Start angle of the arc in degrees.}{Startwinkel des Bogens in Grad.}}
  \SemioTableRow{\Key{end} & \Key{number} & --- & \ApiText{End angle of the arc in degrees.}{Endwinkel des Bogens in Grad.}}
  \SemioTableRow{\Key{rings} & \Key{integer} & --- & \ApiText{Number of concentric rings drawn.}{Anzahl der gezeichneten konzentrischen Ringe.}}
  \SemioTableRow{\Key{grid} & \Key{boolean} & --- & \ApiText{Grid lines drawn behind the marks.}{Gitterlinien hinter den Marken.}}
}

\subsection{\Key{quadrant}}

\SemioTableLong[text-size=8pt]{\Key{quadrant}}{0.20,0.13,0.15,0.52}{\ApiKey & \ApiType & \ApiDefault & \ApiMeaning}{%
  \SemioTableRow{\Key{preset} & \Key{string} & --- & \ApiText{Named preset of captions and thresholds.}{Benannte Voreinstellung aus Beschriftungen und Schwellen.}}
  \SemioTableRow{\Key{xcol} & \Key{string} & --- & \ApiText{Column mapped onto the horizontal position.}{Spalte, die auf die waagerechte Position abgebildet wird.}}
  \SemioTableRow{\Key{ycol} & \Key{string} & --- & \ApiText{Column mapped onto the vertical position.}{Spalte, die auf die senkrechte Position abgebildet wird.}}
  \SemioTableRow{\Key{sizecol} & \Key{string} & --- & \ApiText{Column mapped onto the mark size.}{Spalte, die auf die Markengröße abgebildet wird.}}
  \SemioTableRow{\Key{itemlabel} & \Key{string} & --- & \ApiText{Column holding the label of one item.}{Spalte mit der Beschriftung eines Eintrags.}}
  \SemioTableRow{\Key{split} & \Key{number} & --- & \ApiText{Where the drawing splits into its two halves.}{Wo die Zeichnung in ihre zwei Hälften zerfällt.}}
  \SemioTableRow{\Key{tint} & \Key{boolean} & --- & \ApiText{Tint the four quadrants.}{Die vier Quadranten einfärben.}}
  \SemioTableRow{\Key{cells} & \Key{integer} & --- & \ApiText{Number of cells laid out, or the column naming them.}{Anzahl der angeordneten Zellen oder die Spalte, die sie benennt.}}
}

\subsection{\Key{quantile}}

\SemioTableLong[text-size=8pt]{\Key{quantile}}{0.20,0.13,0.15,0.52}{\ApiKey & \ApiType & \ApiDefault & \ApiMeaning}{%
  \SemioTableRow{\Key{kind} & \Key{string} & --- & \ApiText{Which member of the family is drawn.}{Welches Mitglied der Familie gezeichnet wird.}}
  \SemioTableRow{\Key{steps} & \Key{integer} & --- & \ApiText{Number of integration or iteration steps.}{Anzahl der Integrations- bzw. Iterationsschritte.}}
  \SemioTableRow{\Key{reference} & \Key{boolean} & --- & \ApiText{Reference value the marks are compared against.}{Referenzwert, mit dem die Marken verglichen werden.}}
  \SemioTableRow{\Key{quantiledots} & \Key{integer} & --- & \ApiText{Number of quantile dots drawn.}{Anzahl der gezeichneten Quantilpunkte.}}
}

\subsection{\Key{radar}}

\SemioTableLong[text-size=8pt]{\Key{radar}}{0.20,0.13,0.15,0.52}{\ApiKey & \ApiType & \ApiDefault & \ApiMeaning}{%
  \SemioTableRow{\Key{axis} & \Key{string} & --- & \ApiText{Column holding the radar axis names.}{Spalte mit den Namen der Netzdiagrammachsen.}}
  \SemioTableRow{\Key{series} & \Key{string} & --- & \ApiText{Column that splits the rows into series.}{Spalte, die die Zeilen in Reihen aufteilt.}}
  \SemioTableRow{\Key{shape} & \Key{string} & --- & \ApiText{Shape of the node marks.}{Form der Knotenmarken.}}
  \SemioTableRow{\Key{levels} & \Key{integer} & --- & \ApiText{Number of contour levels.}{Anzahl der Höhenlinienniveaus.}}
  \SemioTableRow{\Key{filled} & \Key{boolean} & --- & \ApiText{Fill the bands between contour levels.}{Die Bänder zwischen den Höhenlinien füllen.}}
  \SemioTableRow{\Key{star} & \Key{boolean} & --- & \ApiText{Draw a star grid instead of concentric polygons.}{Ein Sterngitter statt konzentrischer Vielecke zeichnen.}}
  \SemioTableRow{\Key{labels} & \Key{boolean} & --- & \ApiText{Draw the captions.}{Die Beschriftungen zeichnen.}}
}

\subsection{\Key{radial-partition}}

\SemioTableLong[text-size=8pt]{\Key{radial-partition}}{0.20,0.13,0.15,0.52}{\ApiKey & \ApiType & \ApiDefault & \ApiMeaning}{%
  \SemioTableRow{\Key{node} & \Key{string} & --- & \ApiText{Mark drawn at every node.}{Marke, die an jedem Knoten gezeichnet wird.}}
  \SemioTableRow{\Key{parent} & \Key{string} & --- & \ApiText{Column holding the parent identifier.}{Spalte mit der Kennung des Elternknotens.}}
  \SemioTableRow{\Key{depth} & \Key{string} & --- & \ApiText{First depth that is drawn; the lower levels stay hidden.}{Erste gezeichnete Tiefe; die darüberliegenden Ebenen bleiben verborgen.}}
  \SemioTableRow{\Key{amount} & \Key{string} & --- & \ApiText{Column holding the quantity each cell or segment encodes.}{Spalte mit der Menge, die jede Zelle oder jedes Segment codiert.}}
  \SemioTableRow{\Key{inner} & \Key{number} & --- & \ApiText{Inner radius in millimetres; larger values open a hole.}{Innenradius in Millimetern; größere Werte öffnen ein Loch.}}
  \SemioTableRow{\Key{levels} & \Key{integer} & --- & \ApiText{Number of contour levels.}{Anzahl der Höhenlinienniveaus.}}
  \SemioTableRow{\Key{ringgap} & \Key{number} & --- & \ApiText{Millimetre gap between consecutive rings.}{Millimeterabstand zwischen aufeinanderfolgenden Ringen.}}
}

\subsection{\Key{rank}}

\SemioTableLong[text-size=8pt]{\Key{rank}}{0.20,0.13,0.15,0.52}{\ApiKey & \ApiType & \ApiDefault & \ApiMeaning}{%
  \SemioTableRow{\Key{mode} & \Key{string} & \Key{bump} & \ApiText{How several series share one band.}{Wie sich mehrere Reihen ein Band teilen.}}
  \SemioTableRow{\Key{n} & \Key{integer} & \Key{3} & \ApiText{Ranks kept by the podium and top-n modes.}{Ränge, die die Modi podium und top-n behalten.}}
  \SemioTableRow{\Key{marker} & \Key{boolean} & \Key{True} & \ApiText{Dots on the rank path.}{Punkte auf dem Rangpfad.}}
  \SemioTableRow{\Key{x} & \Key{string} & \Key{cat} & \ApiText{Independent / category column.}{Spalte der unabhängigen Größe bzw. der Kategorie.}}
  \SemioTableRow{\Key{y} & \Key{string} & \Key{val} & \ApiText{Value column.}{Spalte mit dem Wert.}}
  \SemioTableRow{\Key{y2} & \Key{string} & --- & \ApiText{Upper value column of span modes.}{Spalte des oberen Werts der Spannenmodi.}}
  \SemioTableRow{\Key{group} & \Key{string} & --- & \ApiText{Series / stack key column.}{Spalte des Reihen- bzw. Stapelschlüssels.}}
  \SemioTableRow{\Key{series} & \Key{string} & --- & \ApiText{Column that splits the rows into series.}{Spalte, die die Zeilen in Reihen aufteilt.}}
  \SemioTableRow{\Key{label} & \Key{string} & --- & \ApiText{Text column for direct labels.}{Textspalte für Direktbeschriftungen.}}
  \SemioTableRow{\Key{padLeft} & \Key{number} & --- & \ApiText{Plot inset inside the canvas.}{Einrückung der Zeichenfläche innerhalb der Leinwand.}}
  \SemioTableRow{\Key{padRight} & \Key{number} & --- & \ApiText{Plot inset inside the canvas.}{Einrückung der Zeichenfläche innerhalb der Leinwand.}}
  \SemioTableRow{\Key{padTop} & \Key{number} & --- & \ApiText{Plot inset inside the canvas.}{Einrückung der Zeichenfläche innerhalb der Leinwand.}}
  \SemioTableRow{\Key{padBottom} & \Key{number} & --- & \ApiText{Plot inset inside the canvas.}{Einrückung der Zeichenfläche innerhalb der Leinwand.}}
  \SemioTableRow{\Key{yMin} & \Key{string} & \Key{from the data} & \ApiText{Value-axis domain override.}{Vorgabe des Wertebereichs der Werteachse.}}
  \SemioTableRow{\Key{yMax} & \Key{string} & \Key{from the data} & \ApiText{Value-axis domain override.}{Vorgabe des Wertebereichs der Werteachse.}}
  \SemioTableRow{\Key{axes} & \Key{boolean} & \Key{True} & \ApiText{Axis lines, tick labels, legend.}{Achsenlinien, Teilstrichbeschriftungen und Legende.}}
  \SemioTableRow{\Key{grid} & \Key{string} & \Key{none} & \ApiText{Grid lines.}{Gitterlinien hinter den Marken.}}
  \SemioTableRow{\Key{labels} & \Key{string} & \Key{none} & \ApiText{Per-mark labels.}{Beschriftungen an den einzelnen Marken.}}
  \SemioTableRow{\Key{sort} & \Key{string} & \Key{none} & \ApiText{Category order by total value.}{Reihenfolge der Kategorien nach Gesamtwert.}}
  \SemioTableRow{\Key{curve} & \Key{string} & --- & \ApiText{Interpolation used between consecutive points.}{Interpolation zwischen aufeinanderfolgenden Punkten.}}
}

\subsection{\Key{residual}}

\SemioTableLong[text-size=8pt]{\Key{residual}}{0.20,0.13,0.15,0.52}{\ApiKey & \ApiType & \ApiDefault & \ApiMeaning}{%
  \SemioTableRow{\Key{kind} & \Key{string} & --- & \ApiText{Which member of the family is drawn.}{Welches Mitglied der Familie gezeichnet wird.}}
  \SemioTableRow{\Key{fitted} & \Key{string} & --- & \ApiText{Column holding the fitted values.}{Spalte mit den angepassten Werten.}}
  \SemioTableRow{\Key{resid} & \Key{string} & --- & \ApiText{Column holding the residuals.}{Spalte mit den Residuen.}}
  \SemioTableRow{\Key{weight} & \Key{string} & --- & \ApiText{Column holding the edge weight.}{Spalte mit dem Kantengewicht.}}
  \SemioTableRow{\Key{smooth} & \Key{boolean} & --- & \ApiText{Draw the smoothing curve through the residuals.}{Die Glättungskurve durch die Residuen zeichnen.}}
  \SemioTableRow{\Key{bands} & \Key{boolean} & --- & \ApiText{Draw the residual reference bands.}{Die Referenzbänder der Residuen zeichnen.}}
}

\subsection{\Key{scatter}}

\SemioTableLong[text-size=8pt]{\Key{scatter}}{0.20,0.13,0.15,0.52}{\ApiKey & \ApiType & \ApiDefault & \ApiMeaning}{%
  \SemioTableRow{\Key{size} & \Key{string} & \Key{1.4} & \ApiText{Radius, through a sqrt scale when it is a column.}{Radius, über eine Wurzelskala, wenn eine Spalte angegeben ist.}}
  \SemioTableRow{\Key{shape} & \Key{string} & \Key{circle} & \ApiText{Shape of the node marks.}{Form der Knotenmarken.}}
  \SemioTableRow{\Key{color} & \Key{string} & --- & \ApiText{Categorical colour encoding.}{Kategoriale Farbcodierung.}}
  \SemioTableRow{\Key{trend} & \Key{string} & \Key{none} & \ApiText{Trend line fitted through the points.}{Trendlinie, die durch die Punkte gelegt wird.}}
  \SemioTableRow{\Key{confidence} & \Key{boolean} & \Key{False} & \ApiText{Band of one mean absolute residual around a trend.}{Band eines mittleren absoluten Residuums um einen Trend.}}
  \SemioTableRow{\Key{connected} & \Key{boolean} & \Key{False} & \ApiText{Join the points in row order.}{Die Punkte in Zeilenreihenfolge verbinden.}}
  \SemioTableRow{\Key{jitter} & \Key{number} & \Key{0} & \ApiText{Deterministic displacement to separate ties.}{Deterministische Verschiebung, die gleiche Werte trennt.}}
  \SemioTableRow{\Key{x} & \Key{string} & \Key{cat} & \ApiText{Independent / category column.}{Spalte der unabhängigen Größe bzw. der Kategorie.}}
  \SemioTableRow{\Key{y} & \Key{string} & \Key{val} & \ApiText{Value column.}{Spalte mit dem Wert.}}
  \SemioTableRow{\Key{y2} & \Key{string} & --- & \ApiText{Upper value column of span modes.}{Spalte des oberen Werts der Spannenmodi.}}
  \SemioTableRow{\Key{group} & \Key{string} & --- & \ApiText{Series / stack key column.}{Spalte des Reihen- bzw. Stapelschlüssels.}}
  \SemioTableRow{\Key{series} & \Key{string} & --- & \ApiText{Column that splits the rows into series.}{Spalte, die die Zeilen in Reihen aufteilt.}}
  \SemioTableRow{\Key{label} & \Key{string} & --- & \ApiText{Text column for direct labels.}{Textspalte für Direktbeschriftungen.}}
  \SemioTableRow{\Key{padLeft} & \Key{number} & --- & \ApiText{Plot inset inside the canvas.}{Einrückung der Zeichenfläche innerhalb der Leinwand.}}
  \SemioTableRow{\Key{padRight} & \Key{number} & --- & \ApiText{Plot inset inside the canvas.}{Einrückung der Zeichenfläche innerhalb der Leinwand.}}
  \SemioTableRow{\Key{padTop} & \Key{number} & --- & \ApiText{Plot inset inside the canvas.}{Einrückung der Zeichenfläche innerhalb der Leinwand.}}
  \SemioTableRow{\Key{padBottom} & \Key{number} & --- & \ApiText{Plot inset inside the canvas.}{Einrückung der Zeichenfläche innerhalb der Leinwand.}}
  \SemioTableRow{\Key{yMin} & \Key{string} & \Key{from the data} & \ApiText{Value-axis domain override.}{Vorgabe des Wertebereichs der Werteachse.}}
  \SemioTableRow{\Key{yMax} & \Key{string} & \Key{from the data} & \ApiText{Value-axis domain override.}{Vorgabe des Wertebereichs der Werteachse.}}
  \SemioTableRow{\Key{axes} & \Key{boolean} & \Key{True} & \ApiText{Axis lines, tick labels, legend.}{Achsenlinien, Teilstrichbeschriftungen und Legende.}}
  \SemioTableRow{\Key{grid} & \Key{string} & \Key{none} & \ApiText{Grid lines.}{Gitterlinien hinter den Marken.}}
  \SemioTableRow{\Key{labels} & \Key{string} & --- & \ApiText{Per-point text from the label column.}{Text je Punkt aus der Beschriftungsspalte.}}
  \SemioTableRow{\Key{sort} & \Key{string} & \Key{none} & \ApiText{Category order by total value.}{Reihenfolge der Kategorien nach Gesamtwert.}}
  \SemioTableRow{\Key{curve} & \Key{string} & --- & \ApiText{Interpolation used between consecutive points.}{Interpolation zwischen aufeinanderfolgenden Punkten.}}
}

\subsection{\Key{scatter-matrix}}

\SemioTableLong[text-size=8pt]{\Key{scatter-matrix}}{0.20,0.13,0.15,0.52}{\ApiKey & \ApiType & \ApiDefault & \ApiMeaning}{%
  \SemioTableRow{\Key{columns} & \Key{string} & --- & \ApiText{Number of columns of the layout grid.}{Anzahl der Spalten des Anordnungsgitters.}}
  \SemioTableRow{\Key{diagonal} & \Key{string} & --- & \ApiText{What the diagonal cells show.}{Was die Diagonalzellen zeigen.}}
  \SemioTableRow{\Key{radius} & \Key{number} & --- & \ApiText{Radius in millimetres, or the column that drives it.}{Radius in Millimetern oder die Spalte, die ihn steuert.}}
  \SemioTableRow{\Key{labels} & \Key{boolean} & --- & \ApiText{Draw the captions.}{Die Beschriftungen zeichnen.}}
}

\subsection{\Key{schedule}}

\ApiText{This family takes no options beyond \Key{variant} and \Key{data}.}{Diese Familie nimmt außer \Key{variant} und \Key{data} keine Optionen.}

\subsection{\Key{slope}}

\SemioTableLong[text-size=8pt]{\Key{slope}}{0.20,0.13,0.15,0.52}{\ApiKey & \ApiType & \ApiDefault & \ApiMeaning}{%
  \SemioTableRow{\Key{mode} & \Key{string} & \Key{slope} & \ApiText{How several series share one band.}{Wie sich mehrere Reihen ein Band teilen.}}
  \SemioTableRow{\Key{markers} & \Key{boolean} & \Key{True} & \ApiText{Dots at both ends.}{Punkte an beiden Enden.}}
  \SemioTableRow{\Key{x} & \Key{string} & \Key{cat} & \ApiText{Independent / category column.}{Spalte der unabhängigen Größe bzw. der Kategorie.}}
  \SemioTableRow{\Key{y} & \Key{string} & \Key{val} & \ApiText{Value column.}{Spalte mit dem Wert.}}
  \SemioTableRow{\Key{y2} & \Key{string} & --- & \ApiText{Upper value column of span modes.}{Spalte des oberen Werts der Spannenmodi.}}
  \SemioTableRow{\Key{group} & \Key{string} & --- & \ApiText{Series / stack key column.}{Spalte des Reihen- bzw. Stapelschlüssels.}}
  \SemioTableRow{\Key{series} & \Key{string} & --- & \ApiText{Column that splits the rows into series.}{Spalte, die die Zeilen in Reihen aufteilt.}}
  \SemioTableRow{\Key{label} & \Key{string} & --- & \ApiText{Text column for direct labels.}{Textspalte für Direktbeschriftungen.}}
  \SemioTableRow{\Key{padLeft} & \Key{number} & --- & \ApiText{Plot inset inside the canvas.}{Einrückung der Zeichenfläche innerhalb der Leinwand.}}
  \SemioTableRow{\Key{padRight} & \Key{number} & --- & \ApiText{Plot inset inside the canvas.}{Einrückung der Zeichenfläche innerhalb der Leinwand.}}
  \SemioTableRow{\Key{padTop} & \Key{number} & --- & \ApiText{Plot inset inside the canvas.}{Einrückung der Zeichenfläche innerhalb der Leinwand.}}
  \SemioTableRow{\Key{padBottom} & \Key{number} & --- & \ApiText{Plot inset inside the canvas.}{Einrückung der Zeichenfläche innerhalb der Leinwand.}}
  \SemioTableRow{\Key{yMin} & \Key{string} & \Key{from the data} & \ApiText{Value-axis domain override.}{Vorgabe des Wertebereichs der Werteachse.}}
  \SemioTableRow{\Key{yMax} & \Key{string} & \Key{from the data} & \ApiText{Value-axis domain override.}{Vorgabe des Wertebereichs der Werteachse.}}
  \SemioTableRow{\Key{axes} & \Key{boolean} & \Key{True} & \ApiText{Axis lines, tick labels, legend.}{Achsenlinien, Teilstrichbeschriftungen und Legende.}}
  \SemioTableRow{\Key{grid} & \Key{string} & \Key{none} & \ApiText{Grid lines.}{Gitterlinien hinter den Marken.}}
  \SemioTableRow{\Key{labels} & \Key{string} & --- & \ApiText{Per-mark labels.}{Beschriftungen an den einzelnen Marken.}}
  \SemioTableRow{\Key{sort} & \Key{string} & \Key{none} & \ApiText{Category order by total value.}{Reihenfolge der Kategorien nach Gesamtwert.}}
  \SemioTableRow{\Key{curve} & \Key{string} & --- & \ApiText{Interpolation between points.}{Interpolation zwischen den Punkten.}}
}

\subsection{\Key{spark}}

\SemioTableLong[text-size=8pt]{\Key{spark}}{0.20,0.13,0.15,0.52}{\ApiKey & \ApiType & \ApiDefault & \ApiMeaning}{%
  \SemioTableRow{\Key{mode} & \Key{string} & \Key{line} & \ApiText{The spark form.}{Die Form der Sparklinie.}}
  \SemioTableRow{\Key{band} & \Key{boolean} & \Key{False} & \ApiText{Min/max reference band behind the spark.}{Referenzband aus Minimum und Maximum hinter der Sparklinie.}}
  \SemioTableRow{\Key{x} & \Key{string} & \Key{cat} & \ApiText{Independent / category column.}{Spalte der unabhängigen Größe bzw. der Kategorie.}}
  \SemioTableRow{\Key{y} & \Key{string} & \Key{val} & \ApiText{Value column.}{Spalte mit dem Wert.}}
  \SemioTableRow{\Key{y2} & \Key{string} & --- & \ApiText{Upper value column of span modes.}{Spalte des oberen Werts der Spannenmodi.}}
  \SemioTableRow{\Key{group} & \Key{string} & --- & \ApiText{Series / stack key column.}{Spalte des Reihen- bzw. Stapelschlüssels.}}
  \SemioTableRow{\Key{series} & \Key{string} & --- & \ApiText{Column that splits the rows into series.}{Spalte, die die Zeilen in Reihen aufteilt.}}
  \SemioTableRow{\Key{label} & \Key{string} & --- & \ApiText{Text column for direct labels.}{Textspalte für Direktbeschriftungen.}}
  \SemioTableRow{\Key{padLeft} & \Key{number} & --- & \ApiText{Plot inset inside the canvas.}{Einrückung der Zeichenfläche innerhalb der Leinwand.}}
  \SemioTableRow{\Key{padRight} & \Key{number} & --- & \ApiText{Plot inset inside the canvas.}{Einrückung der Zeichenfläche innerhalb der Leinwand.}}
  \SemioTableRow{\Key{padTop} & \Key{number} & --- & \ApiText{Plot inset inside the canvas.}{Einrückung der Zeichenfläche innerhalb der Leinwand.}}
  \SemioTableRow{\Key{padBottom} & \Key{number} & --- & \ApiText{Plot inset inside the canvas.}{Einrückung der Zeichenfläche innerhalb der Leinwand.}}
  \SemioTableRow{\Key{yMin} & \Key{string} & \Key{from the data} & \ApiText{Value-axis domain override.}{Vorgabe des Wertebereichs der Werteachse.}}
  \SemioTableRow{\Key{yMax} & \Key{string} & \Key{from the data} & \ApiText{Value-axis domain override.}{Vorgabe des Wertebereichs der Werteachse.}}
  \SemioTableRow{\Key{axes} & \Key{boolean} & \Key{True} & \ApiText{Axis lines, tick labels, legend.}{Achsenlinien, Teilstrichbeschriftungen und Legende.}}
  \SemioTableRow{\Key{grid} & \Key{string} & \Key{none} & \ApiText{Grid lines.}{Gitterlinien hinter den Marken.}}
  \SemioTableRow{\Key{labels} & \Key{string} & \Key{none} & \ApiText{Per-mark labels.}{Beschriftungen an den einzelnen Marken.}}
  \SemioTableRow{\Key{sort} & \Key{string} & \Key{none} & \ApiText{Category order by total value.}{Reihenfolge der Kategorien nach Gesamtwert.}}
  \SemioTableRow{\Key{curve} & \Key{string} & --- & \ApiText{Interpolation between points.}{Interpolation zwischen den Punkten.}}
}

\subsection{\Key{statepanel}}

\SemioTableLong[text-size=8pt]{\Key{statepanel}}{0.20,0.13,0.15,0.52}{\ApiKey & \ApiType & \ApiDefault & \ApiMeaning}{%
  \SemioTableRow{\Key{mode} & \Key{string} & \Key{overview-detail} & \ApiText{How several series share one band.}{Wie sich mehrere Reihen ein Band teilen.}}
  \SemioTableRow{\Key{frames} & \Key{integer} & \Key{4} & \ApiText{Frames of the frame-based modes.}{Bilder der bildbasierten Modi.}}
  \SemioTableRow{\Key{from} & \Key{number} & --- & \ApiText{The selected data window.}{Das ausgewählte Datenfenster.}}
  \SemioTableRow{\Key{to} & \Key{number} & --- & \ApiText{The selected data window.}{Das ausgewählte Datenfenster.}}
  \SemioTableRow{\Key{x} & \Key{string} & \Key{cat} & \ApiText{Independent / category column.}{Spalte der unabhängigen Größe bzw. der Kategorie.}}
  \SemioTableRow{\Key{y} & \Key{string} & \Key{val} & \ApiText{Value column.}{Spalte mit dem Wert.}}
  \SemioTableRow{\Key{y2} & \Key{string} & --- & \ApiText{Upper value column of span modes.}{Spalte des oberen Werts der Spannenmodi.}}
  \SemioTableRow{\Key{group} & \Key{string} & --- & \ApiText{Series / stack key column.}{Spalte des Reihen- bzw. Stapelschlüssels.}}
  \SemioTableRow{\Key{series} & \Key{string} & --- & \ApiText{Column that splits the rows into series.}{Spalte, die die Zeilen in Reihen aufteilt.}}
  \SemioTableRow{\Key{label} & \Key{string} & --- & \ApiText{Text column for direct labels.}{Textspalte für Direktbeschriftungen.}}
  \SemioTableRow{\Key{padLeft} & \Key{number} & --- & \ApiText{Plot inset inside the canvas.}{Einrückung der Zeichenfläche innerhalb der Leinwand.}}
  \SemioTableRow{\Key{padRight} & \Key{number} & --- & \ApiText{Plot inset inside the canvas.}{Einrückung der Zeichenfläche innerhalb der Leinwand.}}
  \SemioTableRow{\Key{padTop} & \Key{number} & --- & \ApiText{Plot inset inside the canvas.}{Einrückung der Zeichenfläche innerhalb der Leinwand.}}
  \SemioTableRow{\Key{padBottom} & \Key{number} & --- & \ApiText{Plot inset inside the canvas.}{Einrückung der Zeichenfläche innerhalb der Leinwand.}}
  \SemioTableRow{\Key{yMin} & \Key{string} & \Key{from the data} & \ApiText{Value-axis domain override.}{Vorgabe des Wertebereichs der Werteachse.}}
  \SemioTableRow{\Key{yMax} & \Key{string} & \Key{from the data} & \ApiText{Value-axis domain override.}{Vorgabe des Wertebereichs der Werteachse.}}
  \SemioTableRow{\Key{axes} & \Key{boolean} & \Key{True} & \ApiText{Axis lines, tick labels, legend.}{Achsenlinien, Teilstrichbeschriftungen und Legende.}}
  \SemioTableRow{\Key{grid} & \Key{string} & \Key{none} & \ApiText{Grid lines.}{Gitterlinien hinter den Marken.}}
  \SemioTableRow{\Key{labels} & \Key{string} & \Key{none} & \ApiText{Per-mark labels.}{Beschriftungen an den einzelnen Marken.}}
  \SemioTableRow{\Key{sort} & \Key{string} & \Key{none} & \ApiText{Category order by total value.}{Reihenfolge der Kategorien nach Gesamtwert.}}
  \SemioTableRow{\Key{curve} & \Key{string} & --- & \ApiText{Interpolation used between consecutive points.}{Interpolation zwischen aufeinanderfolgenden Punkten.}}
}

\subsection{\Key{stem-leaf}}

\SemioTableLong[text-size=8pt]{\Key{stem-leaf}}{0.20,0.13,0.15,0.52}{\ApiKey & \ApiType & \ApiDefault & \ApiMeaning}{%
  \SemioTableRow{\Key{unit} & \Key{number} & --- & \ApiText{Data units one glyph or dot stands for.}{Dateneinheiten, für die ein Glyph oder Punkt steht.}}
  \SemioTableRow{\Key{stemplot} & \Key{boolean} & --- & \ApiText{Draw the classic stem-and-leaf column.}{Die klassische Stamm-Blatt-Spalte zeichnen.}}
}

\subsection{\Key{strip}}

\SemioTableLong[text-size=8pt]{\Key{strip}}{0.20,0.13,0.15,0.52}{\ApiKey & \ApiType & \ApiDefault & \ApiMeaning}{%
  \SemioTableRow{\Key{layout} & \Key{string} & --- & \ApiText{Layout algorithm that places the nodes.}{Layoutalgorithmus, der die Knoten platziert.}}
  \SemioTableRow{\Key{jitter} & \Key{number} & --- & \ApiText{Deterministic jitter amplitude that separates ties.}{Deterministische Streuweite, die gleiche Werte trennt.}}
  \SemioTableRow{\Key{radius} & \Key{number} & --- & \ApiText{Radius in millimetres, or the column that drives it.}{Radius in Millimetern oder die Spalte, die ihn steuert.}}
  \SemioTableRow{\Key{seed} & \Key{integer} & --- & \ApiText{Seed of the deterministic pseudo-random layout.}{Startwert des deterministischen Pseudozufallslayouts.}}
  \SemioTableRow{\Key{binned} & \Key{number} & --- & \ApiText{Bin width the strip points are snapped to; '0' keeps them continuous.}{Klassenbreite, auf die die Streifenpunkte gerundet werden; '0' lässt sie stetig.}}
}

\subsection{\Key{survey}}

\ApiText{This family takes no options beyond \Key{variant} and \Key{data}.}{Diese Familie nimmt außer \Key{variant} und \Key{data} keine Optionen.}

\subsection{\Key{table}}

\SemioTableLong[text-size=8pt]{\Key{table}}{0.20,0.13,0.15,0.52}{\ApiKey & \ApiType & \ApiDefault & \ApiMeaning}{%
  \SemioTableRow{\Key{preset} & \Key{string} & --- & \ApiText{Named preset of captions and thresholds.}{Benannte Voreinstellung aus Beschriftungen und Schwellen.}}
  \SemioTableRow{\Key{columns} & \Key{string} & --- & \ApiText{Number of columns of the layout grid.}{Anzahl der Spalten des Anordnungsgitters.}}
  \SemioTableRow{\Key{headers} & \Key{string} & --- & \ApiText{Comma list of column captions; empty reuses the column names.}{Kommaliste der Spaltenbeschriftungen; leer übernimmt die Spaltennamen.}}
  \SemioTableRow{\Key{cells} & \Key{string} & --- & \ApiText{Number of cells laid out, or the column naming them.}{Anzahl der angeordneten Zellen oder die Spalte, die sie benennt.}}
  \SemioTableRow{\Key{rows} & \Key{integer} & --- & \ApiText{Number of rows of the layout grid.}{Anzahl der Zeilen des Anordnungsgitters.}}
  \SemioTableRow{\Key{rowheight} & \Key{number} & --- & \ApiText{Millimetre height of one table row.}{Millimeterhöhe einer Tabellenzeile.}}
  \SemioTableRow{\Key{headrule} & \Key{boolean} & --- & \ApiText{Draw the rule under the header row.}{Die Linie unter der Kopfzeile zeichnen.}}
  \SemioTableRow{\Key{zebra} & \Key{boolean} & --- & \ApiText{Shade every other row.}{Jede zweite Zeile einfärben.}}
  \SemioTableRow{\Key{sortby} & \Key{string} & --- & \ApiText{Column the rows are sorted by.}{Spalte, nach der die Zeilen sortiert werden.}}
  \SemioTableRow{\Key{total} & \Key{boolean} & --- & \ApiText{Print the summary row.}{Die Summenzeile setzen.}}
  \SemioTableRow{\Key{decimals} & \Key{integer} & --- & \ApiText{Decimal places numeric cells are printed with.}{Nachkommastellen, mit denen numerische Zellen gesetzt werden.}}
}

\subsection{\Key{ternary}}

\SemioTableLong[text-size=8pt]{\Key{ternary}}{0.20,0.13,0.15,0.52}{\ApiKey & \ApiType & \ApiDefault & \ApiMeaning}{%
  \SemioTableRow{\Key{mark} & \Key{string} & --- & \ApiText{Mark kind the data points are drawn with.}{Markenart, mit der die Datenpunkte gezeichnet werden.}}
  \SemioTableRow{\Key{pa} & \Key{string} & --- & \ApiText{Column of the first ternary component.}{Spalte der ersten ternären Komponente.}}
  \SemioTableRow{\Key{pb} & \Key{string} & --- & \ApiText{Column of the second ternary component.}{Spalte der zweiten ternären Komponente.}}
  \SemioTableRow{\Key{pc} & \Key{string} & --- & \ApiText{Column of the third ternary component.}{Spalte der dritten ternären Komponente.}}
  \SemioTableRow{\Key{grid} & \Key{integer} & --- & \ApiText{Grid lines drawn behind the marks.}{Gitterlinien hinter den Marken.}}
  \SemioTableRow{\Key{radius} & \Key{number} & --- & \ApiText{Radius in millimetres, or the column that drives it.}{Radius in Millimetern oder die Spalte, die ihn steuert.}}
}

\subsection{\Key{text-viz}}

\ApiText{This family takes no options beyond \Key{variant} and \Key{data}.}{Diese Familie nimmt außer \Key{variant} und \Key{data} keine Optionen.}

\subsection{\Key{timeline}}

\SemioTableLong[text-size=8pt]{\Key{timeline}}{0.20,0.13,0.15,0.52}{\ApiKey & \ApiType & \ApiDefault & \ApiMeaning}{%
  \SemioTableRow{\Key{mode} & \Key{string} & \Key{interval} & \ApiText{How several series share one band.}{Wie sich mehrere Reihen ein Band teilen.}}
  \SemioTableRow{\Key{lane} & \Key{string} & \Key{lane} & \ApiText{The swimlane / row key.}{Der Schlüssel der Schwimmbahn bzw. Zeile.}}
  \SemioTableRow{\Key{start} & \Key{string} & \Key{start} & \ApiText{Interval start.}{Beginn des Intervalls.}}
  \SemioTableRow{\Key{end} & \Key{string} & \Key{end} & \ApiText{Interval end.}{Ende des Intervalls.}}
  \SemioTableRow{\Key{state} & \Key{string} & \Key{state} & \ApiText{Categorical state colouring.}{Kategoriale Zustandsfärbung.}}
  \SemioTableRow{\Key{columns} & \Key{integer} & \Key{7} & \ApiText{Cells per row of the calendar mode.}{Zellen je Zeile im Kalendermodus.}}
  \SemioTableRow{\Key{turns} & \Key{number} & \Key{2.5} & \ApiText{Revolutions of the spiral mode.}{Umdrehungen des Spiralmodus.}}
  \SemioTableRow{\Key{split} & \Key{number} & \Key{4} & \ApiText{Where forecast and fan modes leave the observed data.}{Wo die Modi forecast und fan die beobachteten Daten verlassen.}}
  \SemioTableRow{\Key{x} & \Key{string} & \Key{cat} & \ApiText{Independent / category column.}{Spalte der unabhängigen Größe bzw. der Kategorie.}}
  \SemioTableRow{\Key{y} & \Key{string} & \Key{val} & \ApiText{Value column.}{Spalte mit dem Wert.}}
  \SemioTableRow{\Key{y2} & \Key{string} & --- & \ApiText{Upper value column of span modes.}{Spalte des oberen Werts der Spannenmodi.}}
  \SemioTableRow{\Key{group} & \Key{string} & --- & \ApiText{Series / stack key column.}{Spalte des Reihen- bzw. Stapelschlüssels.}}
  \SemioTableRow{\Key{series} & \Key{string} & --- & \ApiText{Column that splits the rows into series.}{Spalte, die die Zeilen in Reihen aufteilt.}}
  \SemioTableRow{\Key{label} & \Key{string} & --- & \ApiText{Text column for direct labels.}{Textspalte für Direktbeschriftungen.}}
  \SemioTableRow{\Key{padLeft} & \Key{number} & --- & \ApiText{Plot inset inside the canvas.}{Einrückung der Zeichenfläche innerhalb der Leinwand.}}
  \SemioTableRow{\Key{padRight} & \Key{number} & --- & \ApiText{Plot inset inside the canvas.}{Einrückung der Zeichenfläche innerhalb der Leinwand.}}
  \SemioTableRow{\Key{padTop} & \Key{number} & --- & \ApiText{Plot inset inside the canvas.}{Einrückung der Zeichenfläche innerhalb der Leinwand.}}
  \SemioTableRow{\Key{padBottom} & \Key{number} & --- & \ApiText{Plot inset inside the canvas.}{Einrückung der Zeichenfläche innerhalb der Leinwand.}}
  \SemioTableRow{\Key{yMin} & \Key{string} & \Key{from the data} & \ApiText{Value-axis domain override.}{Vorgabe des Wertebereichs der Werteachse.}}
  \SemioTableRow{\Key{yMax} & \Key{string} & \Key{from the data} & \ApiText{Value-axis domain override.}{Vorgabe des Wertebereichs der Werteachse.}}
  \SemioTableRow{\Key{axes} & \Key{boolean} & \Key{True} & \ApiText{Axis lines, tick labels, legend.}{Achsenlinien, Teilstrichbeschriftungen und Legende.}}
  \SemioTableRow{\Key{grid} & \Key{string} & \Key{none} & \ApiText{Grid lines.}{Gitterlinien hinter den Marken.}}
  \SemioTableRow{\Key{labels} & \Key{string} & \Key{none} & \ApiText{Per-mark labels.}{Beschriftungen an den einzelnen Marken.}}
  \SemioTableRow{\Key{sort} & \Key{string} & \Key{none} & \ApiText{Category order by total value.}{Reihenfolge der Kategorien nach Gesamtwert.}}
  \SemioTableRow{\Key{curve} & \Key{string} & --- & \ApiText{Interpolation used between consecutive points.}{Interpolation zwischen aufeinanderfolgenden Punkten.}}
}

\subsection{\Key{uncertainty}}

\SemioTableLong[text-size=8pt]{\Key{uncertainty}}{0.20,0.13,0.15,0.52}{\ApiKey & \ApiType & \ApiDefault & \ApiMeaning}{%
  \SemioTableRow{\Key{mark} & \Key{string} & --- & \ApiText{Mark kind the data points are drawn with.}{Markenart, mit der die Datenpunkte gezeichnet werden.}}
  \SemioTableRow{\Key{estimate} & \Key{string} & --- & \ApiText{Column holding the point estimate.}{Spalte mit der Punktschätzung.}}
  \SemioTableRow{\Key{lower} & \Key{string} & --- & \ApiText{Column holding the lower interval bound.}{Spalte mit der unteren Intervallgrenze.}}
  \SemioTableRow{\Key{upper} & \Key{string} & --- & \ApiText{Column holding the upper interval bound.}{Spalte mit der oberen Intervallgrenze.}}
  \SemioTableRow{\Key{weight} & \Key{string} & --- & \ApiText{Column holding the edge weight.}{Spalte mit dem Kantengewicht.}}
  \SemioTableRow{\Key{itemlabel} & \Key{string} & --- & \ApiText{Column holding the label of one item.}{Spalte mit der Beschriftung eines Eintrags.}}
  \SemioTableRow{\Key{orient} & \Key{string} & --- & \ApiText{Orientation of the marks.}{Ausrichtung der Marken.}}
  \SemioTableRow{\Key{reference} & \Key{number} & --- & \ApiText{Reference value the marks are compared against.}{Referenzwert, mit dem die Marken verglichen werden.}}
  \SemioTableRow{\Key{levels} & \Key{integer} & --- & \ApiText{Number of contour levels.}{Anzahl der Höhenlinienniveaus.}}
  \SemioTableRow{\Key{cap} & \Key{number} & --- & \ApiText{Millimetre width of the interval end caps.}{Millimeterbreite der Abschlussstriche des Intervalls.}}
}

\subsection{\Key{unit}}

\SemioTableLong[text-size=8pt]{\Key{unit}}{0.20,0.13,0.15,0.52}{\ApiKey & \ApiType & \ApiDefault & \ApiMeaning}{%
  \SemioTableRow{\Key{mode} & \Key{string} & \Key{unit} & \ApiText{How several series share one band.}{Wie sich mehrere Reihen ein Band teilen.}}
  \SemioTableRow{\Key{columns} & \Key{integer} & \Key{10} & \ApiText{Glyphs per row.}{Glyphen je Zeile.}}
  \SemioTableRow{\Key{unitValue} & \Key{number} & \Key{1} & \ApiText{Data units one glyph stands for.}{Dateneinheiten, für die ein Glyph steht.}}
  \SemioTableRow{\Key{glyph} & \Key{string} & \Key{square} & \ApiText{Glyph one unit of the chart is drawn with.}{Glyph, mit dem eine Einheit des Diagramms gezeichnet wird.}}
  \SemioTableRow{\Key{x} & \Key{string} & \Key{cat} & \ApiText{Independent / category column.}{Spalte der unabhängigen Größe bzw. der Kategorie.}}
  \SemioTableRow{\Key{y} & \Key{string} & \Key{val} & \ApiText{Value column.}{Spalte mit dem Wert.}}
  \SemioTableRow{\Key{y2} & \Key{string} & --- & \ApiText{Upper value column of span modes.}{Spalte des oberen Werts der Spannenmodi.}}
  \SemioTableRow{\Key{group} & \Key{string} & --- & \ApiText{Series / stack key column.}{Spalte des Reihen- bzw. Stapelschlüssels.}}
  \SemioTableRow{\Key{series} & \Key{string} & --- & \ApiText{Column that splits the rows into series.}{Spalte, die die Zeilen in Reihen aufteilt.}}
  \SemioTableRow{\Key{label} & \Key{string} & --- & \ApiText{Text column for direct labels.}{Textspalte für Direktbeschriftungen.}}
  \SemioTableRow{\Key{padLeft} & \Key{number} & --- & \ApiText{Plot inset inside the canvas.}{Einrückung der Zeichenfläche innerhalb der Leinwand.}}
  \SemioTableRow{\Key{padRight} & \Key{number} & --- & \ApiText{Plot inset inside the canvas.}{Einrückung der Zeichenfläche innerhalb der Leinwand.}}
  \SemioTableRow{\Key{padTop} & \Key{number} & --- & \ApiText{Plot inset inside the canvas.}{Einrückung der Zeichenfläche innerhalb der Leinwand.}}
  \SemioTableRow{\Key{padBottom} & \Key{number} & --- & \ApiText{Plot inset inside the canvas.}{Einrückung der Zeichenfläche innerhalb der Leinwand.}}
  \SemioTableRow{\Key{yMin} & \Key{string} & \Key{from the data} & \ApiText{Value-axis domain override.}{Vorgabe des Wertebereichs der Werteachse.}}
  \SemioTableRow{\Key{yMax} & \Key{string} & \Key{from the data} & \ApiText{Value-axis domain override.}{Vorgabe des Wertebereichs der Werteachse.}}
  \SemioTableRow{\Key{axes} & \Key{boolean} & \Key{True} & \ApiText{Axis lines, tick labels, legend.}{Achsenlinien, Teilstrichbeschriftungen und Legende.}}
  \SemioTableRow{\Key{grid} & \Key{string} & \Key{none} & \ApiText{Grid lines.}{Gitterlinien hinter den Marken.}}
  \SemioTableRow{\Key{labels} & \Key{string} & \Key{none} & \ApiText{Per-mark labels.}{Beschriftungen an den einzelnen Marken.}}
  \SemioTableRow{\Key{sort} & \Key{string} & \Key{none} & \ApiText{Category order by total value.}{Reihenfolge der Kategorien nach Gesamtwert.}}
  \SemioTableRow{\Key{curve} & \Key{string} & --- & \ApiText{Interpolation used between consecutive points.}{Interpolation zwischen aufeinanderfolgenden Punkten.}}
}

\subsection{\Key{violin}}

\SemioTableLong[text-size=8pt]{\Key{violin}}{0.20,0.13,0.15,0.52}{\ApiKey & \ApiType & \ApiDefault & \ApiMeaning}{%
  \SemioTableRow{\Key{half} & \Key{boolean} & --- & \ApiText{Draw only one half of the violin.}{Nur eine Hälfte der Violine zeichnen.}}
  \SemioTableRow{\Key{split} & \Key{boolean} & --- & \ApiText{Where the drawing splits into its two halves.}{Wo die Zeichnung in ihre zwei Hälften zerfällt.}}
  \SemioTableRow{\Key{bean} & \Key{boolean} & --- & \ApiText{Add the bean-plot rug of the individual observations.}{Den Bean-Plot-Teppich der Einzelbeobachtungen ergänzen.}}
  \SemioTableRow{\Key{raincloud} & \Key{boolean} & --- & \ApiText{Add the raincloud strip of raw observations.}{Den Raincloud-Streifen der Rohdaten ergänzen.}}
  \SemioTableRow{\Key{box} & \Key{boolean} & --- & \ApiText{Draw a box plot inside the violin.}{Einen Kastenplot in die Violine zeichnen.}}
  \SemioTableRow{\Key{samples} & \Key{integer} & --- & \ApiText{Number of samples the curve or field is evaluated at.}{Anzahl der Stützstellen, an denen Kurve oder Feld ausgewertet werden.}}
  \SemioTableRow{\Key{bandwidth} & \Key{number} & --- & \ApiText{Kernel bandwidth of the density estimate; '0' uses Silverman's rule.}{Bandbreite des Kerns der Dichteschätzung; '0' nutzt Silvermans Regel.}}
}

\subsection{\Key{waffle}}

\SemioTableLong[text-size=8pt]{\Key{waffle}}{0.20,0.13,0.15,0.52}{\ApiKey & \ApiType & \ApiDefault & \ApiMeaning}{%
  \SemioTableRow{\Key{rows} & \Key{integer} & --- & \ApiText{Number of rows of the layout grid.}{Anzahl der Zeilen des Anordnungsgitters.}}
  \SemioTableRow{\Key{columns} & \Key{integer} & --- & \ApiText{Number of columns of the layout grid.}{Anzahl der Spalten des Anordnungsgitters.}}
  \SemioTableRow{\Key{shape} & \Key{string} & --- & \ApiText{Shape of the node marks.}{Form der Knotenmarken.}}
  \SemioTableRow{\Key{unit} & \Key{number} & --- & \ApiText{Data units one glyph or dot stands for.}{Dateneinheiten, für die ein Glyph oder Punkt steht.}}
  \SemioTableRow{\Key{part} & \Key{string} & --- & \ApiText{Column that splits a column into its parts.}{Spalte, die eine Spalte in ihre Teile zerlegt.}}
  \SemioTableRow{\Key{amount} & \Key{string} & --- & \ApiText{Column holding the quantity each cell or segment encodes.}{Spalte mit der Menge, die jede Zelle oder jedes Segment codiert.}}
  \SemioTableRow{\Key{gap} & \Key{number} & --- & \ApiText{Millimetre gap between adjacent parts.}{Millimeterabstand zwischen benachbarten Teilen.}}
}

\subsection{\Key{waterfall}}

\SemioTableLong[text-size=8pt]{\Key{waterfall}}{0.20,0.13,0.15,0.52}{\ApiKey & \ApiType & \ApiDefault & \ApiMeaning}{%
  \SemioTableRow{\Key{mode} & \Key{string} & \Key{waterfall} & \ApiText{How several series share one band.}{Wie sich mehrere Reihen ein Band teilen.}}
  \SemioTableRow{\Key{connector} & \Key{boolean} & \Key{False} & \ApiText{The step rule between consecutive bars.}{Die Stufenregel zwischen aufeinanderfolgenden Balken.}}
  \SemioTableRow{\Key{total} & \Key{boolean} & \Key{True} & \ApiText{A closing total bar from the baseline.}{Ein abschließender Summenbalken von der Grundlinie aus.}}
  \SemioTableRow{\Key{padding} & \Key{number} & \Key{0.3} & \ApiText{Fraction of the band left empty.}{Anteil des Bands, der leer bleibt.}}
  \SemioTableRow{\Key{x} & \Key{string} & \Key{cat} & \ApiText{Independent / category column.}{Spalte der unabhängigen Größe bzw. der Kategorie.}}
  \SemioTableRow{\Key{y} & \Key{string} & \Key{val} & \ApiText{Value column.}{Spalte mit dem Wert.}}
  \SemioTableRow{\Key{y2} & \Key{string} & --- & \ApiText{Upper value column of span modes.}{Spalte des oberen Werts der Spannenmodi.}}
  \SemioTableRow{\Key{group} & \Key{string} & --- & \ApiText{Series / stack key column.}{Spalte des Reihen- bzw. Stapelschlüssels.}}
  \SemioTableRow{\Key{series} & \Key{string} & --- & \ApiText{Column that splits the rows into series.}{Spalte, die die Zeilen in Reihen aufteilt.}}
  \SemioTableRow{\Key{label} & \Key{string} & --- & \ApiText{Text column for direct labels.}{Textspalte für Direktbeschriftungen.}}
  \SemioTableRow{\Key{padLeft} & \Key{number} & --- & \ApiText{Plot inset inside the canvas.}{Einrückung der Zeichenfläche innerhalb der Leinwand.}}
  \SemioTableRow{\Key{padRight} & \Key{number} & --- & \ApiText{Plot inset inside the canvas.}{Einrückung der Zeichenfläche innerhalb der Leinwand.}}
  \SemioTableRow{\Key{padTop} & \Key{number} & --- & \ApiText{Plot inset inside the canvas.}{Einrückung der Zeichenfläche innerhalb der Leinwand.}}
  \SemioTableRow{\Key{padBottom} & \Key{number} & --- & \ApiText{Plot inset inside the canvas.}{Einrückung der Zeichenfläche innerhalb der Leinwand.}}
  \SemioTableRow{\Key{yMin} & \Key{string} & \Key{from the data} & \ApiText{Value-axis domain override.}{Vorgabe des Wertebereichs der Werteachse.}}
  \SemioTableRow{\Key{yMax} & \Key{string} & \Key{from the data} & \ApiText{Value-axis domain override.}{Vorgabe des Wertebereichs der Werteachse.}}
  \SemioTableRow{\Key{axes} & \Key{boolean} & \Key{True} & \ApiText{Axis lines, tick labels, legend.}{Achsenlinien, Teilstrichbeschriftungen und Legende.}}
  \SemioTableRow{\Key{grid} & \Key{string} & \Key{none} & \ApiText{Grid lines.}{Gitterlinien hinter den Marken.}}
  \SemioTableRow{\Key{labels} & \Key{string} & \Key{none} & \ApiText{Per-mark labels.}{Beschriftungen an den einzelnen Marken.}}
  \SemioTableRow{\Key{sort} & \Key{string} & \Key{none} & \ApiText{Category order by total value.}{Reihenfolge der Kategorien nach Gesamtwert.}}
  \SemioTableRow{\Key{curve} & \Key{string} & --- & \ApiText{Interpolation used between consecutive points.}{Interpolation zwischen aufeinanderfolgenden Punkten.}}
}

\section{\ApiText{Hierarchy}{Hierarchie}}

\ApiText{Trees, dendrograms, treemaps, partitions and circle packings — everything that reads a parent-linked table.}{Bäume, Dendrogramme, Kacheldiagramme, Partitionen und Kreispackungen — alles, was eine elternverknüpfte Tabelle liest.}

\subsection{\Key{dendrogram}}

\SemioTableLong[text-size=8pt]{\Key{dendrogram}}{0.20,0.13,0.15,0.52}{\ApiKey & \ApiType & \ApiDefault & \ApiMeaning}{%
  \SemioTableRow{\Key{heightColumn} & \Key{string} & --- & \ApiText{Column with a merge height; empty uses the cluster depth.}{Spalte mit einer Verschmelzungshöhe; leer nutzt die Clustertiefe.}}
  \SemioTableRow{\Key{scale} & \Key{string} & --- & \ApiText{Named scale the values pass through.}{Benannte Skala, durch die die Werte laufen.}}
  \SemioTableRow{\Key{fill} & \Key{string} & --- & \ApiText{Fill colour of the node marks.}{Füllfarbe der Knotenmarken.}}
  \SemioTableRow{\Key{labels} & \Key{string} & --- & \ApiText{Draw the captions.}{Die Beschriftungen zeichnen.}}
  \SemioTableRow{\Key{link} & \Key{string} & --- & \ApiText{Shape of the link between parent and child.}{Form der Verbindung zwischen Eltern- und Kindknoten.}}
  \SemioTableRow{\Key{node} & \Key{string} & --- & \ApiText{Mark drawn at every node.}{Marke, die an jedem Knoten gezeichnet wird.}}
  \SemioTableRow{\Key{nodeRadius} & \Key{string} & --- & \ApiText{Millimetre radius of a node mark.}{Millimeterradius einer Knotenmarke.}}
  \SemioTableRow{\Key{orientation} & \Key{string} & --- & \ApiText{Primary direction of the layout.}{Hauptrichtung der Anordnung.}}
}

\subsection{\Key{pack}}

\SemioTableLong[text-size=8pt]{\Key{pack}}{0.20,0.13,0.15,0.52}{\ApiKey & \ApiType & \ApiDefault & \ApiMeaning}{%
  \SemioTableRow{\Key{depth} & \Key{integer} & --- & \ApiText{First depth that is drawn; the lower levels stay hidden.}{Erste gezeichnete Tiefe; die darüberliegenden Ebenen bleiben verborgen.}}
  \SemioTableRow{\Key{labels} & \Key{string} & --- & \ApiText{Draw the captions.}{Die Beschriftungen zeichnen.}}
  \SemioTableRow{\Key{padding} & \Key{string} & --- & \ApiText{Padding between adjacent elements.}{Abstand zwischen benachbarten Elementen.}}
  \SemioTableRow{\Key{radius} & \Key{string} & --- & \ApiText{Radius in millimetres, or the column that drives it.}{Radius in Millimetern oder die Spalte, die ihn steuert.}}
}

\subsection{\Key{partition}}

\SemioTableLong[text-size=8pt]{\Key{partition}}{0.20,0.13,0.15,0.52}{\ApiKey & \ApiType & \ApiDefault & \ApiMeaning}{%
  \SemioTableRow{\Key{depth} & \Key{integer} & --- & \ApiText{First depth that is drawn; the lower levels stay hidden.}{Erste gezeichnete Tiefe; die darüberliegenden Ebenen bleiben verborgen.}}
  \SemioTableRow{\Key{orientation} & \Key{string} & --- & \ApiText{Primary direction of the layout.}{Hauptrichtung der Anordnung.}}
  \SemioTableRow{\Key{round} & \Key{string} & --- & \ApiText{Round the band positions to whole millimetres.}{Die Bandpositionen auf ganze Millimeter runden.}}
}

\subsection{\Key{phylogram}}

\SemioTableLong[text-size=8pt]{\Key{phylogram}}{0.20,0.13,0.15,0.52}{\ApiKey & \ApiType & \ApiDefault & \ApiMeaning}{%
  \SemioTableRow{\Key{lengthColumn} & \Key{string} & --- & \ApiText{Column with the branch length of a phylogram.}{Spalte mit der Astlänge eines Phylogramms.}}
  \SemioTableRow{\Key{scale} & \Key{string} & --- & \ApiText{Named scale the values pass through.}{Benannte Skala, durch die die Werte laufen.}}
  \SemioTableRow{\Key{link} & \Key{string} & --- & \ApiText{Shape of the link between parent and child.}{Form der Verbindung zwischen Eltern- und Kindknoten.}}
  \SemioTableRow{\Key{orientation} & \Key{string} & --- & \ApiText{Primary direction of the layout.}{Hauptrichtung der Anordnung.}}
}

\subsection{\Key{sunburst}}

\SemioTableLong[text-size=8pt]{\Key{sunburst}}{0.20,0.13,0.15,0.52}{\ApiKey & \ApiType & \ApiDefault & \ApiMeaning}{%
  \SemioTableRow{\Key{depth} & \Key{integer} & --- & \ApiText{First depth that is drawn; the lower levels stay hidden.}{Erste gezeichnete Tiefe; die darüberliegenden Ebenen bleiben verborgen.}}
  \SemioTableRow{\Key{innerRadius} & \Key{string} & --- & \ApiText{Millimetre radius the rings or axes start at.}{Millimeterradius, bei dem die Ringe bzw. Achsen beginnen.}}
}

\subsection{\Key{tree}}

\SemioTableLong[text-size=8pt]{\Key{tree}}{0.20,0.13,0.15,0.52}{\ApiKey & \ApiType & \ApiDefault & \ApiMeaning}{%
  \SemioTableRow{\Key{fill} & \Key{string} & --- & \ApiText{Fill colour of the node marks.}{Füllfarbe der Knotenmarken.}}
  \SemioTableRow{\Key{labels} & \Key{string} & --- & \ApiText{Draw the captions.}{Die Beschriftungen zeichnen.}}
  \SemioTableRow{\Key{link} & \Key{string} & --- & \ApiText{Shape of the link between parent and child.}{Form der Verbindung zwischen Eltern- und Kindknoten.}}
  \SemioTableRow{\Key{node} & \Key{string} & --- & \ApiText{Mark drawn at every node.}{Marke, die an jedem Knoten gezeichnet wird.}}
  \SemioTableRow{\Key{nodeHeight} & \Key{string} & --- & \ApiText{Millimetre height reserved for one node.}{Millimeterhöhe, die für einen Knoten reserviert ist.}}
  \SemioTableRow{\Key{nodeRadius} & \Key{string} & --- & \ApiText{Millimetre radius of a node mark.}{Millimeterradius einer Knotenmarke.}}
  \SemioTableRow{\Key{nodeWidth} & \Key{string} & --- & \ApiText{Millimetre width reserved for one node.}{Millimeterbreite, die für einen Knoten reserviert ist.}}
  \SemioTableRow{\Key{orientation} & \Key{string} & --- & \ApiText{Primary direction of the layout.}{Hauptrichtung der Anordnung.}}
  \SemioTableRow{\Key{path} & \Key{string} & --- & \ApiText{Column holding the slash-joined path of a node.}{Spalte mit dem durch Schrägstriche verbundenen Pfad eines Knotens.}}
  \SemioTableRow{\Key{radialOffset} & \Key{string} & --- & \ApiText{Millimetre offset of the radial root ring.}{Millimeterversatz des radialen Wurzelrings.}}
  \SemioTableRow{\Key{separationDepth} & \Key{string} & --- & \ApiText{Extra separation between tree depths.}{Zusätzlicher Abstand zwischen den Baumtiefen.}}
  \SemioTableRow{\Key{separationOther} & \Key{string} & --- & \ApiText{Extra separation between siblings of different parents.}{Zusätzlicher Abstand zwischen Geschwistern verschiedener Eltern.}}
  \SemioTableRow{\Key{sort} & \Key{string} & --- & \ApiText{Order the entries are laid out in.}{Reihenfolge, in der die Einträge angeordnet werden.}}
}

\subsection{\Key{treemap}}

\SemioTableLong[text-size=8pt]{\Key{treemap}}{0.20,0.13,0.15,0.52}{\ApiKey & \ApiType & \ApiDefault & \ApiMeaning}{%
  \SemioTableRow{\Key{depth} & \Key{integer} & --- & \ApiText{First depth that is drawn; the lower levels stay hidden.}{Erste gezeichnete Tiefe; die darüberliegenden Ebenen bleiben verborgen.}}
  \SemioTableRow{\Key{padding} & \Key{string} & --- & \ApiText{Padding between adjacent elements.}{Abstand zwischen benachbarten Elementen.}}
  \SemioTableRow{\Key{relaxation} & \Key{string} & --- & \ApiText{Relaxation passes that even out the tiles.}{Entspannungsdurchläufe, die die Kacheln ausgleichen.}}
  \SemioTableRow{\Key{tile} & \Key{string} & --- & \ApiText{Tiling algorithm that fills the rectangle.}{Kachelalgorithmus, der das Rechteck füllt.}}
}

\section{\ApiText{Network and flow}{Netzwerk und Fluss}}

\ApiText{Graphs, adjacency matrices, arc and chord diagrams, edge bundling, and the flow namespaces sankey, alluvial and parallel sets.}{Graphen, Adjazenzmatrizen, Bogen- und Sehnendiagramme, Kantenbündelung sowie die Flussnamensräume Sankey, Alluvial und Parallelmengen.}

\subsection{\Key{adjacency-matrix}}

\SemioTableLong[text-size=8pt]{\Key{adjacency-matrix}}{0.20,0.13,0.15,0.52}{\ApiKey & \ApiType & \ApiDefault & \ApiMeaning}{%
  \SemioTableRow{\Key{nodes} & \Key{string} & --- & \ApiText{Name of the node table.}{Name der Knotentabelle.}}
  \SemioTableRow{\Key{id} & \Key{string} & --- & \ApiText{Column holding the row identifier.}{Spalte mit der Zeilenkennung.}}
  \SemioTableRow{\Key{group} & \Key{string} & --- & \ApiText{Column that groups the nodes.}{Spalte, die die Knoten gruppiert.}}
  \SemioTableRow{\Key{source} & \Key{string} & --- & \ApiText{Column holding the source node of an edge.}{Spalte mit dem Quellknoten einer Kante.}}
  \SemioTableRow{\Key{target} & \Key{string} & --- & \ApiText{Column holding the target node of an edge, or the target value.}{Spalte mit dem Zielknoten einer Kante oder mit dem Zielwert.}}
  \SemioTableRow{\Key{weight} & \Key{string} & --- & \ApiText{Column holding the edge weight.}{Spalte mit dem Kantengewicht.}}
  \SemioTableRow{\Key{label} & \Key{string} & --- & \ApiText{Column holding the visible caption.}{Spalte mit der sichtbaren Beschriftung.}}
  \SemioTableRow{\Key{layout} & \Key{string} & --- & \ApiText{Layout algorithm that places the nodes.}{Layoutalgorithmus, der die Knoten platziert.}}
  \SemioTableRow{\Key{iterations} & \Key{integer} & --- & \ApiText{Fixed iteration count of the deterministic layout.}{Feste Iterationszahl des deterministischen Layouts.}}
  \SemioTableRow{\Key{seed} & \Key{number} & --- & \ApiText{Seed of the deterministic pseudo-random layout.}{Startwert des deterministischen Pseudozufallslayouts.}}
  \SemioTableRow{\Key{order} & \Key{string} & \Key{index} & \ApiText{Row and column ordering.}{Reihenfolge von Zeilen und Spalten.}}
  \SemioTableRow{\Key{columns} & \Key{integer} & --- & \ApiText{Number of columns of the layout grid.}{Anzahl der Spalten des Anordnungsgitters.}}
  \SemioTableRow{\Key{root} & \Key{integer} & --- & \ApiText{Node the layout starts from.}{Knoten, bei dem das Layout beginnt.}}
  \SemioTableRow{\Key{sweeps} & \Key{integer} & --- & \ApiText{Ordering sweeps of the crossing reduction.}{Sortierdurchläufe der Kreuzungsreduktion.}}
  \SemioTableRow{\Key{orientation} & \Key{string} & --- & \ApiText{Primary direction of the layout.}{Hauptrichtung der Anordnung.}}
  \SemioTableRow{\Key{shape} & \Key{string} & --- & \ApiText{Shape of the node marks.}{Form der Knotenmarken.}}
  \SemioTableRow{\Key{size} & \Key{number} & --- & \ApiText{Millimetre size of a node mark.}{Millimetergröße einer Knotenmarke.}}
  \SemioTableRow{\Key{sizeBy} & \Key{string} & --- & \ApiText{What the node size encodes.}{Was die Knotengröße codiert.}}
  \SemioTableRow{\Key{directed} & \Key{boolean} & --- & \ApiText{Draw arrow heads, so the edges read as directed.}{Pfeilspitzen zeichnen, sodass die Kanten gerichtet wirken.}}
  \SemioTableRow{\Key{curvature} & \Key{number} & --- & \ApiText{Curvature of a link, as a fraction of its span.}{Krümmung einer Kante als Bruchteil ihrer Spannweite.}}
  \SemioTableRow{\Key{weightWidth} & \Key{boolean} & --- & \ApiText{Derive the link width from the weight.}{Die Kantenbreite aus dem Gewicht ableiten.}}
  \SemioTableRow{\Key{signed} & \Key{boolean} & --- & \ApiText{Colour negative weights in the danger hue.}{Negative Gewichte im Warnfarbton einfärben.}}
  \SemioTableRow{\Key{colorBy} & \Key{string} & --- & \ApiText{What the node or link colour encodes.}{Was die Knoten- bzw. Kantenfarbe codiert.}}
  \SemioTableRow{\Key{labels} & \Key{boolean} & \Key{True} & \ApiText{Row and column captions.}{Zeilen- und Spaltenbeschriftungen.}}
  \SemioTableRow{\Key{linkWidth} & \Key{number} & --- & \ApiText{Millimetre width of a link.}{Millimeterbreite einer Kante.}}
  \SemioTableRow{\Key{linkDistance} & \Key{number} & --- & \ApiText{Rest length of a link in the force layout.}{Ruhelänge einer Kante im Kräftelayout.}}
  \SemioTableRow{\Key{charge} & \Key{number} & --- & \ApiText{Many-body charge of the force layout; negative values repel.}{Ladung der Vielteilchenkraft im Kräftelayout; negative Werte stoßen ab.}}
  \SemioTableRow{\Key{forces} & \Key{string} & --- & \ApiText{Forces the deterministic layout applies.}{Kräfte, die das deterministische Layout anwendet.}}
  \SemioTableRow{\Key{padding} & \Key{string} & \Key{3} & \ApiText{Padding between adjacent elements.}{Abstand zwischen benachbarten Elementen.}}
  \SemioTableRow{\Key{symmetric} & \Key{boolean} & \Key{True} & \ApiText{Mirror every edge so an undirected graph fills both triangles.}{Jede Kante spiegeln, damit ein ungerichteter Graph beide Dreiecke füllt.}}
  \SemioTableRow{\Key{cellGap} & \Key{number} & \Key{0.15} & \ApiText{Gap between cells.}{Abstand zwischen den Zellen.}}
  \SemioTableRow{\Key{scaleBy} & \Key{string} & \Key{weight} & \ApiText{Cell shade from |weight| or from presence alone.}{Zellton aus |weight| oder allein aus dem Vorhandensein.}}
  \SemioTableRow{\Key{diagonal} & \Key{boolean} & \Key{False} & \ApiText{Draw the diagonal cells.}{Die Diagonalzellen zeichnen.}}
  \SemioTableRow{\Key{labelWidth} & \Key{number} & \Key{9} & \ApiText{Millimetre width reserved for the labels.}{Millimeterbreite, die für die Beschriftungen reserviert ist.}}
  \SemioTableRow{\Key{blocks} & \Key{boolean} & \Key{False} & \ApiText{Hairlines between consecutive ordering groups.}{Haarlinien zwischen aufeinanderfolgenden Sortiergruppen.}}
  \SemioTableRow{\Key{split} & \Key{boolean} & --- & \ApiText{Where the drawing splits into its two halves.}{Wo die Zeichnung in ihre zwei Hälften zerfällt.}}
}

\subsection{\Key{alluvial}}

\SemioTableLong[text-size=8pt]{\Key{alluvial}}{0.20,0.13,0.15,0.52}{\ApiKey & \ApiType & \ApiDefault & \ApiMeaning}{%
  \SemioTableRow{\Key{iterations} & \Key{integer} & --- & \ApiText{Relaxation passes.}{Entspannungsdurchläufe.}}
  \SemioTableRow{\Key{curvature} & \Key{number} & --- & \ApiText{Horizontal control-point fraction of a link ribbon.}{Waagerechter Kontrollpunktanteil eines Verbindungsbands.}}
  \SemioTableRow{\Key{nodes} & \Key{string} & --- & \ApiText{Name of the node table.}{Name der Knotentabelle.}}
  \SemioTableRow{\Key{id} & \Key{string} & --- & \ApiText{Column holding the row identifier.}{Spalte mit der Zeilenkennung.}}
  \SemioTableRow{\Key{group} & \Key{string} & --- & \ApiText{Column that groups the nodes.}{Spalte, die die Knoten gruppiert.}}
  \SemioTableRow{\Key{source} & \Key{string} & --- & \ApiText{Column holding the source node of an edge.}{Spalte mit dem Quellknoten einer Kante.}}
  \SemioTableRow{\Key{target} & \Key{string} & --- & \ApiText{Column holding the target node of an edge, or the target value.}{Spalte mit dem Zielknoten einer Kante oder mit dem Zielwert.}}
  \SemioTableRow{\Key{weight} & \Key{string} & --- & \ApiText{Per-call binding override.}{Bindungsvorgabe für einen einzelnen Aufruf.}}
  \SemioTableRow{\Key{label} & \Key{string} & --- & \ApiText{Column holding the visible caption.}{Spalte mit der sichtbaren Beschriftung.}}
  \SemioTableRow{\Key{layout} & \Key{string} & --- & \ApiText{Layout algorithm that places the nodes.}{Layoutalgorithmus, der die Knoten platziert.}}
  \SemioTableRow{\Key{seed} & \Key{number} & --- & \ApiText{Seed of the deterministic pseudo-random layout.}{Startwert des deterministischen Pseudozufallslayouts.}}
  \SemioTableRow{\Key{order} & \Key{string} & --- & \ApiText{Order the nodes are stacked in.}{Reihenfolge, in der die Knoten gestapelt werden.}}
  \SemioTableRow{\Key{columns} & \Key{integer} & --- & \ApiText{Number of columns of the layout grid.}{Anzahl der Spalten des Anordnungsgitters.}}
  \SemioTableRow{\Key{root} & \Key{integer} & --- & \ApiText{Node the layout starts from.}{Knoten, bei dem das Layout beginnt.}}
  \SemioTableRow{\Key{sweeps} & \Key{integer} & --- & \ApiText{Ordering sweeps of the crossing reduction.}{Sortierdurchläufe der Kreuzungsreduktion.}}
  \SemioTableRow{\Key{orientation} & \Key{string} & --- & \ApiText{Primary direction of the layout.}{Hauptrichtung der Anordnung.}}
  \SemioTableRow{\Key{shape} & \Key{string} & --- & \ApiText{Shape of the node marks.}{Form der Knotenmarken.}}
  \SemioTableRow{\Key{size} & \Key{number} & --- & \ApiText{Millimetre size of a node mark.}{Millimetergröße einer Knotenmarke.}}
  \SemioTableRow{\Key{sizeBy} & \Key{string} & --- & \ApiText{What the node size encodes.}{Was die Knotengröße codiert.}}
  \SemioTableRow{\Key{directed} & \Key{boolean} & --- & \ApiText{Draw arrow heads, so the edges read as directed.}{Pfeilspitzen zeichnen, sodass die Kanten gerichtet wirken.}}
  \SemioTableRow{\Key{weightWidth} & \Key{boolean} & --- & \ApiText{Derive the link width from the weight.}{Die Kantenbreite aus dem Gewicht ableiten.}}
  \SemioTableRow{\Key{signed} & \Key{boolean} & --- & \ApiText{Colour negative weights in the danger hue.}{Negative Gewichte im Warnfarbton einfärben.}}
  \SemioTableRow{\Key{colorBy} & \Key{string} & --- & \ApiText{Ribbon hue.}{Farbton der Bänder.}}
  \SemioTableRow{\Key{labels} & \Key{boolean} & --- & \ApiText{Node captions.}{Beschriftungen der Knoten.}}
  \SemioTableRow{\Key{linkWidth} & \Key{number} & --- & \ApiText{Millimetre width of a link.}{Millimeterbreite einer Kante.}}
  \SemioTableRow{\Key{linkDistance} & \Key{number} & --- & \ApiText{Rest length of a link in the force layout.}{Ruhelänge einer Kante im Kräftelayout.}}
  \SemioTableRow{\Key{charge} & \Key{number} & --- & \ApiText{Many-body charge of the force layout; negative values repel.}{Ladung der Vielteilchenkraft im Kräftelayout; negative Werte stoßen ab.}}
  \SemioTableRow{\Key{forces} & \Key{string} & --- & \ApiText{Forces the deterministic layout applies.}{Kräfte, die das deterministische Layout anwendet.}}
  \SemioTableRow{\Key{padding} & \Key{string} & --- & \ApiText{Padding between adjacent elements.}{Abstand zwischen benachbarten Elementen.}}
  \SemioTableRow{\Key{nodeWidth} & \Key{number} & --- & \ApiText{Width of a node rectangle.}{Breite eines Knotenrechtecks.}}
  \SemioTableRow{\Key{nodePadding} & \Key{number} & --- & \ApiText{Vertical gap between the nodes of one column.}{Senkrechter Abstand zwischen den Knoten einer Spalte.}}
  \SemioTableRow{\Key{nodeAlign} & \Key{string} & --- & \ApiText{How nodes are aligned inside their column.}{Wie die Knoten innerhalb ihrer Spalte ausgerichtet werden.}}
  \SemioTableRow{\Key{nodeSort} & \Key{string} & --- & \ApiText{Auto keeps d3's "sort each column by y0".}{'auto' behält d3s Regel „jede Spalte nach y0 sortieren“ bei.}}
  \SemioTableRow{\Key{linkSort} & \Key{string} & --- & \ApiText{Order the links leave and enter a node in.}{Reihenfolge, in der die Kanten einen Knoten verlassen und erreichen.}}
  \SemioTableRow{\Key{labelWidth} & \Key{number} & --- & \ApiText{Room reserved left and right for the captions.}{Platz, der links und rechts für die Beschriftungen freigehalten wird.}}
}

\subsection{\Key{arc-diagram}}

\SemioTableLong[text-size=8pt]{\Key{arc-diagram}}{0.20,0.13,0.15,0.52}{\ApiKey & \ApiType & \ApiDefault & \ApiMeaning}{%
  \SemioTableRow{\Key{nodes} & \Key{string} & --- & \ApiText{Name of the node table.}{Name der Knotentabelle.}}
  \SemioTableRow{\Key{id} & \Key{string} & --- & \ApiText{Column holding the row identifier.}{Spalte mit der Zeilenkennung.}}
  \SemioTableRow{\Key{group} & \Key{string} & --- & \ApiText{Column that groups the nodes.}{Spalte, die die Knoten gruppiert.}}
  \SemioTableRow{\Key{source} & \Key{string} & --- & \ApiText{Column holding the source node of an edge.}{Spalte mit dem Quellknoten einer Kante.}}
  \SemioTableRow{\Key{target} & \Key{string} & --- & \ApiText{Column holding the target node of an edge, or the target value.}{Spalte mit dem Zielknoten einer Kante oder mit dem Zielwert.}}
  \SemioTableRow{\Key{weight} & \Key{string} & --- & \ApiText{Column holding the edge weight.}{Spalte mit dem Kantengewicht.}}
  \SemioTableRow{\Key{label} & \Key{string} & --- & \ApiText{Column holding the visible caption.}{Spalte mit der sichtbaren Beschriftung.}}
  \SemioTableRow{\Key{layout} & \Key{string} & --- & \ApiText{Layout algorithm that places the nodes.}{Layoutalgorithmus, der die Knoten platziert.}}
  \SemioTableRow{\Key{iterations} & \Key{integer} & --- & \ApiText{Fixed iteration count of the deterministic layout.}{Feste Iterationszahl des deterministischen Layouts.}}
  \SemioTableRow{\Key{seed} & \Key{number} & --- & \ApiText{Seed of the deterministic pseudo-random layout.}{Startwert des deterministischen Pseudozufallslayouts.}}
  \SemioTableRow{\Key{order} & \Key{string} & \Key{index} & \ApiText{Position along the axis.}{Position entlang der Achse.}}
  \SemioTableRow{\Key{columns} & \Key{integer} & --- & \ApiText{Number of columns of the layout grid.}{Anzahl der Spalten des Anordnungsgitters.}}
  \SemioTableRow{\Key{root} & \Key{integer} & --- & \ApiText{Node the layout starts from.}{Knoten, bei dem das Layout beginnt.}}
  \SemioTableRow{\Key{sweeps} & \Key{integer} & --- & \ApiText{Ordering sweeps of the crossing reduction.}{Sortierdurchläufe der Kreuzungsreduktion.}}
  \SemioTableRow{\Key{orientation} & \Key{string} & --- & \ApiText{Primary direction of the layout.}{Hauptrichtung der Anordnung.}}
  \SemioTableRow{\Key{shape} & \Key{string} & --- & \ApiText{Shape of the node marks.}{Form der Knotenmarken.}}
  \SemioTableRow{\Key{size} & \Key{number} & --- & \ApiText{As in 'graph'.}{Wie in der Familie 'graph'.}}
  \SemioTableRow{\Key{sizeBy} & \Key{string} & --- & \ApiText{What the node size encodes.}{Was die Knotengröße codiert.}}
  \SemioTableRow{\Key{directed} & \Key{boolean} & --- & \ApiText{As in 'graph'.}{Wie in der Familie 'graph'.}}
  \SemioTableRow{\Key{curvature} & \Key{number} & --- & \ApiText{Curvature of a link, as a fraction of its span.}{Krümmung einer Kante als Bruchteil ihrer Spannweite.}}
  \SemioTableRow{\Key{weightWidth} & \Key{boolean} & --- & \ApiText{As in 'graph'.}{Wie in der Familie 'graph'.}}
  \SemioTableRow{\Key{signed} & \Key{boolean} & --- & \ApiText{As in 'graph'.}{Wie in der Familie 'graph'.}}
  \SemioTableRow{\Key{colorBy} & \Key{string} & --- & \ApiText{As in 'graph'.}{Wie in der Familie 'graph'.}}
  \SemioTableRow{\Key{labels} & \Key{boolean} & --- & \ApiText{As in 'graph'.}{Wie in der Familie 'graph'.}}
  \SemioTableRow{\Key{linkWidth} & \Key{number} & --- & \ApiText{As in 'graph'.}{Wie in der Familie 'graph'.}}
  \SemioTableRow{\Key{linkDistance} & \Key{number} & --- & \ApiText{Rest length of a link in the force layout.}{Ruhelänge einer Kante im Kräftelayout.}}
  \SemioTableRow{\Key{charge} & \Key{number} & --- & \ApiText{Many-body charge of the force layout; negative values repel.}{Ladung der Vielteilchenkraft im Kräftelayout; negative Werte stoßen ab.}}
  \SemioTableRow{\Key{forces} & \Key{string} & --- & \ApiText{Forces the deterministic layout applies.}{Kräfte, die das deterministische Layout anwendet.}}
  \SemioTableRow{\Key{padding} & \Key{string} & --- & \ApiText{As in 'graph'.}{Wie in der Familie 'graph'.}}
  \SemioTableRow{\Key{mode} & \Key{string} & \Key{linear} & \ApiText{Straight axis with bows, or a ring with chords.}{Gerade Achse mit Bögen oder ein Ring mit Sehnen.}}
  \SemioTableRow{\Key{side} & \Key{string} & \Key{above} & \ApiText{Which half the bows of a linear diagram occupy.}{Welche Hälfte die Bögen eines linearen Diagramms einnehmen.}}
  \SemioTableRow{\Key{bow} & \Key{number} & \Key{0.5} & \ApiText{Bow height as a fraction of the span.}{Bogenhöhe als Bruchteil der Spannweite.}}
}

\subsection{\Key{biofabric}}

\SemioTableLong[text-size=8pt]{\Key{biofabric}}{0.20,0.13,0.15,0.52}{\ApiKey & \ApiType & \ApiDefault & \ApiMeaning}{%
  \SemioTableRow{\Key{nodes} & \Key{string} & --- & \ApiText{Name of the node table.}{Name der Knotentabelle.}}
  \SemioTableRow{\Key{id} & \Key{string} & --- & \ApiText{Column holding the row identifier.}{Spalte mit der Zeilenkennung.}}
  \SemioTableRow{\Key{group} & \Key{string} & --- & \ApiText{Column that groups the nodes.}{Spalte, die die Knoten gruppiert.}}
  \SemioTableRow{\Key{source} & \Key{string} & --- & \ApiText{Column holding the source node of an edge.}{Spalte mit dem Quellknoten einer Kante.}}
  \SemioTableRow{\Key{target} & \Key{string} & --- & \ApiText{Column holding the target node of an edge, or the target value.}{Spalte mit dem Zielknoten einer Kante oder mit dem Zielwert.}}
  \SemioTableRow{\Key{weight} & \Key{string} & --- & \ApiText{Column holding the edge weight.}{Spalte mit dem Kantengewicht.}}
  \SemioTableRow{\Key{label} & \Key{string} & --- & \ApiText{Column holding the visible caption.}{Spalte mit der sichtbaren Beschriftung.}}
  \SemioTableRow{\Key{layout} & \Key{string} & --- & \ApiText{Layout algorithm that places the nodes.}{Layoutalgorithmus, der die Knoten platziert.}}
  \SemioTableRow{\Key{iterations} & \Key{integer} & --- & \ApiText{Fixed iteration count of the deterministic layout.}{Feste Iterationszahl des deterministischen Layouts.}}
  \SemioTableRow{\Key{seed} & \Key{number} & --- & \ApiText{Seed of the deterministic pseudo-random layout.}{Startwert des deterministischen Pseudozufallslayouts.}}
  \SemioTableRow{\Key{order} & \Key{string} & --- & \ApiText{Row and column ordering.}{Reihenfolge von Zeilen und Spalten.}}
  \SemioTableRow{\Key{columns} & \Key{integer} & --- & \ApiText{Number of columns of the layout grid.}{Anzahl der Spalten des Anordnungsgitters.}}
  \SemioTableRow{\Key{root} & \Key{integer} & --- & \ApiText{Node the layout starts from.}{Knoten, bei dem das Layout beginnt.}}
  \SemioTableRow{\Key{sweeps} & \Key{integer} & --- & \ApiText{Ordering sweeps of the crossing reduction.}{Sortierdurchläufe der Kreuzungsreduktion.}}
  \SemioTableRow{\Key{orientation} & \Key{string} & --- & \ApiText{Primary direction of the layout.}{Hauptrichtung der Anordnung.}}
  \SemioTableRow{\Key{shape} & \Key{string} & --- & \ApiText{Shape of the node marks.}{Form der Knotenmarken.}}
  \SemioTableRow{\Key{size} & \Key{number} & --- & \ApiText{Millimetre size of a node mark.}{Millimetergröße einer Knotenmarke.}}
  \SemioTableRow{\Key{sizeBy} & \Key{string} & --- & \ApiText{What the node size encodes.}{Was die Knotengröße codiert.}}
  \SemioTableRow{\Key{directed} & \Key{boolean} & --- & \ApiText{Draw arrow heads, so the edges read as directed.}{Pfeilspitzen zeichnen, sodass die Kanten gerichtet wirken.}}
  \SemioTableRow{\Key{curvature} & \Key{number} & --- & \ApiText{Curvature of a link, as a fraction of its span.}{Krümmung einer Kante als Bruchteil ihrer Spannweite.}}
  \SemioTableRow{\Key{weightWidth} & \Key{boolean} & --- & \ApiText{Derive the link width from the weight.}{Die Kantenbreite aus dem Gewicht ableiten.}}
  \SemioTableRow{\Key{signed} & \Key{boolean} & --- & \ApiText{Colour negative weights in the danger hue.}{Negative Gewichte im Warnfarbton einfärben.}}
  \SemioTableRow{\Key{colorBy} & \Key{string} & --- & \ApiText{What the node or link colour encodes.}{Was die Knoten- bzw. Kantenfarbe codiert.}}
  \SemioTableRow{\Key{labels} & \Key{boolean} & --- & \ApiText{Row and column captions.}{Zeilen- und Spaltenbeschriftungen.}}
  \SemioTableRow{\Key{linkWidth} & \Key{number} & --- & \ApiText{Millimetre width of a link.}{Millimeterbreite einer Kante.}}
  \SemioTableRow{\Key{linkDistance} & \Key{number} & --- & \ApiText{Rest length of a link in the force layout.}{Ruhelänge einer Kante im Kräftelayout.}}
  \SemioTableRow{\Key{charge} & \Key{number} & --- & \ApiText{Many-body charge of the force layout; negative values repel.}{Ladung der Vielteilchenkraft im Kräftelayout; negative Werte stoßen ab.}}
  \SemioTableRow{\Key{forces} & \Key{string} & --- & \ApiText{Forces the deterministic layout applies.}{Kräfte, die das deterministische Layout anwendet.}}
  \SemioTableRow{\Key{padding} & \Key{string} & --- & \ApiText{Padding between adjacent elements.}{Abstand zwischen benachbarten Elementen.}}
  \SemioTableRow{\Key{symmetric} & \Key{boolean} & --- & \ApiText{Mirror every edge so an undirected graph fills both triangles.}{Jede Kante spiegeln, damit ein ungerichteter Graph beide Dreiecke füllt.}}
  \SemioTableRow{\Key{cellGap} & \Key{number} & --- & \ApiText{Gap between cells.}{Abstand zwischen den Zellen.}}
  \SemioTableRow{\Key{scaleBy} & \Key{string} & --- & \ApiText{Cell shade from |weight| or from presence alone.}{Zellton aus |weight| oder allein aus dem Vorhandensein.}}
  \SemioTableRow{\Key{diagonal} & \Key{boolean} & --- & \ApiText{Draw the diagonal cells.}{Die Diagonalzellen zeichnen.}}
  \SemioTableRow{\Key{labelWidth} & \Key{number} & --- & \ApiText{Millimetre width reserved for the labels.}{Millimeterbreite, die für die Beschriftungen reserviert ist.}}
  \SemioTableRow{\Key{blocks} & \Key{boolean} & --- & \ApiText{Hairlines between consecutive ordering groups.}{Haarlinien zwischen aufeinanderfolgenden Sortiergruppen.}}
  \SemioTableRow{\Key{split} & \Key{boolean} & --- & \ApiText{Where the drawing splits into its two halves.}{Wo die Zeichnung in ihre zwei Hälften zerfällt.}}
}

\subsection{\Key{chord}}

\SemioTableLong[text-size=8pt]{\Key{chord}}{0.20,0.13,0.15,0.52}{\ApiKey & \ApiType & \ApiDefault & \ApiMeaning}{%
  \SemioTableRow{\Key{nodes} & \Key{string} & --- & \ApiText{Name of the node table.}{Name der Knotentabelle.}}
  \SemioTableRow{\Key{id} & \Key{string} & --- & \ApiText{Column holding the row identifier.}{Spalte mit der Zeilenkennung.}}
  \SemioTableRow{\Key{group} & \Key{string} & --- & \ApiText{Column that groups the nodes.}{Spalte, die die Knoten gruppiert.}}
  \SemioTableRow{\Key{source} & \Key{string} & --- & \ApiText{Column holding the source node of an edge.}{Spalte mit dem Quellknoten einer Kante.}}
  \SemioTableRow{\Key{target} & \Key{string} & --- & \ApiText{Column holding the target node of an edge, or the target value.}{Spalte mit dem Zielknoten einer Kante oder mit dem Zielwert.}}
  \SemioTableRow{\Key{weight} & \Key{string} & --- & \ApiText{Column holding the edge weight.}{Spalte mit dem Kantengewicht.}}
  \SemioTableRow{\Key{label} & \Key{string} & --- & \ApiText{Column holding the visible caption.}{Spalte mit der sichtbaren Beschriftung.}}
  \SemioTableRow{\Key{layout} & \Key{string} & --- & \ApiText{Layout algorithm that places the nodes.}{Layoutalgorithmus, der die Knoten platziert.}}
  \SemioTableRow{\Key{iterations} & \Key{integer} & --- & \ApiText{Fixed iteration count of the deterministic layout.}{Feste Iterationszahl des deterministischen Layouts.}}
  \SemioTableRow{\Key{seed} & \Key{number} & --- & \ApiText{Seed of the deterministic pseudo-random layout.}{Startwert des deterministischen Pseudozufallslayouts.}}
  \SemioTableRow{\Key{order} & \Key{string} & --- & \ApiText{Order the nodes are stacked in.}{Reihenfolge, in der die Knoten gestapelt werden.}}
  \SemioTableRow{\Key{columns} & \Key{integer} & --- & \ApiText{Number of columns of the layout grid.}{Anzahl der Spalten des Anordnungsgitters.}}
  \SemioTableRow{\Key{root} & \Key{integer} & --- & \ApiText{Node the layout starts from.}{Knoten, bei dem das Layout beginnt.}}
  \SemioTableRow{\Key{sweeps} & \Key{integer} & --- & \ApiText{Ordering sweeps of the crossing reduction.}{Sortierdurchläufe der Kreuzungsreduktion.}}
  \SemioTableRow{\Key{orientation} & \Key{string} & --- & \ApiText{Primary direction of the layout.}{Hauptrichtung der Anordnung.}}
  \SemioTableRow{\Key{shape} & \Key{string} & --- & \ApiText{Shape of the node marks.}{Form der Knotenmarken.}}
  \SemioTableRow{\Key{size} & \Key{number} & --- & \ApiText{Millimetre size of a node mark.}{Millimetergröße einer Knotenmarke.}}
  \SemioTableRow{\Key{sizeBy} & \Key{string} & --- & \ApiText{What the node size encodes.}{Was die Knotengröße codiert.}}
  \SemioTableRow{\Key{directed} & \Key{boolean} & --- & \ApiText{Draw arrow heads, so the edges read as directed.}{Pfeilspitzen zeichnen, sodass die Kanten gerichtet wirken.}}
  \SemioTableRow{\Key{curvature} & \Key{number} & --- & \ApiText{Curvature of a link, as a fraction of its span.}{Krümmung einer Kante als Bruchteil ihrer Spannweite.}}
  \SemioTableRow{\Key{weightWidth} & \Key{boolean} & --- & \ApiText{Derive the link width from the weight.}{Die Kantenbreite aus dem Gewicht ableiten.}}
  \SemioTableRow{\Key{signed} & \Key{boolean} & --- & \ApiText{Colour negative weights in the danger hue.}{Negative Gewichte im Warnfarbton einfärben.}}
  \SemioTableRow{\Key{colorBy} & \Key{string} & --- & \ApiText{What the node or link colour encodes.}{Was die Knoten- bzw. Kantenfarbe codiert.}}
  \SemioTableRow{\Key{labels} & \Key{boolean} & \Key{True} & \ApiText{Group captions outside the band.}{Gruppenbeschriftungen außerhalb des Bands.}}
  \SemioTableRow{\Key{linkWidth} & \Key{number} & --- & \ApiText{Millimetre width of a link.}{Millimeterbreite einer Kante.}}
  \SemioTableRow{\Key{linkDistance} & \Key{number} & --- & \ApiText{Rest length of a link in the force layout.}{Ruhelänge einer Kante im Kräftelayout.}}
  \SemioTableRow{\Key{charge} & \Key{number} & --- & \ApiText{Many-body charge of the force layout; negative values repel.}{Ladung der Vielteilchenkraft im Kräftelayout; negative Werte stoßen ab.}}
  \SemioTableRow{\Key{forces} & \Key{string} & --- & \ApiText{Forces the deterministic layout applies.}{Kräfte, die das deterministische Layout anwendet.}}
  \SemioTableRow{\Key{padding} & \Key{string} & \Key{3} & \ApiText{Padding between adjacent elements.}{Abstand zwischen benachbarten Elementen.}}
  \SemioTableRow{\Key{padAngle} & \Key{number} & \Key{0.04} & \ApiText{Gap between the group arcs.}{Abstand zwischen den Gruppenbögen.}}
  \SemioTableRow{\Key{sortGroups} & \Key{string} & \Key{none} & \ApiText{Order the groups appear in around the ring.}{Reihenfolge der Gruppen rund um den Ring.}}
  \SemioTableRow{\Key{sortSubgroups} & \Key{string} & \Key{none} & \ApiText{Order the subgroups appear in inside a group.}{Reihenfolge der Untergruppen innerhalb einer Gruppe.}}
  \SemioTableRow{\Key{symmetric} & \Key{boolean} & \Key{True} & \ApiText{Mirror every edge into the transposed cell.}{Jede Kante in die transponierte Zelle spiegeln.}}
  \SemioTableRow{\Key{innerRadius} & \Key{number} & \Key{0.86} & \ApiText{Inner edge of the group band.}{Innenkante des Gruppenbands.}}
  \SemioTableRow{\Key{ribbonInset} & \Key{number} & \Key{0.86} & \ApiText{Radius the ribbons start from.}{Radius, an dem die Bänder beginnen.}}
  \SemioTableRow{\Key{tension} & \Key{number} & \Key{0.55} & \ApiText{How strongly a ribbon is pulled towards the centre.}{Wie stark ein Band zur Mitte gezogen wird.}}
}

\subsection{\Key{commit-graph}}

\SemioTableLong[text-size=8pt]{\Key{commit-graph}}{0.20,0.13,0.15,0.52}{\ApiKey & \ApiType & \ApiDefault & \ApiMeaning}{%
  \SemioTableRow{\Key{nodes} & \Key{string} & --- & \ApiText{Per-call binding overrides, passed to the kernel.}{Bindungsvorgaben für einen einzelnen Aufruf, an den Kern weitergereicht.}}
  \SemioTableRow{\Key{id} & \Key{string} & --- & \ApiText{Per-call binding overrides, passed to the kernel.}{Bindungsvorgaben für einen einzelnen Aufruf, an den Kern weitergereicht.}}
  \SemioTableRow{\Key{group} & \Key{string} & --- & \ApiText{Per-call binding overrides, passed to the kernel.}{Bindungsvorgaben für einen einzelnen Aufruf, an den Kern weitergereicht.}}
  \SemioTableRow{\Key{source} & \Key{string} & --- & \ApiText{Per-call binding overrides, passed to the kernel.}{Bindungsvorgaben für einen einzelnen Aufruf, an den Kern weitergereicht.}}
  \SemioTableRow{\Key{target} & \Key{string} & --- & \ApiText{Per-call binding overrides, passed to the kernel.}{Bindungsvorgaben für einen einzelnen Aufruf, an den Kern weitergereicht.}}
  \SemioTableRow{\Key{weight} & \Key{string} & --- & \ApiText{Per-call binding overrides, passed to the kernel.}{Bindungsvorgaben für einen einzelnen Aufruf, an den Kern weitergereicht.}}
  \SemioTableRow{\Key{label} & \Key{string} & --- & \ApiText{Column holding the visible caption.}{Spalte mit der sichtbaren Beschriftung.}}
  \SemioTableRow{\Key{layout} & \Key{string} & --- & \ApiText{Layout algorithm that places the nodes.}{Layoutalgorithmus, der die Knoten platziert.}}
  \SemioTableRow{\Key{iterations} & \Key{integer} & --- & \ApiText{Fixed iteration count of the deterministic layout.}{Feste Iterationszahl des deterministischen Layouts.}}
  \SemioTableRow{\Key{seed} & \Key{number} & --- & \ApiText{Seed of the deterministic LCG.}{Startwert des deterministischen linearen Kongruenzgenerators.}}
  \SemioTableRow{\Key{order} & \Key{string} & --- & \ApiText{Closed-form layout ordering.}{Geschlossene Anordnungsreihenfolge des Layouts.}}
  \SemioTableRow{\Key{columns} & \Key{integer} & --- & \ApiText{Grid columns.}{Spalten des Gitters.}}
  \SemioTableRow{\Key{root} & \Key{integer} & --- & \ApiText{Radial/BFS root.}{Wurzel der radialen Anordnung bzw. der Breitensuche.}}
  \SemioTableRow{\Key{sweeps} & \Key{integer} & --- & \ApiText{Sugiyama ordering passes.}{Sortierdurchläufe des Sugiyama-Verfahrens.}}
  \SemioTableRow{\Key{orientation} & \Key{string} & --- & \ApiText{Primary direction of the layout.}{Hauptrichtung der Anordnung.}}
  \SemioTableRow{\Key{shape} & \Key{string} & --- & \ApiText{Shape of the node marks.}{Form der Knotenmarken.}}
  \SemioTableRow{\Key{size} & \Key{number} & --- & \ApiText{Node radius.}{Radius eines Knotens.}}
  \SemioTableRow{\Key{sizeBy} & \Key{string} & --- & \ApiText{Node radius encoding.}{Codierung des Knotenradius.}}
  \SemioTableRow{\Key{directed} & \Key{boolean} & --- & \ApiText{Arrow heads.}{Pfeilspitzen an den Kanten.}}
  \SemioTableRow{\Key{curvature} & \Key{number} & --- & \ApiText{Base link bend; parallel edges add to it.}{Grundkrümmung einer Kante; parallele Kanten addieren sich darauf.}}
  \SemioTableRow{\Key{weightWidth} & \Key{boolean} & --- & \ApiText{Link width from |weight|.}{Kantenbreite aus |weight|.}}
  \SemioTableRow{\Key{signed} & \Key{boolean} & --- & \ApiText{Negative weights in the danger hue.}{Negative Gewichte im Warnfarbton.}}
  \SemioTableRow{\Key{colorBy} & \Key{string} & --- & \ApiText{What the node or link colour encodes.}{Was die Knoten- bzw. Kantenfarbe codiert.}}
  \SemioTableRow{\Key{labels} & \Key{boolean} & --- & \ApiText{Node ids next to the glyphs.}{Knotenkennungen neben den Glyphen.}}
  \SemioTableRow{\Key{linkWidth} & \Key{number} & --- & \ApiText{Millimetre width of a link.}{Millimeterbreite einer Kante.}}
  \SemioTableRow{\Key{linkDistance} & \Key{number} & --- & \ApiText{Rest length of a link in the force layout.}{Ruhelänge einer Kante im Kräftelayout.}}
  \SemioTableRow{\Key{charge} & \Key{number} & --- & \ApiText{Many-body charge of the force layout; negative values repel.}{Ladung der Vielteilchenkraft im Kräftelayout; negative Werte stoßen ab.}}
  \SemioTableRow{\Key{forces} & \Key{string} & --- & \ApiText{Forces the deterministic layout applies.}{Kräfte, die das deterministische Layout anwendet.}}
  \SemioTableRow{\Key{padding} & \Key{string} & --- & \ApiText{Same inset on all four sides.}{Gleiche Einrückung auf allen vier Seiten.}}
}

\subsection{\Key{data-structure}}

\SemioTableLong[text-size=8pt]{\Key{data-structure}}{0.20,0.13,0.15,0.52}{\ApiKey & \ApiType & \ApiDefault & \ApiMeaning}{%
  \SemioTableRow{\Key{mode} & \Key{string} & --- & \ApiText{Which member of the family the preset renders.}{Welches Mitglied der Familie die Voreinstellung zeichnet.}}
  \SemioTableRow{\Key{idColumn} & \Key{string} & --- & \ApiText{Column holding the node identifier.}{Spalte mit der Knotenkennung.}}
  \SemioTableRow{\Key{parentColumn} & \Key{string} & --- & \ApiText{Column holding the parent node identifier.}{Spalte mit der Kennung des übergeordneten Knotens.}}
  \SemioTableRow{\Key{valueColumn} & \Key{string} & --- & \ApiText{Column holding the stored value of a node.}{Spalte mit dem gespeicherten Wert eines Knotens.}}
  \SemioTableRow{\Key{nodes} & \Key{string} & --- & \ApiText{Per-call binding overrides, passed to the kernel.}{Bindungsvorgaben für einen einzelnen Aufruf, an den Kern weitergereicht.}}
  \SemioTableRow{\Key{id} & \Key{string} & --- & \ApiText{Per-call binding overrides, passed to the kernel.}{Bindungsvorgaben für einen einzelnen Aufruf, an den Kern weitergereicht.}}
  \SemioTableRow{\Key{group} & \Key{string} & --- & \ApiText{Per-call binding overrides, passed to the kernel.}{Bindungsvorgaben für einen einzelnen Aufruf, an den Kern weitergereicht.}}
  \SemioTableRow{\Key{source} & \Key{string} & --- & \ApiText{Per-call binding overrides, passed to the kernel.}{Bindungsvorgaben für einen einzelnen Aufruf, an den Kern weitergereicht.}}
  \SemioTableRow{\Key{target} & \Key{string} & --- & \ApiText{Per-call binding overrides, passed to the kernel.}{Bindungsvorgaben für einen einzelnen Aufruf, an den Kern weitergereicht.}}
  \SemioTableRow{\Key{weight} & \Key{string} & --- & \ApiText{Per-call binding overrides, passed to the kernel.}{Bindungsvorgaben für einen einzelnen Aufruf, an den Kern weitergereicht.}}
  \SemioTableRow{\Key{label} & \Key{string} & --- & \ApiText{Column holding the visible caption.}{Spalte mit der sichtbaren Beschriftung.}}
  \SemioTableRow{\Key{layout} & \Key{string} & --- & \ApiText{Layout algorithm that places the nodes.}{Layoutalgorithmus, der die Knoten platziert.}}
  \SemioTableRow{\Key{iterations} & \Key{integer} & --- & \ApiText{Fixed iteration count of the deterministic layout.}{Feste Iterationszahl des deterministischen Layouts.}}
  \SemioTableRow{\Key{seed} & \Key{number} & --- & \ApiText{Seed of the deterministic LCG.}{Startwert des deterministischen linearen Kongruenzgenerators.}}
  \SemioTableRow{\Key{order} & \Key{string} & --- & \ApiText{Closed-form layout ordering.}{Geschlossene Anordnungsreihenfolge des Layouts.}}
  \SemioTableRow{\Key{columns} & \Key{integer} & --- & \ApiText{Grid columns.}{Spalten des Gitters.}}
  \SemioTableRow{\Key{root} & \Key{integer} & --- & \ApiText{Radial/BFS root.}{Wurzel der radialen Anordnung bzw. der Breitensuche.}}
  \SemioTableRow{\Key{sweeps} & \Key{integer} & --- & \ApiText{Sugiyama ordering passes.}{Sortierdurchläufe des Sugiyama-Verfahrens.}}
  \SemioTableRow{\Key{orientation} & \Key{string} & --- & \ApiText{Primary direction of the layout.}{Hauptrichtung der Anordnung.}}
  \SemioTableRow{\Key{shape} & \Key{string} & --- & \ApiText{Shape of the node marks.}{Form der Knotenmarken.}}
  \SemioTableRow{\Key{size} & \Key{number} & --- & \ApiText{Node radius.}{Radius eines Knotens.}}
  \SemioTableRow{\Key{sizeBy} & \Key{string} & --- & \ApiText{Node radius encoding.}{Codierung des Knotenradius.}}
  \SemioTableRow{\Key{directed} & \Key{boolean} & --- & \ApiText{Arrow heads.}{Pfeilspitzen an den Kanten.}}
  \SemioTableRow{\Key{curvature} & \Key{number} & --- & \ApiText{Base link bend; parallel edges add to it.}{Grundkrümmung einer Kante; parallele Kanten addieren sich darauf.}}
  \SemioTableRow{\Key{weightWidth} & \Key{boolean} & --- & \ApiText{Link width from |weight|.}{Kantenbreite aus |weight|.}}
  \SemioTableRow{\Key{signed} & \Key{boolean} & --- & \ApiText{Negative weights in the danger hue.}{Negative Gewichte im Warnfarbton.}}
  \SemioTableRow{\Key{colorBy} & \Key{string} & --- & \ApiText{What the node or link colour encodes.}{Was die Knoten- bzw. Kantenfarbe codiert.}}
  \SemioTableRow{\Key{labels} & \Key{boolean} & --- & \ApiText{Node ids next to the glyphs.}{Knotenkennungen neben den Glyphen.}}
  \SemioTableRow{\Key{linkWidth} & \Key{number} & --- & \ApiText{Millimetre width of a link.}{Millimeterbreite einer Kante.}}
  \SemioTableRow{\Key{linkDistance} & \Key{number} & --- & \ApiText{Rest length of a link in the force layout.}{Ruhelänge einer Kante im Kräftelayout.}}
  \SemioTableRow{\Key{charge} & \Key{number} & --- & \ApiText{Many-body charge of the force layout; negative values repel.}{Ladung der Vielteilchenkraft im Kräftelayout; negative Werte stoßen ab.}}
  \SemioTableRow{\Key{forces} & \Key{string} & --- & \ApiText{Forces the deterministic layout applies.}{Kräfte, die das deterministische Layout anwendet.}}
  \SemioTableRow{\Key{padding} & \Key{string} & --- & \ApiText{Same inset on all four sides.}{Gleiche Einrückung auf allen vier Seiten.}}
}

\subsection{\Key{dependency-wheel}}

\SemioTableLong[text-size=8pt]{\Key{dependency-wheel}}{0.20,0.13,0.15,0.52}{\ApiKey & \ApiType & \ApiDefault & \ApiMeaning}{%
  \SemioTableRow{\Key{nodes} & \Key{string} & --- & \ApiText{Name of the node table.}{Name der Knotentabelle.}}
  \SemioTableRow{\Key{id} & \Key{string} & --- & \ApiText{Column holding the row identifier.}{Spalte mit der Zeilenkennung.}}
  \SemioTableRow{\Key{group} & \Key{string} & --- & \ApiText{Column that groups the nodes.}{Spalte, die die Knoten gruppiert.}}
  \SemioTableRow{\Key{source} & \Key{string} & --- & \ApiText{Column holding the source node of an edge.}{Spalte mit dem Quellknoten einer Kante.}}
  \SemioTableRow{\Key{target} & \Key{string} & --- & \ApiText{Column holding the target node of an edge, or the target value.}{Spalte mit dem Zielknoten einer Kante oder mit dem Zielwert.}}
  \SemioTableRow{\Key{weight} & \Key{string} & --- & \ApiText{Column holding the edge weight.}{Spalte mit dem Kantengewicht.}}
  \SemioTableRow{\Key{label} & \Key{string} & --- & \ApiText{Column holding the visible caption.}{Spalte mit der sichtbaren Beschriftung.}}
  \SemioTableRow{\Key{layout} & \Key{string} & --- & \ApiText{Layout algorithm that places the nodes.}{Layoutalgorithmus, der die Knoten platziert.}}
  \SemioTableRow{\Key{iterations} & \Key{integer} & --- & \ApiText{Fixed iteration count of the deterministic layout.}{Feste Iterationszahl des deterministischen Layouts.}}
  \SemioTableRow{\Key{seed} & \Key{number} & --- & \ApiText{Seed of the deterministic pseudo-random layout.}{Startwert des deterministischen Pseudozufallslayouts.}}
  \SemioTableRow{\Key{order} & \Key{string} & --- & \ApiText{Order the nodes are stacked in.}{Reihenfolge, in der die Knoten gestapelt werden.}}
  \SemioTableRow{\Key{columns} & \Key{integer} & --- & \ApiText{Number of columns of the layout grid.}{Anzahl der Spalten des Anordnungsgitters.}}
  \SemioTableRow{\Key{root} & \Key{integer} & --- & \ApiText{Node the layout starts from.}{Knoten, bei dem das Layout beginnt.}}
  \SemioTableRow{\Key{sweeps} & \Key{integer} & --- & \ApiText{Ordering sweeps of the crossing reduction.}{Sortierdurchläufe der Kreuzungsreduktion.}}
  \SemioTableRow{\Key{orientation} & \Key{string} & --- & \ApiText{Primary direction of the layout.}{Hauptrichtung der Anordnung.}}
  \SemioTableRow{\Key{shape} & \Key{string} & --- & \ApiText{Shape of the node marks.}{Form der Knotenmarken.}}
  \SemioTableRow{\Key{size} & \Key{number} & --- & \ApiText{Millimetre size of a node mark.}{Millimetergröße einer Knotenmarke.}}
  \SemioTableRow{\Key{sizeBy} & \Key{string} & --- & \ApiText{What the node size encodes.}{Was die Knotengröße codiert.}}
  \SemioTableRow{\Key{directed} & \Key{boolean} & --- & \ApiText{Draw arrow heads, so the edges read as directed.}{Pfeilspitzen zeichnen, sodass die Kanten gerichtet wirken.}}
  \SemioTableRow{\Key{curvature} & \Key{number} & --- & \ApiText{Curvature of a link, as a fraction of its span.}{Krümmung einer Kante als Bruchteil ihrer Spannweite.}}
  \SemioTableRow{\Key{weightWidth} & \Key{boolean} & --- & \ApiText{Derive the link width from the weight.}{Die Kantenbreite aus dem Gewicht ableiten.}}
  \SemioTableRow{\Key{signed} & \Key{boolean} & --- & \ApiText{Colour negative weights in the danger hue.}{Negative Gewichte im Warnfarbton einfärben.}}
  \SemioTableRow{\Key{colorBy} & \Key{string} & --- & \ApiText{What the node or link colour encodes.}{Was die Knoten- bzw. Kantenfarbe codiert.}}
  \SemioTableRow{\Key{labels} & \Key{boolean} & \Key{True} & \ApiText{Group captions outside the band.}{Gruppenbeschriftungen außerhalb des Bands.}}
  \SemioTableRow{\Key{linkWidth} & \Key{number} & --- & \ApiText{Millimetre width of a link.}{Millimeterbreite einer Kante.}}
  \SemioTableRow{\Key{linkDistance} & \Key{number} & --- & \ApiText{Rest length of a link in the force layout.}{Ruhelänge einer Kante im Kräftelayout.}}
  \SemioTableRow{\Key{charge} & \Key{number} & --- & \ApiText{Many-body charge of the force layout; negative values repel.}{Ladung der Vielteilchenkraft im Kräftelayout; negative Werte stoßen ab.}}
  \SemioTableRow{\Key{forces} & \Key{string} & --- & \ApiText{Forces the deterministic layout applies.}{Kräfte, die das deterministische Layout anwendet.}}
  \SemioTableRow{\Key{padding} & \Key{string} & \Key{3} & \ApiText{Padding between adjacent elements.}{Abstand zwischen benachbarten Elementen.}}
  \SemioTableRow{\Key{padAngle} & \Key{number} & \Key{0.04} & \ApiText{Gap between the group arcs.}{Abstand zwischen den Gruppenbögen.}}
  \SemioTableRow{\Key{sortGroups} & \Key{string} & \Key{none} & \ApiText{Order the groups appear in around the ring.}{Reihenfolge der Gruppen rund um den Ring.}}
  \SemioTableRow{\Key{sortSubgroups} & \Key{string} & \Key{none} & \ApiText{Order the subgroups appear in inside a group.}{Reihenfolge der Untergruppen innerhalb einer Gruppe.}}
  \SemioTableRow{\Key{symmetric} & \Key{boolean} & \Key{True} & \ApiText{Mirror every edge into the transposed cell.}{Jede Kante in die transponierte Zelle spiegeln.}}
  \SemioTableRow{\Key{innerRadius} & \Key{number} & \Key{0.86} & \ApiText{Inner edge of the group band.}{Innenkante des Gruppenbands.}}
  \SemioTableRow{\Key{ribbonInset} & \Key{number} & \Key{0.86} & \ApiText{Radius the ribbons start from.}{Radius, an dem die Bänder beginnen.}}
  \SemioTableRow{\Key{tension} & \Key{number} & \Key{0.55} & \ApiText{How strongly a ribbon is pulled towards the centre.}{Wie stark ein Band zur Mitte gezogen wird.}}
}

\subsection{\Key{edge-bundled}}

\SemioTableLong[text-size=8pt]{\Key{edge-bundled}}{0.20,0.13,0.15,0.52}{\ApiKey & \ApiType & \ApiDefault & \ApiMeaning}{%
  \SemioTableRow{\Key{nodes} & \Key{string} & --- & \ApiText{Name of the node table.}{Name der Knotentabelle.}}
  \SemioTableRow{\Key{id} & \Key{string} & --- & \ApiText{Column holding the row identifier.}{Spalte mit der Zeilenkennung.}}
  \SemioTableRow{\Key{group} & \Key{string} & --- & \ApiText{Column that groups the nodes.}{Spalte, die die Knoten gruppiert.}}
  \SemioTableRow{\Key{source} & \Key{string} & --- & \ApiText{Column holding the source node of an edge.}{Spalte mit dem Quellknoten einer Kante.}}
  \SemioTableRow{\Key{target} & \Key{string} & --- & \ApiText{Column holding the target node of an edge, or the target value.}{Spalte mit dem Zielknoten einer Kante oder mit dem Zielwert.}}
  \SemioTableRow{\Key{weight} & \Key{string} & --- & \ApiText{Column holding the edge weight.}{Spalte mit dem Kantengewicht.}}
  \SemioTableRow{\Key{label} & \Key{string} & --- & \ApiText{Column holding the visible caption.}{Spalte mit der sichtbaren Beschriftung.}}
  \SemioTableRow{\Key{layout} & \Key{string} & --- & \ApiText{Layout algorithm that places the nodes.}{Layoutalgorithmus, der die Knoten platziert.}}
  \SemioTableRow{\Key{iterations} & \Key{integer} & --- & \ApiText{Fixed iteration count of the deterministic layout.}{Feste Iterationszahl des deterministischen Layouts.}}
  \SemioTableRow{\Key{seed} & \Key{number} & --- & \ApiText{Seed of the deterministic pseudo-random layout.}{Startwert des deterministischen Pseudozufallslayouts.}}
  \SemioTableRow{\Key{order} & \Key{string} & \Key{group} & \ApiText{Order of the leaves on the ring.}{Reihenfolge der Blätter auf dem Ring.}}
  \SemioTableRow{\Key{columns} & \Key{integer} & --- & \ApiText{Number of columns of the layout grid.}{Anzahl der Spalten des Anordnungsgitters.}}
  \SemioTableRow{\Key{root} & \Key{integer} & --- & \ApiText{Node the layout starts from.}{Knoten, bei dem das Layout beginnt.}}
  \SemioTableRow{\Key{sweeps} & \Key{integer} & --- & \ApiText{Ordering sweeps of the crossing reduction.}{Sortierdurchläufe der Kreuzungsreduktion.}}
  \SemioTableRow{\Key{orientation} & \Key{string} & --- & \ApiText{Primary direction of the layout.}{Hauptrichtung der Anordnung.}}
  \SemioTableRow{\Key{shape} & \Key{string} & --- & \ApiText{Shape of the node marks.}{Form der Knotenmarken.}}
  \SemioTableRow{\Key{size} & \Key{number} & --- & \ApiText{As in 'graph'.}{Wie in der Familie 'graph'.}}
  \SemioTableRow{\Key{sizeBy} & \Key{string} & --- & \ApiText{What the node size encodes.}{Was die Knotengröße codiert.}}
  \SemioTableRow{\Key{directed} & \Key{boolean} & --- & \ApiText{Draw arrow heads, so the edges read as directed.}{Pfeilspitzen zeichnen, sodass die Kanten gerichtet wirken.}}
  \SemioTableRow{\Key{curvature} & \Key{number} & --- & \ApiText{Curvature of a link, as a fraction of its span.}{Krümmung einer Kante als Bruchteil ihrer Spannweite.}}
  \SemioTableRow{\Key{weightWidth} & \Key{boolean} & --- & \ApiText{As in 'graph'.}{Wie in der Familie 'graph'.}}
  \SemioTableRow{\Key{signed} & \Key{boolean} & --- & \ApiText{Colour negative weights in the danger hue.}{Negative Gewichte im Warnfarbton einfärben.}}
  \SemioTableRow{\Key{colorBy} & \Key{string} & --- & \ApiText{As in 'graph'.}{Wie in der Familie 'graph'.}}
  \SemioTableRow{\Key{labels} & \Key{boolean} & \Key{True} & \ApiText{Leaf captions outside the ring.}{Blattbeschriftungen außerhalb des Rings.}}
  \SemioTableRow{\Key{linkWidth} & \Key{number} & --- & \ApiText{As in 'graph'.}{Wie in der Familie 'graph'.}}
  \SemioTableRow{\Key{linkDistance} & \Key{number} & --- & \ApiText{Rest length of a link in the force layout.}{Ruhelänge einer Kante im Kräftelayout.}}
  \SemioTableRow{\Key{charge} & \Key{number} & --- & \ApiText{Many-body charge of the force layout; negative values repel.}{Ladung der Vielteilchenkraft im Kräftelayout; negative Werte stoßen ab.}}
  \SemioTableRow{\Key{forces} & \Key{string} & --- & \ApiText{Forces the deterministic layout applies.}{Kräfte, die das deterministische Layout anwendet.}}
  \SemioTableRow{\Key{padding} & \Key{string} & --- & \ApiText{As in 'graph'.}{Wie in der Familie 'graph'.}}
  \SemioTableRow{\Key{beta} & \Key{number} & \Key{0.85} & \ApiText{Bundle tension: 1 follows the hierarchy, 0 is the straight chord.}{Bündelspannung: 1 folgt der Hierarchie, 0 ist die gerade Sehne.}}
  \SemioTableRow{\Key{hubRadius} & \Key{number} & \Key{0.45} & \ApiText{Radius of the group hubs.}{Radius der Gruppenknoten.}}
  \SemioTableRow{\Key{samples} & \Key{integer} & \Key{24} & \ApiText{Points sampled on the B-spline through the control polygon.}{Punkte, die auf dem B-Spline durch das Kontrollpolygon abgetastet werden.}}
}

\subsection{\Key{graph}}

\SemioTableLong[text-size=8pt]{\Key{graph}}{0.20,0.13,0.15,0.52}{\ApiKey & \ApiType & \ApiDefault & \ApiMeaning}{%
  \SemioTableRow{\Key{nodes} & \Key{string} & --- & \ApiText{Per-call binding overrides, passed to the kernel.}{Bindungsvorgaben für einen einzelnen Aufruf, an den Kern weitergereicht.}}
  \SemioTableRow{\Key{id} & \Key{string} & --- & \ApiText{Per-call binding overrides, passed to the kernel.}{Bindungsvorgaben für einen einzelnen Aufruf, an den Kern weitergereicht.}}
  \SemioTableRow{\Key{group} & \Key{string} & --- & \ApiText{Per-call binding overrides, passed to the kernel.}{Bindungsvorgaben für einen einzelnen Aufruf, an den Kern weitergereicht.}}
  \SemioTableRow{\Key{source} & \Key{string} & --- & \ApiText{Per-call binding overrides, passed to the kernel.}{Bindungsvorgaben für einen einzelnen Aufruf, an den Kern weitergereicht.}}
  \SemioTableRow{\Key{target} & \Key{string} & --- & \ApiText{Per-call binding overrides, passed to the kernel.}{Bindungsvorgaben für einen einzelnen Aufruf, an den Kern weitergereicht.}}
  \SemioTableRow{\Key{weight} & \Key{string} & --- & \ApiText{Per-call binding overrides, passed to the kernel.}{Bindungsvorgaben für einen einzelnen Aufruf, an den Kern weitergereicht.}}
  \SemioTableRow{\Key{label} & \Key{string} & --- & \ApiText{Column holding the visible caption.}{Spalte mit der sichtbaren Beschriftung.}}
  \SemioTableRow{\Key{layout} & \Key{string} & \Key{force} & \ApiText{Layout algorithm that places the nodes.}{Layoutalgorithmus, der die Knoten platziert.}}
  \SemioTableRow{\Key{iterations} & \Key{integer} & \Key{60} & \ApiText{Fixed iteration count of the deterministic layout.}{Feste Iterationszahl des deterministischen Layouts.}}
  \SemioTableRow{\Key{seed} & \Key{number} & \Key{1} & \ApiText{Seed of the deterministic LCG.}{Startwert des deterministischen linearen Kongruenzgenerators.}}
  \SemioTableRow{\Key{order} & \Key{string} & \Key{index} & \ApiText{Closed-form layout ordering.}{Geschlossene Anordnungsreihenfolge des Layouts.}}
  \SemioTableRow{\Key{columns} & \Key{integer} & --- & \ApiText{Grid columns.}{Spalten des Gitters.}}
  \SemioTableRow{\Key{root} & \Key{integer} & \Key{1} & \ApiText{Radial/BFS root.}{Wurzel der radialen Anordnung bzw. der Breitensuche.}}
  \SemioTableRow{\Key{sweeps} & \Key{integer} & \Key{4} & \ApiText{Sugiyama ordering passes.}{Sortierdurchläufe des Sugiyama-Verfahrens.}}
  \SemioTableRow{\Key{orientation} & \Key{string} & \Key{horizontal} & \ApiText{Primary direction of the layout.}{Hauptrichtung der Anordnung.}}
  \SemioTableRow{\Key{shape} & \Key{string} & \Key{circle} & \ApiText{Shape of the node marks.}{Form der Knotenmarken.}}
  \SemioTableRow{\Key{size} & \Key{number} & \Key{1.5} & \ApiText{Node radius.}{Radius eines Knotens.}}
  \SemioTableRow{\Key{sizeBy} & \Key{string} & \Key{none} & \ApiText{Node radius encoding.}{Codierung des Knotenradius.}}
  \SemioTableRow{\Key{directed} & \Key{boolean} & \Key{False} & \ApiText{Arrow heads.}{Pfeilspitzen an den Kanten.}}
  \SemioTableRow{\Key{curvature} & \Key{number} & \Key{0} & \ApiText{Base link bend; parallel edges add to it.}{Grundkrümmung einer Kante; parallele Kanten addieren sich darauf.}}
  \SemioTableRow{\Key{weightWidth} & \Key{boolean} & \Key{False} & \ApiText{Link width from |weight|.}{Kantenbreite aus |weight|.}}
  \SemioTableRow{\Key{signed} & \Key{boolean} & \Key{False} & \ApiText{Negative weights in the danger hue.}{Negative Gewichte im Warnfarbton.}}
  \SemioTableRow{\Key{colorBy} & \Key{string} & \Key{group} & \ApiText{What the node or link colour encodes.}{Was die Knoten- bzw. Kantenfarbe codiert.}}
  \SemioTableRow{\Key{labels} & \Key{boolean} & \Key{False} & \ApiText{Node ids next to the glyphs.}{Knotenkennungen neben den Glyphen.}}
  \SemioTableRow{\Key{linkWidth} & \Key{number} & \Key{0.16} & \ApiText{Millimetre width of a link.}{Millimeterbreite einer Kante.}}
  \SemioTableRow{\Key{linkDistance} & \Key{number} & --- & \ApiText{Rest length of a link in the force layout.}{Ruhelänge einer Kante im Kräftelayout.}}
  \SemioTableRow{\Key{charge} & \Key{number} & --- & \ApiText{Many-body charge of the force layout; negative values repel.}{Ladung der Vielteilchenkraft im Kräftelayout; negative Werte stoßen ab.}}
  \SemioTableRow{\Key{forces} & \Key{string} & --- & \ApiText{Forces the deterministic layout applies.}{Kräfte, die das deterministische Layout anwendet.}}
  \SemioTableRow{\Key{padding} & \Key{string} & \Key{3} & \ApiText{Same inset on all four sides.}{Gleiche Einrückung auf allen vier Seiten.}}
}

\subsection{\Key{hive-plot}}

\SemioTableLong[text-size=8pt]{\Key{hive-plot}}{0.20,0.13,0.15,0.52}{\ApiKey & \ApiType & \ApiDefault & \ApiMeaning}{%
  \SemioTableRow{\Key{innerRadius} & \Key{number} & --- & \ApiText{Millimetre radius the rings or axes start at.}{Millimeterradius, bei dem die Ringe bzw. Achsen beginnen.}}
  \SemioTableRow{\Key{outerRadius} & \Key{number} & --- & \ApiText{Millimetre radius the axes end at.}{Millimeterradius, an dem die Achsen enden.}}
  \SemioTableRow{\Key{nodes} & \Key{string} & --- & \ApiText{Name of the node table.}{Name der Knotentabelle.}}
  \SemioTableRow{\Key{id} & \Key{string} & --- & \ApiText{Column holding the row identifier.}{Spalte mit der Zeilenkennung.}}
  \SemioTableRow{\Key{group} & \Key{string} & --- & \ApiText{Column that groups the nodes.}{Spalte, die die Knoten gruppiert.}}
  \SemioTableRow{\Key{source} & \Key{string} & --- & \ApiText{Column holding the source node of an edge.}{Spalte mit dem Quellknoten einer Kante.}}
  \SemioTableRow{\Key{target} & \Key{string} & --- & \ApiText{Column holding the target node of an edge, or the target value.}{Spalte mit dem Zielknoten einer Kante oder mit dem Zielwert.}}
  \SemioTableRow{\Key{weight} & \Key{string} & --- & \ApiText{Column holding the edge weight.}{Spalte mit dem Kantengewicht.}}
  \SemioTableRow{\Key{label} & \Key{string} & --- & \ApiText{Column holding the visible caption.}{Spalte mit der sichtbaren Beschriftung.}}
  \SemioTableRow{\Key{layout} & \Key{string} & --- & \ApiText{Layout algorithm that places the nodes.}{Layoutalgorithmus, der die Knoten platziert.}}
  \SemioTableRow{\Key{iterations} & \Key{integer} & --- & \ApiText{Fixed iteration count of the deterministic layout.}{Feste Iterationszahl des deterministischen Layouts.}}
  \SemioTableRow{\Key{seed} & \Key{number} & --- & \ApiText{Seed of the deterministic pseudo-random layout.}{Startwert des deterministischen Pseudozufallslayouts.}}
  \SemioTableRow{\Key{order} & \Key{string} & --- & \ApiText{Position along the axis.}{Position entlang der Achse.}}
  \SemioTableRow{\Key{columns} & \Key{integer} & --- & \ApiText{Number of columns of the layout grid.}{Anzahl der Spalten des Anordnungsgitters.}}
  \SemioTableRow{\Key{root} & \Key{integer} & --- & \ApiText{Node the layout starts from.}{Knoten, bei dem das Layout beginnt.}}
  \SemioTableRow{\Key{sweeps} & \Key{integer} & --- & \ApiText{Ordering sweeps of the crossing reduction.}{Sortierdurchläufe der Kreuzungsreduktion.}}
  \SemioTableRow{\Key{orientation} & \Key{string} & --- & \ApiText{Primary direction of the layout.}{Hauptrichtung der Anordnung.}}
  \SemioTableRow{\Key{shape} & \Key{string} & --- & \ApiText{Shape of the node marks.}{Form der Knotenmarken.}}
  \SemioTableRow{\Key{size} & \Key{number} & --- & \ApiText{As in 'graph'.}{Wie in der Familie 'graph'.}}
  \SemioTableRow{\Key{sizeBy} & \Key{string} & --- & \ApiText{What the node size encodes.}{Was die Knotengröße codiert.}}
  \SemioTableRow{\Key{directed} & \Key{boolean} & --- & \ApiText{As in 'graph'.}{Wie in der Familie 'graph'.}}
  \SemioTableRow{\Key{curvature} & \Key{number} & --- & \ApiText{Curvature of a link, as a fraction of its span.}{Krümmung einer Kante als Bruchteil ihrer Spannweite.}}
  \SemioTableRow{\Key{weightWidth} & \Key{boolean} & --- & \ApiText{As in 'graph'.}{Wie in der Familie 'graph'.}}
  \SemioTableRow{\Key{signed} & \Key{boolean} & --- & \ApiText{As in 'graph'.}{Wie in der Familie 'graph'.}}
  \SemioTableRow{\Key{colorBy} & \Key{string} & --- & \ApiText{As in 'graph'.}{Wie in der Familie 'graph'.}}
  \SemioTableRow{\Key{labels} & \Key{boolean} & --- & \ApiText{As in 'graph'.}{Wie in der Familie 'graph'.}}
  \SemioTableRow{\Key{linkWidth} & \Key{number} & --- & \ApiText{As in 'graph'.}{Wie in der Familie 'graph'.}}
  \SemioTableRow{\Key{linkDistance} & \Key{number} & --- & \ApiText{Rest length of a link in the force layout.}{Ruhelänge einer Kante im Kräftelayout.}}
  \SemioTableRow{\Key{charge} & \Key{number} & --- & \ApiText{Many-body charge of the force layout; negative values repel.}{Ladung der Vielteilchenkraft im Kräftelayout; negative Werte stoßen ab.}}
  \SemioTableRow{\Key{forces} & \Key{string} & --- & \ApiText{Forces the deterministic layout applies.}{Kräfte, die das deterministische Layout anwendet.}}
  \SemioTableRow{\Key{padding} & \Key{string} & --- & \ApiText{As in 'graph'.}{Wie in der Familie 'graph'.}}
  \SemioTableRow{\Key{mode} & \Key{string} & --- & \ApiText{Straight axis with bows, or a ring with chords.}{Gerade Achse mit Bögen oder ein Ring mit Sehnen.}}
  \SemioTableRow{\Key{side} & \Key{string} & --- & \ApiText{Which half the bows of a linear diagram occupy.}{Welche Hälfte die Bögen eines linearen Diagramms einnehmen.}}
  \SemioTableRow{\Key{bow} & \Key{number} & --- & \ApiText{Bow height as a fraction of the span.}{Bogenhöhe als Bruchteil der Spannweite.}}
}

\subsection{\Key{neural}}

\ApiText{This family takes no options beyond \Key{variant} and \Key{data}.}{Diese Familie nimmt außer \Key{variant} und \Key{data} keine Optionen.}

\subsection{\Key{neural-network}}

\SemioTableLong[text-size=8pt]{\Key{neural-network}}{0.20,0.13,0.15,0.52}{\ApiKey & \ApiType & \ApiDefault & \ApiMeaning}{%
  \SemioTableRow{\Key{mode} & \Key{string} & --- & \ApiText{Which member of the family the preset renders.}{Welches Mitglied der Familie die Voreinstellung zeichnet.}}
  \SemioTableRow{\Key{render} & \Key{string} & --- & \ApiText{How the cells or units are rendered.}{Wie die Zellen bzw. Einheiten gezeichnet werden.}}
  \SemioTableRow{\Key{connections} & \Key{string} & --- & \ApiText{How consecutive layers are wired.}{Wie aufeinanderfolgende Schichten verdrahtet werden.}}
  \SemioTableRow{\Key{layerColumn} & \Key{string} & --- & \ApiText{Column holding the layer index.}{Spalte mit dem Schichtindex.}}
  \SemioTableRow{\Key{unitsColumn} & \Key{string} & --- & \ApiText{Column holding the unit count of a layer.}{Spalte mit der Einheitenzahl einer Schicht.}}
  \SemioTableRow{\Key{kindColumn} & \Key{string} & --- & \ApiText{Column holding the layer kind.}{Spalte mit der Schichtart.}}
  \SemioTableRow{\Key{maxUnits} & \Key{integer} & --- & \ApiText{Largest number of units drawn per layer.}{Größte Anzahl der je Schicht gezeichneten Einheiten.}}
  \SemioTableRow{\Key{unitSize} & \Key{number} & --- & \ApiText{Millimetre size of one unit glyph.}{Millimetergröße eines Einheitenglyphs.}}
  \SemioTableRow{\Key{nodes} & \Key{string} & --- & \ApiText{Per-call binding overrides, passed to the kernel.}{Bindungsvorgaben für einen einzelnen Aufruf, an den Kern weitergereicht.}}
  \SemioTableRow{\Key{id} & \Key{string} & --- & \ApiText{Per-call binding overrides, passed to the kernel.}{Bindungsvorgaben für einen einzelnen Aufruf, an den Kern weitergereicht.}}
  \SemioTableRow{\Key{group} & \Key{string} & --- & \ApiText{Per-call binding overrides, passed to the kernel.}{Bindungsvorgaben für einen einzelnen Aufruf, an den Kern weitergereicht.}}
  \SemioTableRow{\Key{source} & \Key{string} & --- & \ApiText{Per-call binding overrides, passed to the kernel.}{Bindungsvorgaben für einen einzelnen Aufruf, an den Kern weitergereicht.}}
  \SemioTableRow{\Key{target} & \Key{string} & --- & \ApiText{Per-call binding overrides, passed to the kernel.}{Bindungsvorgaben für einen einzelnen Aufruf, an den Kern weitergereicht.}}
  \SemioTableRow{\Key{weight} & \Key{string} & --- & \ApiText{Per-call binding overrides, passed to the kernel.}{Bindungsvorgaben für einen einzelnen Aufruf, an den Kern weitergereicht.}}
  \SemioTableRow{\Key{label} & \Key{string} & --- & \ApiText{Column holding the visible caption.}{Spalte mit der sichtbaren Beschriftung.}}
  \SemioTableRow{\Key{layout} & \Key{string} & --- & \ApiText{Layout algorithm that places the nodes.}{Layoutalgorithmus, der die Knoten platziert.}}
  \SemioTableRow{\Key{iterations} & \Key{integer} & --- & \ApiText{Fixed iteration count of the deterministic layout.}{Feste Iterationszahl des deterministischen Layouts.}}
  \SemioTableRow{\Key{seed} & \Key{number} & --- & \ApiText{Seed of the deterministic LCG.}{Startwert des deterministischen linearen Kongruenzgenerators.}}
  \SemioTableRow{\Key{order} & \Key{string} & --- & \ApiText{Closed-form layout ordering.}{Geschlossene Anordnungsreihenfolge des Layouts.}}
  \SemioTableRow{\Key{columns} & \Key{integer} & --- & \ApiText{Grid columns.}{Spalten des Gitters.}}
  \SemioTableRow{\Key{root} & \Key{integer} & --- & \ApiText{Radial/BFS root.}{Wurzel der radialen Anordnung bzw. der Breitensuche.}}
  \SemioTableRow{\Key{sweeps} & \Key{integer} & --- & \ApiText{Sugiyama ordering passes.}{Sortierdurchläufe des Sugiyama-Verfahrens.}}
  \SemioTableRow{\Key{orientation} & \Key{string} & --- & \ApiText{Primary direction of the layout.}{Hauptrichtung der Anordnung.}}
  \SemioTableRow{\Key{shape} & \Key{string} & --- & \ApiText{Shape of the node marks.}{Form der Knotenmarken.}}
  \SemioTableRow{\Key{size} & \Key{number} & --- & \ApiText{Node radius.}{Radius eines Knotens.}}
  \SemioTableRow{\Key{sizeBy} & \Key{string} & --- & \ApiText{Node radius encoding.}{Codierung des Knotenradius.}}
  \SemioTableRow{\Key{directed} & \Key{boolean} & --- & \ApiText{Arrow heads.}{Pfeilspitzen an den Kanten.}}
  \SemioTableRow{\Key{curvature} & \Key{number} & --- & \ApiText{Base link bend; parallel edges add to it.}{Grundkrümmung einer Kante; parallele Kanten addieren sich darauf.}}
  \SemioTableRow{\Key{weightWidth} & \Key{boolean} & --- & \ApiText{Link width from |weight|.}{Kantenbreite aus |weight|.}}
  \SemioTableRow{\Key{signed} & \Key{boolean} & --- & \ApiText{Negative weights in the danger hue.}{Negative Gewichte im Warnfarbton.}}
  \SemioTableRow{\Key{colorBy} & \Key{string} & --- & \ApiText{What the node or link colour encodes.}{Was die Knoten- bzw. Kantenfarbe codiert.}}
  \SemioTableRow{\Key{labels} & \Key{boolean} & --- & \ApiText{Node ids next to the glyphs.}{Knotenkennungen neben den Glyphen.}}
  \SemioTableRow{\Key{linkWidth} & \Key{number} & --- & \ApiText{Millimetre width of a link.}{Millimeterbreite einer Kante.}}
  \SemioTableRow{\Key{linkDistance} & \Key{number} & --- & \ApiText{Rest length of a link in the force layout.}{Ruhelänge einer Kante im Kräftelayout.}}
  \SemioTableRow{\Key{charge} & \Key{number} & --- & \ApiText{Many-body charge of the force layout; negative values repel.}{Ladung der Vielteilchenkraft im Kräftelayout; negative Werte stoßen ab.}}
  \SemioTableRow{\Key{forces} & \Key{string} & --- & \ApiText{Forces the deterministic layout applies.}{Kräfte, die das deterministische Layout anwendet.}}
  \SemioTableRow{\Key{padding} & \Key{string} & --- & \ApiText{Same inset on all four sides.}{Gleiche Einrückung auf allen vier Seiten.}}
}

\subsection{\Key{node-link-matrix}}

\SemioTableLong[text-size=8pt]{\Key{node-link-matrix}}{0.20,0.13,0.15,0.52}{\ApiKey & \ApiType & \ApiDefault & \ApiMeaning}{%
  \SemioTableRow{\Key{nodes} & \Key{string} & --- & \ApiText{Name of the node table.}{Name der Knotentabelle.}}
  \SemioTableRow{\Key{id} & \Key{string} & --- & \ApiText{Column holding the row identifier.}{Spalte mit der Zeilenkennung.}}
  \SemioTableRow{\Key{group} & \Key{string} & --- & \ApiText{Column that groups the nodes.}{Spalte, die die Knoten gruppiert.}}
  \SemioTableRow{\Key{source} & \Key{string} & --- & \ApiText{Column holding the source node of an edge.}{Spalte mit dem Quellknoten einer Kante.}}
  \SemioTableRow{\Key{target} & \Key{string} & --- & \ApiText{Column holding the target node of an edge, or the target value.}{Spalte mit dem Zielknoten einer Kante oder mit dem Zielwert.}}
  \SemioTableRow{\Key{weight} & \Key{string} & --- & \ApiText{Column holding the edge weight.}{Spalte mit dem Kantengewicht.}}
  \SemioTableRow{\Key{label} & \Key{string} & --- & \ApiText{Column holding the visible caption.}{Spalte mit der sichtbaren Beschriftung.}}
  \SemioTableRow{\Key{layout} & \Key{string} & --- & \ApiText{Layout algorithm that places the nodes.}{Layoutalgorithmus, der die Knoten platziert.}}
  \SemioTableRow{\Key{iterations} & \Key{integer} & --- & \ApiText{Fixed iteration count of the deterministic layout.}{Feste Iterationszahl des deterministischen Layouts.}}
  \SemioTableRow{\Key{seed} & \Key{number} & --- & \ApiText{Seed of the deterministic pseudo-random layout.}{Startwert des deterministischen Pseudozufallslayouts.}}
  \SemioTableRow{\Key{order} & \Key{string} & --- & \ApiText{Row and column ordering.}{Reihenfolge von Zeilen und Spalten.}}
  \SemioTableRow{\Key{columns} & \Key{integer} & --- & \ApiText{Number of columns of the layout grid.}{Anzahl der Spalten des Anordnungsgitters.}}
  \SemioTableRow{\Key{root} & \Key{integer} & --- & \ApiText{Node the layout starts from.}{Knoten, bei dem das Layout beginnt.}}
  \SemioTableRow{\Key{sweeps} & \Key{integer} & --- & \ApiText{Ordering sweeps of the crossing reduction.}{Sortierdurchläufe der Kreuzungsreduktion.}}
  \SemioTableRow{\Key{orientation} & \Key{string} & --- & \ApiText{Primary direction of the layout.}{Hauptrichtung der Anordnung.}}
  \SemioTableRow{\Key{shape} & \Key{string} & --- & \ApiText{Shape of the node marks.}{Form der Knotenmarken.}}
  \SemioTableRow{\Key{size} & \Key{number} & --- & \ApiText{Millimetre size of a node mark.}{Millimetergröße einer Knotenmarke.}}
  \SemioTableRow{\Key{sizeBy} & \Key{string} & --- & \ApiText{What the node size encodes.}{Was die Knotengröße codiert.}}
  \SemioTableRow{\Key{directed} & \Key{boolean} & --- & \ApiText{Draw arrow heads, so the edges read as directed.}{Pfeilspitzen zeichnen, sodass die Kanten gerichtet wirken.}}
  \SemioTableRow{\Key{curvature} & \Key{number} & --- & \ApiText{Curvature of a link, as a fraction of its span.}{Krümmung einer Kante als Bruchteil ihrer Spannweite.}}
  \SemioTableRow{\Key{weightWidth} & \Key{boolean} & --- & \ApiText{Derive the link width from the weight.}{Die Kantenbreite aus dem Gewicht ableiten.}}
  \SemioTableRow{\Key{signed} & \Key{boolean} & --- & \ApiText{Colour negative weights in the danger hue.}{Negative Gewichte im Warnfarbton einfärben.}}
  \SemioTableRow{\Key{colorBy} & \Key{string} & --- & \ApiText{What the node or link colour encodes.}{Was die Knoten- bzw. Kantenfarbe codiert.}}
  \SemioTableRow{\Key{labels} & \Key{boolean} & --- & \ApiText{Row and column captions.}{Zeilen- und Spaltenbeschriftungen.}}
  \SemioTableRow{\Key{linkWidth} & \Key{number} & --- & \ApiText{Millimetre width of a link.}{Millimeterbreite einer Kante.}}
  \SemioTableRow{\Key{linkDistance} & \Key{number} & --- & \ApiText{Rest length of a link in the force layout.}{Ruhelänge einer Kante im Kräftelayout.}}
  \SemioTableRow{\Key{charge} & \Key{number} & --- & \ApiText{Many-body charge of the force layout; negative values repel.}{Ladung der Vielteilchenkraft im Kräftelayout; negative Werte stoßen ab.}}
  \SemioTableRow{\Key{forces} & \Key{string} & --- & \ApiText{Forces the deterministic layout applies.}{Kräfte, die das deterministische Layout anwendet.}}
  \SemioTableRow{\Key{padding} & \Key{string} & --- & \ApiText{Padding between adjacent elements.}{Abstand zwischen benachbarten Elementen.}}
  \SemioTableRow{\Key{symmetric} & \Key{boolean} & --- & \ApiText{Mirror every edge so an undirected graph fills both triangles.}{Jede Kante spiegeln, damit ein ungerichteter Graph beide Dreiecke füllt.}}
  \SemioTableRow{\Key{cellGap} & \Key{number} & --- & \ApiText{Gap between cells.}{Abstand zwischen den Zellen.}}
  \SemioTableRow{\Key{scaleBy} & \Key{string} & --- & \ApiText{Cell shade from |weight| or from presence alone.}{Zellton aus |weight| oder allein aus dem Vorhandensein.}}
  \SemioTableRow{\Key{diagonal} & \Key{boolean} & --- & \ApiText{Draw the diagonal cells.}{Die Diagonalzellen zeichnen.}}
  \SemioTableRow{\Key{labelWidth} & \Key{number} & --- & \ApiText{Millimetre width reserved for the labels.}{Millimeterbreite, die für die Beschriftungen reserviert ist.}}
  \SemioTableRow{\Key{blocks} & \Key{boolean} & --- & \ApiText{Hairlines between consecutive ordering groups.}{Haarlinien zwischen aufeinanderfolgenden Sortiergruppen.}}
  \SemioTableRow{\Key{split} & \Key{boolean} & --- & \ApiText{Where the drawing splits into its two halves.}{Wo die Zeichnung in ihre zwei Hälften zerfällt.}}
}

\subsection{\Key{parallel-sets}}

\SemioTableLong[text-size=8pt]{\Key{parallel-sets}}{0.20,0.13,0.15,0.52}{\ApiKey & \ApiType & \ApiDefault & \ApiMeaning}{%
  \SemioTableRow{\Key{iterations} & \Key{integer} & --- & \ApiText{Relaxation passes.}{Entspannungsdurchläufe.}}
  \SemioTableRow{\Key{curvature} & \Key{number} & --- & \ApiText{Horizontal control-point fraction of a link ribbon.}{Waagerechter Kontrollpunktanteil eines Verbindungsbands.}}
  \SemioTableRow{\Key{nodes} & \Key{string} & --- & \ApiText{Name of the node table.}{Name der Knotentabelle.}}
  \SemioTableRow{\Key{id} & \Key{string} & --- & \ApiText{Column holding the row identifier.}{Spalte mit der Zeilenkennung.}}
  \SemioTableRow{\Key{group} & \Key{string} & --- & \ApiText{Column that groups the nodes.}{Spalte, die die Knoten gruppiert.}}
  \SemioTableRow{\Key{source} & \Key{string} & --- & \ApiText{Column holding the source node of an edge.}{Spalte mit dem Quellknoten einer Kante.}}
  \SemioTableRow{\Key{target} & \Key{string} & --- & \ApiText{Column holding the target node of an edge, or the target value.}{Spalte mit dem Zielknoten einer Kante oder mit dem Zielwert.}}
  \SemioTableRow{\Key{weight} & \Key{string} & --- & \ApiText{Per-call binding override.}{Bindungsvorgabe für einen einzelnen Aufruf.}}
  \SemioTableRow{\Key{label} & \Key{string} & --- & \ApiText{Column holding the visible caption.}{Spalte mit der sichtbaren Beschriftung.}}
  \SemioTableRow{\Key{layout} & \Key{string} & --- & \ApiText{Layout algorithm that places the nodes.}{Layoutalgorithmus, der die Knoten platziert.}}
  \SemioTableRow{\Key{seed} & \Key{number} & --- & \ApiText{Seed of the deterministic pseudo-random layout.}{Startwert des deterministischen Pseudozufallslayouts.}}
  \SemioTableRow{\Key{order} & \Key{string} & --- & \ApiText{Order the nodes are stacked in.}{Reihenfolge, in der die Knoten gestapelt werden.}}
  \SemioTableRow{\Key{columns} & \Key{integer} & --- & \ApiText{Number of columns of the layout grid.}{Anzahl der Spalten des Anordnungsgitters.}}
  \SemioTableRow{\Key{root} & \Key{integer} & --- & \ApiText{Node the layout starts from.}{Knoten, bei dem das Layout beginnt.}}
  \SemioTableRow{\Key{sweeps} & \Key{integer} & --- & \ApiText{Ordering sweeps of the crossing reduction.}{Sortierdurchläufe der Kreuzungsreduktion.}}
  \SemioTableRow{\Key{orientation} & \Key{string} & --- & \ApiText{Primary direction of the layout.}{Hauptrichtung der Anordnung.}}
  \SemioTableRow{\Key{shape} & \Key{string} & --- & \ApiText{Shape of the node marks.}{Form der Knotenmarken.}}
  \SemioTableRow{\Key{size} & \Key{number} & --- & \ApiText{Millimetre size of a node mark.}{Millimetergröße einer Knotenmarke.}}
  \SemioTableRow{\Key{sizeBy} & \Key{string} & --- & \ApiText{What the node size encodes.}{Was die Knotengröße codiert.}}
  \SemioTableRow{\Key{directed} & \Key{boolean} & --- & \ApiText{Draw arrow heads, so the edges read as directed.}{Pfeilspitzen zeichnen, sodass die Kanten gerichtet wirken.}}
  \SemioTableRow{\Key{weightWidth} & \Key{boolean} & --- & \ApiText{Derive the link width from the weight.}{Die Kantenbreite aus dem Gewicht ableiten.}}
  \SemioTableRow{\Key{signed} & \Key{boolean} & --- & \ApiText{Colour negative weights in the danger hue.}{Negative Gewichte im Warnfarbton einfärben.}}
  \SemioTableRow{\Key{colorBy} & \Key{string} & --- & \ApiText{Ribbon hue.}{Farbton der Bänder.}}
  \SemioTableRow{\Key{labels} & \Key{boolean} & --- & \ApiText{Node captions.}{Beschriftungen der Knoten.}}
  \SemioTableRow{\Key{linkWidth} & \Key{number} & --- & \ApiText{Millimetre width of a link.}{Millimeterbreite einer Kante.}}
  \SemioTableRow{\Key{linkDistance} & \Key{number} & --- & \ApiText{Rest length of a link in the force layout.}{Ruhelänge einer Kante im Kräftelayout.}}
  \SemioTableRow{\Key{charge} & \Key{number} & --- & \ApiText{Many-body charge of the force layout; negative values repel.}{Ladung der Vielteilchenkraft im Kräftelayout; negative Werte stoßen ab.}}
  \SemioTableRow{\Key{forces} & \Key{string} & --- & \ApiText{Forces the deterministic layout applies.}{Kräfte, die das deterministische Layout anwendet.}}
  \SemioTableRow{\Key{padding} & \Key{string} & --- & \ApiText{Padding between adjacent elements.}{Abstand zwischen benachbarten Elementen.}}
  \SemioTableRow{\Key{nodeWidth} & \Key{number} & --- & \ApiText{Width of a node rectangle.}{Breite eines Knotenrechtecks.}}
  \SemioTableRow{\Key{nodePadding} & \Key{number} & --- & \ApiText{Vertical gap between the nodes of one column.}{Senkrechter Abstand zwischen den Knoten einer Spalte.}}
  \SemioTableRow{\Key{nodeAlign} & \Key{string} & --- & \ApiText{How nodes are aligned inside their column.}{Wie die Knoten innerhalb ihrer Spalte ausgerichtet werden.}}
  \SemioTableRow{\Key{nodeSort} & \Key{string} & --- & \ApiText{Auto keeps d3's "sort each column by y0".}{'auto' behält d3s Regel „jede Spalte nach y0 sortieren“ bei.}}
  \SemioTableRow{\Key{linkSort} & \Key{string} & --- & \ApiText{Order the links leave and enter a node in.}{Reihenfolge, in der die Kanten einen Knoten verlassen und erreichen.}}
  \SemioTableRow{\Key{labelWidth} & \Key{number} & --- & \ApiText{Room reserved left and right for the captions.}{Platz, der links und rechts für die Beschriftungen freigehalten wird.}}
}

\subsection{\Key{sankey}}

\SemioTableLong[text-size=8pt]{\Key{sankey}}{0.20,0.13,0.15,0.52}{\ApiKey & \ApiType & \ApiDefault & \ApiMeaning}{%
  \SemioTableRow{\Key{iterations} & \Key{integer} & \Key{6} & \ApiText{Relaxation passes.}{Entspannungsdurchläufe.}}
  \SemioTableRow{\Key{curvature} & \Key{number} & \Key{0.5} & \ApiText{Horizontal control-point fraction of a link ribbon.}{Waagerechter Kontrollpunktanteil eines Verbindungsbands.}}
  \SemioTableRow{\Key{nodes} & \Key{string} & --- & \ApiText{Name of the node table.}{Name der Knotentabelle.}}
  \SemioTableRow{\Key{id} & \Key{string} & --- & \ApiText{Column holding the row identifier.}{Spalte mit der Zeilenkennung.}}
  \SemioTableRow{\Key{group} & \Key{string} & --- & \ApiText{Column that groups the nodes.}{Spalte, die die Knoten gruppiert.}}
  \SemioTableRow{\Key{source} & \Key{string} & --- & \ApiText{Column holding the source node of an edge.}{Spalte mit dem Quellknoten einer Kante.}}
  \SemioTableRow{\Key{target} & \Key{string} & --- & \ApiText{Column holding the target node of an edge, or the target value.}{Spalte mit dem Zielknoten einer Kante oder mit dem Zielwert.}}
  \SemioTableRow{\Key{weight} & \Key{string} & \Key{value} & \ApiText{Per-call binding override.}{Bindungsvorgabe für einen einzelnen Aufruf.}}
  \SemioTableRow{\Key{label} & \Key{string} & --- & \ApiText{Column holding the visible caption.}{Spalte mit der sichtbaren Beschriftung.}}
  \SemioTableRow{\Key{layout} & \Key{string} & --- & \ApiText{Layout algorithm that places the nodes.}{Layoutalgorithmus, der die Knoten platziert.}}
  \SemioTableRow{\Key{seed} & \Key{number} & --- & \ApiText{Seed of the deterministic pseudo-random layout.}{Startwert des deterministischen Pseudozufallslayouts.}}
  \SemioTableRow{\Key{order} & \Key{string} & --- & \ApiText{Order the nodes are stacked in.}{Reihenfolge, in der die Knoten gestapelt werden.}}
  \SemioTableRow{\Key{columns} & \Key{integer} & --- & \ApiText{Number of columns of the layout grid.}{Anzahl der Spalten des Anordnungsgitters.}}
  \SemioTableRow{\Key{root} & \Key{integer} & --- & \ApiText{Node the layout starts from.}{Knoten, bei dem das Layout beginnt.}}
  \SemioTableRow{\Key{sweeps} & \Key{integer} & --- & \ApiText{Ordering sweeps of the crossing reduction.}{Sortierdurchläufe der Kreuzungsreduktion.}}
  \SemioTableRow{\Key{orientation} & \Key{string} & --- & \ApiText{Primary direction of the layout.}{Hauptrichtung der Anordnung.}}
  \SemioTableRow{\Key{shape} & \Key{string} & --- & \ApiText{Shape of the node marks.}{Form der Knotenmarken.}}
  \SemioTableRow{\Key{size} & \Key{number} & --- & \ApiText{Millimetre size of a node mark.}{Millimetergröße einer Knotenmarke.}}
  \SemioTableRow{\Key{sizeBy} & \Key{string} & --- & \ApiText{What the node size encodes.}{Was die Knotengröße codiert.}}
  \SemioTableRow{\Key{directed} & \Key{boolean} & --- & \ApiText{Draw arrow heads, so the edges read as directed.}{Pfeilspitzen zeichnen, sodass die Kanten gerichtet wirken.}}
  \SemioTableRow{\Key{weightWidth} & \Key{boolean} & --- & \ApiText{Derive the link width from the weight.}{Die Kantenbreite aus dem Gewicht ableiten.}}
  \SemioTableRow{\Key{signed} & \Key{boolean} & --- & \ApiText{Colour negative weights in the danger hue.}{Negative Gewichte im Warnfarbton einfärben.}}
  \SemioTableRow{\Key{colorBy} & \Key{string} & \Key{source} & \ApiText{Ribbon hue.}{Farbton der Bänder.}}
  \SemioTableRow{\Key{labels} & \Key{boolean} & \Key{True} & \ApiText{Node captions.}{Beschriftungen der Knoten.}}
  \SemioTableRow{\Key{linkWidth} & \Key{number} & --- & \ApiText{Millimetre width of a link.}{Millimeterbreite einer Kante.}}
  \SemioTableRow{\Key{linkDistance} & \Key{number} & --- & \ApiText{Rest length of a link in the force layout.}{Ruhelänge einer Kante im Kräftelayout.}}
  \SemioTableRow{\Key{charge} & \Key{number} & --- & \ApiText{Many-body charge of the force layout; negative values repel.}{Ladung der Vielteilchenkraft im Kräftelayout; negative Werte stoßen ab.}}
  \SemioTableRow{\Key{forces} & \Key{string} & --- & \ApiText{Forces the deterministic layout applies.}{Kräfte, die das deterministische Layout anwendet.}}
  \SemioTableRow{\Key{padding} & \Key{string} & \Key{3} & \ApiText{Padding between adjacent elements.}{Abstand zwischen benachbarten Elementen.}}
  \SemioTableRow{\Key{nodeWidth} & \Key{number} & \Key{4} & \ApiText{Width of a node rectangle.}{Breite eines Knotenrechtecks.}}
  \SemioTableRow{\Key{nodePadding} & \Key{number} & \Key{2.4} & \ApiText{Vertical gap between the nodes of one column.}{Senkrechter Abstand zwischen den Knoten einer Spalte.}}
  \SemioTableRow{\Key{nodeAlign} & \Key{string} & \Key{justify} & \ApiText{How nodes are aligned inside their column.}{Wie die Knoten innerhalb ihrer Spalte ausgerichtet werden.}}
  \SemioTableRow{\Key{nodeSort} & \Key{string} & \Key{auto} & \ApiText{Auto keeps d3's "sort each column by y0".}{'auto' behält d3s Regel „jede Spalte nach y0 sortieren“ bei.}}
  \SemioTableRow{\Key{linkSort} & \Key{string} & \Key{auto} & \ApiText{Order the links leave and enter a node in.}{Reihenfolge, in der die Kanten einen Knoten verlassen und erreichen.}}
  \SemioTableRow{\Key{labelWidth} & \Key{number} & \Key{11} & \ApiText{Room reserved left and right for the captions.}{Platz, der links und rechts für die Beschriftungen freigehalten wird.}}
}

\subsection{\Key{schema-graph}}

\SemioTableLong[text-size=8pt]{\Key{schema-graph}}{0.20,0.13,0.15,0.52}{\ApiKey & \ApiType & \ApiDefault & \ApiMeaning}{%
  \SemioTableRow{\Key{mode} & \Key{string} & --- & \ApiText{Which member of the family the preset renders.}{Welches Mitglied der Familie die Voreinstellung zeichnet.}}
  \SemioTableRow{\Key{nodes} & \Key{string} & --- & \ApiText{Per-call binding overrides, passed to the kernel.}{Bindungsvorgaben für einen einzelnen Aufruf, an den Kern weitergereicht.}}
  \SemioTableRow{\Key{id} & \Key{string} & --- & \ApiText{Per-call binding overrides, passed to the kernel.}{Bindungsvorgaben für einen einzelnen Aufruf, an den Kern weitergereicht.}}
  \SemioTableRow{\Key{group} & \Key{string} & --- & \ApiText{Per-call binding overrides, passed to the kernel.}{Bindungsvorgaben für einen einzelnen Aufruf, an den Kern weitergereicht.}}
  \SemioTableRow{\Key{source} & \Key{string} & --- & \ApiText{Per-call binding overrides, passed to the kernel.}{Bindungsvorgaben für einen einzelnen Aufruf, an den Kern weitergereicht.}}
  \SemioTableRow{\Key{target} & \Key{string} & --- & \ApiText{Per-call binding overrides, passed to the kernel.}{Bindungsvorgaben für einen einzelnen Aufruf, an den Kern weitergereicht.}}
  \SemioTableRow{\Key{weight} & \Key{string} & --- & \ApiText{Per-call binding overrides, passed to the kernel.}{Bindungsvorgaben für einen einzelnen Aufruf, an den Kern weitergereicht.}}
  \SemioTableRow{\Key{label} & \Key{string} & --- & \ApiText{Column holding the visible caption.}{Spalte mit der sichtbaren Beschriftung.}}
  \SemioTableRow{\Key{layout} & \Key{string} & --- & \ApiText{Layout algorithm that places the nodes.}{Layoutalgorithmus, der die Knoten platziert.}}
  \SemioTableRow{\Key{iterations} & \Key{integer} & --- & \ApiText{Fixed iteration count of the deterministic layout.}{Feste Iterationszahl des deterministischen Layouts.}}
  \SemioTableRow{\Key{seed} & \Key{number} & --- & \ApiText{Seed of the deterministic LCG.}{Startwert des deterministischen linearen Kongruenzgenerators.}}
  \SemioTableRow{\Key{order} & \Key{string} & --- & \ApiText{Closed-form layout ordering.}{Geschlossene Anordnungsreihenfolge des Layouts.}}
  \SemioTableRow{\Key{columns} & \Key{integer} & --- & \ApiText{Grid columns.}{Spalten des Gitters.}}
  \SemioTableRow{\Key{root} & \Key{integer} & --- & \ApiText{Radial/BFS root.}{Wurzel der radialen Anordnung bzw. der Breitensuche.}}
  \SemioTableRow{\Key{sweeps} & \Key{integer} & --- & \ApiText{Sugiyama ordering passes.}{Sortierdurchläufe des Sugiyama-Verfahrens.}}
  \SemioTableRow{\Key{orientation} & \Key{string} & --- & \ApiText{Primary direction of the layout.}{Hauptrichtung der Anordnung.}}
  \SemioTableRow{\Key{shape} & \Key{string} & --- & \ApiText{Shape of the node marks.}{Form der Knotenmarken.}}
  \SemioTableRow{\Key{size} & \Key{number} & --- & \ApiText{Node radius.}{Radius eines Knotens.}}
  \SemioTableRow{\Key{sizeBy} & \Key{string} & --- & \ApiText{Node radius encoding.}{Codierung des Knotenradius.}}
  \SemioTableRow{\Key{directed} & \Key{boolean} & --- & \ApiText{Arrow heads.}{Pfeilspitzen an den Kanten.}}
  \SemioTableRow{\Key{curvature} & \Key{number} & --- & \ApiText{Base link bend; parallel edges add to it.}{Grundkrümmung einer Kante; parallele Kanten addieren sich darauf.}}
  \SemioTableRow{\Key{weightWidth} & \Key{boolean} & --- & \ApiText{Link width from |weight|.}{Kantenbreite aus |weight|.}}
  \SemioTableRow{\Key{signed} & \Key{boolean} & --- & \ApiText{Negative weights in the danger hue.}{Negative Gewichte im Warnfarbton.}}
  \SemioTableRow{\Key{colorBy} & \Key{string} & --- & \ApiText{What the node or link colour encodes.}{Was die Knoten- bzw. Kantenfarbe codiert.}}
  \SemioTableRow{\Key{labels} & \Key{boolean} & --- & \ApiText{Node ids next to the glyphs.}{Knotenkennungen neben den Glyphen.}}
  \SemioTableRow{\Key{linkWidth} & \Key{number} & --- & \ApiText{Millimetre width of a link.}{Millimeterbreite einer Kante.}}
  \SemioTableRow{\Key{linkDistance} & \Key{number} & --- & \ApiText{Rest length of a link in the force layout.}{Ruhelänge einer Kante im Kräftelayout.}}
  \SemioTableRow{\Key{charge} & \Key{number} & --- & \ApiText{Many-body charge of the force layout; negative values repel.}{Ladung der Vielteilchenkraft im Kräftelayout; negative Werte stoßen ab.}}
  \SemioTableRow{\Key{forces} & \Key{string} & --- & \ApiText{Forces the deterministic layout applies.}{Kräfte, die das deterministische Layout anwendet.}}
  \SemioTableRow{\Key{padding} & \Key{string} & --- & \ApiText{Same inset on all four sides.}{Gleiche Einrückung auf allen vier Seiten.}}
}

\subsection{\Key{state-machine}}

\SemioTableLong[text-size=8pt]{\Key{state-machine}}{0.20,0.13,0.15,0.52}{\ApiKey & \ApiType & \ApiDefault & \ApiMeaning}{%
  \SemioTableRow{\Key{mode} & \Key{string} & --- & \ApiText{Which member of the family the preset renders.}{Welches Mitglied der Familie die Voreinstellung zeichnet.}}
  \SemioTableRow{\Key{role} & \Key{string} & --- & \ApiText{Column holding the role of a state.}{Spalte mit der Rolle eines Zustands.}}
  \SemioTableRow{\Key{edgeLabels} & \Key{boolean} & --- & \ApiText{Print the transition labels on the edges.}{Die Übergangsbeschriftungen an die Kanten schreiben.}}
  \SemioTableRow{\Key{nodes} & \Key{string} & --- & \ApiText{Per-call binding overrides, passed to the kernel.}{Bindungsvorgaben für einen einzelnen Aufruf, an den Kern weitergereicht.}}
  \SemioTableRow{\Key{id} & \Key{string} & --- & \ApiText{Per-call binding overrides, passed to the kernel.}{Bindungsvorgaben für einen einzelnen Aufruf, an den Kern weitergereicht.}}
  \SemioTableRow{\Key{group} & \Key{string} & --- & \ApiText{Per-call binding overrides, passed to the kernel.}{Bindungsvorgaben für einen einzelnen Aufruf, an den Kern weitergereicht.}}
  \SemioTableRow{\Key{source} & \Key{string} & --- & \ApiText{Per-call binding overrides, passed to the kernel.}{Bindungsvorgaben für einen einzelnen Aufruf, an den Kern weitergereicht.}}
  \SemioTableRow{\Key{target} & \Key{string} & --- & \ApiText{Per-call binding overrides, passed to the kernel.}{Bindungsvorgaben für einen einzelnen Aufruf, an den Kern weitergereicht.}}
  \SemioTableRow{\Key{weight} & \Key{string} & --- & \ApiText{Per-call binding overrides, passed to the kernel.}{Bindungsvorgaben für einen einzelnen Aufruf, an den Kern weitergereicht.}}
  \SemioTableRow{\Key{label} & \Key{string} & --- & \ApiText{Column holding the visible caption.}{Spalte mit der sichtbaren Beschriftung.}}
  \SemioTableRow{\Key{layout} & \Key{string} & --- & \ApiText{Layout algorithm that places the nodes.}{Layoutalgorithmus, der die Knoten platziert.}}
  \SemioTableRow{\Key{iterations} & \Key{integer} & --- & \ApiText{Fixed iteration count of the deterministic layout.}{Feste Iterationszahl des deterministischen Layouts.}}
  \SemioTableRow{\Key{seed} & \Key{number} & --- & \ApiText{Seed of the deterministic LCG.}{Startwert des deterministischen linearen Kongruenzgenerators.}}
  \SemioTableRow{\Key{order} & \Key{string} & --- & \ApiText{Closed-form layout ordering.}{Geschlossene Anordnungsreihenfolge des Layouts.}}
  \SemioTableRow{\Key{columns} & \Key{integer} & --- & \ApiText{Grid columns.}{Spalten des Gitters.}}
  \SemioTableRow{\Key{root} & \Key{integer} & --- & \ApiText{Radial/BFS root.}{Wurzel der radialen Anordnung bzw. der Breitensuche.}}
  \SemioTableRow{\Key{sweeps} & \Key{integer} & --- & \ApiText{Sugiyama ordering passes.}{Sortierdurchläufe des Sugiyama-Verfahrens.}}
  \SemioTableRow{\Key{orientation} & \Key{string} & --- & \ApiText{Primary direction of the layout.}{Hauptrichtung der Anordnung.}}
  \SemioTableRow{\Key{shape} & \Key{string} & --- & \ApiText{Shape of the node marks.}{Form der Knotenmarken.}}
  \SemioTableRow{\Key{size} & \Key{number} & --- & \ApiText{Node radius.}{Radius eines Knotens.}}
  \SemioTableRow{\Key{sizeBy} & \Key{string} & --- & \ApiText{Node radius encoding.}{Codierung des Knotenradius.}}
  \SemioTableRow{\Key{directed} & \Key{boolean} & --- & \ApiText{Arrow heads.}{Pfeilspitzen an den Kanten.}}
  \SemioTableRow{\Key{curvature} & \Key{number} & --- & \ApiText{Base link bend; parallel edges add to it.}{Grundkrümmung einer Kante; parallele Kanten addieren sich darauf.}}
  \SemioTableRow{\Key{weightWidth} & \Key{boolean} & --- & \ApiText{Link width from |weight|.}{Kantenbreite aus |weight|.}}
  \SemioTableRow{\Key{signed} & \Key{boolean} & --- & \ApiText{Negative weights in the danger hue.}{Negative Gewichte im Warnfarbton.}}
  \SemioTableRow{\Key{colorBy} & \Key{string} & --- & \ApiText{What the node or link colour encodes.}{Was die Knoten- bzw. Kantenfarbe codiert.}}
  \SemioTableRow{\Key{labels} & \Key{boolean} & --- & \ApiText{Node ids next to the glyphs.}{Knotenkennungen neben den Glyphen.}}
  \SemioTableRow{\Key{linkWidth} & \Key{number} & --- & \ApiText{Millimetre width of a link.}{Millimeterbreite einer Kante.}}
  \SemioTableRow{\Key{linkDistance} & \Key{number} & --- & \ApiText{Rest length of a link in the force layout.}{Ruhelänge einer Kante im Kräftelayout.}}
  \SemioTableRow{\Key{charge} & \Key{number} & --- & \ApiText{Many-body charge of the force layout; negative values repel.}{Ladung der Vielteilchenkraft im Kräftelayout; negative Werte stoßen ab.}}
  \SemioTableRow{\Key{forces} & \Key{string} & --- & \ApiText{Forces the deterministic layout applies.}{Kräfte, die das deterministische Layout anwendet.}}
  \SemioTableRow{\Key{padding} & \Key{string} & --- & \ApiText{Same inset on all four sides.}{Gleiche Einrückung auf allen vier Seiten.}}
}

\subsection{\Key{version-control}}

\ApiText{This family takes no options beyond \Key{variant} and \Key{data}.}{Diese Familie nimmt außer \Key{variant} und \Key{data} keine Optionen.}

\section{\ApiText{Geo and spatial}{Geo und Raum}}

\ApiText{Maps, projections, choropleths, symbol maps, contours and routes, plus the planar spatial kernel: voronoi, delaunay, hull, hexbin, contour and density.}{Karten, Projektionen, Choroplethen, Symbolkarten, Höhenlinien und Routen sowie der ebene Raumkern: Voronoi, Delaunay, Hülle, Hexbin, Höhenlinie und Dichte.}

\subsection{\Key{geo-basemap}}

\SemioTableLong[text-size=8pt]{\Key{geo-basemap}}{0.20,0.13,0.15,0.52}{\ApiKey & \ApiType & \ApiDefault & \ApiMeaning}{%
  \SemioTableRow{\Key{relief} & \Key{boolean} & --- & \ApiText{Draw the shaded relief.}{Das Schummerungsrelief zeichnen.}}
  \SemioTableRow{\Key{hydrology} & \Key{boolean} & --- & \ApiText{Draw rivers and lakes.}{Flüsse und Seen zeichnen.}}
  \SemioTableRow{\Key{grid} & \Key{boolean} & --- & \ApiText{Grid lines drawn behind the marks.}{Gitterlinien hinter den Marken.}}
  \SemioTableRow{\Key{labels} & \Key{boolean} & --- & \ApiText{Draw the captions.}{Die Beschriftungen zeichnen.}}
  \SemioTableRow{\Key{settlements} & \Key{boolean} & --- & \ApiText{Draw the settlement symbols.}{Die Siedlungssymbole zeichnen.}}
  \SemioTableRow{\Key{water} & \Key{string} & --- & \ApiText{Water bodies drawn on the base map.}{Gewässer, die auf der Grundkarte gezeichnet werden.}}
  \SemioTableRow{\Key{places} & \Key{string} & --- & \ApiText{Draw the place labels.}{Die Ortsbeschriftungen zeichnen.}}
  \SemioTableRow{\Key{geometry} & \Key{string} & --- & \ApiText{Name of the geometry table the map draws.}{Name der Geometrietabelle, die die Karte zeichnet.}}
  \SemioTableRow{\Key{height} & \Key{string} & --- & \ApiText{Frame height in millimetres.}{Rahmenhöhe in Millimetern.}}
  \SemioTableRow{\Key{opacity} & \Key{string} & --- & \ApiText{Draw opacity of the layer.}{Zeichendeckkraft der Ebene.}}
  \SemioTableRow{\Key{palette} & \Key{string} & --- & \ApiText{Named colour palette resolved by semio-viz-theme.}{Benannte Farbpalette, die semio-viz-theme auflöst.}}
  \SemioTableRow{\Key{steps} & \Key{string} & --- & \ApiText{Number of integration or iteration steps.}{Anzahl der Integrations- bzw. Iterationsschritte.}}
  \SemioTableRow{\Key{width} & \Key{string} & --- & \ApiText{Frame width in millimetres.}{Rahmenbreite in Millimetern.}}
}

\subsection{\Key{geo-cartogram}}

\SemioTableLong[text-size=8pt]{\Key{geo-cartogram}}{0.20,0.13,0.15,0.52}{\ApiKey & \ApiType & \ApiDefault & \ApiMeaning}{%
  \SemioTableRow{\Key{kind} & \Key{string} & --- & \ApiText{Which member of the family is drawn.}{Welches Mitglied der Familie gezeichnet wird.}}
  \SemioTableRow{\Key{maxSize} & \Key{number} & --- & \ApiText{Millimetre size of the largest symbol.}{Millimetergröße des größten Symbols.}}
  \SemioTableRow{\Key{iterations} & \Key{integer} & --- & \ApiText{Fixed iteration count of the deterministic layout.}{Feste Iterationszahl des deterministischen Layouts.}}
  \SemioTableRow{\Key{geometry} & \Key{string} & --- & \ApiText{Name of the geometry table the map draws.}{Name der Geometrietabelle, die die Karte zeichnet.}}
  \SemioTableRow{\Key{palette} & \Key{string} & --- & \ApiText{Named colour palette resolved by semio-viz-theme.}{Benannte Farbpalette, die semio-viz-theme auflöst.}}
}

\subsection{\Key{geo-choropleth}}

\SemioTableLong[text-size=8pt]{\Key{geo-choropleth}}{0.20,0.13,0.15,0.52}{\ApiKey & \ApiType & \ApiDefault & \ApiMeaning}{%
  \SemioTableRow{\Key{scale} & \Key{string} & --- & \ApiText{Named scale the values pass through.}{Benannte Skala, durch die die Werte laufen.}}
  \SemioTableRow{\Key{breaks} & \Key{string} & --- & \ApiText{Class breaks of the choropleth scale.}{Klassengrenzen der Choroplethenskala.}}
  \SemioTableRow{\Key{bivariate} & \Key{boolean} & --- & \ApiText{Shade the regions from two variables at once.}{Die Regionen aus zwei Variablen zugleich einfärben.}}
  \SemioTableRow{\Key{legend} & \Key{boolean} & --- & \ApiText{Draw the class legend.}{Die Klassenlegende zeichnen.}}
  \SemioTableRow{\Key{dasymetric} & \Key{boolean} & --- & \ApiText{Restrict the shading to the populated areas.}{Die Einfärbung auf die besiedelten Flächen beschränken.}}
  \SemioTableRow{\Key{weight} & \Key{number} & --- & \ApiText{Column holding the edge weight.}{Spalte mit dem Kantengewicht.}}
  \SemioTableRow{\Key{geometry} & \Key{string} & --- & \ApiText{Name of the geometry table the map draws.}{Name der Geometrietabelle, die die Karte zeichnet.}}
  \SemioTableRow{\Key{palette} & \Key{string} & --- & \ApiText{Named colour palette resolved by semio-viz-theme.}{Benannte Farbpalette, die semio-viz-theme auflöst.}}
  \SemioTableRow{\Key{steps} & \Key{string} & --- & \ApiText{Number of integration or iteration steps.}{Anzahl der Integrations- bzw. Iterationsschritte.}}
}

\subsection{\Key{geo-dotdensity}}

\SemioTableLong[text-size=8pt]{\Key{geo-dotdensity}}{0.20,0.13,0.15,0.52}{\ApiKey & \ApiType & \ApiDefault & \ApiMeaning}{%
  \SemioTableRow{\Key{unit} & \Key{number} & --- & \ApiText{Data units one glyph or dot stands for.}{Dateneinheiten, für die ein Glyph oder Punkt steht.}}
  \SemioTableRow{\Key{dotSize} & \Key{number} & --- & \ApiText{Millimetre radius of one density dot.}{Millimeterradius eines Dichtepunkts.}}
  \SemioTableRow{\Key{spread} & \Key{number} & --- & \ApiText{Deterministic spread of the placed marks.}{Deterministische Streuung der platzierten Marken.}}
  \SemioTableRow{\Key{geometry} & \Key{string} & --- & \ApiText{Name of the geometry table the map draws.}{Name der Geometrietabelle, die die Karte zeichnet.}}
}

\subsection{\Key{geo-field}}

\SemioTableLong[text-size=8pt]{\Key{geo-field}}{0.20,0.13,0.15,0.52}{\ApiKey & \ApiType & \ApiDefault & \ApiMeaning}{%
  \SemioTableRow{\Key{grid} & \Key{string} & --- & \ApiText{Grid lines drawn behind the marks.}{Gitterlinien hinter den Marken.}}
  \SemioTableRow{\Key{cell} & \Key{number} & --- & \ApiText{Cell size in millimetres, or the column that carries it.}{Zellgröße in Millimetern oder die Spalte, die sie trägt.}}
  \SemioTableRow{\Key{levels} & \Key{integer} & --- & \ApiText{Number of contour levels.}{Anzahl der Höhenlinienniveaus.}}
  \SemioTableRow{\Key{filled} & \Key{boolean} & --- & \ApiText{Fill the bands between contour levels.}{Die Bänder zwischen den Höhenlinien füllen.}}
  \SemioTableRow{\Key{labels} & \Key{boolean} & --- & \ApiText{Draw the captions.}{Die Beschriftungen zeichnen.}}
  \SemioTableRow{\Key{raster} & \Key{boolean} & --- & \ApiText{Draw the field as a filled raster.}{Das Feld als gefülltes Raster zeichnen.}}
  \SemioTableRow{\Key{palette} & \Key{string} & --- & \ApiText{Named colour palette resolved by semio-viz-theme.}{Benannte Farbpalette, die semio-viz-theme auflöst.}}
}

\subsection{\Key{geo-flow}}

\SemioTableLong[text-size=8pt]{\Key{geo-flow}}{0.20,0.13,0.15,0.52}{\ApiKey & \ApiType & \ApiDefault & \ApiMeaning}{%
  \SemioTableRow{\Key{flow} & \Key{string} & --- & \ApiText{Table of origin-destination flows.}{Tabelle der Quelle-Ziel-Ströme.}}
  \SemioTableRow{\Key{curvature} & \Key{number} & --- & \ApiText{Curvature of a link, as a fraction of its span.}{Krümmung einer Kante als Bruchteil ihrer Spannweite.}}
  \SemioTableRow{\Key{origins} & \Key{boolean} & --- & \ApiText{Mark the origin of every flow.}{Den Ursprung jedes Stroms markieren.}}
  \SemioTableRow{\Key{routes} & \Key{string} & --- & \ApiText{Table of routes drawn on the map.}{Tabelle der Routen, die auf der Karte gezeichnet werden.}}
}

\subsection{\Key{geo-geodesic}}

\SemioTableLong[text-size=8pt]{\Key{geo-geodesic}}{0.20,0.13,0.15,0.52}{\ApiKey & \ApiType & \ApiDefault & \ApiMeaning}{%
  \SemioTableRow{\Key{arcs} & \Key{string} & --- & \ApiText{Great-circle arcs drawn, as 'from/to' pairs.}{Gezeichnete Großkreisbögen als 'from/to'-Paare.}}
  \SemioTableRow{\Key{samples} & \Key{integer} & --- & \ApiText{Number of samples the curve or field is evaluated at.}{Anzahl der Stützstellen, an denen Kurve oder Feld ausgewertet werden.}}
  \SemioTableRow{\Key{tissot} & \Key{boolean} & --- & \ApiText{Draw Tissot's indicatrices.}{Die Tissotschen Indikatrizen zeichnen.}}
  \SemioTableRow{\Key{radius} & \Key{number} & --- & \ApiText{Radius in millimetres, or the column that drives it.}{Radius in Millimetern oder die Spalte, die ihn steuert.}}
  \SemioTableRow{\Key{projection} & \Key{string} & --- & \ApiText{Map or camera projection used.}{Verwendete Karten- bzw. Kameraprojektion.}}
}

\subsection{\Key{geo-graticule}}

\SemioTableLong[text-size=8pt]{\Key{geo-graticule}}{0.20,0.13,0.15,0.52}{\ApiKey & \ApiType & \ApiDefault & \ApiMeaning}{%
  \SemioTableRow{\Key{step} & \Key{string} & --- & \ApiText{Spacing between consecutive steps.}{Abstand zwischen aufeinanderfolgenden Schritten.}}
  \SemioTableRow{\Key{stepMajor} & \Key{string} & --- & \ApiText{Degree spacing of the major graticule lines.}{Gradabstand der Hauptgradnetzlinien.}}
  \SemioTableRow{\Key{outline} & \Key{boolean} & --- & \ApiText{Draw the outline of the projected sphere.}{Den Umriss der projizierten Kugel zeichnen.}}
  \SemioTableRow{\Key{sphere} & \Key{boolean} & --- & \ApiText{Draw the sphere outline behind the graticule.}{Den Kugelumriss hinter dem Gradnetz zeichnen.}}
  \SemioTableRow{\Key{labels} & \Key{boolean} & --- & \ApiText{Draw the captions.}{Die Beschriftungen zeichnen.}}
  \SemioTableRow{\Key{geometry} & \Key{string} & --- & \ApiText{Name of the geometry table the map draws.}{Name der Geometrietabelle, die die Karte zeichnet.}}
  \SemioTableRow{\Key{projection} & \Key{string} & --- & \ApiText{Map or camera projection used.}{Verwendete Karten- bzw. Kameraprojektion.}}
}

\subsection{\Key{geo-hexbin}}

\SemioTableLong[text-size=8pt]{\Key{geo-hexbin}}{0.20,0.13,0.15,0.52}{\ApiKey & \ApiType & \ApiDefault & \ApiMeaning}{%
  \SemioTableRow{\Key{points} & \Key{string} & --- & \ApiText{Draw the individual data points.}{Die einzelnen Datenpunkte zeichnen.}}
  \SemioTableRow{\Key{radius} & \Key{number} & --- & \ApiText{Radius in millimetres, or the column that drives it.}{Radius in Millimetern oder die Spalte, die ihn steuert.}}
  \SemioTableRow{\Key{render} & \Key{string} & --- & \ApiText{How the cells or units are rendered.}{Wie die Zellen bzw. Einheiten gezeichnet werden.}}
  \SemioTableRow{\Key{bandwidth} & \Key{number} & --- & \ApiText{Kernel bandwidth of the density estimate; '0' uses Silverman's rule.}{Bandbreite des Kerns der Dichteschätzung; '0' nutzt Silvermans Regel.}}
  \SemioTableRow{\Key{levels} & \Key{integer} & --- & \ApiText{Number of contour levels.}{Anzahl der Höhenlinienniveaus.}}
}

\subsection{\Key{geo-route}}

\SemioTableLong[text-size=8pt]{\Key{geo-route}}{0.20,0.13,0.15,0.52}{\ApiKey & \ApiType & \ApiDefault & \ApiMeaning}{%
  \SemioTableRow{\Key{routes} & \Key{string} & --- & \ApiText{Table of routes drawn on the map.}{Tabelle der Routen, die auf der Karte gezeichnet werden.}}
  \SemioTableRow{\Key{mode} & \Key{string} & --- & \ApiText{Which member of the family the preset renders.}{Welches Mitglied der Familie die Voreinstellung zeichnet.}}
  \SemioTableRow{\Key{samples} & \Key{integer} & --- & \ApiText{Number of samples the curve or field is evaluated at.}{Anzahl der Stützstellen, an denen Kurve oder Feld ausgewertet werden.}}
  \SemioTableRow{\Key{stations} & \Key{boolean} & --- & \ApiText{Draw the stations along the route.}{Die Stationen entlang der Route zeichnen.}}
  \SemioTableRow{\Key{basemap} & \Key{boolean} & --- & \ApiText{Draw the base map behind the overlay.}{Die Grundkarte hinter der Überlagerung zeichnen.}}
  \SemioTableRow{\Key{routeWidth} & \Key{number} & --- & \ApiText{Millimetre width of a route line.}{Millimeterbreite einer Routenlinie.}}
  \SemioTableRow{\Key{palette} & \Key{string} & --- & \ApiText{Named colour palette resolved by semio-viz-theme.}{Benannte Farbpalette, die semio-viz-theme auflöst.}}
  \SemioTableRow{\Key{projection} & \Key{string} & --- & \ApiText{Map or camera projection used.}{Verwendete Karten- bzw. Kameraprojektion.}}
}

\subsection{\Key{geo-symbol}}

\SemioTableLong[text-size=8pt]{\Key{geo-symbol}}{0.20,0.13,0.15,0.52}{\ApiKey & \ApiType & \ApiDefault & \ApiMeaning}{%
  \SemioTableRow{\Key{symbols} & \Key{string} & --- & \ApiText{Symbol vocabulary used on the map.}{Symbolvorrat, der auf der Karte verwendet wird.}}
  \SemioTableRow{\Key{shape} & \Key{string} & --- & \ApiText{Shape of the node marks.}{Form der Knotenmarken.}}
  \SemioTableRow{\Key{sizing} & \Key{string} & --- & \ApiText{How symbol sizes are derived from the values.}{Wie die Symbolgrößen aus den Werten abgeleitet werden.}}
  \SemioTableRow{\Key{maxSize} & \Key{number} & --- & \ApiText{Millimetre size of the largest symbol.}{Millimetergröße des größten Symbols.}}
  \SemioTableRow{\Key{classes} & \Key{integer} & --- & \ApiText{Number of symbol size classes in the legend.}{Anzahl der Symbolgrößenklassen in der Legende.}}
  \SemioTableRow{\Key{basemap} & \Key{boolean} & --- & \ApiText{Draw the base map behind the overlay.}{Die Grundkarte hinter der Überlagerung zeichnen.}}
  \SemioTableRow{\Key{opacity} & \Key{string} & --- & \ApiText{Draw opacity of the layer.}{Zeichendeckkraft der Ebene.}}
}

\subsection{\Key{geo-terrain}}

\SemioTableLong[text-size=8pt]{\Key{geo-terrain}}{0.20,0.13,0.15,0.52}{\ApiKey & \ApiType & \ApiDefault & \ApiMeaning}{%
  \SemioTableRow{\Key{azimuth} & \Key{number} & --- & \ApiText{Camera or light azimuth in degrees.}{Azimut der Kamera bzw. des Lichts in Grad.}}
  \SemioTableRow{\Key{altitude} & \Key{number} & --- & \ApiText{Sun altitude in degrees used for the hillshade.}{Sonnenhöhe in Grad für die Schummerung.}}
  \SemioTableRow{\Key{mode} & \Key{string} & --- & \ApiText{Which member of the family the preset renders.}{Welches Mitglied der Familie die Voreinstellung zeichnet.}}
  \SemioTableRow{\Key{row} & \Key{integer} & --- & \ApiText{Row index on the layout grid.}{Zeilenindex im Anordnungsgitter.}}
  \SemioTableRow{\Key{cell} & \Key{string} & --- & \ApiText{Cell size in millimetres, or the column that carries it.}{Zellgröße in Millimetern oder die Spalte, die sie trägt.}}
  \SemioTableRow{\Key{grid} & \Key{string} & --- & \ApiText{Grid lines drawn behind the marks.}{Gitterlinien hinter den Marken.}}
  \SemioTableRow{\Key{levels} & \Key{string} & --- & \ApiText{Number of contour levels.}{Anzahl der Höhenlinienniveaus.}}
  \SemioTableRow{\Key{palette} & \Key{string} & --- & \ApiText{Named colour palette resolved by semio-viz-theme.}{Benannte Farbpalette, die semio-viz-theme auflöst.}}
  \SemioTableRow{\Key{zScale} & \Key{string} & --- & \ApiText{Vertical exaggeration of the height values.}{Überhöhung der Höhenwerte.}}
}

\subsection{\Key{geo-tilegrid}}

\SemioTableLong[text-size=8pt]{\Key{geo-tilegrid}}{0.20,0.13,0.15,0.52}{\ApiKey & \ApiType & \ApiDefault & \ApiMeaning}{%
  \SemioTableRow{\Key{tile} & \Key{string} & --- & \ApiText{Tiling algorithm that fills the rectangle.}{Kachelalgorithmus, der das Rechteck füllt.}}
  \SemioTableRow{\Key{columns} & \Key{integer} & --- & \ApiText{Number of columns of the layout grid.}{Anzahl der Spalten des Anordnungsgitters.}}
  \SemioTableRow{\Key{cell} & \Key{number} & --- & \ApiText{Cell size in millimetres, or the column that carries it.}{Zellgröße in Millimetern oder die Spalte, die sie trägt.}}
  \SemioTableRow{\Key{labels} & \Key{boolean} & --- & \ApiText{Draw the captions.}{Die Beschriftungen zeichnen.}}
  \SemioTableRow{\Key{geometry} & \Key{string} & --- & \ApiText{Name of the geometry table the map draws.}{Name der Geometrietabelle, die die Karte zeichnet.}}
}

\subsection{\Key{geo-weather}}

\SemioTableLong[text-size=8pt]{\Key{geo-weather}}{0.20,0.13,0.15,0.52}{\ApiKey & \ApiType & \ApiDefault & \ApiMeaning}{%
  \SemioTableRow{\Key{barbs} & \Key{boolean} & --- & \ApiText{Draw wind barbs instead of plain arrows.}{Windfahnen statt einfacher Pfeile zeichnen.}}
  \SemioTableRow{\Key{fronts} & \Key{boolean} & --- & \ApiText{Draw the weather fronts.}{Die Wetterfronten zeichnen.}}
  \SemioTableRow{\Key{barbScale} & \Key{number} & --- & \ApiText{Millimetre length of a full wind barb.}{Millimeterlänge einer vollen Windfahne.}}
  \SemioTableRow{\Key{filled} & \Key{string} & --- & \ApiText{Fill the bands between contour levels.}{Die Bänder zwischen den Höhenlinien füllen.}}
  \SemioTableRow{\Key{grid} & \Key{string} & --- & \ApiText{Grid lines drawn behind the marks.}{Gitterlinien hinter den Marken.}}
  \SemioTableRow{\Key{labels} & \Key{string} & --- & \ApiText{Draw the captions.}{Die Beschriftungen zeichnen.}}
  \SemioTableRow{\Key{levels} & \Key{string} & --- & \ApiText{Number of contour levels.}{Anzahl der Höhenlinienniveaus.}}
  \SemioTableRow{\Key{palette} & \Key{string} & --- & \ApiText{Named colour palette resolved by semio-viz-theme.}{Benannte Farbpalette, die semio-viz-theme auflöst.}}
  \SemioTableRow{\Key{raster} & \Key{string} & --- & \ApiText{Draw the field as a filled raster.}{Das Feld als gefülltes Raster zeichnen.}}
}

\subsection{\Key{logistics}}

\ApiText{This family takes no options beyond \Key{variant} and \Key{data}.}{Diese Familie nimmt außer \Key{variant} und \Key{data} keine Optionen.}

\subsection{\Key{spatial-layout}}

\ApiText{This family takes no options beyond \Key{variant} and \Key{data}.}{Diese Familie nimmt außer \Key{variant} und \Key{data} keine Optionen.}

\subsection{\Key{spatial-scalar-field}}

\SemioTableLong[text-size=8pt]{\Key{spatial-scalar-field}}{0.20,0.13,0.15,0.52}{\ApiKey & \ApiType & \ApiDefault & \ApiMeaning}{%
  \SemioTableRow{\Key{surface} & \Key{string} & --- & \ApiText{Expression or table that defines the surface.}{Ausdruck oder Tabelle, die die Fläche definiert.}}
  \SemioTableRow{\Key{tilt} & \Key{number} & --- & \ApiText{Tilt of the projected surface in degrees.}{Neigung der projizierten Fläche in Grad.}}
  \SemioTableRow{\Key{turn} & \Key{number} & --- & \ApiText{Rotation of the sampled domain in degrees.}{Drehung des abgetasteten Bereichs in Grad.}}
  \SemioTableRow{\Key{zScale} & \Key{number} & --- & \ApiText{Vertical exaggeration of the height values.}{Überhöhung der Höhenwerte.}}
  \SemioTableRow{\Key{cell} & \Key{string} & --- & \ApiText{Cell size in millimetres, or the column that carries it.}{Zellgröße in Millimetern oder die Spalte, die sie trägt.}}
  \SemioTableRow{\Key{grid} & \Key{string} & --- & \ApiText{Grid lines drawn behind the marks.}{Gitterlinien hinter den Marken.}}
  \SemioTableRow{\Key{labels} & \Key{string} & --- & \ApiText{Draw the captions.}{Die Beschriftungen zeichnen.}}
  \SemioTableRow{\Key{levels} & \Key{string} & --- & \ApiText{Number of contour levels.}{Anzahl der Höhenlinienniveaus.}}
  \SemioTableRow{\Key{opacity} & \Key{string} & --- & \ApiText{Draw opacity of the layer.}{Zeichendeckkraft der Ebene.}}
  \SemioTableRow{\Key{palette} & \Key{string} & --- & \ApiText{Named colour palette resolved by semio-viz-theme.}{Benannte Farbpalette, die semio-viz-theme auflöst.}}
  \SemioTableRow{\Key{row} & \Key{string} & --- & \ApiText{Row index on the layout grid.}{Zeilenindex im Anordnungsgitter.}}
  \SemioTableRow{\Key{steps} & \Key{string} & --- & \ApiText{Number of integration or iteration steps.}{Anzahl der Integrations- bzw. Iterationsschritte.}}
}

\subsection{\Key{spatial-tessellation}}

\SemioTableLong[text-size=8pt]{\Key{spatial-tessellation}}{0.20,0.13,0.15,0.52}{\ApiKey & \ApiType & \ApiDefault & \ApiMeaning}{%
  \SemioTableRow{\Key{points} & \Key{string} & --- & \ApiText{Draw the individual data points.}{Die einzelnen Datenpunkte zeichnen.}}
  \SemioTableRow{\Key{kind} & \Key{string} & --- & \ApiText{Which member of the family is drawn.}{Welches Mitglied der Familie gezeichnet wird.}}
  \SemioTableRow{\Key{sites} & \Key{boolean} & --- & \ApiText{Draw the generating sites of the tessellation.}{Die erzeugenden Punkte der Parkettierung zeichnen.}}
  \SemioTableRow{\Key{extent} & \Key{string} & --- & \ApiText{Bounding box the tessellation is clipped to.}{Begrenzungsrechteck, auf das die Parkettierung beschnitten wird.}}
  \SemioTableRow{\Key{palette} & \Key{string} & --- & \ApiText{Named colour palette resolved by semio-viz-theme.}{Benannte Farbpalette, die semio-viz-theme auflöst.}}
}

\subsection{\Key{spatial-vector-field}}

\SemioTableLong[text-size=8pt]{\Key{spatial-vector-field}}{0.20,0.13,0.15,0.52}{\ApiKey & \ApiType & \ApiDefault & \ApiMeaning}{%
  \SemioTableRow{\Key{render} & \Key{string} & --- & \ApiText{How the cells or units are rendered.}{Wie die Zellen bzw. Einheiten gezeichnet werden.}}
  \SemioTableRow{\Key{field} & \Key{string} & --- & \ApiText{Expression or table that defines the vector field.}{Ausdruck oder Tabelle, die das Vektorfeld definiert.}}
  \SemioTableRow{\Key{columns} & \Key{integer} & --- & \ApiText{Number of columns of the layout grid.}{Anzahl der Spalten des Anordnungsgitters.}}
  \SemioTableRow{\Key{rows} & \Key{integer} & --- & \ApiText{Number of rows of the layout grid.}{Anzahl der Zeilen des Anordnungsgitters.}}
  \SemioTableRow{\Key{cell} & \Key{number} & --- & \ApiText{Cell size in millimetres, or the column that carries it.}{Zellgröße in Millimetern oder die Spalte, die sie trägt.}}
  \SemioTableRow{\Key{arrow} & \Key{number} & --- & \ApiText{Arrow head drawn at the end of a connector.}{Pfeilspitze am Ende eines Verbinders.}}
  \SemioTableRow{\Key{steps} & \Key{integer} & --- & \ApiText{Number of integration or iteration steps.}{Anzahl der Integrations- bzw. Iterationsschritte.}}
  \SemioTableRow{\Key{dt} & \Key{number} & --- & \ApiText{Integration step of the streamline tracer.}{Integrationsschritt des Stromlinienverfolgers.}}
  \SemioTableRow{\Key{seeds} & \Key{integer} & --- & \ApiText{Number of streamline seed points.}{Anzahl der Startpunkte der Stromlinien.}}
  \SemioTableRow{\Key{points} & \Key{string} & --- & \ApiText{Draw the individual data points.}{Die einzelnen Datenpunkte zeichnen.}}
}

\section{\ApiText{Diagram}{Diagramm}}

\ApiText{Flowcharts, UML, architecture, process, concept diagrams, infographics and interaction states.}{Flussdiagramme, UML, Architektur, Prozess, Begriffsdiagramme, Infografiken und Interaktionszustände.}

\subsection{\Key{arch-axonometric}}

\SemioTableLong[text-size=8pt]{\Key{arch-axonometric}}{0.20,0.13,0.15,0.52}{\ApiKey & \ApiType & \ApiDefault & \ApiMeaning}{%
  \SemioTableRow{\Key{label} & \Key{string} & --- & \ApiText{Column holding the visible caption.}{Spalte mit der sichtbaren Beschriftung.}}
  \SemioTableRow{\Key{tone} & \Key{string} & --- & \ApiText{Column holding the colour tone of a row.}{Spalte mit dem Farbton einer Zeile.}}
  \SemioTableRow{\Key{explode} & \Key{number} & \Key{0} & \ApiText{Factor applied to each box z-offset for an exploded view.}{Faktor auf den z-Versatz jedes Kastens für eine Explosionsdarstellung.}}
  \SemioTableRow{\Key{unit} & \Key{number} & \Key{0} & \ApiText{Millimetres per model unit, 0 fits the frame.}{Millimeter je Modelleinheit, 0 passt sich dem Rahmen an.}}
  \SemioTableRow{\Key{labels} & \Key{boolean} & --- & \ApiText{Draw the room captions.}{Die Raumbeschriftungen zeichnen.}}
  \SemioTableRow{\Key{width} & \Key{number} & --- & \ApiText{Frame size in millimetres.}{Rahmengröße in Millimetern.}}
  \SemioTableRow{\Key{height} & \Key{number} & --- & \ApiText{Frame size in millimetres.}{Rahmengröße in Millimetern.}}
  \SemioTableRow{\Key{scale-z} & \Key{string} & \Key{1} & \ApiText{Vertical foreshortening.}{Senkrechte Verkürzung.}}
}

\subsection{\Key{arch-bubble}}

\SemioTableLong[text-size=8pt]{\Key{arch-bubble}}{0.20,0.13,0.15,0.52}{\ApiKey & \ApiType & \ApiDefault & \ApiMeaning}{%
  \SemioTableRow{\Key{edges} & \Key{string} & --- & \ApiText{Adjacency table (default demo-arch-bubble-edges), columns from/to/kind.}{Nachbarschaftstabelle (Vorgabe demo-arch-bubble-edges), Spalten from/to/kind.}}
  \SemioTableRow{\Key{id} & \Key{string} & --- & \ApiText{Column holding the row identifier.}{Spalte mit der Zeilenkennung.}}
  \SemioTableRow{\Key{label} & \Key{string} & --- & \ApiText{Column holding the visible caption.}{Spalte mit der sichtbaren Beschriftung.}}
  \SemioTableRow{\Key{tone} & \Key{string} & --- & \ApiText{Column holding the colour tone of a row.}{Spalte mit dem Farbton einer Zeile.}}
  \SemioTableRow{\Key{value} & \Key{string} & \Key{value} & \ApiText{Area column.}{Spalte mit der Fläche.}}
  \SemioTableRow{\Key{radius} & \Key{number} & \Key{0} & \ApiText{Millimetre radius of the placement circle, 0 fits the frame.}{Millimeterradius des Platzierungskreises, 0 passt sich dem Rahmen an.}}
  \SemioTableRow{\Key{scale} & \Key{number} & \Key{1.6} & \ApiText{Millimetres per square root of value.}{Millimeter je Quadratwurzel des Werts.}}
  \SemioTableRow{\Key{start} & \Key{number} & \Key{90} & \ApiText{Degrees of the first bubble.}{Winkellage der ersten Blase in Grad.}}
  \SemioTableRow{\Key{routing} & \Key{string} & --- & \ApiText{Orthogonal|straight|curved.}{Orthogonal|straight|curved — Führung der Verbinder.}}
  \SemioTableRow{\Key{arrow} & \Key{string} & --- & \ApiText{Arrow head drawn at the end of a connector.}{Pfeilspitze am Ende eines Verbinders.}}
  \SemioTableRow{\Key{width} & \Key{number} & --- & \ApiText{Frame size in millimetres.}{Rahmengröße in Millimetern.}}
  \SemioTableRow{\Key{height} & \Key{number} & --- & \ApiText{Frame size in millimetres.}{Rahmengröße in Millimetern.}}
}

\subsection{\Key{arch-layer}}

\SemioTableLong[text-size=8pt]{\Key{arch-layer}}{0.20,0.13,0.15,0.52}{\ApiKey & \ApiType & \ApiDefault & \ApiMeaning}{%
  \SemioTableRow{\Key{edges} & \Key{string} & --- & \ApiText{Dependency table (default demo-arch-layer-edges), columns from/to/kind.}{Abhängigkeitstabelle (Vorgabe demo-arch-layer-edges), Spalten from/to/kind.}}
  \SemioTableRow{\Key{id} & \Key{string} & --- & \ApiText{Column holding the row identifier.}{Spalte mit der Zeilenkennung.}}
  \SemioTableRow{\Key{label} & \Key{string} & --- & \ApiText{Column holding the visible caption.}{Spalte mit der sichtbaren Beschriftung.}}
  \SemioTableRow{\Key{tone} & \Key{string} & --- & \ApiText{Column holding the colour tone of a row.}{Spalte mit dem Farbton einer Zeile.}}
  \SemioTableRow{\Key{tier} & \Key{string} & \Key{tier} & \ApiText{Tier column name.}{Name der Ebenenspalte.}}
  \SemioTableRow{\Key{tiers} & \Key{string} & --- & \ApiText{Tier names top to bottom; empty derives them in first-seen order.}{Ebenennamen von oben nach unten; leer leitet sie in der Reihenfolge des ersten Auftretens ab.}}
  \SemioTableRow{\Key{shape} & \Key{string} & \Key{process} & \ApiText{Box shape.}{Form des Kastens.}}
  \SemioTableRow{\Key{gap} & \Key{number} & \Key{2} & \ApiText{Millimetres between tier bands.}{Millimeter zwischen den Ebenenbändern.}}
  \SemioTableRow{\Key{routing} & \Key{string} & --- & \ApiText{Orthogonal|straight|curved.}{Orthogonal|straight|curved — Führung der Verbinder.}}
  \SemioTableRow{\Key{arrow} & \Key{string} & --- & \ApiText{Arrow head drawn at the end of a connector.}{Pfeilspitze am Ende eines Verbinders.}}
  \SemioTableRow{\Key{width} & \Key{number} & --- & \ApiText{Frame size in millimetres.}{Rahmengröße in Millimetern.}}
  \SemioTableRow{\Key{height} & \Key{number} & --- & \ApiText{Frame size in millimetres.}{Rahmengröße in Millimetern.}}
  \SemioTableRow{\Key{tier-title} & \Key{string} & \Key{true} & \ApiText{Draw the tier caption on the left.}{Die Ebenenbeschriftung links zeichnen.}}
}

\subsection{\Key{arch-plan}}

\SemioTableLong[text-size=8pt]{\Key{arch-plan}}{0.20,0.13,0.15,0.52}{\ApiKey & \ApiType & \ApiDefault & \ApiMeaning}{%
  \SemioTableRow{\Key{id} & \Key{string} & --- & \ApiText{Column holding the row identifier.}{Spalte mit der Zeilenkennung.}}
  \SemioTableRow{\Key{label} & \Key{string} & --- & \ApiText{Column holding the visible caption.}{Spalte mit der sichtbaren Beschriftung.}}
  \SemioTableRow{\Key{tone} & \Key{string} & --- & \ApiText{Column holding the colour tone of a row.}{Spalte mit dem Farbton einer Zeile.}}
  \SemioTableRow{\Key{points} & \Key{string} & \Key{points} & \ApiText{Vertex column, 'x y x y …' in plan units.}{Spalte der Eckpunkte, 'x y x y …' in Planeinheiten.}}
  \SemioTableRow{\Key{hatch} & \Key{string} & \Key{hatch} & \ApiText{Hatch column, values solid|none.}{Schraffurspalte, Werte solid|none.}}
  \SemioTableRow{\Key{domain} & \Key{string} & --- & \ApiText{Plan bounding box 'xmin,ymin,xmax,ymax'.}{Begrenzungsrechteck des Plans 'xmin,ymin,xmax,ymax'.}}
  \SemioTableRow{\Key{grid} & \Key{number} & \Key{0} & \ApiText{Plan-unit spacing of a background grid, 0 for none.}{Abstand eines Hintergrundgitters in Planeinheiten, 0 für keines.}}
  \SemioTableRow{\Key{labels} & \Key{boolean} & \Key{True} & \ApiText{Draw the room captions.}{Die Raumbeschriftungen zeichnen.}}
  \SemioTableRow{\Key{width} & \Key{number} & --- & \ApiText{Frame size in millimetres.}{Rahmengröße in Millimetern.}}
  \SemioTableRow{\Key{height} & \Key{number} & --- & \ApiText{Frame size in millimetres.}{Rahmengröße in Millimetern.}}
}

\subsection{\Key{arch-profile}}

\SemioTableLong[text-size=8pt]{\Key{arch-profile}}{0.20,0.13,0.15,0.52}{\ApiKey & \ApiType & \ApiDefault & \ApiMeaning}{%
  \SemioTableRow{\Key{label} & \Key{string} & --- & \ApiText{Column holding the visible caption.}{Spalte mit der sichtbaren Beschriftung.}}
  \SemioTableRow{\Key{tone} & \Key{string} & --- & \ApiText{Column holding the colour tone of a row.}{Spalte mit dem Farbton einer Zeile.}}
  \SemioTableRow{\Key{base} & \Key{string} & \Key{base/height/width/offset} & \ApiText{Column names in plan units.}{Spaltennamen in Planeinheiten.}}
  \SemioTableRow{\Key{height} & \Key{string} & \Key{base/height/width/offset} & \ApiText{Column names in plan units.}{Spaltennamen in Planeinheiten.}}
  \SemioTableRow{\Key{span} & \Key{string} & --- & \ApiText{Horizontal span of the element in model units.}{Waagerechte Spannweite des Elements in Modelleinheiten.}}
  \SemioTableRow{\Key{offset} & \Key{string} & \Key{base/height/width/offset} & \ApiText{Column names in plan units.}{Spaltennamen in Planeinheiten.}}
  \SemioTableRow{\Key{domain} & \Key{string} & --- & \ApiText{Section extent 'xmin,ymin,xmax,ymax'.}{Ausdehnung des Schnitts 'xmin,ymin,xmax,ymax'.}}
  \SemioTableRow{\Key{ground} & \Key{boolean} & \Key{True} & \ApiText{Draw a ground line at ymin.}{Eine Grundlinie bei ymin zeichnen.}}
}

\subsection{\Key{arch-schematic}}

\SemioTableLong[text-size=8pt]{\Key{arch-schematic}}{0.20,0.13,0.15,0.52}{\ApiKey & \ApiType & \ApiDefault & \ApiMeaning}{%
  \SemioTableRow{\Key{edges} & \Key{string} & --- & \ApiText{Line table (default demo-arch-schematic-edges), columns from/to/kind.}{Leitungstabelle (Vorgabe demo-arch-schematic-edges), Spalten from/to/kind.}}
  \SemioTableRow{\Key{id} & \Key{string} & --- & \ApiText{Column holding the row identifier.}{Spalte mit der Zeilenkennung.}}
  \SemioTableRow{\Key{label} & \Key{string} & --- & \ApiText{Column holding the visible caption.}{Spalte mit der sichtbaren Beschriftung.}}
  \SemioTableRow{\Key{row} & \Key{string} & --- & \ApiText{Row index on the layout grid.}{Zeilenindex im Anordnungsgitter.}}
  \SemioTableRow{\Key{col} & \Key{string} & --- & \ApiText{Column index of a component on the schematic grid.}{Spaltenindex eines Bauteils im Schemagitter.}}
  \SemioTableRow{\Key{symbol} & \Key{string} & --- & \ApiText{Symbol column; pump|valve|tank|heater|cylinder|filter|sensor|controller|actuator|junction.}{Symbolspalte; pump|valve|tank|heater|cylinder|filter|sensor|controller|actuator|junction.}}
  \SemioTableRow{\Key{routing} & \Key{string} & \Key{orthogonal} & \ApiText{Orthogonal|straight|curved.}{Orthogonal|straight|curved — Führung der Verbinder.}}
  \SemioTableRow{\Key{arrow} & \Key{string} & --- & \ApiText{Arrow head drawn at the end of a connector.}{Pfeilspitze am Ende eines Verbinders.}}
  \SemioTableRow{\Key{width} & \Key{number} & --- & \ApiText{Frame size in millimetres.}{Rahmengröße in Millimetern.}}
  \SemioTableRow{\Key{height} & \Key{number} & --- & \ApiText{Frame size in millimetres.}{Rahmengröße in Millimetern.}}
  \SemioTableRow{\Key{col-gap} & \Key{string} & --- & \ApiText{Millimetre gap between grid columns.}{Millimeterabstand zwischen den Gitterspalten.}}
  \SemioTableRow{\Key{row-gap} & \Key{string} & --- & \ApiText{Millimetre gap between grid rows.}{Millimeterabstand zwischen den Gitterzeilen.}}
}

\subsection{\Key{arch-sunpath}}

\SemioTableLong[text-size=8pt]{\Key{arch-sunpath}}{0.20,0.13,0.15,0.52}{\ApiKey & \ApiType & \ApiDefault & \ApiMeaning}{%
  \SemioTableRow{\Key{latitude} & \Key{number} & \Key{52.37} & \ApiText{Degrees north, negative south.}{Grad nördlicher Breite, negativ südlich.}}
  \SemioTableRow{\Key{longitude} & \Key{number} & \Key{9.73} & \ApiText{Degrees east, negative west.}{Grad östlicher Länge, negativ westlich.}}
  \SemioTableRow{\Key{date} & \Key{string} & \Key{2026-06-21} & \ApiText{ISO calendar day 'YYYY-MM-DD' of the arc.}{ISO-Kalendertag 'YYYY-MM-DD' des Bogens.}}
  \SemioTableRow{\Key{hours} & \Key{string} & --- & \ApiText{Local solar hours in UTC to sample.}{Abzutastende wahre Ortsstunden in UTC.}}
  \SemioTableRow{\Key{mode} & \Key{string} & \Key{stereographic} & \ApiText{Stereographic|shadow  polar sun chart or a shadow fan.}{Stereographic|shadow — polares Sonnenbahndiagramm oder ein Schattenfächer.}}
  \SemioTableRow{\Key{pole} & \Key{number} & \Key{0} & \ApiText{Millimetre radius of the horizon circle, 0 fits the frame.}{Millimeterradius des Horizontkreises, 0 passt sich dem Rahmen an.}}
  \SemioTableRow{\Key{object} & \Key{number} & \Key{3} & \ApiText{Metres of object height used by mode=shadow.}{Objekthöhe in Metern, die mode=shadow verwendet.}}
  \SemioTableRow{\Key{width} & \Key{number} & --- & \ApiText{Frame size in millimetres.}{Rahmengröße in Millimetern.}}
  \SemioTableRow{\Key{height} & \Key{number} & --- & \ApiText{Frame size in millimetres.}{Rahmengröße in Millimetern.}}
}

\subsection{\Key{concept-map}}

\SemioTableLong[text-size=8pt]{\Key{concept-map}}{0.20,0.13,0.15,0.52}{\ApiKey & \ApiType & \ApiDefault & \ApiMeaning}{%
  \SemioTableRow{\Key{id} & \Key{string} & --- & \ApiText{Column holding the row identifier.}{Spalte mit der Zeilenkennung.}}
  \SemioTableRow{\Key{label} & \Key{string} & --- & \ApiText{Column holding the visible caption.}{Spalte mit der sichtbaren Beschriftung.}}
  \SemioTableRow{\Key{tone} & \Key{string} & --- & \ApiText{Column holding the colour tone of a row.}{Spalte mit dem Farbton einer Zeile.}}
  \SemioTableRow{\Key{parent} & \Key{string} & \Key{parent} & \ApiText{Parent column, empty for a root.}{Spalte des Elternknotens, leer für eine Wurzel.}}
  \SemioTableRow{\Key{layout} & \Key{string} & \Key{radial} & \ApiText{Radial|balanced|tree  radial rings, left/right mind map, or a top-down tree.}{Radial|balanced|tree — radiale Ringe, links/rechts verzweigte Mindmap oder ein Baum von oben nach unten.}}
  \SemioTableRow{\Key{start} & \Key{number} & \Key{90} & \ApiText{Degrees of the first first-level branch.}{Winkellage des ersten Zweigs der obersten Ebene in Grad.}}
  \SemioTableRow{\Key{sweep} & \Key{number} & \Key{360} & \ApiText{Degrees covered by the first level.}{Gradmaß, das die erste Ebene einnimmt.}}
  \SemioTableRow{\Key{ring} & \Key{number} & \Key{0} & \ApiText{Millimetre distance between depth rings, 0 fits the frame.}{Millimeterabstand zwischen den Tiefenringen, 0 passt sich dem Rahmen an.}}
  \SemioTableRow{\Key{node} & \Key{string} & \Key{pill} & \ApiText{Pill|box|plain  branch node rendering.}{Pill|box|plain — Darstellung eines Zweigknotens.}}
  \SemioTableRow{\Key{link} & \Key{string} & \Key{curved} & \ApiText{Curved|straight|orthogonal.}{Curved|straight|orthogonal — Führung der Verbinder.}}
  \SemioTableRow{\Key{arrow} & \Key{string} & --- & \ApiText{Arrow head drawn at the end of a connector.}{Pfeilspitze am Ende eines Verbinders.}}
  \SemioTableRow{\Key{width} & \Key{number} & --- & \ApiText{Frame size in millimetres.}{Rahmengröße in Millimetern.}}
  \SemioTableRow{\Key{height} & \Key{number} & --- & \ApiText{Frame size in millimetres.}{Rahmengröße in Millimetern.}}
}

\subsection{\Key{concept-shape}}

\SemioTableLong[text-size=8pt]{\Key{concept-shape}}{0.20,0.13,0.15,0.52}{\ApiKey & \ApiType & \ApiDefault & \ApiMeaning}{%
  \SemioTableRow{\Key{label} & \Key{string} & --- & \ApiText{Column holding the visible caption.}{Spalte mit der sichtbaren Beschriftung.}}
  \SemioTableRow{\Key{tone} & \Key{string} & --- & \ApiText{Column holding the colour tone of a row.}{Spalte mit dem Farbton einer Zeile.}}
  \SemioTableRow{\Key{value} & \Key{string} & \Key{value} & \ApiText{Optional weight column driving the level width.}{Optionale Gewichtsspalte, die die Ebenenbreite steuert.}}
  \SemioTableRow{\Key{shape} & \Key{string} & --- & \ApiText{Pyramid|inverted-pyramid|layer|onion|stair|continuum|spectrum|comparison.}{Pyramid|inverted-pyramid|layer|onion|stair|continuum|spectrum|comparison — Aufbau des Diagramms.}}
  \SemioTableRow{\Key{arrows} & \Key{boolean} & \Key{False} & \ApiText{Draw the direction arrow along the sequence.}{Den Richtungspfeil entlang der Folge zeichnen.}}
  \SemioTableRow{\Key{width} & \Key{number} & --- & \ApiText{Frame size in millimetres.}{Rahmengröße in Millimetern.}}
  \SemioTableRow{\Key{height} & \Key{number} & --- & \ApiText{Frame size in millimetres.}{Rahmengröße in Millimetern.}}
}

\subsection{\Key{cycle}}

\SemioTableLong[text-size=8pt]{\Key{cycle}}{0.20,0.13,0.15,0.52}{\ApiKey & \ApiType & \ApiDefault & \ApiMeaning}{%
  \SemioTableRow{\Key{label} & \Key{string} & \Key{label} & \ApiText{Label column.}{Spalte mit der Beschriftung.}}
  \SemioTableRow{\Key{sign} & \Key{string} & \Key{empty = no polarity} & \ApiText{Polarity column for causal loops, values '+'/'-'.}{Polaritätsspalte für Wirkungskreise, Werte '+'/'-'.}}
  \SemioTableRow{\Key{shape} & \Key{string} & \Key{terminator} & \ApiText{Node shape.}{Form eines Knotens.}}
  \SemioTableRow{\Key{start} & \Key{number} & \Key{90} & \ApiText{Degrees of the first stage.}{Winkellage der ersten Stufe in Grad.}}
  \SemioTableRow{\Key{sweep} & \Key{number} & \Key{-360} & \ApiText{Total degrees covered, negative for clockwise.}{Insgesamt überstrichenes Gradmaß, negativ im Uhrzeigersinn.}}
  \SemioTableRow{\Key{closed} & \Key{boolean} & \Key{True} & \ApiText{Close the loop back to the first stage.}{Den Kreis zur ersten Stufe zurück schließen.}}
  \SemioTableRow{\Key{radius} & \Key{number} & \Key{0} & \ApiText{Millimetre radius, '0' fits the frame.}{Millimeterradius, '0' passt sich dem Rahmen an.}}
  \SemioTableRow{\Key{hub} & \Key{string} & --- & \ApiText{Text of a centre hub, empty for none.}{Text einer Mittelnabe, leer für keine.}}
  \SemioTableRow{\Key{width} & \Key{number} & --- & \ApiText{Frame size override in millimetres.}{Vorgabe der Rahmengröße in Millimetern.}}
  \SemioTableRow{\Key{height} & \Key{number} & --- & \ApiText{Frame size override in millimetres.}{Vorgabe der Rahmengröße in Millimetern.}}
}

\subsection{\Key{decision-table}}

\SemioTableLong[text-size=8pt]{\Key{decision-table}}{0.20,0.13,0.15,0.52}{\ApiKey & \ApiType & \ApiDefault & \ApiMeaning}{%
  \SemioTableRow{\Key{label} & \Key{string} & \Key{label} & \ApiText{Rule label column.}{Spalte mit der Beschriftung der Regel.}}
  \SemioTableRow{\Key{columns} & \Key{string} & --- & \ApiText{Comma list of value columns rendered as grid columns.}{Kommaliste der Wertspalten, die als Gitterspalten gezeichnet werden.}}
  \SemioTableRow{\Key{headers} & \Key{string} & --- & \ApiText{Comma list of column captions, empty reuses the column names.}{Kommaliste der Spaltenbeschriftungen; leer übernimmt die Spaltennamen.}}
  \SemioTableRow{\Key{width} & \Key{number} & --- & \ApiText{Frame size override in millimetres.}{Vorgabe der Rahmengröße in Millimetern.}}
  \SemioTableRow{\Key{height} & \Key{number} & --- & \ApiText{Frame size override in millimetres.}{Vorgabe der Rahmengröße in Millimetern.}}
}

\subsection{\Key{engineering-diagram}}

\ApiText{This family takes no options beyond \Key{variant} and \Key{data}.}{Diese Familie nimmt außer \Key{variant} und \Key{data} keine Optionen.}

\subsection{\Key{fishbone}}

\SemioTableLong[text-size=8pt]{\Key{fishbone}}{0.20,0.13,0.15,0.52}{\ApiKey & \ApiType & \ApiDefault & \ApiMeaning}{%
  \SemioTableRow{\Key{label} & \Key{string} & \Key{label} & \ApiText{Cause label column.}{Spalte mit der Beschriftung der Ursache.}}
  \SemioTableRow{\Key{group} & \Key{string} & \Key{grp} & \ApiText{Bone/category column.}{Spalte der Gräte bzw. Kategorie.}}
  \SemioTableRow{\Key{effect} & \Key{string} & \Key{the table name} & \ApiText{Text of the head box.}{Text des Kopfkastens.}}
  \SemioTableRow{\Key{angle} & \Key{number} & \Key{32} & \ApiText{Degrees of bone inclination.}{Neigung der Gräten in Grad.}}
  \SemioTableRow{\Key{width} & \Key{number} & --- & \ApiText{Frame size override in millimetres.}{Vorgabe der Rahmengröße in Millimetern.}}
  \SemioTableRow{\Key{height} & \Key{number} & --- & \ApiText{Frame size override in millimetres.}{Vorgabe der Rahmengröße in Millimetern.}}
}

\subsection{\Key{flow}}

\SemioTableLong[text-size=8pt]{\Key{flow}}{0.20,0.13,0.15,0.52}{\ApiKey & \ApiType & \ApiDefault & \ApiMeaning}{%
  \SemioTableRow{\Key{arrow} & \Key{string} & --- & \ApiText{Arrow head drawn at the end of a connector.}{Pfeilspitze am Ende eines Verbinders.}}
  \SemioTableRow{\Key{bend} & \Key{string} & --- & \ApiText{Bend style of the flow connectors.}{Krümmungsstil der Ablaufverbinder.}}
  \SemioTableRow{\Key{default-shape} & \Key{string} & --- & \ApiText{Node shape used when a row carries none.}{Knotenform, wenn eine Zeile keine mitbringt.}}
  \SemioTableRow{\Key{direction} & \Key{string} & --- & \ApiText{Primary direction the flow grows in.}{Hauptrichtung, in die der Ablauf wächst.}}
  \SemioTableRow{\Key{groups} & \Key{string} & --- & \ApiText{Group boxes drawn behind the members.}{Gruppenrahmen hinter den Mitgliedern.}}
  \SemioTableRow{\Key{routing} & \Key{string} & --- & \ApiText{How the connectors are routed between nodes.}{Wie die Verbinder zwischen den Knoten geführt werden.}}
}

\subsection{\Key{infographic-icon}}

\SemioTableLong[text-size=8pt]{\Key{infographic-icon}}{0.20,0.13,0.15,0.52}{\ApiKey & \ApiType & \ApiDefault & \ApiMeaning}{%
  \SemioTableRow{\Key{label} & \Key{string} & --- & \ApiText{Column holding the visible caption.}{Spalte mit der sichtbaren Beschriftung.}}
  \SemioTableRow{\Key{tone} & \Key{string} & --- & \ApiText{Column holding the colour tone of a row.}{Spalte mit dem Farbton einer Zeile.}}
  \SemioTableRow{\Key{value} & \Key{string} & \Key{value} & \ApiText{Count column.}{Spalte mit der Anzahl.}}
  \SemioTableRow{\Key{columns} & \Key{integer} & \Key{10} & \ApiText{Glyphs per row.}{Glyphen je Zeile.}}
  \SemioTableRow{\Key{glyph} & \Key{string} & \Key{square} & \ApiText{Square|circle|person.}{Square|circle|person — Form eines Sitzes.}}
  \SemioTableRow{\Key{width} & \Key{number} & --- & \ApiText{Frame size in millimetres.}{Rahmengröße in Millimetern.}}
  \SemioTableRow{\Key{height} & \Key{number} & --- & \ApiText{Frame size in millimetres.}{Rahmengröße in Millimetern.}}
  \SemioTableRow{\Key{unit-value} & \Key{string} & \Key{1} & \ApiText{Units represented by one glyph.}{Einheiten, die ein Glyph darstellt.}}
}

\subsection{\Key{infographic-illustration}}

\SemioTableLong[text-size=8pt]{\Key{infographic-illustration}}{0.20,0.13,0.15,0.52}{\ApiKey & \ApiType & \ApiDefault & \ApiMeaning}{%
  \SemioTableRow{\Key{label} & \Key{string} & --- & \ApiText{Column holding the visible caption.}{Spalte mit der sichtbaren Beschriftung.}}
  \SemioTableRow{\Key{tone} & \Key{string} & --- & \ApiText{Column holding the colour tone of a row.}{Spalte mit dem Farbton einer Zeile.}}
  \SemioTableRow{\Key{side} & \Key{string} & --- & \ApiText{Which side the secondary element sits on.}{Auf welcher Seite das Nebenelement sitzt.}}
  \SemioTableRow{\Key{outline} & \Key{string} & \Key{empty = an ellipse} & \ApiText{Silhouette vertex column list 'x y x y …' in model units.}{Spaltenliste der Silhouettenpunkte 'x y x y …' in Modelleinheiten.}}
  \SemioTableRow{\Key{cutaway} & \Key{number} & \Key{0} & \ApiText{Fraction of the silhouette drawn as a cut wedge, 0 for none.}{Anteil der Silhouette, der als Schnittkeil gezeichnet wird, 0 für keinen.}}
  \SemioTableRow{\Key{explode} & \Key{number} & \Key{0} & \ApiText{Radial offset factor of the callout targets.}{Radialer Versatzfaktor der Sprechblasenziele.}}
  \SemioTableRow{\Key{unit} & \Key{number} & \Key{0} & \ApiText{Millimetres per model unit, 0 fits the frame.}{Millimeter je Modelleinheit, 0 passt sich dem Rahmen an.}}
  \SemioTableRow{\Key{width} & \Key{number} & --- & \ApiText{Frame size in millimetres.}{Rahmengröße in Millimetern.}}
  \SemioTableRow{\Key{height} & \Key{number} & --- & \ApiText{Frame size in millimetres.}{Rahmengröße in Millimetern.}}
}

\subsection{\Key{infographic-list}}

\SemioTableLong[text-size=8pt]{\Key{infographic-list}}{0.20,0.13,0.15,0.52}{\ApiKey & \ApiType & \ApiDefault & \ApiMeaning}{%
  \SemioTableRow{\Key{label} & \Key{string} & --- & \ApiText{Column holding the visible caption.}{Spalte mit der sichtbaren Beschriftung.}}
  \SemioTableRow{\Key{tone} & \Key{string} & --- & \ApiText{Column holding the colour tone of a row.}{Spalte mit dem Farbton einer Zeile.}}
  \SemioTableRow{\Key{badge} & \Key{string} & \Key{number} & \ApiText{Number|bullet|icon  leading marker.}{Number|bullet|icon — führende Marke.}}
  \SemioTableRow{\Key{detail} & \Key{string} & \Key{value} & \ApiText{Optional second-line column.}{Optionale Spalte für die zweite Zeile.}}
  \SemioTableRow{\Key{width} & \Key{number} & --- & \ApiText{Frame size in millimetres.}{Rahmengröße in Millimetern.}}
  \SemioTableRow{\Key{height} & \Key{number} & --- & \ApiText{Frame size in millimetres.}{Rahmengröße in Millimetern.}}
}

\subsection{\Key{infographic-number}}

\SemioTableLong[text-size=8pt]{\Key{infographic-number}}{0.20,0.13,0.15,0.52}{\ApiKey & \ApiType & \ApiDefault & \ApiMeaning}{%
  \SemioTableRow{\Key{label} & \Key{string} & --- & \ApiText{Column holding the visible caption.}{Spalte mit der sichtbaren Beschriftung.}}
  \SemioTableRow{\Key{tone} & \Key{string} & --- & \ApiText{Column holding the colour tone of a row.}{Spalte mit dem Farbton einer Zeile.}}
  \SemioTableRow{\Key{value} & \Key{string} & \Key{value} & \ApiText{Headline number column.}{Spalte mit der Leitzahl.}}
  \SemioTableRow{\Key{unit} & \Key{string} & \Key{unit} & \ApiText{Suffix column drawn next to the number.}{Spalte des Zusatzes, neben der Zahl gezeichnet.}}
  \SemioTableRow{\Key{columns} & \Key{integer} & \Key{3} & \ApiText{Cards per row.}{Karten je Zeile.}}
  \SemioTableRow{\Key{rule} & \Key{boolean} & \Key{True} & \ApiText{Draw a separating rule under each card.}{Unter jeder Karte eine Trennlinie zeichnen.}}
  \SemioTableRow{\Key{align} & \Key{string} & \Key{center} & \ApiText{Center|left.}{Center|left — Ausrichtung der Zahl.}}
  \SemioTableRow{\Key{width} & \Key{number} & --- & \ApiText{Frame size in millimetres.}{Rahmengröße in Millimetern.}}
  \SemioTableRow{\Key{height} & \Key{number} & --- & \ApiText{Frame size in millimetres.}{Rahmengröße in Millimetern.}}
}

\subsection{\Key{logic-tree}}

\SemioTableLong[text-size=8pt]{\Key{logic-tree}}{0.20,0.13,0.15,0.52}{\ApiKey & \ApiType & \ApiDefault & \ApiMeaning}{%
  \SemioTableRow{\Key{id} & \Key{string} & \Key{id/parent/label/shape} & \ApiText{Column names.}{Namen der gebundenen Spalten.}}
  \SemioTableRow{\Key{parent} & \Key{string} & \Key{id/parent/label/shape} & \ApiText{Column names.}{Namen der gebundenen Spalten.}}
  \SemioTableRow{\Key{label} & \Key{string} & \Key{id/parent/label/shape} & \ApiText{Column names.}{Namen der gebundenen Spalten.}}
  \SemioTableRow{\Key{shape} & \Key{string} & \Key{id/parent/label/shape} & \ApiText{Column names.}{Namen der gebundenen Spalten.}}
  \SemioTableRow{\Key{gate} & \Key{string} & \Key{none} & \ApiText{Shape drawn on a node that has more than one child.}{Form, die an einem Knoten mit mehr als einem Kind gezeichnet wird.}}
  \SemioTableRow{\Key{orientation} & \Key{string} & \Key{down} & \ApiText{Down|right.}{Down|right — Wachstumsrichtung.}}
  \SemioTableRow{\Key{routing} & \Key{string} & --- & \ApiText{Orthogonal|straight|curved.}{Orthogonal|straight|curved — Führung der Verbinder.}}
  \SemioTableRow{\Key{width} & \Key{number} & --- & \ApiText{Frame size override in millimetres.}{Vorgabe der Rahmengröße in Millimetern.}}
  \SemioTableRow{\Key{height} & \Key{number} & --- & \ApiText{Frame size override in millimetres.}{Vorgabe der Rahmengröße in Millimetern.}}
  \SemioTableRow{\Key{default-shape} & \Key{string} & \Key{process} & \ApiText{Shape for rows without one.}{Form für Zeilen ohne eigene Angabe.}}
}

\subsection{\Key{matrix-grid}}

\SemioTableLong[text-size=8pt]{\Key{matrix-grid}}{0.20,0.13,0.15,0.52}{\ApiKey & \ApiType & \ApiDefault & \ApiMeaning}{%
  \SemioTableRow{\Key{label} & \Key{string} & --- & \ApiText{As in 'pm-gantt'.}{Wie in der Familie 'pm-gantt'.}}
  \SemioTableRow{\Key{tone} & \Key{string} & --- & \ApiText{Column holding the colour tone of a row.}{Spalte mit dem Farbton einer Zeile.}}
  \SemioTableRow{\Key{row} & \Key{string} & --- & \ApiText{Grid position columns (default row/col); missing columns place by CPM start.}{Spalten der Gitterposition (Vorgabe row/col); fehlende Spalten platzieren nach dem CPM-Start.}}
  \SemioTableRow{\Key{col} & \Key{string} & --- & \ApiText{Grid position columns (default row/col); missing columns place by CPM start.}{Spalten der Gitterposition (Vorgabe row/col); fehlende Spalten platzieren nach dem CPM-Start.}}
  \SemioTableRow{\Key{rows} & \Key{string} & --- & \ApiText{Row captions top to bottom.}{Zeilenbeschriftungen von oben nach unten.}}
  \SemioTableRow{\Key{columns} & \Key{string} & --- & \ApiText{Column captions left to right.}{Spaltenbeschriftungen von links nach rechts.}}
  \SemioTableRow{\Key{diagonal} & \Key{boolean} & \Key{False} & \ApiText{Draw the leading diagonal, e.g. for a compatibility matrix.}{Die Hauptdiagonale zeichnen, etwa für eine Verträglichkeitsmatrix.}}
  \SemioTableRow{\Key{roof} & \Key{boolean} & \Key{False} & \ApiText{Draw the QFD correlation roof above the columns.}{Das QFD-Korrelationsdach über den Spalten zeichnen.}}
  \SemioTableRow{\Key{width} & \Key{number} & --- & \ApiText{Frame size in millimetres.}{Rahmengröße in Millimetern.}}
  \SemioTableRow{\Key{height} & \Key{number} & --- & \ApiText{Frame size in millimetres.}{Rahmengröße in Millimetern.}}
  \SemioTableRow{\Key{axis-x} & \Key{string} & --- & \ApiText{Axis captions, empty for none.}{Achsenbeschriftungen, leer für keine.}}
  \SemioTableRow{\Key{axis-y} & \Key{string} & --- & \ApiText{Axis captions, empty for none.}{Achsenbeschriftungen, leer für keine.}}
}

\subsection{\Key{notation}}

\ApiText{This family takes no options beyond \Key{variant} and \Key{data}.}{Diese Familie nimmt außer \Key{variant} und \Key{data} keine Optionen.}

\subsection{\Key{notation-board}}

\SemioTableLong[text-size=8pt]{\Key{notation-board}}{0.20,0.13,0.15,0.52}{\ApiKey & \ApiType & \ApiDefault & \ApiMeaning}{%
  \SemioTableRow{\Key{files} & \Key{integer} & --- & \ApiText{Board size.}{Größe des Spielbretts.}}
  \SemioTableRow{\Key{ranks} & \Key{integer} & --- & \ApiText{Board size.}{Größe des Spielbretts.}}
  \SemioTableRow{\Key{checker} & \Key{boolean} & \Key{True} & \ApiText{Alternate the cell fill, off for a Go grid.}{Die Zellfüllung abwechseln; ausgeschaltet für ein Go-Gitter.}}
  \SemioTableRow{\Key{lines} & \Key{boolean} & \Key{False} & \ApiText{Draw pieces on the intersections instead of in the cells.}{Die Steine auf die Schnittpunkte statt in die Felder setzen.}}
  \SemioTableRow{\Key{coords} & \Key{boolean} & \Key{True} & \ApiText{Draw the file and rank captions.}{Die Linien- und Reihenbeschriftungen zeichnen.}}
  \SemioTableRow{\Key{width} & \Key{number} & --- & \ApiText{Frame size in millimetres.}{Rahmengröße in Millimetern.}}
  \SemioTableRow{\Key{height} & \Key{number} & --- & \ApiText{Frame size in millimetres.}{Rahmengröße in Millimetern.}}
}

\subsection{\Key{notation-railroad}}

\SemioTableLong[text-size=8pt]{\Key{notation-railroad}}{0.20,0.13,0.15,0.52}{\ApiKey & \ApiType & \ApiDefault & \ApiMeaning}{%
  \SemioTableRow{\Key{label} & \Key{string} & --- & \ApiText{As in 'pm-gantt'.}{Wie in der Familie 'pm-gantt'.}}
  \SemioTableRow{\Key{kind} & \Key{string} & \Key{kind} & \ApiText{Term kind column; terminal|nonterminal|skip.}{Spalte der Termart; terminal|nonterminal|skip.}}
  \SemioTableRow{\Key{branch} & \Key{string} & \Key{branch} & \ApiText{Branch index column; '0' is the main line.}{Spalte des Zweigindex; '0' ist die Hauptlinie.}}
  \SemioTableRow{\Key{width} & \Key{number} & --- & \ApiText{Frame size in millimetres.}{Rahmengröße in Millimetern.}}
  \SemioTableRow{\Key{height} & \Key{number} & --- & \ApiText{Frame size in millimetres.}{Rahmengröße in Millimetern.}}
}

\subsection{\Key{pm-board}}

\SemioTableLong[text-size=8pt]{\Key{pm-board}}{0.20,0.13,0.15,0.52}{\ApiKey & \ApiType & \ApiDefault & \ApiMeaning}{%
  \SemioTableRow{\Key{label} & \Key{string} & --- & \ApiText{As in 'pm-gantt'.}{Wie in der Familie 'pm-gantt'.}}
  \SemioTableRow{\Key{tone} & \Key{string} & --- & \ApiText{Column holding the colour tone of a row.}{Spalte mit dem Farbton einer Zeile.}}
  \SemioTableRow{\Key{status} & \Key{string} & \Key{status} & \ApiText{Column key.}{Spalte mit dem Schlüssel.}}
  \SemioTableRow{\Key{columns} & \Key{string} & --- & \ApiText{Column order; empty derives it in first-seen order.}{Spaltenreihenfolge; leer leitet sie in der Reihenfolge des ersten Auftretens ab.}}
  \SemioTableRow{\Key{wip} & \Key{string} & --- & \ApiText{Per-column limits drawn in the header, empty for none.}{Grenzwerte je Spalte, im Kopf gezeichnet, leer für keine.}}
  \SemioTableRow{\Key{size} & \Key{string} & \Key{size} & \ApiText{Optional effort column shown as a chip.}{Optionale Aufwandsspalte, als Chip dargestellt.}}
  \SemioTableRow{\Key{width} & \Key{number} & --- & \ApiText{Frame size in millimetres.}{Rahmengröße in Millimetern.}}
  \SemioTableRow{\Key{height} & \Key{number} & --- & \ApiText{Frame size in millimetres.}{Rahmengröße in Millimetern.}}
}

\subsection{\Key{pm-gantt}}

\SemioTableLong[text-size=8pt]{\Key{pm-gantt}}{0.20,0.13,0.15,0.52}{\ApiKey & \ApiType & \ApiDefault & \ApiMeaning}{%
  \SemioTableRow{\Key{id} & \Key{string} & --- & \ApiText{As in 'pm-gantt'.}{Wie in der Familie 'pm-gantt'.}}
  \SemioTableRow{\Key{label} & \Key{string} & --- & \ApiText{As in 'pm-gantt'.}{Wie in der Familie 'pm-gantt'.}}
  \SemioTableRow{\Key{tone} & \Key{string} & --- & \ApiText{Column holding the colour tone of a row.}{Spalte mit dem Farbton einer Zeile.}}
  \SemioTableRow{\Key{lane} & \Key{string} & \Key{lane} & \ApiText{Optional swimlane column.}{Optionale Spalte für die Schwimmbahn.}}
  \SemioTableRow{\Key{duration} & \Key{string} & \Key{duration} & \ApiText{Duration column.}{Spalte mit der Dauer.}}
  \SemioTableRow{\Key{deps} & \Key{string} & \Key{deps} & \ApiText{Space-separated predecessor id column.}{Spalte der Vorgängerkennungen, durch Leerzeichen getrennt.}}
  \SemioTableRow{\Key{start} & \Key{string} & \Key{start} & \ApiText{Explicit start column used by schedule=explicit.}{Ausdrückliche Startspalte, die schedule=explicit verwendet.}}
  \SemioTableRow{\Key{schedule} & \Key{string} & \Key{cpm} & \ApiText{Cpm|explicit  bar start from the CPM forward pass or from a 'start' column.}{Cpm|explicit — Balkenbeginn aus dem CPM-Vorwärtsdurchlauf oder aus einer 'start'-Spalte.}}
  \SemioTableRow{\Key{domain} & \Key{string} & --- & \ApiText{Time domain 'min,max'; empty derives it from the schedule.}{Zeitbereich 'min,max'; leer leitet ihn aus dem Terminplan ab.}}
  \SemioTableRow{\Key{links} & \Key{boolean} & \Key{True} & \ApiText{Draw the dependency arrows.}{Die Abhängigkeitspfeile zeichnen.}}
  \SemioTableRow{\Key{critical} & \Key{boolean} & \Key{True} & \ApiText{Highlight the zero-float chain.}{Die Kette ohne Pufferzeit hervorheben.}}
  \SemioTableRow{\Key{milestone} & \Key{string} & \Key{milestone} & \ApiText{Column marking zero-duration milestones.}{Spalte, die Meilensteine ohne Dauer kennzeichnet.}}
  \SemioTableRow{\Key{axis} & \Key{number} & \Key{0} & \ApiText{Time-axis tick spacing, 0 for no axis.}{Teilstrichabstand der Zeitachse, 0 für keine Achse.}}
  \SemioTableRow{\Key{gutter} & \Key{number} & --- & \ApiText{Millimetre gutter between the label column and the bars.}{Millimeterabstand zwischen Beschriftungsspalte und Balken.}}
  \SemioTableRow{\Key{width} & \Key{number} & --- & \ApiText{Frame size in millimetres.}{Rahmengröße in Millimetern.}}
  \SemioTableRow{\Key{height} & \Key{number} & --- & \ApiText{Frame size in millimetres.}{Rahmengröße in Millimetern.}}
}

\subsection{\Key{pm-network}}

\SemioTableLong[text-size=8pt]{\Key{pm-network}}{0.20,0.13,0.15,0.52}{\ApiKey & \ApiType & \ApiDefault & \ApiMeaning}{%
  \SemioTableRow{\Key{id} & \Key{string} & --- & \ApiText{As in 'pm-gantt'.}{Wie in der Familie 'pm-gantt'.}}
  \SemioTableRow{\Key{label} & \Key{string} & --- & \ApiText{As in 'pm-gantt'.}{Wie in der Familie 'pm-gantt'.}}
  \SemioTableRow{\Key{row} & \Key{string} & --- & \ApiText{Grid position columns (default row/col); missing columns place by CPM start.}{Spalten der Gitterposition (Vorgabe row/col); fehlende Spalten platzieren nach dem CPM-Start.}}
  \SemioTableRow{\Key{col} & \Key{string} & --- & \ApiText{Grid position columns (default row/col); missing columns place by CPM start.}{Spalten der Gitterposition (Vorgabe row/col); fehlende Spalten platzieren nach dem CPM-Start.}}
  \SemioTableRow{\Key{duration} & \Key{string} & --- & \ApiText{As in 'pm-gantt'.}{Wie in der Familie 'pm-gantt'.}}
  \SemioTableRow{\Key{deps} & \Key{string} & --- & \ApiText{As in 'pm-gantt'.}{Wie in der Familie 'pm-gantt'.}}
  \SemioTableRow{\Key{times} & \Key{boolean} & \Key{True} & \ApiText{Print the four schedule numbers inside the node.}{Die vier Terminzahlen in den Knoten schreiben.}}
  \SemioTableRow{\Key{critical} & \Key{boolean} & \Key{True} & \ApiText{Highlight the zero-float chain.}{Die Kette ohne Pufferzeit hervorheben.}}
  \SemioTableRow{\Key{routing} & \Key{string} & --- & \ApiText{How the connectors are routed between nodes.}{Wie die Verbinder zwischen den Knoten geführt werden.}}
  \SemioTableRow{\Key{width} & \Key{number} & --- & \ApiText{Frame size in millimetres.}{Rahmengröße in Millimetern.}}
  \SemioTableRow{\Key{height} & \Key{number} & --- & \ApiText{Frame size in millimetres.}{Rahmengröße in Millimetern.}}
}

\subsection{\Key{schedule-calendar}}

\SemioTableLong[text-size=8pt]{\Key{schedule-calendar}}{0.20,0.13,0.15,0.52}{\ApiKey & \ApiType & \ApiDefault & \ApiMeaning}{%
  \SemioTableRow{\Key{label} & \Key{string} & --- & \ApiText{As in 'pm-gantt'.}{Wie in der Familie 'pm-gantt'.}}
  \SemioTableRow{\Key{tone} & \Key{string} & --- & \ApiText{Column holding the colour tone of a row.}{Spalte mit dem Farbton einer Zeile.}}
  \SemioTableRow{\Key{resource} & \Key{string} & \Key{resource} & \ApiText{Lane column.}{Spalte mit der Bahn.}}
  \SemioTableRow{\Key{lanes} & \Key{string} & --- & \ApiText{Lane order; empty derives it in first-seen order.}{Reihenfolge der Bahnen; leer leitet sie in der Reihenfolge des ersten Auftretens ab.}}
  \SemioTableRow{\Key{start} & \Key{string} & \Key{start/end} & \ApiText{Numeric interval columns in domain units.}{Numerische Intervallspalten in Bereichseinheiten.}}
  \SemioTableRow{\Key{end} & \Key{string} & \Key{start/end} & \ApiText{Numeric interval columns in domain units.}{Numerische Intervallspalten in Bereichseinheiten.}}
  \SemioTableRow{\Key{domain} & \Key{string} & --- & \ApiText{'min,max' of the time axis.}{'min,max' der Zeitachse.}}
  \SemioTableRow{\Key{ticks} & \Key{number} & \Key{6} & \ApiText{Axis tick spacing, 0 for none.}{Abstand der Achsenteilstriche, 0 für keine.}}
  \SemioTableRow{\Key{orientation} & \Key{string} & \Key{horizontal} & \ApiText{Horizontal|vertical  lanes as rows or as columns.}{Horizontal|vertical — Bahnen als Zeilen oder als Spalten.}}
  \SemioTableRow{\Key{width} & \Key{number} & --- & \ApiText{Frame size in millimetres.}{Rahmengröße in Millimetern.}}
  \SemioTableRow{\Key{height} & \Key{number} & --- & \ApiText{Frame size in millimetres.}{Rahmengröße in Millimetern.}}
}

\subsection{\Key{schedule-timedistance}}

\SemioTableLong[text-size=8pt]{\Key{schedule-timedistance}}{0.20,0.13,0.15,0.52}{\ApiKey & \ApiType & \ApiDefault & \ApiMeaning}{%
  \SemioTableRow{\Key{service} & \Key{string} & \Key{service} & \ApiText{Service column.}{Spalte mit der Verbindung.}}
  \SemioTableRow{\Key{distance} & \Key{string} & \Key{distance} & \ApiText{Distance column.}{Spalte mit der Entfernung.}}
  \SemioTableRow{\Key{time} & \Key{string} & \Key{time} & \ApiText{Time column.}{Spalte mit der Zeit.}}
  \SemioTableRow{\Key{stops} & \Key{string} & \Key{stop} & \ApiText{Stop-name column drawn as horizontal guides.}{Spalte der Haltestellennamen, als waagerechte Hilfslinien gezeichnet.}}
  \SemioTableRow{\Key{width} & \Key{number} & --- & \ApiText{Frame size in millimetres.}{Rahmengröße in Millimetern.}}
  \SemioTableRow{\Key{height} & \Key{number} & --- & \ApiText{Frame size in millimetres.}{Rahmengröße in Millimetern.}}
  \SemioTableRow{\Key{domain-x} & \Key{string} & --- & \ApiText{Time domain 'min,max'.}{Zeitbereich 'min,max'.}}
  \SemioTableRow{\Key{domain-y} & \Key{string} & --- & \ApiText{Distance domain 'min,max'.}{Entfernungsbereich 'min,max'.}}
}

\subsection{\Key{seating}}

\SemioTableLong[text-size=8pt]{\Key{seating}}{0.20,0.13,0.15,0.52}{\ApiKey & \ApiType & \ApiDefault & \ApiMeaning}{%
  \SemioTableRow{\Key{label} & \Key{string} & --- & \ApiText{As in 'pm-gantt'.}{Wie in der Familie 'pm-gantt'.}}
  \SemioTableRow{\Key{tone} & \Key{string} & --- & \ApiText{Column holding the colour tone of a row.}{Spalte mit dem Farbton einer Zeile.}}
  \SemioTableRow{\Key{count} & \Key{string} & \Key{count} & \ApiText{Seat-count column.}{Spalte mit der Sitzzahl.}}
  \SemioTableRow{\Key{shape} & \Key{string} & \Key{hemicycle} & \ApiText{Hemicycle|rows|circle|rectangle|theatre.}{Hemicycle|rows|circle|rectangle|theatre — Sitzanordnung.}}
  \SemioTableRow{\Key{rows} & \Key{integer} & \Key{4} & \ApiText{Seat rows for hemicycle/theatre/rectangle.}{Sitzreihen für hemicycle, theatre und rectangle.}}
  \SemioTableRow{\Key{seat} & \Key{number} & \Key{1.1} & \ApiText{Millimetre seat radius.}{Millimeterradius eines Sitzes.}}
  \SemioTableRow{\Key{start} & \Key{number} & --- & \ApiText{Degrees of the arc for hemicycle and circle.}{Gradmaß des Bogens für Halbrund und Kreis.}}
  \SemioTableRow{\Key{sweep} & \Key{number} & --- & \ApiText{Degrees of the arc for hemicycle and circle.}{Gradmaß des Bogens für Halbrund und Kreis.}}
  \SemioTableRow{\Key{width} & \Key{number} & --- & \ApiText{Frame size in millimetres.}{Rahmengröße in Millimetern.}}
  \SemioTableRow{\Key{height} & \Key{number} & --- & \ApiText{Frame size in millimetres.}{Rahmengröße in Millimetern.}}
}

\subsection{\Key{stage}}

\SemioTableLong[text-size=8pt]{\Key{stage}}{0.20,0.13,0.15,0.52}{\ApiKey & \ApiType & \ApiDefault & \ApiMeaning}{%
  \SemioTableRow{\Key{label} & \Key{string} & \Key{label} & \ApiText{Label column.}{Spalte mit der Beschriftung.}}
  \SemioTableRow{\Key{value} & \Key{string} & \Key{empty = equal bands} & \ApiText{Optional numeric column driving band height.}{Optionale numerische Spalte, die die Bandhöhe steuert.}}
  \SemioTableRow{\Key{shape} & \Key{string} & \Key{chevron} & \ApiText{Band shape: chevron|box|trapezoid|step.}{Bandform: chevron|box|trapezoid|step.}}
  \SemioTableRow{\Key{orientation} & \Key{string} & \Key{right} & \ApiText{Right|down.}{Right|down — Wachstumsrichtung.}}
  \SemioTableRow{\Key{gate} & \Key{boolean} & \Key{False} & \ApiText{Draw a gate marker between stages.}{Eine Torkennzeichnung zwischen den Stufen zeichnen.}}
  \SemioTableRow{\Key{taper} & \Key{number} & \Key{1} & \ApiText{Fraction the last band narrows to, 1 = no taper.}{Anteil, auf den sich das letzte Band verjüngt, 1 = keine Verjüngung.}}
  \SemioTableRow{\Key{width} & \Key{number} & --- & \ApiText{Frame size override in millimetres.}{Vorgabe der Rahmengröße in Millimetern.}}
  \SemioTableRow{\Key{height} & \Key{number} & --- & \ApiText{Frame size override in millimetres.}{Vorgabe der Rahmengröße in Millimetern.}}
}

\subsection{\Key{state-overview}}

\SemioTableLong[text-size=8pt]{\Key{state-overview}}{0.20,0.13,0.15,0.52}{\ApiKey & \ApiType & \ApiDefault & \ApiMeaning}{%
  \SemioTableRow{\Key{brush} & \Key{string} & --- & \ApiText{Selection window 't0,t1' in data units.}{Auswahlfenster 't0,t1' in Dateneinheiten.}}
  \SemioTableRow{\Key{domain} & \Key{string} & --- & \ApiText{Data domain 'min,max' of t.}{Datenbereich 'min,max' von t.}}
  \SemioTableRow{\Key{arrangement} & \Key{string} & \Key{side-by-side} & \ApiText{Side-by-side|inset|stacked.}{Side-by-side|inset|stacked — Anordnung der beiden Felder.}}
  \SemioTableRow{\Key{zoom} & \Key{boolean} & \Key{True} & \ApiText{Draw the leader lines from the selection to the detail panel.}{Die Führungslinien von der Auswahl zum Detailfeld zeichnen.}}
  \SemioTableRow{\Key{width} & \Key{number} & --- & \ApiText{Frame size in millimetres.}{Rahmengröße in Millimetern.}}
  \SemioTableRow{\Key{height} & \Key{number} & --- & \ApiText{Frame size in millimetres.}{Rahmengröße in Millimetern.}}
  \SemioTableRow{\Key{detail-share} & \Key{string} & \Key{0.5} & \ApiText{Fraction of the frame taken by the detail panel.}{Anteil des Rahmens, den das Detailfeld einnimmt.}}
}

\subsection{\Key{state-selection}}

\SemioTableLong[text-size=8pt]{\Key{state-selection}}{0.20,0.13,0.15,0.52}{\ApiKey & \ApiType & \ApiDefault & \ApiMeaning}{%
  \SemioTableRow{\Key{nodes} & \Key{string} & --- & \ApiText{As in the 'flow' kernel.}{Wie im 'flow'-Kern.}}
  \SemioTableRow{\Key{edges} & \Key{string} & --- & \ApiText{As in the 'flow' kernel.}{Wie im 'flow'-Kern.}}
  \SemioTableRow{\Key{id} & \Key{string} & --- & \ApiText{As in the 'flow' kernel.}{Wie im 'flow'-Kern.}}
  \SemioTableRow{\Key{label} & \Key{string} & --- & \ApiText{As in the 'flow' kernel.}{Wie im 'flow'-Kern.}}
  \SemioTableRow{\Key{row} & \Key{string} & --- & \ApiText{As in the 'flow' kernel.}{Wie im 'flow'-Kern.}}
  \SemioTableRow{\Key{col} & \Key{string} & --- & \ApiText{As in the 'flow' kernel.}{Wie im 'flow'-Kern.}}
  \SemioTableRow{\Key{selected} & \Key{string} & \Key{check} & \ApiText{Id of the selected node.}{Kennung des ausgewählten Knotens.}}
  \SemioTableRow{\Key{mode} & \Key{string} & \Key{highlight} & \ApiText{Highlight|neighbourhood|callout|filtered.}{Highlight|neighbourhood|callout|filtered — Art der Hervorhebung.}}
  \SemioTableRow{\Key{dim} & \Key{number} & \Key{0.25} & \ApiText{Opacity of everything outside the selection.}{Deckkraft von allem außerhalb der Auswahl.}}
  \SemioTableRow{\Key{callout} & \Key{string} & --- & \ApiText{Text of the hover callout, empty uses the selected node label.}{Text der Sprechblase; leer nutzt die Beschriftung des ausgewählten Knotens.}}
  \SemioTableRow{\Key{routing} & \Key{string} & --- & \ApiText{How the connectors are routed between nodes.}{Wie die Verbinder zwischen den Knoten geführt werden.}}
  \SemioTableRow{\Key{width} & \Key{number} & --- & \ApiText{Frame size in millimetres.}{Rahmengröße in Millimetern.}}
  \SemioTableRow{\Key{height} & \Key{number} & --- & \ApiText{Frame size in millimetres.}{Rahmengröße in Millimetern.}}
}

\subsection{\Key{state-sequence}}

\SemioTableLong[text-size=8pt]{\Key{state-sequence}}{0.20,0.13,0.15,0.52}{\ApiKey & \ApiType & \ApiDefault & \ApiMeaning}{%
  \SemioTableRow{\Key{label} & \Key{string} & --- & \ApiText{As in the 'flow' kernel.}{Wie im 'flow'-Kern.}}
  \SemioTableRow{\Key{tone} & \Key{string} & --- & \ApiText{Column holding the colour tone of a row.}{Spalte mit dem Farbton einer Zeile.}}
  \SemioTableRow{\Key{frames} & \Key{integer} & \Key{0} & \ApiText{Number of frames drawn; 0 uses the table length.}{Anzahl der gezeichneten Bilder; 0 nutzt die Tabellenlänge.}}
  \SemioTableRow{\Key{columns} & \Key{integer} & \Key{4} & \ApiText{Frames per row.}{Bilder je Zeile.}}
  \SemioTableRow{\Key{depiction} & \Key{string} & \Key{bar} & \ApiText{Bar|position|radius|path  what the frame value drives.}{Bar|position|radius|path — was der Bildwert steuert.}}
  \SemioTableRow{\Key{arrows} & \Key{boolean} & \Key{True} & \ApiText{Draw the transition arrow between frames.}{Den Übergangspfeil zwischen den Bildern zeichnen.}}
  \SemioTableRow{\Key{ghost} & \Key{boolean} & \Key{False} & \ApiText{Draw the previous frame as a faint ghost.}{Das vorige Bild als blasses Nachbild zeichnen.}}
  \SemioTableRow{\Key{width} & \Key{number} & --- & \ApiText{Frame size in millimetres.}{Rahmengröße in Millimetern.}}
  \SemioTableRow{\Key{height} & \Key{number} & --- & \ApiText{Frame size in millimetres.}{Rahmengröße in Millimetern.}}
}

\subsection{\Key{swimlane}}

\SemioTableLong[text-size=8pt]{\Key{swimlane}}{0.20,0.13,0.15,0.52}{\ApiKey & \ApiType & \ApiDefault & \ApiMeaning}{%
  \SemioTableRow{\Key{arrow} & \Key{string} & --- & \ApiText{Arrow head drawn at the end of a connector.}{Pfeilspitze am Ende eines Verbinders.}}
  \SemioTableRow{\Key{default-shape} & \Key{string} & --- & \ApiText{Node shape used when a row carries none.}{Knotenform, wenn eine Zeile keine mitbringt.}}
  \SemioTableRow{\Key{groups} & \Key{string} & --- & \ApiText{Group boxes drawn behind the members.}{Gruppenrahmen hinter den Mitgliedern.}}
  \SemioTableRow{\Key{lanes} & \Key{string} & --- & \ApiText{Lane names, in drawing order.}{Bahnnamen in Zeichenreihenfolge.}}
  \SemioTableRow{\Key{routing} & \Key{string} & --- & \ApiText{How the connectors are routed between nodes.}{Wie die Verbinder zwischen den Knoten geführt werden.}}
}

\subsection{\Key{uml-class}}

\SemioTableLong[text-size=8pt]{\Key{uml-class}}{0.20,0.13,0.15,0.52}{\ApiKey & \ApiType & \ApiDefault & \ApiMeaning}{%
  \SemioTableRow{\Key{members} & \Key{string} & --- & \ApiText{Member table (default demo-uml-class-members), columns owner/section/text.}{Mitgliedertabelle (Vorgabe demo-uml-class-members), Spalten owner/section/text.}}
  \SemioTableRow{\Key{edges} & \Key{string} & --- & \ApiText{Relation table (default demo-uml-class-edges), columns from/to/kind/label.}{Beziehungstabelle (Vorgabe demo-uml-class-edges), Spalten from/to/kind/label.}}
  \SemioTableRow{\Key{id} & \Key{string} & --- & \ApiText{Classifier column names.}{Spaltennamen der Klassifizierer.}}
  \SemioTableRow{\Key{label} & \Key{string} & --- & \ApiText{Classifier column names.}{Spaltennamen der Klassifizierer.}}
  \SemioTableRow{\Key{row} & \Key{string} & --- & \ApiText{Classifier column names.}{Spaltennamen der Klassifizierer.}}
  \SemioTableRow{\Key{col} & \Key{string} & --- & \ApiText{Classifier column names.}{Spaltennamen der Klassifizierer.}}
  \SemioTableRow{\Key{tone} & \Key{string} & --- & \ApiText{Column holding the colour tone of a row.}{Spalte mit dem Farbton einer Zeile.}}
  \SemioTableRow{\Key{stereotype} & \Key{string} & --- & \ApiText{Classifier column names.}{Spaltennamen der Klassifizierer.}}
  \SemioTableRow{\Key{owner} & \Key{string} & --- & \ApiText{Member column names.}{Spaltennamen der Mitglieder.}}
  \SemioTableRow{\Key{section} & \Key{string} & --- & \ApiText{Member column names.}{Spaltennamen der Mitglieder.}}
  \SemioTableRow{\Key{text} & \Key{string} & --- & \ApiText{Member column names.}{Spaltennamen der Mitglieder.}}
  \SemioTableRow{\Key{sections} & \Key{string} & --- & \ApiText{Compartment keys in drawing order.}{Schlüssel der Fächer in Zeichenreihenfolge.}}
  \SemioTableRow{\Key{shape} & \Key{string} & \Key{process} & \ApiText{Classifier outline shape.}{Umrissform eines Klassifizierers.}}
  \SemioTableRow{\Key{routing} & \Key{string} & \Key{straight} & \ApiText{Orthogonal|straight|curved.}{Orthogonal|straight|curved — Führung der Verbinder.}}
  \SemioTableRow{\Key{groups} & \Key{boolean} & --- & \ApiText{Group boxes drawn behind the members.}{Gruppenrahmen hinter den Mitgliedern.}}
  \SemioTableRow{\Key{group} & \Key{string} & --- & \ApiText{Column that groups the nodes.}{Spalte, die die Knoten gruppiert.}}
  \SemioTableRow{\Key{width} & \Key{number} & --- & \ApiText{Frame size in millimetres.}{Rahmengröße in Millimetern.}}
  \SemioTableRow{\Key{height} & \Key{number} & --- & \ApiText{Frame size in millimetres.}{Rahmengröße in Millimetern.}}
  \SemioTableRow{\Key{col-gap} & \Key{string} & --- & \ApiText{Millimetre gap between grid columns.}{Millimeterabstand zwischen den Gitterspalten.}}
  \SemioTableRow{\Key{title-style} & \Key{string} & \Key{plain} & \ApiText{Plain|underline|stereotype.}{Plain|underline|stereotype — Schriftauszeichnung.}}
}

\subsection{\Key{uml-sequence}}

\SemioTableLong[text-size=8pt]{\Key{uml-sequence}}{0.20,0.13,0.15,0.52}{\ApiKey & \ApiType & \ApiDefault & \ApiMeaning}{%
  \SemioTableRow{\Key{messages} & \Key{string} & --- & \ApiText{Message table (default demo-uml-sequence-messages), columns from/to/label/kind/activate.}{Nachrichtentabelle (Vorgabe demo-uml-sequence-messages), Spalten from/to/label/kind/activate.}}
  \SemioTableRow{\Key{fragments} & \Key{string} & --- & \ApiText{Fragment table (default empty), columns label/first/last.}{Fragmenttabelle (Vorgabe leer), Spalten label/first/last.}}
  \SemioTableRow{\Key{id} & \Key{string} & --- & \ApiText{Classifier column names.}{Spaltennamen der Klassifizierer.}}
  \SemioTableRow{\Key{label} & \Key{string} & --- & \ApiText{Classifier column names.}{Spaltennamen der Klassifizierer.}}
  \SemioTableRow{\Key{shape} & \Key{string} & --- & \ApiText{Classifier outline shape.}{Umrissform eines Klassifizierers.}}
  \SemioTableRow{\Key{kind} & \Key{string} & \Key{kind} & \ApiText{Message kind column; values sync|async|return|create|destroy|self.}{Spalte der Nachrichtenart; Werte sync|async|return|create|destroy|self.}}
  \SemioTableRow{\Key{activate} & \Key{string} & \Key{activate} & \ApiText{Activation column; '1' opens a bar on the target, '0' closes it.}{Aktivierungsspalte; '1' öffnet einen Balken am Ziel, '0' schließt ihn.}}
  \SemioTableRow{\Key{head} & \Key{number} & \Key{6} & \ApiText{Participant head height in millimetres.}{Kopfhöhe eines Beteiligten in Millimetern.}}
  \SemioTableRow{\Key{step} & \Key{number} & \Key{0} & \ApiText{Vertical distance between messages, '0' fits the frame.}{Senkrechter Abstand zwischen den Nachrichten, '0' passt sich dem Rahmen an.}}
  \SemioTableRow{\Key{activation} & \Key{number} & \Key{1.6} & \ApiText{Width of an activation bar.}{Breite eines Aktivierungsbalkens.}}
  \SemioTableRow{\Key{width} & \Key{number} & --- & \ApiText{Frame size in millimetres.}{Rahmengröße in Millimetern.}}
  \SemioTableRow{\Key{height} & \Key{number} & --- & \ApiText{Frame size in millimetres.}{Rahmengröße in Millimetern.}}
}

\subsection{\Key{uml-timing}}

\SemioTableLong[text-size=8pt]{\Key{uml-timing}}{0.20,0.13,0.15,0.52}{\ApiKey & \ApiType & \ApiDefault & \ApiMeaning}{%
  \SemioTableRow{\Key{segments} & \Key{string} & --- & \ApiText{Segment table (default demo-uml-timing-segments), columns owner/state/start/end.}{Segmenttabelle (Vorgabe demo-uml-timing-segments), Spalten owner/state/start/end.}}
  \SemioTableRow{\Key{states} & \Key{string} & --- & \ApiText{State levels bottom to top.}{Zustandsniveaus von unten nach oben.}}
  \SemioTableRow{\Key{domain} & \Key{string} & --- & \ApiText{Time domain 'min,max'.}{Zeitbereich 'min,max'.}}
  \SemioTableRow{\Key{width} & \Key{number} & --- & \ApiText{Frame size in millimetres.}{Rahmengröße in Millimetern.}}
  \SemioTableRow{\Key{height} & \Key{number} & --- & \ApiText{Frame size in millimetres.}{Rahmengröße in Millimetern.}}
}

\subsection{\Key{uml-usecase}}

\SemioTableLong[text-size=8pt]{\Key{uml-usecase}}{0.20,0.13,0.15,0.52}{\ApiKey & \ApiType & \ApiDefault & \ApiMeaning}{%
  \SemioTableRow{\Key{actors} & \Key{string} & --- & \ApiText{Actor table (default demo-uml-usecase-actors), columns id/label/side.}{Akteurtabelle (Vorgabe demo-uml-usecase-actors), Spalten id/label/side.}}
  \SemioTableRow{\Key{edges} & \Key{string} & --- & \ApiText{Association table (default demo-uml-usecase-edges), columns from/to/kind.}{Assoziationstabelle (Vorgabe demo-uml-usecase-edges), Spalten from/to/kind.}}
  \SemioTableRow{\Key{id} & \Key{string} & --- & \ApiText{Classifier column names.}{Spaltennamen der Klassifizierer.}}
  \SemioTableRow{\Key{label} & \Key{string} & --- & \ApiText{Classifier column names.}{Spaltennamen der Klassifizierer.}}
  \SemioTableRow{\Key{side} & \Key{string} & \Key{side} & \ApiText{Actor side column, values left|right.}{Spalte der Akteurseite, Werte left|right.}}
  \SemioTableRow{\Key{system} & \Key{string} & --- & \ApiText{System boundary caption, empty for none.}{Beschriftung der Systemgrenze, leer für keine.}}
  \SemioTableRow{\Key{width} & \Key{number} & --- & \ApiText{Frame size in millimetres.}{Rahmengröße in Millimetern.}}
  \SemioTableRow{\Key{height} & \Key{number} & --- & \ApiText{Frame size in millimetres.}{Rahmengröße in Millimetern.}}
}

\section{\ApiText{Scientific}{Wissenschaft}}

\ApiText{Fields, surfaces, signals, physics, chemistry, biology, engineering, mathematics, geometry and three-dimensional projection.}{Felder, Flächen, Signale, Physik, Chemie, Biologie, Ingenieurwesen, Mathematik, Geometrie und dreidimensionale Projektion.}

\subsection{\Key{biology}}

\ApiText{This family takes no options beyond \Key{variant} and \Key{data}.}{Diese Familie nimmt außer \Key{variant} und \Key{data} keine Optionen.}

\subsection{\Key{optimization}}

\ApiText{This family takes no options beyond \Key{variant} and \Key{data}.}{Diese Familie nimmt außer \Key{variant} und \Key{data} keine Optionen.}

\subsection{\Key{sci-3d}}

\SemioTableLong[text-size=8pt]{\Key{sci-3d}}{0.20,0.13,0.15,0.52}{\ApiKey & \ApiType & \ApiDefault & \ApiMeaning}{%
  \SemioTableRow{\Key{mode} & \Key{string} & --- & \ApiText{Which member of the family the preset renders.}{Welches Mitglied der Familie die Voreinstellung zeichnet.}}
  \SemioTableRow{\Key{expr} & \Key{string} & --- & \ApiText{The surface height, also the volume field of the voxel modes.}{Die Flächenhöhe, zugleich das Volumenfeld der Voxelmodi.}}
  \SemioTableRow{\Key{solid} & \Key{string} & --- & \ApiText{The polyhedron drawn in the solid modes.}{Das Polyeder der Körpermodi.}}
  \SemioTableRow{\Key{points} & \Key{string} & --- & \ApiText{Draw the individual data points.}{Die einzelnen Datenpunkte zeichnen.}}
  \SemioTableRow{\Key{edges} & \Key{string} & --- & \ApiText{The discrete geometry.}{Die diskrete Geometrie.}}
  \SemioTableRow{\Key{columns} & \Key{integer} & --- & \ApiText{Surface lattice.}{Gitter der Fläche.}}
  \SemioTableRow{\Key{rows} & \Key{integer} & --- & \ApiText{Surface lattice.}{Gitter der Fläche.}}
  \SemioTableRow{\Key{levels} & \Key{integer} & --- & \ApiText{Contour or isosurface levels.}{Höhenlinien- bzw. Isoflächenniveaus.}}
  \SemioTableRow{\Key{extent} & \Key{number} & --- & \ApiText{World half-extent of the sampled domain.}{Halbe Weltausdehnung des abgetasteten Bereichs.}}
  \SemioTableRow{\Key{explode} & \Key{number} & --- & \ApiText{Radial offset of the exploded assembly.}{Radialer Versatz der Explosionsdarstellung.}}
  \SemioTableRow{\Key{faces} & \Key{boolean} & --- & \ApiText{Fill facets; false leaves a wireframe.}{Facetten füllen; false lässt ein Drahtgitter stehen.}}
  \SemioTableRow{\Key{azimuth} & \Key{string} & --- & \ApiText{Camera or light azimuth in degrees.}{Azimut der Kamera bzw. des Lichts in Grad.}}
  \SemioTableRow{\Key{distance} & \Key{string} & --- & \ApiText{Camera distance from the scene centre.}{Kameraabstand zum Szenenmittelpunkt.}}
  \SemioTableRow{\Key{elevation} & \Key{string} & --- & \ApiText{Camera elevation above the horizon in degrees.}{Kamerahöhe über dem Horizont in Grad.}}
  \SemioTableRow{\Key{projection} & \Key{string} & --- & \ApiText{Map or camera projection used.}{Verwendete Karten- bzw. Kameraprojektion.}}
  \SemioTableRow{\Key{zoom} & \Key{string} & --- & \ApiText{Zoom factor of the camera.}{Zoomfaktor der Kamera.}}
}

\subsection{\Key{sci-abstract-diagram}}

\SemioTableLong[text-size=8pt]{\Key{sci-abstract-diagram}}{0.20,0.13,0.15,0.52}{\ApiKey & \ApiType & \ApiDefault & \ApiMeaning}{%
  \SemioTableRow{\Key{mode} & \Key{string} & --- & \ApiText{What the expressions mean.}{Was die Ausdrücke bedeuten.}}
  \SemioTableRow{\Key{nodes} & \Key{string} & --- & \ApiText{Name of the node table.}{Name der Knotentabelle.}}
  \SemioTableRow{\Key{edges} & \Key{string} & --- & \ApiText{Name of the edge table.}{Name der Kantentabelle.}}
  \SemioTableRow{\Key{nodeShape} & \Key{string} & --- & \ApiText{Decoration drawn behind each node label.}{Dekoration hinter jeder Knotenbeschriftung.}}
  \SemioTableRow{\Key{unit} & \Key{number} & --- & \ApiText{Millimetres per data unit of the node coordinates.}{Millimeter je Dateneinheit der Knotenkoordinaten.}}
}

\subsection{\Key{sci-approximation}}

\SemioTableLong[text-size=8pt]{\Key{sci-approximation}}{0.20,0.13,0.15,0.52}{\ApiKey & \ApiType & \ApiDefault & \ApiMeaning}{%
  \SemioTableRow{\Key{mode} & \Key{string} & --- & \ApiText{What the expressions mean.}{Was die Ausdrücke bedeuten.}}
  \SemioTableRow{\Key{expr} & \Key{string} & --- & \ApiText{The function, or the x component in parametric modes.}{Die Funktion oder die x-Komponente in parametrischen Modi.}}
  \SemioTableRow{\Key{order} & \Key{integer} & --- & \ApiText{Highest term.}{Höchstes Glied.}}
  \SemioTableRow{\Key{centre} & \Key{number} & --- & \ApiText{Taylor expansion point.}{Entwicklungspunkt der Taylorreihe.}}
  \SemioTableRow{\Key{period} & \Key{number} & --- & \ApiText{Fourier fundamental period.}{Grundperiode der Fourierreihe.}}
  \SemioTableRow{\Key{partials} & \Key{boolean} & --- & \ApiText{Draw every partial sum, not only the highest.}{Jede Partialsumme zeichnen, nicht nur die höchste.}}
  \SemioTableRow{\Key{samples} & \Key{integer} & --- & \ApiText{Uniform sample count.}{Anzahl gleichmäßiger Stützstellen.}}
  \SemioTableRow{\Key{domain} & \Key{string} & --- & \ApiText{Input domain 'min,max' of the scale or function.}{Eingabebereich 'min,max' der Skala bzw. Funktion.}}
  \SemioTableRow{\Key{range} & \Key{string} & --- & \ApiText{Output range 'min,max' of the scale or sweep.}{Ausgabebereich 'min,max' der Skala bzw. des Durchlaufs.}}
}

\subsection{\Key{sci-beam}}

\SemioTableLong[text-size=8pt]{\Key{sci-beam}}{0.20,0.13,0.15,0.52}{\ApiKey & \ApiType & \ApiDefault & \ApiMeaning}{%
  \SemioTableRow{\Key{mode} & \Key{string} & --- & \ApiText{Which symbol vocabulary is used.}{Welcher Symbolvorrat verwendet wird.}}
  \SemioTableRow{\Key{length} & \Key{number} & --- & \ApiText{Span in millimetres, also the length unit of the load positions.}{Spannweite in Millimetern, zugleich Längeneinheit der Lastpositionen.}}
  \SemioTableRow{\Key{supports} & \Key{string} & --- & \ApiText{Kinds pin, roller, fixed.}{Arten pin, roller, fixed.}}
  \SemioTableRow{\Key{loads} & \Key{string} & --- & \ApiText{Downward point loads.}{Nach unten gerichtete Einzellasten.}}
  \SemioTableRow{\Key{distributed} & \Key{string} & --- & \ApiText{Uniformly distributed loads.}{Gleichstreckenlasten.}}
  \SemioTableRow{\Key{samples} & \Key{integer} & --- & \ApiText{Number of samples the curve or field is evaluated at.}{Anzahl der Stützstellen, an denen Kurve oder Feld ausgewertet werden.}}
}

\subsection{\Key{sci-bracket}}

\SemioTableLong[text-size=8pt]{\Key{sci-bracket}}{0.20,0.13,0.15,0.52}{\ApiKey & \ApiType & \ApiDefault & \ApiMeaning}{%
  \SemioTableRow{\Key{rounds} & \Key{integer} & --- & \ApiText{Knockout rounds.}{Runden des K.-o.-Turniers.}}
  \SemioTableRow{\Key{entries} & \Key{string} & --- & \ApiText{The first-round competitors.}{Die Teilnehmer der ersten Runde.}}
  \SemioTableRow{\Key{results} & \Key{string} & --- & \ApiText{The winner label placed in each later slot.}{Die Siegerbeschriftung in jedem späteren Platz.}}
  \SemioTableRow{\Key{slotHeight} & \Key{number} & --- & \ApiText{Millimetre height of one bracket slot.}{Millimeterhöhe eines Turnierplatzes.}}
  \SemioTableRow{\Key{roundWidth} & \Key{number} & --- & \ApiText{Millimetre width of one tournament round.}{Millimeterbreite einer Turnierrunde.}}
}

\subsection{\Key{sci-braid}}

\SemioTableLong[text-size=8pt]{\Key{sci-braid}}{0.20,0.13,0.15,0.52}{\ApiKey & \ApiType & \ApiDefault & \ApiMeaning}{%
  \SemioTableRow{\Key{mode} & \Key{string} & --- & \ApiText{What the expressions mean.}{Was die Ausdrücke bedeuten.}}
  \SemioTableRow{\Key{strands} & \Key{integer} & --- & \ApiText{Number of strands of the braid.}{Anzahl der Stränge des Zopfs.}}
  \SemioTableRow{\Key{crossings} & \Key{string} & --- & \ApiText{The braid word, sign ±1.}{Das Zopfwort, Vorzeichen ±1.}}
  \SemioTableRow{\Key{stride} & \Key{number} & --- & \ApiText{Millimetres per crossing.}{Millimeter je Kreuzung.}}
  \SemioTableRow{\Key{gauge} & \Key{number} & --- & \ApiText{Millimetres between strands.}{Millimeter zwischen den Strängen.}}
}

\subsection{\Key{sci-circos}}

\SemioTableLong[text-size=8pt]{\Key{sci-circos}}{0.20,0.13,0.15,0.52}{\ApiKey & \ApiType & \ApiDefault & \ApiMeaning}{%
  \SemioTableRow{\Key{segments} & \Key{string} & --- & \ApiText{The chromosomes, drawn as arcs proportional to their length.}{Die Chromosomen, als längenproportionale Bögen gezeichnet.}}
  \SemioTableRow{\Key{tracks} & \Key{string} & --- & \ApiText{The ring data; track 1 is the outermost ring.}{Die Ringdaten; Spur 1 ist der äußerste Ring.}}
  \SemioTableRow{\Key{links} & \Key{string} & --- & \ApiText{The chords inside the ring.}{Die Sehnen im Inneren des Rings.}}
  \SemioTableRow{\Key{radius} & \Key{number} & --- & \ApiText{Ring radius in millimetres.}{Ringradius in Millimetern.}}
  \SemioTableRow{\Key{ringWidth} & \Key{number} & --- & \ApiText{Millimetre width of one circos ring.}{Millimeterbreite eines Circos-Rings.}}
  \SemioTableRow{\Key{gap} & \Key{number} & --- & \ApiText{The ring geometry in millimetres and degrees.}{Die Ringgeometrie in Millimetern und Grad.}}
  \SemioTableRow{\Key{showTracks} & \Key{boolean} & --- & \ApiText{Draw the data tracks of the rings.}{Die Datenspuren der Ringe zeichnen.}}
  \SemioTableRow{\Key{showLinks} & \Key{boolean} & --- & \ApiText{Draw the links inside the circos ring.}{Die Verbindungen im Circos-Ring zeichnen.}}
}

\subsection{\Key{sci-circuit}}

\SemioTableLong[text-size=8pt]{\Key{sci-circuit}}{0.20,0.13,0.15,0.52}{\ApiKey & \ApiType & \ApiDefault & \ApiMeaning}{%
  \SemioTableRow{\Key{mode} & \Key{string} & --- & \ApiText{Which symbol vocabulary is used.}{Welcher Symbolvorrat verwendet wird.}}
  \SemioTableRow{\Key{elements} & \Key{string} & --- & \ApiText{The netlist; coordinates in millimetres.}{Die Netzliste; Koordinaten in Millimetern.}}
  \SemioTableRow{\Key{symbolSize} & \Key{number} & --- & \ApiText{Body length of a two-terminal symbol in millimetres.}{Körperlänge eines zweipoligen Symbols in Millimetern.}}
  \SemioTableRow{\Key{labels} & \Key{boolean} & --- & \ApiText{Print element labels.}{Die Elementbeschriftungen setzen.}}
}

\subsection{\Key{sci-climate-series}}

\SemioTableLong[text-size=8pt]{\Key{sci-climate-series}}{0.20,0.13,0.15,0.52}{\ApiKey & \ApiType & \ApiDefault & \ApiMeaning}{%
  \SemioTableRow{\Key{mode} & \Key{string} & --- & \ApiText{The coordinate skew of the diagram.}{Die Koordinatenscherung des Diagramms.}}
  \SemioTableRow{\Key{values} & \Key{string} & --- & \ApiText{The anomaly or observation series.}{Die Anomalie- oder Beobachtungsreihe.}}
  \SemioTableRow{\Key{baseline} & \Key{number} & --- & \ApiText{The reference the anomaly is measured against.}{Die Referenz, gegen die die Anomalie gemessen wird.}}
  \SemioTableRow{\Key{panels} & \Key{string} & --- & \ApiText{Where series is a '+'-joined 'time:value' list, for a meteogram.}{Wobei series eine mit '+' verbundene 'time:value'-Liste ist, für ein Meteogramm.}}
}

\subsection{\Key{sci-climograph}}

\SemioTableLong[text-size=8pt]{\Key{sci-climograph}}{0.20,0.13,0.15,0.52}{\ApiKey & \ApiType & \ApiDefault & \ApiMeaning}{%
  \SemioTableRow{\Key{mode} & \Key{string} & --- & \ApiText{The coordinate skew of the diagram.}{Die Koordinatenscherung des Diagramms.}}
  \SemioTableRow{\Key{months} & \Key{string} & --- & \ApiText{The twelve monthly normals.}{Die zwölf Monatsmittel.}}
  \SemioTableRow{\Key{discharge} & \Key{string} & --- & \ApiText{The hydrograph series.}{Die Ganglinienreihe.}}
  \SemioTableRow{\Key{rainfall} & \Key{string} & --- & \ApiText{The hyetograph bars.}{Die Balken des Niederschlagsdiagramms.}}
  \SemioTableRow{\Key{tempScale} & \Key{number} & --- & \ApiText{Millimetres of precipitation per degree, the Walter–Lieth 2:1 convention.}{Millimeter Niederschlag je Grad, nach der Walter-Lieth-Konvention 2:1.}}
}

\subsection{\Key{sci-clinical-timeline}}

\SemioTableLong[text-size=8pt]{\Key{sci-clinical-timeline}}{0.20,0.13,0.15,0.52}{\ApiKey & \ApiType & \ApiDefault & \ApiMeaning}{%
  \SemioTableRow{\Key{events} & \Key{string} & --- & \ApiText{Kinds point, interval-start, interval-end, marker.}{Arten point, interval-start, interval-end, marker.}}
  \SemioTableRow{\Key{lanes} & \Key{string} & --- & \ApiText{The lane names, top to bottom.}{Die Bahnnamen von oben nach unten.}}
  \SemioTableRow{\Key{laneHeight} & \Key{number} & --- & \ApiText{Track height in millimetres.}{Spurhöhe in Millimetern.}}
}

\subsection{\Key{sci-complex}}

\SemioTableLong[text-size=8pt]{\Key{sci-complex}}{0.20,0.13,0.15,0.52}{\ApiKey & \ApiType & \ApiDefault & \ApiMeaning}{%
  \SemioTableRow{\Key{mode} & \Key{string} & --- & \ApiText{What the expressions mean.}{Was die Ausdrücke bedeuten.}}
  \SemioTableRow{\Key{values} & \Key{string} & --- & \ApiText{The numbers of an Argand diagram.}{Die Zahlen eines Argand-Diagramms.}}
  \SemioTableRow{\Key{exprRe} & \Key{string} & --- & \ApiText{The mapped function for the other modes.}{Die abgebildete Funktion der übrigen Modi.}}
  \SemioTableRow{\Key{exprIm} & \Key{string} & --- & \ApiText{The mapped function for the other modes.}{Die abgebildete Funktion der übrigen Modi.}}
  \SemioTableRow{\Key{columns} & \Key{integer} & --- & \ApiText{Implicit-curve lattice.}{Gitter der impliziten Kurve.}}
  \SemioTableRow{\Key{rows} & \Key{integer} & --- & \ApiText{Implicit-curve lattice.}{Gitter der impliziten Kurve.}}
  \SemioTableRow{\Key{sheets} & \Key{integer} & --- & \ApiText{Riemann-surface sheets.}{Blätter der Riemannschen Fläche.}}
}

\subsection{\Key{sci-conic}}

\SemioTableLong[text-size=8pt]{\Key{sci-conic}}{0.20,0.13,0.15,0.52}{\ApiKey & \ApiType & \ApiDefault & \ApiMeaning}{%
  \SemioTableRow{\Key{kind} & \Key{string} & --- & \ApiText{Which conic.}{Welcher Kegelschnitt.}}
  \SemioTableRow{\Key{a} & \Key{number} & --- & \ApiText{The semi-axes, or the focal parameter of a parabola.}{Die Halbachsen oder der Halbparameter einer Parabel.}}
  \SemioTableRow{\Key{b} & \Key{number} & --- & \ApiText{The semi-axes, or the focal parameter of a parabola.}{Die Halbachsen oder der Halbparameter einer Parabel.}}
  \SemioTableRow{\Key{foci} & \Key{boolean} & --- & \ApiText{Mark the foci of the conic.}{Die Brennpunkte des Kegelschnitts markieren.}}
  \SemioTableRow{\Key{directrix} & \Key{boolean} & --- & \ApiText{Draw the directrix of the conic.}{Die Leitlinie des Kegelschnitts zeichnen.}}
  \SemioTableRow{\Key{asymptotes} & \Key{boolean} & --- & \ApiText{Construction elements to draw.}{Zu zeichnende Konstruktionselemente.}}
  \SemioTableRow{\Key{samples} & \Key{integer} & --- & \ApiText{Number of samples the curve or field is evaluated at.}{Anzahl der Stützstellen, an denen Kurve oder Feld ausgewertet werden.}}
}

\subsection{\Key{sci-constellation}}

\SemioTableLong[text-size=8pt]{\Key{sci-constellation}}{0.20,0.13,0.15,0.52}{\ApiKey & \ApiType & \ApiDefault & \ApiMeaning}{%
  \SemioTableRow{\Key{mode} & \Key{string} & --- & \ApiText{Which symbol vocabulary is used.}{Welcher Symbolvorrat verwendet wird.}}
  \SemioTableRow{\Key{order} & \Key{integer} & --- & \ApiText{The QAM order; the lattice is its square root per axis.}{Die QAM-Ordnung; das Gitter ist ihre Quadratwurzel je Achse.}}
  \SemioTableRow{\Key{noise} & \Key{number} & --- & \ApiText{The standard deviation of the seeded symbol noise.}{Die Standardabweichung des deterministisch erzeugten Symbolrauschens.}}
  \SemioTableRow{\Key{symbols} & \Key{integer} & --- & \ApiText{Symbol vocabulary used on the map.}{Symbolvorrat, der auf der Karte verwendet wird.}}
  \SemioTableRow{\Key{span} & \Key{integer} & --- & \ApiText{Horizontal span of the element in model units.}{Waagerechte Spannweite des Elements in Modelleinheiten.}}
  \SemioTableRow{\Key{samples} & \Key{integer} & --- & \ApiText{Number of samples the curve or field is evaluated at.}{Anzahl der Stützstellen, an denen Kurve oder Feld ausgewertet werden.}}
}

\subsection{\Key{sci-construction}}

\SemioTableLong[text-size=8pt]{\Key{sci-construction}}{0.20,0.13,0.15,0.52}{\ApiKey & \ApiType & \ApiDefault & \ApiMeaning}{%
  \SemioTableRow{\Key{mode} & \Key{string} & --- & \ApiText{Which member of the family the preset renders.}{Welches Mitglied der Familie die Voreinstellung zeichnet.}}
  \SemioTableRow{\Key{points} & \Key{string} & --- & \ApiText{The vertices, in millimetres of the figure.}{Die Eckpunkte in Millimetern der Abbildung.}}
  \SemioTableRow{\Key{centre} & \Key{string} & --- & \ApiText{Circle centre.}{Mittelpunkt des Kreises.}}
  \SemioTableRow{\Key{through} & \Key{string} & --- & \ApiText{The external point of a tangent or secant.}{Der äußere Punkt einer Tangente oder Sekante.}}
  \SemioTableRow{\Key{radius} & \Key{number} & --- & \ApiText{Circle radius for the circle modes.}{Kreisradius der Kreismodi.}}
  \SemioTableRow{\Key{marks} & \Key{boolean} & --- & \ApiText{Draw angle arcs and side ticks.}{Winkelbögen und Seitenmarken zeichnen.}}
  \SemioTableRow{\Key{labels} & \Key{boolean} & --- & \ApiText{Print vertex labels.}{Die Knotenbeschriftungen setzen.}}
}

\subsection{\Key{sci-contact-map}}

\SemioTableLong[text-size=8pt]{\Key{sci-contact-map}}{0.20,0.13,0.15,0.52}{\ApiKey & \ApiType & \ApiDefault & \ApiMeaning}{%
  \SemioTableRow{\Key{bins} & \Key{integer} & --- & \ApiText{Number of histogram bins.}{Anzahl der Histogrammklassen.}}
  \SemioTableRow{\Key{expr} & \Key{string} & --- & \ApiText{The contact frequency.}{Die Kontakthäufigkeit.}}
  \SemioTableRow{\Key{cell} & \Key{number} & --- & \ApiText{Cell size in millimetres.}{Zellgröße in Millimetern.}}
  \SemioTableRow{\Key{logScale} & \Key{boolean} & --- & \ApiText{Plot the estimates on a log axis.}{Die Schätzwerte auf einer logarithmischen Achse zeichnen.}}
}

\subsection{\Key{sci-control}}

\SemioTableLong[text-size=8pt]{\Key{sci-control}}{0.20,0.13,0.15,0.52}{\ApiKey & \ApiType & \ApiDefault & \ApiMeaning}{%
  \SemioTableRow{\Key{mode} & \Key{string} & --- & \ApiText{Which symbol vocabulary is used.}{Welcher Symbolvorrat verwendet wird.}}
  \SemioTableRow{\Key{numerator} & \Key{string} & --- & \ApiText{Column holding the numerator of the ratio.}{Spalte mit dem Zähler des Verhältnisses.}}
  \SemioTableRow{\Key{denominator} & \Key{string} & --- & \ApiText{Column holding the denominator of the ratio.}{Spalte mit dem Nenner des Verhältnisses.}}
  \SemioTableRow{\Key{decades} & \Key{string} & --- & \ApiText{The log frequency sweep.}{Der logarithmische Frequenzdurchlauf.}}
  \SemioTableRow{\Key{samples} & \Key{integer} & --- & \ApiText{Number of samples the curve or field is evaluated at.}{Anzahl der Stützstellen, an denen Kurve oder Feld ausgewertet werden.}}
  \SemioTableRow{\Key{gains} & \Key{integer} & --- & \ApiText{Root-locus gain steps.}{Verstärkungsschritte der Wurzelortskurve.}}
  \SemioTableRow{\Key{gainMax} & \Key{number} & --- & \ApiText{The largest gain traced.}{Die größte verfolgte Verstärkung.}}
  \SemioTableRow{\Key{steps} & \Key{integer} & --- & \ApiText{Number of integration or iteration steps.}{Anzahl der Integrations- bzw. Iterationsschritte.}}
  \SemioTableRow{\Key{step} & \Key{number} & --- & \ApiText{Spacing between consecutive steps.}{Abstand zwischen aufeinanderfolgenden Schritten.}}
}

\subsection{\Key{sci-crystal}}

\SemioTableLong[text-size=8pt]{\Key{sci-crystal}}{0.20,0.13,0.15,0.52}{\ApiKey & \ApiType & \ApiDefault & \ApiMeaning}{%
  \SemioTableRow{\Key{lattice} & \Key{string} & --- & \ApiText{Which sites are occupied.}{Welche Plätze besetzt sind.}}
  \SemioTableRow{\Key{cell} & \Key{number} & --- & \ApiText{Edge length in millimetres.}{Kantenlänge in Millimetern.}}
  \SemioTableRow{\Key{azimuth} & \Key{number} & --- & \ApiText{The isometric camera.}{Die isometrische Kamera.}}
  \SemioTableRow{\Key{elevation} & \Key{number} & --- & \ApiText{The isometric camera.}{Die isometrische Kamera.}}
  \SemioTableRow{\Key{radius} & \Key{number} & --- & \ApiText{Atom radius.}{Radius eines Atoms.}}
  \SemioTableRow{\Key{bonds} & \Key{boolean} & --- & \ApiText{Order 1–3, style plain, wedge, dash, aromatic.}{Ordnung 1–3, Stil plain, wedge, dash, aromatic.}}
}

\subsection{\Key{sci-density}}

\SemioTableLong[text-size=8pt]{\Key{sci-density}}{0.20,0.13,0.15,0.52}{\ApiKey & \ApiType & \ApiDefault & \ApiMeaning}{%
  \SemioTableRow{\Key{mode} & \Key{string} & --- & \ApiText{What the expressions mean.}{Was die Ausdrücke bedeuten.}}
  \SemioTableRow{\Key{points} & \Key{string} & --- & \ApiText{The candidate set a Pareto frontier is taken from.}{Die Kandidatenmenge, aus der eine Pareto-Front gebildet wird.}}
  \SemioTableRow{\Key{bandwidth} & \Key{number} & --- & \ApiText{Gaussian kernel width.}{Breite des Gauß-Kerns.}}
  \SemioTableRow{\Key{columns} & \Key{integer} & --- & \ApiText{Implicit-curve lattice.}{Gitter der impliziten Kurve.}}
  \SemioTableRow{\Key{rows} & \Key{integer} & --- & \ApiText{Implicit-curve lattice.}{Gitter der impliziten Kurve.}}
  \SemioTableRow{\Key{levels} & \Key{integer} & --- & \ApiText{Objective contour count.}{Anzahl der Höhenlinien der Zielfunktion.}}
  \SemioTableRow{\Key{at} & \Key{number} & --- & \ApiText{The conditioning value of x in 'conditional' mode.}{Der Bedingungswert von x im Modus 'conditional'.}}
}

\subsection{\Key{sci-dimensioned}}

\SemioTableLong[text-size=8pt]{\Key{sci-dimensioned}}{0.20,0.13,0.15,0.52}{\ApiKey & \ApiType & \ApiDefault & \ApiMeaning}{%
  \SemioTableRow{\Key{outline} & \Key{string} & --- & \ApiText{The part outline, closed automatically.}{Der Umriss des Teils, automatisch geschlossen.}}
  \SemioTableRow{\Key{dimensions} & \Key{string} & --- & \ApiText{Witness lines and their dimension text.}{Maßhilfslinien und ihr Maßtext.}}
  \SemioTableRow{\Key{hatch} & \Key{boolean} & --- & \ApiText{Fill the outline with a section hatch.}{Den Umriss mit einer Schnittschraffur füllen.}}
}

\subsection{\Key{sci-distribution}}

\SemioTableLong[text-size=8pt]{\Key{sci-distribution}}{0.20,0.13,0.15,0.52}{\ApiKey & \ApiType & \ApiDefault & \ApiMeaning}{%
  \SemioTableRow{\Key{law} & \Key{string} & --- & \ApiText{Distribution law that is plotted.}{Verteilungsgesetz, das gezeichnet wird.}}
  \SemioTableRow{\Key{curve} & \Key{string} & --- & \ApiText{Which function of the law is drawn.}{Welche Funktion des Gesetzes gezeichnet wird.}}
  \SemioTableRow{\Key{p1} & \Key{number} & --- & \ApiText{The two law parameters.}{Die beiden Parameter des Gesetzes.}}
  \SemioTableRow{\Key{p2} & \Key{number} & --- & \ApiText{The two law parameters.}{Die beiden Parameter des Gesetzes.}}
  \SemioTableRow{\Key{discrete} & \Key{boolean} & --- & \ApiText{Draw stems instead of a curve.}{Stiele statt einer Kurve zeichnen.}}
  \SemioTableRow{\Key{samples} & \Key{integer} & --- & \ApiText{Uniform sample count.}{Anzahl gleichmäßiger Stützstellen.}}
  \SemioTableRow{\Key{shade} & \Key{string} & --- & \ApiText{A shaded sub-interval of the density.}{Ein schattiertes Teilintervall der Dichte.}}
  \SemioTableRow{\Key{domain} & \Key{string} & --- & \ApiText{Input domain 'min,max' of the scale or function.}{Eingabebereich 'min,max' der Skala bzw. Funktion.}}
  \SemioTableRow{\Key{range} & \Key{string} & --- & \ApiText{Output range 'min,max' of the scale or sweep.}{Ausgabebereich 'min,max' der Skala bzw. des Durchlaufs.}}
}

\subsection{\Key{sci-dose-response}}

\SemioTableLong[text-size=8pt]{\Key{sci-dose-response}}{0.20,0.13,0.15,0.52}{\ApiKey & \ApiType & \ApiDefault & \ApiMeaning}{%
  \SemioTableRow{\Key{bottom} & \Key{number} & --- & \ApiText{The four-parameter logistic curve.}{Die vierparametrige logistische Kurve.}}
  \SemioTableRow{\Key{top} & \Key{number} & --- & \ApiText{The four-parameter logistic curve.}{Die vierparametrige logistische Kurve.}}
  \SemioTableRow{\Key{ec50} & \Key{number} & --- & \ApiText{The four-parameter logistic curve.}{Die vierparametrige logistische Kurve.}}
  \SemioTableRow{\Key{hill} & \Key{number} & --- & \ApiText{The four-parameter logistic curve.}{Die vierparametrige logistische Kurve.}}
  \SemioTableRow{\Key{points} & \Key{string} & --- & \ApiText{The measured data, doses on a decadic log axis.}{Die Messdaten, Dosen auf einer dekadisch logarithmischen Achse.}}
  \SemioTableRow{\Key{decades} & \Key{string} & --- & \ApiText{The log dose sweep.}{Der logarithmische Dosisdurchlauf.}}
  \SemioTableRow{\Key{samples} & \Key{integer} & --- & \ApiText{Number of samples the curve or field is evaluated at.}{Anzahl der Stützstellen, an denen Kurve oder Feld ausgewertet werden.}}
  \SemioTableRow{\Key{markEC} & \Key{boolean} & --- & \ApiText{Mark the EC50 on the curve.}{Den EC50 auf der Kurve markieren.}}
}

\subsection{\Key{sci-dynamics}}

\SemioTableLong[text-size=8pt]{\Key{sci-dynamics}}{0.20,0.13,0.15,0.52}{\ApiKey & \ApiType & \ApiDefault & \ApiMeaning}{%
  \SemioTableRow{\Key{mode} & \Key{string} & --- & \ApiText{What the expressions mean.}{Was die Ausdrücke bedeuten.}}
  \SemioTableRow{\Key{u} & \Key{string} & --- & \ApiText{The vector field.}{Das Vektorfeld.}}
  \SemioTableRow{\Key{v} & \Key{string} & --- & \ApiText{The vector field.}{Das Vektorfeld.}}
  \SemioTableRow{\Key{columns} & \Key{integer} & --- & \ApiText{Implicit-curve lattice.}{Gitter der impliziten Kurve.}}
  \SemioTableRow{\Key{rows} & \Key{integer} & --- & \ApiText{Implicit-curve lattice.}{Gitter der impliziten Kurve.}}
  \SemioTableRow{\Key{seeds} & \Key{string} & --- & \ApiText{Trajectory start points.}{Startpunkte der Bahnen.}}
  \SemioTableRow{\Key{steps} & \Key{integer} & --- & \ApiText{Number of integration or iteration steps.}{Anzahl der Integrations- bzw. Iterationsschritte.}}
  \SemioTableRow{\Key{step} & \Key{number} & --- & \ApiText{Spacing between consecutive steps.}{Abstand zwischen aufeinanderfolgenden Schritten.}}
  \SemioTableRow{\Key{scale} & \Key{number} & --- & \ApiText{Glyph length factor.}{Längenfaktor des Glyphs.}}
  \SemioTableRow{\Key{normalize} & \Key{boolean} & --- & \ApiText{Unit-length direction glyphs.}{Richtungsglyphen von Einheitslänge.}}
}

\subsection{\Key{sci-electrochem}}

\SemioTableLong[text-size=8pt]{\Key{sci-electrochem}}{0.20,0.13,0.15,0.52}{\ApiKey & \ApiType & \ApiDefault & \ApiMeaning}{%
  \SemioTableRow{\Key{mode} & \Key{string} & --- & \ApiText{The coordinate system of the diagram.}{Das Koordinatensystem des Diagramms.}}
  \SemioTableRow{\Key{species} & \Key{string} & --- & \ApiText{Oxidation state, nE° value and species name for the Frost diagram.}{Oxidationsstufe, nE°-Wert und Speziesname für das Frost-Diagramm.}}
  \SemioTableRow{\Key{couples} & \Key{string} & --- & \ApiText{The Latimer chain.}{Die Latimer-Kette.}}
  \SemioTableRow{\Key{regions} & \Key{string} & --- & \ApiText{Region labels.}{Beschriftungen der Bereiche.}}
}

\subsection{\Key{sci-energy-profile}}

\SemioTableLong[text-size=8pt]{\Key{sci-energy-profile}}{0.20,0.13,0.15,0.52}{\ApiKey & \ApiType & \ApiDefault & \ApiMeaning}{%
  \SemioTableRow{\Key{states} & \Key{string} & --- & \ApiText{Minima and maxima along the reaction coordinate.}{Minima und Maxima entlang der Reaktionskoordinate.}}
  \SemioTableRow{\Key{smooth} & \Key{boolean} & --- & \ApiText{Interpolate the profile with a cosine ramp between consecutive states.}{Das Profil mit einer Kosinusrampe zwischen aufeinanderfolgenden Zuständen interpolieren.}}
  \SemioTableRow{\Key{labels} & \Key{boolean} & --- & \ApiText{Draw the captions.}{Die Beschriftungen zeichnen.}}
  \SemioTableRow{\Key{barriers} & \Key{boolean} & --- & \ApiText{Print state labels and the activation barriers.}{Die Zustandsbeschriftungen und die Aktivierungsbarrieren setzen.}}
}

\subsection{\Key{sci-feynman}}

\SemioTableLong[text-size=8pt]{\Key{sci-feynman}}{0.20,0.13,0.15,0.52}{\ApiKey & \ApiType & \ApiDefault & \ApiMeaning}{%
  \SemioTableRow{\Key{vertices} & \Key{string} & --- & \ApiText{Interaction points, in millimetres.}{Wechselwirkungspunkte in Millimetern.}}
  \SemioTableRow{\Key{propagators} & \Key{string} & --- & \ApiText{Styles fermion, antifermion, photon, gluon, scalar, ghost.}{Stile fermion, antifermion, photon, gluon, scalar, ghost.}}
  \SemioTableRow{\Key{wiggle} & \Key{number} & --- & \ApiText{Amplitude of the wavy propagator decorations.}{Amplitude der welligen Propagatordekorationen.}}
  \SemioTableRow{\Key{coils} & \Key{integer} & --- & \ApiText{The amplitude and count of the photon and gluon decorations.}{Amplitude und Anzahl der Photonen- und Gluondekorationen.}}
}

\subsection{\Key{sci-fieldlines}}

\SemioTableLong[text-size=8pt]{\Key{sci-fieldlines}}{0.20,0.13,0.15,0.52}{\ApiKey & \ApiType & \ApiDefault & \ApiMeaning}{%
  \SemioTableRow{\Key{mode} & \Key{string} & --- & \ApiText{The optical element traced.}{Das durchgerechnete optische Element.}}
  \SemioTableRow{\Key{charges} & \Key{string} & --- & \ApiText{The sources; q is the signed strength.}{Die Quellen; q ist die vorzeichenbehaftete Stärke.}}
  \SemioTableRow{\Key{lines} & \Key{integer} & --- & \ApiText{Field lines seeded per positive charge.}{Feldlinien je positiver Ladung.}}
  \SemioTableRow{\Key{steps} & \Key{integer} & --- & \ApiText{Number of integration or iteration steps.}{Anzahl der Integrations- bzw. Iterationsschritte.}}
  \SemioTableRow{\Key{step} & \Key{number} & --- & \ApiText{The integration budget of the computed trajectories.}{Das Integrationsbudget der berechneten Bahnen.}}
  \SemioTableRow{\Key{levels} & \Key{integer} & --- & \ApiText{The level ladder, energies in the plotted unit.}{Die Niveauleiter, Energien in der gezeichneten Einheit.}}
}

\subsection{\Key{sci-force}}

\SemioTableLong[text-size=8pt]{\Key{sci-force}}{0.20,0.13,0.15,0.52}{\ApiKey & \ApiType & \ApiDefault & \ApiMeaning}{%
  \SemioTableRow{\Key{mode} & \Key{string} & --- & \ApiText{The optical element traced.}{Das durchgerechnete optische Element.}}
  \SemioTableRow{\Key{forces} & \Key{string} & --- & \ApiText{Angle in degrees from the positive x axis.}{Winkel in Grad ab der positiven x-Achse.}}
  \SemioTableRow{\Key{frames} & \Key{string} & --- & \ApiText{The strobe positions of a motion diagram.}{Die Stroboskoppositionen eines Bewegungsdiagramms.}}
  \SemioTableRow{\Key{body} & \Key{string} & --- & \ApiText{The glyph the forces act on.}{Der Glyph, an dem die Kräfte angreifen.}}
  \SemioTableRow{\Key{bodySize} & \Key{number} & --- & \ApiText{Millimetre size of the body glyph the forces act on.}{Millimetergröße des Körperglyphs, an dem die Kräfte angreifen.}}
  \SemioTableRow{\Key{scale} & \Key{number} & --- & \ApiText{Millimetres per unit of force.}{Millimeter je Krafteinheit.}}
}

\subsection{\Key{sci-fractal}}

\SemioTableLong[text-size=8pt]{\Key{sci-fractal}}{0.20,0.13,0.15,0.52}{\ApiKey & \ApiType & \ApiDefault & \ApiMeaning}{%
  \SemioTableRow{\Key{kind} & \Key{string} & --- & \ApiText{Which conic.}{Welcher Kegelschnitt.}}
  \SemioTableRow{\Key{depth} & \Key{integer} & --- & \ApiText{Penrose substitution steps, or hyperbolic recursion depth.}{Penrose-Substitutionsschritte oder hyperbolische Rekursionstiefe.}}
  \SemioTableRow{\Key{size} & \Key{number} & --- & \ApiText{Tile size in millimetres.}{Kachelgröße in Millimetern.}}
  \SemioTableRow{\Key{axiom} & \Key{string} & --- & \ApiText{Start word the L-system rewrites.}{Startwort, das das L-System umschreibt.}}
  \SemioTableRow{\Key{rules} & \Key{string} & --- & \ApiText{Rewrite rules of the L-system.}{Ersetzungsregeln des L-Systems.}}
  \SemioTableRow{\Key{angle} & \Key{number} & --- & \ApiText{Rotation or reflection-axis angle in degrees.}{Winkel der Drehung bzw. der Spiegelachse in Grad.}}
}

\subsection{\Key{sci-function}}

\SemioTableLong[text-size=8pt]{\Key{sci-function}}{0.20,0.13,0.15,0.52}{\ApiKey & \ApiType & \ApiDefault & \ApiMeaning}{%
  \SemioTableRow{\Key{mode} & \Key{string} & --- & \ApiText{What the expressions mean.}{Was die Ausdrücke bedeuten.}}
  \SemioTableRow{\Key{expr} & \Key{string} & --- & \ApiText{The function, or the x component in parametric modes.}{Die Funktion oder die x-Komponente in parametrischen Modi.}}
  \SemioTableRow{\Key{exprY} & \Key{string} & --- & \ApiText{The y component in parametric and vector-valued modes.}{Die y-Komponente in parametrischen und vektorwertigen Modi.}}
  \SemioTableRow{\Key{pieces} & \Key{string} & --- & \ApiText{The segments of a piecewise function.}{Die Abschnitte einer stückweise definierten Funktion.}}
  \SemioTableRow{\Key{samples} & \Key{integer} & --- & \ApiText{Uniform sample count.}{Anzahl gleichmäßiger Stützstellen.}}
  \SemioTableRow{\Key{refine} & \Key{integer} & --- & \ApiText{Adaptive refinement passes of the sampled curve.}{Adaptive Verfeinerungsdurchläufe der abgetasteten Kurve.}}
  \SemioTableRow{\Key{tolerance} & \Key{number} & --- & \ApiText{Adaptive midpoint refinement budget.}{Budget der adaptiven Mittelpunktverfeinerung.}}
  \SemioTableRow{\Key{columns} & \Key{integer} & --- & \ApiText{Implicit-curve lattice.}{Gitter der impliziten Kurve.}}
  \SemioTableRow{\Key{rows} & \Key{integer} & --- & \ApiText{Implicit-curve lattice.}{Gitter der impliziten Kurve.}}
  \SemioTableRow{\Key{level} & \Key{number} & --- & \ApiText{The implicit level set f=level.}{Die implizite Niveaumenge f=level.}}
  \SemioTableRow{\Key{fill} & \Key{string} & --- & \ApiText{Shade between curve and the baseline.}{Zwischen Kurve und Grundlinie schattieren.}}
  \SemioTableRow{\Key{baseline} & \Key{number} & --- & \ApiText{Value the marks grow from.}{Wert, von dem aus die Marken wachsen.}}
  \SemioTableRow{\Key{glyphs} & \Key{integer} & --- & \ApiText{Number of arrow glyphs drawn along a vector-valued curve.}{Anzahl der Pfeilglyphen entlang einer vektorwertigen Kurve.}}
  \SemioTableRow{\Key{domain} & \Key{string} & --- & \ApiText{Input domain 'min,max' of the scale or function.}{Eingabebereich 'min,max' der Skala bzw. Funktion.}}
  \SemioTableRow{\Key{range} & \Key{string} & --- & \ApiText{Output range 'min,max' of the scale or sweep.}{Ausgabebereich 'min,max' der Skala bzw. des Durchlaufs.}}
}

\subsection{\Key{sci-genome-track}}

\SemioTableLong[text-size=8pt]{\Key{sci-genome-track}}{0.20,0.13,0.15,0.52}{\ApiKey & \ApiType & \ApiDefault & \ApiMeaning}{%
  \SemioTableRow{\Key{mode} & \Key{string} & --- & \ApiText{Which member of the family the preset renders.}{Welches Mitglied der Familie die Voreinstellung zeichnet.}}
  \SemioTableRow{\Key{features} & \Key{string} & --- & \ApiText{The differential-expression table.}{Die Tabelle der differenziellen Expression.}}
  \SemioTableRow{\Key{coverage} & \Key{string} & --- & \ApiText{The per-bin depth of the coverage track.}{Die Abdeckungstiefe je Klasse der Spur.}}
  \SemioTableRow{\Key{span} & \Key{string} & --- & \ApiText{The genomic window.}{Das Genomfenster.}}
  \SemioTableRow{\Key{laneHeight} & \Key{number} & --- & \ApiText{Track height in millimetres.}{Spurhöhe in Millimetern.}}
}

\subsection{\Key{sci-geo-section}}

\SemioTableLong[text-size=8pt]{\Key{sci-geo-section}}{0.20,0.13,0.15,0.52}{\ApiKey & \ApiType & \ApiDefault & \ApiMeaning}{%
  \SemioTableRow{\Key{mode} & \Key{string} & --- & \ApiText{The coordinate skew of the diagram.}{Die Koordinatenscherung des Diagramms.}}
  \SemioTableRow{\Key{layers} & \Key{string} & --- & \ApiText{The stratigraphic pile, top layer first.}{Die stratigraphische Abfolge, oberste Schicht zuerst.}}
  \SemioTableRow{\Key{horizons} & \Key{string} & --- & \ApiText{The layer boundaries of a cross-section, as l3fp over \textbackslash{}SemioX.}{Die Schichtgrenzen eines Querschnitts, als l3fp über \textbackslash{}SemioX.}}
  \SemioTableRow{\Key{profile} & \Key{string} & --- & \ApiText{The sounding, pressure in hPa.}{Die Radiosondierung, Druck in hPa.}}
  \SemioTableRow{\Key{traces} & \Key{integer} & --- & \ApiText{The seismic trace count.}{Die Anzahl der seismischen Spuren.}}
  \SemioTableRow{\Key{expr} & \Key{string} & --- & \ApiText{The scalar field.}{Das Skalarfeld.}}
  \SemioTableRow{\Key{columnWidth} & \Key{number} & --- & \ApiText{The width of a stratigraphic column in millimetres.}{Die Breite einer stratigraphischen Säule in Millimetern.}}
}

\subsection{\Key{sci-grid-diagram}}

\SemioTableLong[text-size=8pt]{\Key{sci-grid-diagram}}{0.20,0.13,0.15,0.52}{\ApiKey & \ApiType & \ApiDefault & \ApiMeaning}{%
  \SemioTableRow{\Key{mode} & \Key{string} & --- & \ApiText{What the expressions mean.}{Was die Ausdrücke bedeuten.}}
  \SemioTableRow{\Key{partition} & \Key{string} & --- & \ApiText{Young-diagram row lengths.}{Zeilenlängen des Young-Diagramms.}}
  \SemioTableRow{\Key{cells} & \Key{string} & --- & \ApiText{The filled entries of the other modes.}{Die gefüllten Einträge der übrigen Modi.}}
  \SemioTableRow{\Key{rowLabels} & \Key{string} & --- & \ApiText{Header text.}{Text der Kopfzeile.}}
  \SemioTableRow{\Key{columnLabels} & \Key{string} & --- & \ApiText{Header text.}{Text der Kopfzeile.}}
  \SemioTableRow{\Key{cell} & \Key{number} & --- & \ApiText{Cell size in millimetres.}{Zellgröße in Millimetern.}}
}

\subsection{\Key{sci-growth-chart}}

\SemioTableLong[text-size=8pt]{\Key{sci-growth-chart}}{0.20,0.13,0.15,0.52}{\ApiKey & \ApiType & \ApiDefault & \ApiMeaning}{%
  \SemioTableRow{\Key{percentiles} & \Key{string} & --- & \ApiText{Reference percentile curves, each given as 'label/a/b/c'.}{Referenzperzentilkurven, jeweils als 'label/a/b/c' angegeben.}}
  \SemioTableRow{\Key{subject} & \Key{string} & --- & \ApiText{The individual measurements plotted on the reference.}{Die Einzelmessungen, auf die Referenz gezeichnet.}}
  \SemioTableRow{\Key{samples} & \Key{integer} & --- & \ApiText{Number of samples the curve or field is evaluated at.}{Anzahl der Stützstellen, an denen Kurve oder Feld ausgewertet werden.}}
}

\subsection{\Key{sci-ideogram}}

\SemioTableLong[text-size=8pt]{\Key{sci-ideogram}}{0.20,0.13,0.15,0.52}{\ApiKey & \ApiType & \ApiDefault & \ApiMeaning}{%
  \SemioTableRow{\Key{chromosomes} & \Key{string} & --- & \ApiText{Lengths and centromere positions in the same unit.}{Längen und Zentromerpositionen in derselben Einheit.}}
  \SemioTableRow{\Key{bands} & \Key{string} & --- & \ApiText{Stains gneg, gpos25, gpos50, gpos75, gpos100, acen.}{Färbungen gneg, gpos25, gpos50, gpos75, gpos100, acen.}}
  \SemioTableRow{\Key{marks} & \Key{string} & --- & \ApiText{Loci to annotate.}{Zu beschriftende Loci.}}
  \SemioTableRow{\Key{columnGap} & \Key{number} & --- & \ApiText{Millimetre gap between ideogram columns.}{Millimeterabstand zwischen den Ideogrammspalten.}}
  \SemioTableRow{\Key{scale} & \Key{number} & --- & \ApiText{Millimetres between chromosomes and per length unit.}{Millimeter zwischen den Chromosomen und je Längeneinheit.}}
  \SemioTableRow{\Key{showBands} & \Key{boolean} & --- & \ApiText{Draw the chromosome bands.}{Die Chromosomenbanden zeichnen.}}
  \SemioTableRow{\Key{showMarks} & \Key{boolean} & --- & \ApiText{Draw the annotated loci.}{Die annotierten Loci zeichnen.}}
}

\subsection{\Key{sci-integral}}

\SemioTableLong[text-size=8pt]{\Key{sci-integral}}{0.20,0.13,0.15,0.52}{\ApiKey & \ApiType & \ApiDefault & \ApiMeaning}{%
  \SemioTableRow{\Key{expr} & \Key{string} & --- & \ApiText{The function, or the x component in parametric modes.}{Die Funktion oder die x-Komponente in parametrischen Modi.}}
  \SemioTableRow{\Key{from} & \Key{number} & --- & \ApiText{Column holding the start node of an edge.}{Spalte mit dem Startknoten einer Kante.}}
  \SemioTableRow{\Key{to} & \Key{number} & --- & \ApiText{Column holding the end node of an edge.}{Spalte mit dem Endknoten einer Kante.}}
  \SemioTableRow{\Key{bins} & \Key{integer} & --- & \ApiText{Riemann partition size.}{Größe der Riemannschen Zerlegung.}}
  \SemioTableRow{\Key{rule} & \Key{string} & --- & \ApiText{The sampling rule.}{Die Abtastregel.}}
  \SemioTableRow{\Key{bars} & \Key{boolean} & --- & \ApiText{Draw the partition rectangles as well as the shaded area.}{Neben der schattierten Fläche auch die Zerlegungsrechtecke zeichnen.}}
  \SemioTableRow{\Key{samples} & \Key{integer} & --- & \ApiText{Uniform sample count.}{Anzahl gleichmäßiger Stützstellen.}}
}

\subsection{\Key{sci-iteration}}

\SemioTableLong[text-size=8pt]{\Key{sci-iteration}}{0.20,0.13,0.15,0.52}{\ApiKey & \ApiType & \ApiDefault & \ApiMeaning}{%
  \SemioTableRow{\Key{mode} & \Key{string} & --- & \ApiText{What the expressions mean.}{Was die Ausdrücke bedeuten.}}
  \SemioTableRow{\Key{expr} & \Key{string} & --- & \ApiText{The function, or the x component in parametric modes.}{Die Funktion oder die x-Komponente in parametrischen Modi.}}
  \SemioTableRow{\Key{start} & \Key{number} & --- & \ApiText{The initial value.}{Der Anfangswert.}}
  \SemioTableRow{\Key{steps} & \Key{integer} & --- & \ApiText{Number of integration or iteration steps.}{Anzahl der Integrations- bzw. Iterationsschritte.}}
  \SemioTableRow{\Key{parameter} & \Key{number} & --- & \ApiText{The map parameter, available as \textbackslash{}SemioY.}{Der Parameter der Abbildung, verfügbar als \textbackslash{}SemioY.}}
  \SemioTableRow{\Key{from} & \Key{number} & --- & \ApiText{Column holding the start node of an edge.}{Spalte mit dem Startknoten einer Kante.}}
  \SemioTableRow{\Key{to} & \Key{number} & --- & \ApiText{Column holding the end node of an edge.}{Spalte mit dem Endknoten einer Kante.}}
  \SemioTableRow{\Key{columns} & \Key{integer} & --- & \ApiText{Implicit-curve lattice.}{Gitter der impliziten Kurve.}}
  \SemioTableRow{\Key{settle} & \Key{integer} & --- & \ApiText{Iterations discarded before the orbit is drawn.}{Iterationen, die vor dem Zeichnen des Orbits verworfen werden.}}
  \SemioTableRow{\Key{keep} & \Key{integer} & --- & \ApiText{Bifurcation transient and plotted tail.}{Einschwingphase der Bifurkation und gezeichneter Schwanz.}}
  \SemioTableRow{\Key{samples} & \Key{integer} & --- & \ApiText{Uniform sample count.}{Anzahl gleichmäßiger Stützstellen.}}
}

\subsection{\Key{sci-kinematics}}

\SemioTableLong[text-size=8pt]{\Key{sci-kinematics}}{0.20,0.13,0.15,0.52}{\ApiKey & \ApiType & \ApiDefault & \ApiMeaning}{%
  \SemioTableRow{\Key{mode} & \Key{string} & --- & \ApiText{The optical element traced.}{Das durchgerechnete optische Element.}}
  \SemioTableRow{\Key{expr} & \Key{string} & --- & \ApiText{The position law of the graph modes.}{Das Positionsgesetz der Graphmodi.}}
  \SemioTableRow{\Key{speed} & \Key{number} & --- & \ApiText{Initial speed of the motion.}{Anfangsgeschwindigkeit der Bewegung.}}
  \SemioTableRow{\Key{angle} & \Key{number} & --- & \ApiText{Angular position, in degrees or as the column that drives it.}{Winkelposition, in Grad oder als die Spalte, die sie steuert.}}
  \SemioTableRow{\Key{gravity} & \Key{number} & --- & \ApiText{The projectile launch.}{Der Abschuss des Wurfkörpers.}}
  \SemioTableRow{\Key{radius} & \Key{number} & --- & \ApiText{Sphere radius in millimetres.}{Kugelradius in Millimetern.}}
  \SemioTableRow{\Key{eccentricity} & \Key{number} & --- & \ApiText{The orbital element set.}{Der Satz der Bahnelemente.}}
  \SemioTableRow{\Key{steps} & \Key{integer} & --- & \ApiText{Number of integration or iteration steps.}{Anzahl der Integrations- bzw. Iterationsschritte.}}
  \SemioTableRow{\Key{step} & \Key{number} & --- & \ApiText{The integration budget of the computed trajectories.}{Das Integrationsbudget der berechneten Bahnen.}}
}

\subsection{\Key{sci-levels}}

\SemioTableLong[text-size=8pt]{\Key{sci-levels}}{0.20,0.13,0.15,0.52}{\ApiKey & \ApiType & \ApiDefault & \ApiMeaning}{%
  \SemioTableRow{\Key{mode} & \Key{string} & --- & \ApiText{The optical element traced.}{Das durchgerechnete optische Element.}}
  \SemioTableRow{\Key{levels} & \Key{string} & --- & \ApiText{The level ladder, energies in the plotted unit.}{Die Niveauleiter, Energien in der gezeichneten Einheit.}}
  \SemioTableRow{\Key{transitions} & \Key{string} & --- & \ApiText{Arrows between level indices.}{Pfeile zwischen den Niveauindizes.}}
  \SemioTableRow{\Key{bands} & \Key{string} & --- & \ApiText{The band edges of the 'band' mode.}{Die Bandgrenzen des Modus 'band'.}}
  \SemioTableRow{\Key{expr} & \Key{string} & --- & \ApiText{The position law of the graph modes.}{Das Positionsgesetz der Graphmodi.}}
}

\subsection{\Key{sci-lifecycle}}

\SemioTableLong[text-size=8pt]{\Key{sci-lifecycle}}{0.20,0.13,0.15,0.52}{\ApiKey & \ApiType & \ApiDefault & \ApiMeaning}{%
  \SemioTableRow{\Key{stages} & \Key{string} & --- & \ApiText{The stages, placed evenly around the ring.}{Die Stufen, gleichmäßig auf dem Ring verteilt.}}
  \SemioTableRow{\Key{radius} & \Key{number} & --- & \ApiText{Ring radius in millimetres.}{Ringradius in Millimetern.}}
  \SemioTableRow{\Key{arrows} & \Key{boolean} & --- & \ApiText{Draw the arcs between stages.}{Die Bögen zwischen den Stufen zeichnen.}}
}

\subsection{\Key{sci-manhattan}}

\SemioTableLong[text-size=8pt]{\Key{sci-manhattan}}{0.20,0.13,0.15,0.52}{\ApiKey & \ApiType & \ApiDefault & \ApiMeaning}{%
  \SemioTableRow{\Key{mode} & \Key{string} & --- & \ApiText{Which member of the family the preset renders.}{Welches Mitglied der Familie die Voreinstellung zeichnet.}}
  \SemioTableRow{\Key{points} & \Key{string} & --- & \ApiText{The measured data, doses on a decadic log axis.}{Die Messdaten, Dosen auf einer dekadisch logarithmischen Achse.}}
  \SemioTableRow{\Key{chromosomes} & \Key{string} & --- & \ApiText{Lengths and centromere positions in the same unit.}{Längen und Zentromerpositionen in derselben Einheit.}}
  \SemioTableRow{\Key{threshold} & \Key{number} & --- & \ApiText{The genome-wide significance line.}{Die genomweite Signifikanzlinie.}}
  \SemioTableRow{\Key{suggestive} & \Key{number} & --- & \ApiText{Suggestive significance line of the plot.}{Linie der suggestiven Signifikanz im Diagramm.}}
}

\subsection{\Key{sci-material}}

\SemioTableLong[text-size=8pt]{\Key{sci-material}}{0.20,0.13,0.15,0.52}{\ApiKey & \ApiType & \ApiDefault & \ApiMeaning}{%
  \SemioTableRow{\Key{curve} & \Key{string} & --- & \ApiText{The measured response.}{Die gemessene Antwort.}}
  \SemioTableRow{\Key{yield} & \Key{number} & --- & \ApiText{Strains to mark; 0 leaves the marker out.}{Zu markierende Dehnungen; 0 lässt die Marke weg.}}
  \SemioTableRow{\Key{ultimate} & \Key{number} & --- & \ApiText{Strains to mark; 0 leaves the marker out.}{Zu markierende Dehnungen; 0 lässt die Marke weg.}}
  \SemioTableRow{\Key{modulus} & \Key{boolean} & --- & \ApiText{Draw the elastic tangent through the origin.}{Die elastische Tangente durch den Ursprung zeichnen.}}
}

\subsection{\Key{sci-mesh}}

\SemioTableLong[text-size=8pt]{\Key{sci-mesh}}{0.20,0.13,0.15,0.52}{\ApiKey & \ApiType & \ApiDefault & \ApiMeaning}{%
  \SemioTableRow{\Key{mode} & \Key{string} & --- & \ApiText{Which member of the family the preset renders.}{Welches Mitglied der Familie die Voreinstellung zeichnet.}}
  \SemioTableRow{\Key{columns} & \Key{integer} & --- & \ApiText{Lattice extent of the regular tilings.}{Gitterausdehnung der regelmäßigen Parkettierungen.}}
  \SemioTableRow{\Key{rows} & \Key{integer} & --- & \ApiText{Lattice extent of the regular tilings.}{Gitterausdehnung der regelmäßigen Parkettierungen.}}
  \SemioTableRow{\Key{size} & \Key{number} & --- & \ApiText{Tile size in millimetres.}{Kachelgröße in Millimetern.}}
  \SemioTableRow{\Key{warp} & \Key{string} & --- & \ApiText{A displacement applied to every node, so a mesh can follow a shape.}{Eine Verschiebung, die auf jeden Knoten wirkt, damit ein Netz einer Form folgen kann.}}
  \SemioTableRow{\Key{fillCells} & \Key{boolean} & --- & \ApiText{Shade every cell by its own index.}{Jede Zelle nach ihrem eigenen Index einfärben.}}
}

\subsection{\Key{sci-mohr}}

\SemioTableLong[text-size=8pt]{\Key{sci-mohr}}{0.20,0.13,0.15,0.52}{\ApiKey & \ApiType & \ApiDefault & \ApiMeaning}{%
  \SemioTableRow{\Key{sigmaX} & \Key{number} & --- & \ApiText{The plane stress state.}{Der ebene Spannungszustand.}}
  \SemioTableRow{\Key{sigmaY} & \Key{number} & --- & \ApiText{The plane stress state.}{Der ebene Spannungszustand.}}
  \SemioTableRow{\Key{tauXY} & \Key{number} & --- & \ApiText{The plane stress state.}{Der ebene Spannungszustand.}}
  \SemioTableRow{\Key{principal} & \Key{boolean} & --- & \ApiText{Mark the principal stresses and the maximum shear.}{Die Hauptspannungen und die maximale Schubspannung markieren.}}
  \SemioTableRow{\Key{pole} & \Key{boolean} & --- & \ApiText{Draw the pole construction.}{Die Polkonstruktion zeichnen.}}
}

\subsection{\Key{sci-molecule}}

\SemioTableLong[text-size=8pt]{\Key{sci-molecule}}{0.20,0.13,0.15,0.52}{\ApiKey & \ApiType & \ApiDefault & \ApiMeaning}{%
  \SemioTableRow{\Key{mode} & \Key{string} & --- & \ApiText{The coordinate system of the diagram.}{Das Koordinatensystem des Diagramms.}}
  \SemioTableRow{\Key{atoms} & \Key{string} & --- & \ApiText{An empty element is an implicit carbon vertex.}{Ein leeres Element ist ein impliziter Kohlenstoffknoten.}}
  \SemioTableRow{\Key{bonds} & \Key{string} & --- & \ApiText{Order 1–3, style plain, wedge, dash, aromatic.}{Ordnung 1–3, Stil plain, wedge, dash, aromatic.}}
  \SemioTableRow{\Key{lonePairs} & \Key{string} & --- & \ApiText{The Lewis dots drawn around an atom.}{Die Lewis-Punkte um ein Atom.}}
  \SemioTableRow{\Key{arrow} & \Key{string} & --- & \ApiText{The reaction arrow of the reaction mode.}{Der Reaktionspfeil des Reaktionsmodus.}}
  \SemioTableRow{\Key{bondGap} & \Key{number} & --- & \ApiText{The offset of the second and third bond lines in millimetres.}{Der Versatz der zweiten und dritten Bindungslinie in Millimetern.}}
}

\subsection{\Key{sci-nomogram}}

\SemioTableLong[text-size=8pt]{\Key{sci-nomogram}}{0.20,0.13,0.15,0.52}{\ApiKey & \ApiType & \ApiDefault & \ApiMeaning}{%
  \SemioTableRow{\Key{scales} & \Key{string} & --- & \ApiText{The parallel scales, x in millimetres.}{Die parallelen Skalen, x in Millimetern.}}
  \SemioTableRow{\Key{isopleth} & \Key{string} & --- & \ApiText{The reading line drawn through the given scale values.}{Die Ablesegerade durch die angegebenen Skalenwerte.}}
}

\subsection{\Key{sci-optics}}

\SemioTableLong[text-size=8pt]{\Key{sci-optics}}{0.20,0.13,0.15,0.52}{\ApiKey & \ApiType & \ApiDefault & \ApiMeaning}{%
  \SemioTableRow{\Key{mode} & \Key{string} & --- & \ApiText{The optical element traced.}{Das durchgerechnete optische Element.}}
  \SemioTableRow{\Key{focal} & \Key{number} & --- & \ApiText{The focal length in millimetres, negative for a diverging element.}{Die Brennweite in Millimetern, negativ für ein zerstreuendes Element.}}
  \SemioTableRow{\Key{objectDistance} & \Key{number} & --- & \ApiText{Object distance from the optical element.}{Gegenstandsweite zum optischen Element.}}
  \SemioTableRow{\Key{objectHeight} & \Key{number} & --- & \ApiText{The object placement.}{Die Platzierung des Gegenstands.}}
  \SemioTableRow{\Key{rays} & \Key{boolean} & --- & \ApiText{Draw the principal rays that construct the image.}{Die Hauptstrahlen zeichnen, die das Bild konstruieren.}}
}

\subsection{\Key{sci-optimization}}

\SemioTableLong[text-size=8pt]{\Key{sci-optimization}}{0.20,0.13,0.15,0.52}{\ApiKey & \ApiType & \ApiDefault & \ApiMeaning}{%
  \SemioTableRow{\Key{mode} & \Key{string} & --- & \ApiText{What the expressions mean.}{Was die Ausdrücke bedeuten.}}
  \SemioTableRow{\Key{constraints} & \Key{string} & --- & \ApiText{The half-planes a·x + b·y <= c.}{Die Halbebenen a·x + b·y <= c.}}
  \SemioTableRow{\Key{objective} & \Key{string} & --- & \ApiText{The objective gradient.}{Der Gradient der Zielfunktion.}}
  \SemioTableRow{\Key{levels} & \Key{integer} & --- & \ApiText{Objective contour count.}{Anzahl der Höhenlinien der Zielfunktion.}}
  \SemioTableRow{\Key{points} & \Key{string} & --- & \ApiText{The candidate set a Pareto frontier is taken from.}{Die Kandidatenmenge, aus der eine Pareto-Front gebildet wird.}}
  \SemioTableRow{\Key{stations} & \Key{string} & --- & \ApiText{Draw the stations along the route.}{Die Stationen entlang der Route zeichnen.}}
  \SemioTableRow{\Key{routes} & \Key{string} & --- & \ApiText{The queueing network.}{Das Warteschlangennetz.}}
  \SemioTableRow{\Key{unit} & \Key{number} & --- & \ApiText{Millimetres per data unit of the node coordinates.}{Millimeter je Dateneinheit der Knotenkoordinaten.}}
  \SemioTableRow{\Key{domain} & \Key{string} & --- & \ApiText{Input domain 'min,max' of the scale or function.}{Eingabebereich 'min,max' der Skala bzw. Funktion.}}
  \SemioTableRow{\Key{range} & \Key{string} & --- & \ApiText{Output range 'min,max' of the scale or sweep.}{Ausgabebereich 'min,max' der Skala bzw. des Durchlaufs.}}
}

\subsection{\Key{sci-orbital}}

\SemioTableLong[text-size=8pt]{\Key{sci-orbital}}{0.20,0.13,0.15,0.52}{\ApiKey & \ApiType & \ApiDefault & \ApiMeaning}{%
  \SemioTableRow{\Key{mode} & \Key{string} & --- & \ApiText{The coordinate system of the diagram.}{Das Koordinatensystem des Diagramms.}}
  \SemioTableRow{\Key{levels} & \Key{string} & --- & \ApiText{Side -1 left fragment, 0 molecular, 1 right fragment.}{Seite -1 linkes Fragment, 0 molekular, 1 rechtes Fragment.}}
  \SemioTableRow{\Key{shells} & \Key{string} & --- & \ApiText{The configuration boxes.}{Die Konfigurationskästen.}}
  \SemioTableRow{\Key{lobes} & \Key{integer} & --- & \ApiText{Lobe count of the sketched orbital.}{Anzahl der Keulen des skizzierten Orbitals.}}
}

\subsection{\Key{sci-pathway}}

\SemioTableLong[text-size=8pt]{\Key{sci-pathway}}{0.20,0.13,0.15,0.52}{\ApiKey & \ApiType & \ApiDefault & \ApiMeaning}{%
  \SemioTableRow{\Key{mode} & \Key{string} & --- & \ApiText{Which member of the family the preset renders.}{Welches Mitglied der Familie die Voreinstellung zeichnet.}}
  \SemioTableRow{\Key{nodes} & \Key{string} & --- & \ApiText{Name of the node table.}{Name der Knotentabelle.}}
  \SemioTableRow{\Key{edges} & \Key{string} & --- & \ApiText{Name of the edge table.}{Name der Kantentabelle.}}
  \SemioTableRow{\Key{unit} & \Key{number} & --- & \ApiText{Millimetres per node coordinate unit.}{Millimeter je Einheit der Knotenkoordinaten.}}
  \SemioTableRow{\Key{nodeShape} & \Key{string} & --- & \ApiText{Shape of a pathway node.}{Form eines Pfadknotens.}}
}

\subsection{\Key{sci-pedigree}}

\SemioTableLong[text-size=8pt]{\Key{sci-pedigree}}{0.20,0.13,0.15,0.52}{\ApiKey & \ApiType & \ApiDefault & \ApiMeaning}{%
  \SemioTableRow{\Key{individuals} & \Key{string} & --- & \ApiText{Sex m or f, affected 0 or 1.}{Geschlecht m oder f, betroffen 0 oder 1.}}
  \SemioTableRow{\Key{relations} & \Key{string} & --- & \ApiText{Where children is a '+'-joined id list.}{Wobei children eine mit '+' verbundene Kennungsliste ist.}}
  \SemioTableRow{\Key{generationGap} & \Key{number} & --- & \ApiText{Millimetre distance between pedigree generations.}{Millimeterabstand zwischen den Generationen des Stammbaums.}}
  \SemioTableRow{\Key{siblingGap} & \Key{number} & --- & \ApiText{Millimetre spacing between siblings.}{Millimeterabstand zwischen Geschwistern.}}
}

\subsection{\Key{sci-phase-diagram}}

\SemioTableLong[text-size=8pt]{\Key{sci-phase-diagram}}{0.20,0.13,0.15,0.52}{\ApiKey & \ApiType & \ApiDefault & \ApiMeaning}{%
  \SemioTableRow{\Key{mode} & \Key{string} & --- & \ApiText{The coordinate system of the diagram.}{Das Koordinatensystem des Diagramms.}}
  \SemioTableRow{\Key{boundaries} & \Key{string} & --- & \ApiText{The phase boundaries, in the diagram's own units.}{Die Phasengrenzen in den eigenen Einheiten des Diagramms.}}
  \SemioTableRow{\Key{points} & \Key{string} & --- & \ApiText{Invariant points such as triple or eutectic points.}{Invariante Punkte wie Tripel- oder eutektische Punkte.}}
  \SemioTableRow{\Key{regions} & \Key{string} & --- & \ApiText{Region labels.}{Beschriftungen der Bereiche.}}
  \SemioTableRow{\Key{side} & \Key{number} & --- & \ApiText{Triangle side of the ternary mode.}{Dreiecksseite des ternären Modus.}}
}

\subsection{\Key{sci-pianoroll}}

\SemioTableLong[text-size=8pt]{\Key{sci-pianoroll}}{0.20,0.13,0.15,0.52}{\ApiKey & \ApiType & \ApiDefault & \ApiMeaning}{%
  \SemioTableRow{\Key{mode} & \Key{string} & --- & \ApiText{Which member of the family the preset renders.}{Welches Mitglied der Familie die Voreinstellung zeichnet.}}
  \SemioTableRow{\Key{notes} & \Key{string} & --- & \ApiText{The event table, pitch as a MIDI number.}{Die Ereignistabelle, Tonhöhe als MIDI-Nummer.}}
  \SemioTableRow{\Key{curve} & \Key{string} & --- & \ApiText{The tempo, dynamics or pitch curve of the envelope modes.}{Die Tempo-, Dynamik- oder Tonhöhenkurve der Hüllkurvenmodi.}}
  \SemioTableRow{\Key{lowPitch} & \Key{integer} & --- & \ApiText{The drawn pitch window.}{Das gezeichnete Tonhöhenfenster.}}
  \SemioTableRow{\Key{highPitch} & \Key{integer} & --- & \ApiText{The drawn pitch window.}{Das gezeichnete Tonhöhenfenster.}}
}

\subsection{\Key{sci-pitch}}

\SemioTableLong[text-size=8pt]{\Key{sci-pitch}}{0.20,0.13,0.15,0.52}{\ApiKey & \ApiType & \ApiDefault & \ApiMeaning}{%
  \SemioTableRow{\Key{mode} & \Key{string} & --- & \ApiText{Which member of the family the preset renders.}{Welches Mitglied der Familie die Voreinstellung zeichnet.}}
  \SemioTableRow{\Key{surface} & \Key{string} & --- & \ApiText{The template drawn behind the data.}{Die Vorlage, hinter den Daten gezeichnet.}}
  \SemioTableRow{\Key{events} & \Key{string} & --- & \ApiText{Shots or touches, coordinates normalised to 0..1 of the surface.}{Schüsse oder Ballkontakte, Koordinaten auf 0..1 der Spielfläche normiert.}}
  \SemioTableRow{\Key{nodes} & \Key{string} & --- & \ApiText{Name of the node table.}{Name der Knotentabelle.}}
  \SemioTableRow{\Key{edges} & \Key{string} & --- & \ApiText{The pass network.}{Das Passnetzwerk.}}
  \SemioTableRow{\Key{columns} & \Key{integer} & --- & \ApiText{Lattice extent of the regular tilings.}{Gitterausdehnung der regelmäßigen Parkettierungen.}}
  \SemioTableRow{\Key{rows} & \Key{integer} & --- & \ApiText{Lattice extent of the regular tilings.}{Gitterausdehnung der regelmäßigen Parkettierungen.}}
}

\subsection{\Key{sci-poincare}}

\SemioTableLong[text-size=8pt]{\Key{sci-poincare}}{0.20,0.13,0.15,0.52}{\ApiKey & \ApiType & \ApiDefault & \ApiMeaning}{%
  \SemioTableRow{\Key{states} & \Key{string} & --- & \ApiText{Stokes vectors, normalised before projection.}{Stokes-Vektoren, vor der Projektion normiert.}}
  \SemioTableRow{\Key{radius} & \Key{number} & --- & \ApiText{Sphere radius in millimetres.}{Kugelradius in Millimetern.}}
  \SemioTableRow{\Key{azimuth} & \Key{number} & --- & \ApiText{Camera or light azimuth in degrees.}{Azimut der Kamera bzw. des Lichts in Grad.}}
  \SemioTableRow{\Key{elevation} & \Key{number} & --- & \ApiText{Camera elevation above the horizon in degrees.}{Kamerahöhe über dem Horizont in Grad.}}
  \SemioTableRow{\Key{meridians} & \Key{integer} & --- & \ApiText{The wireframe density.}{Die Dichte des Drahtgitters.}}
  \SemioTableRow{\Key{parallels} & \Key{integer} & --- & \ApiText{The wireframe density.}{Die Dichte des Drahtgitters.}}
}

\subsection{\Key{sci-population}}

\SemioTableLong[text-size=8pt]{\Key{sci-population}}{0.20,0.13,0.15,0.52}{\ApiKey & \ApiType & \ApiDefault & \ApiMeaning}{%
  \SemioTableRow{\Key{mode} & \Key{string} & --- & \ApiText{Which member of the family the preset renders.}{Welches Mitglied der Familie die Voreinstellung zeichnet.}}
  \SemioTableRow{\Key{r} & \Key{number} & --- & \ApiText{The logistic growth parameters.}{Die Parameter des logistischen Wachstums.}}
  \SemioTableRow{\Key{capacity} & \Key{number} & --- & \ApiText{The logistic growth parameters.}{Die Parameter des logistischen Wachstums.}}
  \SemioTableRow{\Key{alpha} & \Key{number} & --- & \ApiText{The Lotka–Volterra coefficients.}{Die Lotka-Volterra-Koeffizienten.}}
  \SemioTableRow{\Key{beta} & \Key{number} & --- & \ApiText{The Lotka–Volterra coefficients.}{Die Lotka-Volterra-Koeffizienten.}}
  \SemioTableRow{\Key{gamma} & \Key{number} & --- & \ApiText{The Lotka–Volterra coefficients.}{Die Lotka-Volterra-Koeffizienten.}}
  \SemioTableRow{\Key{delta} & \Key{number} & --- & \ApiText{The Lotka–Volterra coefficients.}{Die Lotka-Volterra-Koeffizienten.}}
  \SemioTableRow{\Key{prey} & \Key{number} & --- & \ApiText{The initial state.}{Der Anfangszustand.}}
  \SemioTableRow{\Key{predator} & \Key{number} & --- & \ApiText{The initial state.}{Der Anfangszustand.}}
  \SemioTableRow{\Key{steps} & \Key{integer} & --- & \ApiText{Number of integration or iteration steps.}{Anzahl der Integrations- bzw. Iterationsschritte.}}
  \SemioTableRow{\Key{step} & \Key{number} & --- & \ApiText{Spacing between consecutive steps.}{Abstand zwischen aufeinanderfolgenden Schritten.}}
  \SemioTableRow{\Key{counts} & \Key{string} & --- & \ApiText{The ranked abundances or the rarefaction sample.}{Die geordneten Abundanzen oder die Rarefaktionsstichprobe.}}
}

\subsection{\Key{sci-probability-tree}}

\SemioTableLong[text-size=8pt]{\Key{sci-probability-tree}}{0.20,0.13,0.15,0.52}{\ApiKey & \ApiType & \ApiDefault & \ApiMeaning}{%
  \SemioTableRow{\Key{branches} & \Key{string} & --- & \ApiText{One record per node of the tree.}{Ein Datensatz je Knoten des Baums.}}
  \SemioTableRow{\Key{levelWidth} & \Key{number} & --- & \ApiText{Millimetre width of one tree level.}{Millimeterbreite einer Baumebene.}}
  \SemioTableRow{\Key{levelGap} & \Key{number} & --- & \ApiText{Millimetre spacing of depths and siblings.}{Millimeterabstand von Tiefen und Geschwistern.}}
  \SemioTableRow{\Key{showProbabilities} & \Key{boolean} & --- & \ApiText{Print the branch probability on each edge.}{Die Zweigwahrscheinlichkeit an jede Kante schreiben.}}
}

\subsection{\Key{sci-punnett}}

\SemioTableLong[text-size=8pt]{\Key{sci-punnett}}{0.20,0.13,0.15,0.52}{\ApiKey & \ApiType & \ApiDefault & \ApiMeaning}{%
  \SemioTableRow{\Key{parentA} & \Key{string} & --- & \ApiText{The gametes of each parent.}{Die Gameten jedes Elternteils.}}
  \SemioTableRow{\Key{parentB} & \Key{string} & --- & \ApiText{The gametes of each parent.}{Die Gameten jedes Elternteils.}}
  \SemioTableRow{\Key{cell} & \Key{number} & --- & \ApiText{Cell size in millimetres.}{Zellgröße in Millimetern.}}
  \SemioTableRow{\Key{highlight} & \Key{string} & --- & \ApiText{Genotypes to shade.}{Einzufärbende Genotypen.}}
}

\subsection{\Key{sci-scatter-diagnostic}}

\SemioTableLong[text-size=8pt]{\Key{sci-scatter-diagnostic}}{0.20,0.13,0.15,0.52}{\ApiKey & \ApiType & \ApiDefault & \ApiMeaning}{%
  \SemioTableRow{\Key{mode} & \Key{string} & --- & \ApiText{Which member of the family the preset renders.}{Welches Mitglied der Familie die Voreinstellung zeichnet.}}
  \SemioTableRow{\Key{pairs} & \Key{string} & --- & \ApiText{The two measurements of a Bland–Altman plot.}{Die beiden Messungen eines Bland-Altman-Diagramms.}}
  \SemioTableRow{\Key{features} & \Key{string} & --- & \ApiText{The differential-expression table.}{Die Tabelle der differenziellen Expression.}}
  \SemioTableRow{\Key{threshold} & \Key{number} & --- & \ApiText{The genome-wide significance line.}{Die genomweite Signifikanzlinie.}}
  \SemioTableRow{\Key{effect} & \Key{number} & --- & \ApiText{The significance and effect-size cut-offs.}{Die Schwellen für Signifikanz und Effektstärke.}}
}

\subsection{\Key{sci-schematic}}

\SemioTableLong[text-size=8pt]{\Key{sci-schematic}}{0.20,0.13,0.15,0.52}{\ApiKey & \ApiType & \ApiDefault & \ApiMeaning}{%
  \SemioTableRow{\Key{shapes} & \Key{string} & --- & \ApiText{Kinds ellipse, rect, rounded, blob.}{Arten ellipse, rect, rounded, blob.}}
  \SemioTableRow{\Key{leaders} & \Key{string} & --- & \ApiText{The label leaders drawn from a part to its caption.}{Die Führungslinien von einem Teil zu seiner Beschriftung.}}
  \SemioTableRow{\Key{outline} & \Key{string} & --- & \ApiText{An enclosing membrane drawn behind the parts.}{Eine umschließende Membran, hinter den Teilen gezeichnet.}}
}

\subsection{\Key{sci-seqlogo}}

\SemioTableLong[text-size=8pt]{\Key{sci-seqlogo}}{0.20,0.13,0.15,0.52}{\ApiKey & \ApiType & \ApiDefault & \ApiMeaning}{%
  \SemioTableRow{\Key{columns} & \Key{string} & --- & \ApiText{One entry per position.}{Ein Eintrag je Position.}}
  \SemioTableRow{\Key{alphabetSize} & \Key{integer} & --- & \ApiText{The alphabet size that sets the maximum information content.}{Die Alphabetgröße, die den maximalen Informationsgehalt festlegt.}}
  \SemioTableRow{\Key{columnWidth} & \Key{number} & --- & \ApiText{Millimetre width of one sequence-logo position.}{Millimeterbreite einer Position im Sequenzlogo.}}
}

\subsection{\Key{sci-set}}

\SemioTableLong[text-size=8pt]{\Key{sci-set}}{0.20,0.13,0.15,0.52}{\ApiKey & \ApiType & \ApiDefault & \ApiMeaning}{%
  \SemioTableRow{\Key{mode} & \Key{string} & --- & \ApiText{What the expressions mean.}{Was die Ausdrücke bedeuten.}}
  \SemioTableRow{\Key{sets} & \Key{string} & --- & \ApiText{Table of sets and their sizes.}{Tabelle der Mengen und ihrer Größen.}}
  \SemioTableRow{\Key{intersections} & \Key{string} & --- & \ApiText{Where members are set labels joined by '+'.}{Wobei members Mengenbezeichner sind, mit '+' verbunden.}}
  \SemioTableRow{\Key{radius} & \Key{number} & --- & \ApiText{Base circle radius in millimetres; area-proportional scales it by the set sizes.}{Grundkreisradius in Millimetern; flächenproportional skaliert ihn nach den Mengengrößen.}}
  \SemioTableRow{\Key{labels} & \Key{boolean} & --- & \ApiText{Draw region counts.}{Die Anzahlen der Bereiche schreiben.}}
}

\subsection{\Key{sci-signal}}

\SemioTableLong[text-size=8pt]{\Key{sci-signal}}{0.20,0.13,0.15,0.52}{\ApiKey & \ApiType & \ApiDefault & \ApiMeaning}{%
  \SemioTableRow{\Key{mode} & \Key{string} & --- & \ApiText{Which member of the family the preset renders.}{Welches Mitglied der Familie die Voreinstellung zeichnet.}}
  \SemioTableRow{\Key{expr} & \Key{string} & --- & \ApiText{The signal and, for the two-signal modes, its partner.}{Das Signal und, in den Zweisignalmodi, sein Partner.}}
  \SemioTableRow{\Key{expr2} & \Key{string} & --- & \ApiText{The signal and, for the two-signal modes, its partner.}{Das Signal und, in den Zweisignalmodi, sein Partner.}}
  \SemioTableRow{\Key{samples} & \Key{integer} & --- & \ApiText{Number of samples the curve or field is evaluated at.}{Anzahl der Stützstellen, an denen Kurve oder Feld ausgewertet werden.}}
  \SemioTableRow{\Key{rate} & \Key{number} & --- & \ApiText{Buffer length and sampling rate.}{Pufferlänge und Abtastrate.}}
  \SemioTableRow{\Key{numerator} & \Key{string} & --- & \ApiText{The filter transfer function of 'filter-response'.}{Die Filterübertragungsfunktion von 'filter-response'.}}
  \SemioTableRow{\Key{denominator} & \Key{string} & --- & \ApiText{The filter transfer function of 'filter-response'.}{Die Filterübertragungsfunktion von 'filter-response'.}}
  \SemioTableRow{\Key{stems} & \Key{boolean} & --- & \ApiText{Draw a discrete spectrum as stems instead of a curve.}{Ein diskretes Spektrum als Stiele statt als Kurve zeichnen.}}
  \SemioTableRow{\Key{decibel} & \Key{boolean} & --- & \ApiText{Log magnitude axis.}{Logarithmische Betragsachse.}}
}

\subsection{\Key{sci-simplex}}

\SemioTableLong[text-size=8pt]{\Key{sci-simplex}}{0.20,0.13,0.15,0.52}{\ApiKey & \ApiType & \ApiDefault & \ApiMeaning}{%
  \SemioTableRow{\Key{points} & \Key{string} & --- & \ApiText{The candidate set a Pareto frontier is taken from.}{Die Kandidatenmenge, aus der eine Pareto-Front gebildet wird.}}
  \SemioTableRow{\Key{side} & \Key{number} & --- & \ApiText{Triangle side length in millimetres.}{Seitenlänge des Dreiecks in Millimetern.}}
  \SemioTableRow{\Key{gridLines} & \Key{integer} & --- & \ApiText{Interior guide count.}{Anzahl der inneren Hilfslinien.}}
}

\subsection{\Key{sci-skewt}}

\SemioTableLong[text-size=8pt]{\Key{sci-skewt}}{0.20,0.13,0.15,0.52}{\ApiKey & \ApiType & \ApiDefault & \ApiMeaning}{%
  \SemioTableRow{\Key{mode} & \Key{string} & --- & \ApiText{The coordinate skew of the diagram.}{Die Koordinatenscherung des Diagramms.}}
  \SemioTableRow{\Key{profile} & \Key{string} & --- & \ApiText{The sounding, pressure in hPa.}{Die Radiosondierung, Druck in hPa.}}
  \SemioTableRow{\Key{skew} & \Key{number} & --- & \ApiText{The isotherm slope in degrees.}{Die Neigung der Isothermen in Grad.}}
  \SemioTableRow{\Key{top} & \Key{number} & --- & \ApiText{The pressure window.}{Das Druckfenster.}}
  \SemioTableRow{\Key{bottom} & \Key{number} & --- & \ApiText{The pressure window.}{Das Druckfenster.}}
  \SemioTableRow{\Key{isotherms} & \Key{boolean} & --- & \ApiText{Draw the isotherms of the sounding diagram.}{Die Isothermen des Sondierungsdiagramms zeichnen.}}
  \SemioTableRow{\Key{adiabats} & \Key{boolean} & --- & \ApiText{The background line families.}{Die Linienscharen im Hintergrund.}}
}

\subsection{\Key{sci-smith}}

\SemioTableLong[text-size=8pt]{\Key{sci-smith}}{0.20,0.13,0.15,0.52}{\ApiKey & \ApiType & \ApiDefault & \ApiMeaning}{%
  \SemioTableRow{\Key{resistances} & \Key{string} & --- & \ApiText{The constant-resistance circles drawn.}{Die gezeichneten Kreise konstanten Widerstands.}}
  \SemioTableRow{\Key{reactances} & \Key{string} & --- & \ApiText{The constant-reactance arcs drawn.}{Die gezeichneten Bögen konstanter Reaktanz.}}
  \SemioTableRow{\Key{points} & \Key{string} & --- & \ApiText{Normalised impedances marked on the chart.}{Normierte Impedanzen, die im Diagramm markiert werden.}}
  \SemioTableRow{\Key{radius} & \Key{number} & --- & \ApiText{Chart radius in millimetres.}{Radius des Diagramms in Millimetern.}}
}

\subsection{\Key{sci-spacetime}}

\SemioTableLong[text-size=8pt]{\Key{sci-spacetime}}{0.20,0.13,0.15,0.52}{\ApiKey & \ApiType & \ApiDefault & \ApiMeaning}{%
  \SemioTableRow{\Key{mode} & \Key{string} & --- & \ApiText{The optical element traced.}{Das durchgerechnete optische Element.}}
  \SemioTableRow{\Key{beta} & \Key{number} & --- & \ApiText{The boost velocity as a fraction of c, which sets the primed axes.}{Die Boostgeschwindigkeit als Bruchteil von c, die die gestrichenen Achsen festlegt.}}
  \SemioTableRow{\Key{worldlines} & \Key{string} & --- & \ApiText{Further world lines through the origin.}{Weitere Weltlinien durch den Ursprung.}}
  \SemioTableRow{\Key{events} & \Key{string} & --- & \ApiText{Marked events.}{Markierte Ereignisse.}}
  \SemioTableRow{\Key{gridLines} & \Key{integer} & --- & \ApiText{Primed grid density.}{Dichte des gestrichenen Gitters.}}
}

\subsection{\Key{sci-spectrum-peaks}}

\SemioTableLong[text-size=8pt]{\Key{sci-spectrum-peaks}}{0.20,0.13,0.15,0.52}{\ApiKey & \ApiType & \ApiDefault & \ApiMeaning}{%
  \SemioTableRow{\Key{kind} & \Key{string} & --- & \ApiText{Sets the axis direction and the peak shape.}{Legt die Achsenrichtung und die Form der Spitze fest.}}
  \SemioTableRow{\Key{peaks} & \Key{string} & --- & \ApiText{The peak table; width in axis units.}{Die Peaktabelle; Breite in Achseneinheiten.}}
  \SemioTableRow{\Key{reversed} & \Key{boolean} & --- & \ApiText{Draw the x axis right to left.}{Die x-Achse von rechts nach links zeichnen.}}
  \SemioTableRow{\Key{samples} & \Key{integer} & --- & \ApiText{Number of samples the curve or field is evaluated at.}{Anzahl der Stützstellen, an denen Kurve oder Feld ausgewertet werden.}}
  \SemioTableRow{\Key{domain} & \Key{string} & --- & \ApiText{Input domain 'min,max' of the scale or function.}{Eingabebereich 'min,max' der Skala bzw. Funktion.}}
}

\subsection{\Key{sci-sports-series}}

\SemioTableLong[text-size=8pt]{\Key{sci-sports-series}}{0.20,0.13,0.15,0.52}{\ApiKey & \ApiType & \ApiDefault & \ApiMeaning}{%
  \SemioTableRow{\Key{mode} & \Key{string} & --- & \ApiText{Which member of the family the preset renders.}{Welches Mitglied der Familie die Voreinstellung zeichnet.}}
  \SemioTableRow{\Key{series} & \Key{string} & --- & \ApiText{Where points is a '+'-joined 'x:y' list.}{Wobei points eine mit '+' verbundene 'x:y'-Liste ist.}}
  \SemioTableRow{\Key{events} & \Key{string} & --- & \ApiText{Shots or touches, coordinates normalised to 0..1 of the surface.}{Schüsse oder Ballkontakte, Koordinaten auf 0..1 der Spielfläche normiert.}}
  \SemioTableRow{\Key{axesLabels} & \Key{string} & --- & \ApiText{The radar axis names.}{Die Namen der Netzdiagrammachsen.}}
  \SemioTableRow{\Key{reference} & \Key{number} & --- & \ApiText{The win-probability reference line.}{Die Referenzlinie der Siegwahrscheinlichkeit.}}
}

\subsection{\Key{sci-staff}}

\SemioTableLong[text-size=8pt]{\Key{sci-staff}}{0.20,0.13,0.15,0.52}{\ApiKey & \ApiType & \ApiDefault & \ApiMeaning}{%
  \SemioTableRow{\Key{clef} & \Key{string} & --- & \ApiText{The clef drawn at the start of the system.}{Der Schlüssel am Anfang des Systems.}}
  \SemioTableRow{\Key{events} & \Key{string} & --- & \ApiText{Step 0 is the bottom staff line, duration in quarter notes.}{Schritt 0 ist die unterste Notenlinie, Dauer in Vierteln.}}
  \SemioTableRow{\Key{spacing} & \Key{number} & --- & \ApiText{Millimetres per staff step pair.}{Millimeter je Notenlinienschrittpaar.}}
  \SemioTableRow{\Key{stride} & \Key{number} & --- & \ApiText{Millimetres per quarter note.}{Millimeter je Viertelnote.}}
  \SemioTableRow{\Key{beams} & \Key{boolean} & --- & \ApiText{Join consecutive eighth notes.}{Aufeinanderfolgende Achtelnoten verbalken.}}
}

\subsection{\Key{sci-stochastic}}

\SemioTableLong[text-size=8pt]{\Key{sci-stochastic}}{0.20,0.13,0.15,0.52}{\ApiKey & \ApiType & \ApiDefault & \ApiMeaning}{%
  \SemioTableRow{\Key{mode} & \Key{string} & --- & \ApiText{What the expressions mean.}{Was die Ausdrücke bedeuten.}}
  \SemioTableRow{\Key{paths} & \Key{integer} & --- & \ApiText{Number of sample paths drawn.}{Anzahl der gezeichneten Pfade.}}
  \SemioTableRow{\Key{steps} & \Key{integer} & --- & \ApiText{Number of integration or iteration steps.}{Anzahl der Integrations- bzw. Iterationsschritte.}}
  \SemioTableRow{\Key{drift} & \Key{number} & --- & \ApiText{Drift term of the stochastic process.}{Driftterm des stochastischen Prozesses.}}
  \SemioTableRow{\Key{volatility} & \Key{number} & --- & \ApiText{The increment law.}{Das Zuwachsgesetz.}}
  \SemioTableRow{\Key{samples} & \Key{integer} & --- & \ApiText{Uniform sample count.}{Anzahl gleichmäßiger Stützstellen.}}
  \SemioTableRow{\Key{accept} & \Key{string} & --- & \ApiText{The Monte-Carlo acceptance region.}{Der Annahmebereich der Monte-Carlo-Simulation.}}
}

\subsection{\Key{sci-survival}}

\SemioTableLong[text-size=8pt]{\Key{sci-survival}}{0.20,0.13,0.15,0.52}{\ApiKey & \ApiType & \ApiDefault & \ApiMeaning}{%
  \SemioTableRow{\Key{mode} & \Key{string} & --- & \ApiText{Which member of the family the preset renders.}{Welches Mitglied der Familie die Voreinstellung zeichnet.}}
  \SemioTableRow{\Key{groups} & \Key{string} & --- & \ApiText{Where events is a '+'-joined list of 'time:status'.}{Wobei events eine mit '+' verbundene Liste aus 'time:status' ist.}}
  \SemioTableRow{\Key{confidence} & \Key{boolean} & --- & \ApiText{Draw the Greenwood pointwise band.}{Das punktweise Greenwood-Band zeichnen.}}
  \SemioTableRow{\Key{censorMarks} & \Key{boolean} & --- & \ApiText{Tick the censored times.}{Die zensierten Zeiten markieren.}}
}

\subsection{\Key{sci-tiling}}

\SemioTableLong[text-size=8pt]{\Key{sci-tiling}}{0.20,0.13,0.15,0.52}{\ApiKey & \ApiType & \ApiDefault & \ApiMeaning}{%
  \SemioTableRow{\Key{kind} & \Key{string} & --- & \ApiText{Which conic.}{Welcher Kegelschnitt.}}
  \SemioTableRow{\Key{columns} & \Key{integer} & --- & \ApiText{Lattice extent of the regular tilings.}{Gitterausdehnung der regelmäßigen Parkettierungen.}}
  \SemioTableRow{\Key{rows} & \Key{integer} & --- & \ApiText{Lattice extent of the regular tilings.}{Gitterausdehnung der regelmäßigen Parkettierungen.}}
  \SemioTableRow{\Key{size} & \Key{number} & --- & \ApiText{Tile size in millimetres.}{Kachelgröße in Millimetern.}}
  \SemioTableRow{\Key{depth} & \Key{integer} & --- & \ApiText{Penrose substitution steps, or hyperbolic recursion depth.}{Penrose-Substitutionsschritte oder hyperbolische Rekursionstiefe.}}
  \SemioTableRow{\Key{p} & \Key{integer} & --- & \ApiText{The Schläfli symbol \{p,q\} of the hyperbolic tiling.}{Das Schläfli-Symbol \{p,q\} der hyperbolischen Parkettierung.}}
  \SemioTableRow{\Key{q} & \Key{integer} & --- & \ApiText{The Schläfli symbol \{p,q\} of the hyperbolic tiling.}{Das Schläfli-Symbol \{p,q\} der hyperbolischen Parkettierung.}}
  \SemioTableRow{\Key{discRadius} & \Key{number} & --- & \ApiText{Poincaré disc radius.}{Radius der Poincaré-Scheibe.}}
}

\subsection{\Key{sci-timefreq}}

\SemioTableLong[text-size=8pt]{\Key{sci-timefreq}}{0.20,0.13,0.15,0.52}{\ApiKey & \ApiType & \ApiDefault & \ApiMeaning}{%
  \SemioTableRow{\Key{mode} & \Key{string} & --- & \ApiText{Which member of the family the preset renders.}{Welches Mitglied der Familie die Voreinstellung zeichnet.}}
  \SemioTableRow{\Key{expr} & \Key{string} & --- & \ApiText{The signal and, for the two-signal modes, its partner.}{Das Signal und, in den Zweisignalmodi, sein Partner.}}
  \SemioTableRow{\Key{samples} & \Key{integer} & --- & \ApiText{Number of samples the curve or field is evaluated at.}{Anzahl der Stützstellen, an denen Kurve oder Feld ausgewertet werden.}}
  \SemioTableRow{\Key{rate} & \Key{number} & --- & \ApiText{Buffer length and sampling rate.}{Pufferlänge und Abtastrate.}}
  \SemioTableRow{\Key{window} & \Key{integer} & --- & \ApiText{Length of the analysis window in samples.}{Länge des Analysefensters in Abtastwerten.}}
  \SemioTableRow{\Key{hop} & \Key{integer} & --- & \ApiText{The short-time Fourier window and its advance.}{Das Kurzzeit-Fourier-Fenster und sein Vorschub.}}
  \SemioTableRow{\Key{scales} & \Key{integer} & --- & \ApiText{Morlet scale count of the scalogram.}{Anzahl der Morlet-Skalen des Skalogramms.}}
  \SemioTableRow{\Key{decibel} & \Key{boolean} & --- & \ApiText{Log magnitude axis.}{Logarithmische Betragsachse.}}
}

\subsection{\Key{sci-tonnetz}}

\SemioTableLong[text-size=8pt]{\Key{sci-tonnetz}}{0.20,0.13,0.15,0.52}{\ApiKey & \ApiType & \ApiDefault & \ApiMeaning}{%
  \SemioTableRow{\Key{mode} & \Key{string} & --- & \ApiText{Which member of the family the preset renders.}{Welches Mitglied der Familie die Voreinstellung zeichnet.}}
  \SemioTableRow{\Key{columns} & \Key{integer} & --- & \ApiText{Tonnetz extent.}{Ausdehnung des Tonnetzes.}}
  \SemioTableRow{\Key{rows} & \Key{integer} & --- & \ApiText{Tonnetz extent.}{Ausdehnung des Tonnetzes.}}
  \SemioTableRow{\Key{size} & \Key{number} & --- & \ApiText{Lattice spacing in millimetres.}{Gitterabstand in Millimetern.}}
  \SemioTableRow{\Key{highlight} & \Key{string} & --- & \ApiText{Pitch classes to mark.}{Zu markierende Tonklassen.}}
  \SemioTableRow{\Key{frets} & \Key{integer} & --- & \ApiText{Number of frets of the fretboard grid.}{Anzahl der Bünde des Griffbrettgitters.}}
  \SemioTableRow{\Key{strings} & \Key{integer} & --- & \ApiText{The guitar grid.}{Das Griffbrettgitter.}}
  \SemioTableRow{\Key{fingering} & \Key{string} & --- & \ApiText{The chord shape.}{Die Akkordform.}}
}

\subsection{\Key{sci-transform}}

\SemioTableLong[text-size=8pt]{\Key{sci-transform}}{0.20,0.13,0.15,0.52}{\ApiKey & \ApiType & \ApiDefault & \ApiMeaning}{%
  \SemioTableRow{\Key{kind} & \Key{string} & --- & \ApiText{Which conic.}{Welcher Kegelschnitt.}}
  \SemioTableRow{\Key{points} & \Key{string} & --- & \ApiText{The vertices, in millimetres of the figure.}{Die Eckpunkte in Millimetern der Abbildung.}}
  \SemioTableRow{\Key{dx} & \Key{number} & --- & \ApiText{Column holding the horizontal extent of a box.}{Spalte mit der waagerechten Ausdehnung eines Kastens.}}
  \SemioTableRow{\Key{dy} & \Key{number} & --- & \ApiText{Column holding the vertical extent of a box.}{Spalte mit der senkrechten Ausdehnung eines Kastens.}}
  \SemioTableRow{\Key{angle} & \Key{number} & --- & \ApiText{Rotation or reflection-axis angle in degrees.}{Winkel der Drehung bzw. der Spiegelachse in Grad.}}
  \SemioTableRow{\Key{sx} & \Key{number} & --- & \ApiText{Scale factors.}{Skalierungsfaktoren.}}
  \SemioTableRow{\Key{sy} & \Key{number} & --- & \ApiText{Scale factors.}{Skalierungsfaktoren.}}
  \SemioTableRow{\Key{shear} & \Key{number} & --- & \ApiText{Shear factor.}{Scherfaktor.}}
  \SemioTableRow{\Key{about} & \Key{string} & --- & \ApiText{The fixed point.}{Der Fixpunkt.}}
  \SemioTableRow{\Key{labels} & \Key{boolean} & --- & \ApiText{Print vertex labels.}{Die Knotenbeschriftungen setzen.}}
}

\subsection{\Key{sci-transition}}

\SemioTableLong[text-size=8pt]{\Key{sci-transition}}{0.20,0.13,0.15,0.52}{\ApiKey & \ApiType & \ApiDefault & \ApiMeaning}{%
  \SemioTableRow{\Key{mode} & \Key{string} & --- & \ApiText{What the expressions mean.}{Was die Ausdrücke bedeuten.}}
  \SemioTableRow{\Key{nodes} & \Key{string} & --- & \ApiText{Name of the node table.}{Name der Knotentabelle.}}
  \SemioTableRow{\Key{edges} & \Key{string} & --- & \ApiText{Name of the edge table.}{Name der Kantentabelle.}}
  \SemioTableRow{\Key{matrix} & \Key{string} & --- & \ApiText{Name of the matrix table the family reads.}{Name der Matrixtabelle, die die Familie liest.}}
  \SemioTableRow{\Key{cell} & \Key{number} & --- & \ApiText{Cell size in millimetres.}{Zellgröße in Millimetern.}}
  \SemioTableRow{\Key{unit} & \Key{number} & --- & \ApiText{Millimetres per data unit of the node coordinates.}{Millimeter je Dateneinheit der Knotenkoordinaten.}}
}

\subsection{\Key{sci-truss}}

\SemioTableLong[text-size=8pt]{\Key{sci-truss}}{0.20,0.13,0.15,0.52}{\ApiKey & \ApiType & \ApiDefault & \ApiMeaning}{%
  \SemioTableRow{\Key{mode} & \Key{string} & --- & \ApiText{Which symbol vocabulary is used.}{Welcher Symbolvorrat verwendet wird.}}
  \SemioTableRow{\Key{nodes} & \Key{string} & --- & \ApiText{Name of the node table.}{Name der Knotentabelle.}}
  \SemioTableRow{\Key{members} & \Key{string} & --- & \ApiText{The geometry and its member forces.}{Die Geometrie und ihre Stabkräfte.}}
  \SemioTableRow{\Key{columns} & \Key{integer} & --- & \ApiText{The finite-element or CFD lattice of the 'mesh' mode.}{Das Finite-Elemente- oder CFD-Gitter des Modus 'mesh'.}}
  \SemioTableRow{\Key{rows} & \Key{integer} & --- & \ApiText{The finite-element or CFD lattice of the 'mesh' mode.}{Das Finite-Elemente- oder CFD-Gitter des Modus 'mesh'.}}
  \SemioTableRow{\Key{field} & \Key{string} & --- & \ApiText{The element value that colours the mesh.}{Der Elementwert, der das Netz einfärbt.}}
}

\subsection{\Key{sci-upset}}

\SemioTableLong[text-size=8pt]{\Key{sci-upset}}{0.20,0.13,0.15,0.52}{\ApiKey & \ApiType & \ApiDefault & \ApiMeaning}{%
  \SemioTableRow{\Key{mode} & \Key{string} & --- & \ApiText{What the expressions mean.}{Was die Ausdrücke bedeuten.}}
  \SemioTableRow{\Key{sets} & \Key{string} & --- & \ApiText{Table of sets and their sizes.}{Tabelle der Mengen und ihrer Größen.}}
  \SemioTableRow{\Key{intersections} & \Key{string} & --- & \ApiText{Where members are set labels joined by '+'.}{Wobei members Mengenbezeichner sind, mit '+' verbunden.}}
  \SemioTableRow{\Key{cell} & \Key{number} & --- & \ApiText{Cell size in millimetres.}{Zellgröße in Millimetern.}}
  \SemioTableRow{\Key{barHeight} & \Key{number} & --- & \ApiText{Bar panel height in millimetres.}{Höhe des Balkenfelds in Millimetern.}}
}

\subsection{\Key{sci-wave}}

\SemioTableLong[text-size=8pt]{\Key{sci-wave}}{0.20,0.13,0.15,0.52}{\ApiKey & \ApiType & \ApiDefault & \ApiMeaning}{%
  \SemioTableRow{\Key{mode} & \Key{string} & --- & \ApiText{The optical element traced.}{Das durchgerechnete optische Element.}}
  \SemioTableRow{\Key{wavelength} & \Key{number} & --- & \ApiText{Wavelength of the plotted wave.}{Wellenlänge der gezeichneten Welle.}}
  \SemioTableRow{\Key{amplitude} & \Key{number} & --- & \ApiText{Amplitude of the plotted wave.}{Amplitude der gezeichneten Welle.}}
  \SemioTableRow{\Key{phase} & \Key{number} & --- & \ApiText{The wave parameters.}{Die Wellenparameter.}}
  \SemioTableRow{\Key{harmonics} & \Key{integer} & --- & \ApiText{Standing-wave modes drawn.}{Gezeichnete Moden der stehenden Welle.}}
  \SemioTableRow{\Key{sources} & \Key{string} & --- & \ApiText{The interfering sources.}{Die interferierenden Quellen.}}
  \SemioTableRow{\Key{slits} & \Key{integer} & --- & \ApiText{Number of single slits in the diffraction aperture.}{Anzahl der Einzelspalte in der Beugungsblende.}}
  \SemioTableRow{\Key{slitWidth} & \Key{number} & --- & \ApiText{The diffraction aperture.}{Die Beugungsblende.}}
  \SemioTableRow{\Key{samples} & \Key{integer} & --- & \ApiText{Number of samples the curve or field is evaluated at.}{Anzahl der Stützstellen, an denen Kurve oder Feld ausgewertet werden.}}
  \SemioTableRow{\Key{ellipticity} & \Key{number} & --- & \ApiText{Ellipticity of the polarisation ellipse.}{Elliptizität der Polarisationsellipse.}}
  \SemioTableRow{\Key{tilt} & \Key{number} & --- & \ApiText{The polarisation ellipse.}{Die Polarisationsellipse.}}
}

%endregion 🔖️GeneratedNamespaces
\chapter{\ApiText{Extending the library}{Die Bibliothek erweitern}}

\section{\ApiText{How to add a chart kind}{Wie man eine Diagrammart hinzufügt}}

\ApiText{%
  A chart kind is data, not code. Add one entry to the handcrafted catalogue
  \Key{\E{🖼}assets/\E{🔣}viz-catalog.json} and run \Key{bun ./\E{📜}script.ts generate viz}; the generator
  writes \Key{semio-viz-catalog.sty}, \Key{semio-viz-catalog-labels.sty} and the per-section gallery
  documents. Never edit a generated file.%
}{%
  Eine Diagrammart ist Daten, kein Code. Fügen Sie dem handgefertigten Katalog
  \Key{\E{🖼}assets/\E{🔣}viz-catalog.json} einen Eintrag hinzu und starten Sie
  \Key{bun ./\E{📜}script.ts generate viz}; der Generator schreibt \Key{semio-viz-catalog.sty},
  \Key{semio-viz-catalog-labels.sty} und die Galeriedokumente je Abschnitt. Bearbeiten Sie nie eine
  erzeugte Datei.%
}

\SemioTableLong[text-size=8pt]{\ApiText{Fields of a catalogue entry}{Felder eines Katalogeintrags}}{0.16,0.84}{\ApiText{Field}{Feld} & \ApiMeaning}{%
  \SemioTableRow{\Key{id} & \ApiText{Section number and slug, e.g. \Key{7/treemap} --- the sort key of the file.}{Abschnittsnummer und Slug, z.\,B. \Key{7/treemap} --- der Sortierschlüssel der Datei.}}
  \SemioTableRow{\Key{slug} & \ApiText{The global chart-kind name, without a section suffix. Distinct semantics get distinct handcrafted slugs.}{Der globale Name der Diagrammart, ohne Abschnittssuffix. Unterschiedliche Bedeutungen bekommen unterschiedliche, handgemachte Slugs.}}
  \SemioTableRow{\Key{title} & \ApiText{The title in \Key{en} and \Key{de}; both are required.}{Der Titel auf \Key{en} und \Key{de}; beide sind Pflicht.}}
  \SemioTableRow{\Key{kind} & \ApiText{\Key{chart}, or \Key{scale}, \Key{axis}, \Key{layout} for a kernel capability.}{\Key{chart}, oder \Key{scale}, \Key{axis}, \Key{layout} für eine Kernfähigkeit.}}
  \SemioTableRow{\Key{namespace} & \ApiText{The taxonomy section 76 namespace path, e.g. \Key{hierarchy/treemap}.}{Der Namensraumpfad aus Taxonomieabschnitt 76, z.\,B. \Key{hierarchy/treemap}.}}
  \SemioTableRow{\Key{family} & \ApiText{The registered family that renders it.}{Die registrierte Familie, die sie zeichnet.}}
  \SemioTableRow{\Key{options} & \ApiText{The family options this preset sets, including \Key{variant}.}{Die Familienoptionen, die diese Voreinstellung setzt, einschließlich \Key{variant}.}}
  \SemioTableRow{\Key{data} & \ApiText{The demo table the kind draws from.}{Die Demotabelle, aus der die Art zeichnet.}}
  \SemioTableRow{\Key{covers} & \ApiText{The taxonomy leaves this one kind stands for.}{Die Taxonomieblätter, für die diese eine Art steht.}}
}

\ApiText{%
  Five rules are enforced by \Key{test fundamental} and will fail the build if broken: every taxonomy
  leaf is covered exactly once; every option key is declared for that family in
  \Key{\E{🧬}schema/\E{🔣}.json}; every entry carries \Key{variant}; no two kinds of one family share the
  same family and option signature; and the named demo table exists. Both titles must be non-empty.%
}{%
  Fünf Regeln setzt \Key{test fundamental} durch, ihre Verletzung lässt den Bau scheitern: Jedes
  Taxonomieblatt ist genau einmal abgedeckt; jeder Optionsschlüssel ist für diese Familie in
  \Key{\E{🧬}schema/\E{🔣}.json} deklariert; jeder Eintrag trägt \Key{variant}; keine zwei Arten einer
  Familie teilen dieselbe Familien- und Optionssignatur; und die genannte Demotabelle existiert.
  Beide Titel müssen gefüllt sein.%
}

\section{\ApiText{How to add a family}{Wie man eine Familie hinzufügt}}

\ApiText{%
  A family is a renderer registered with \Cs{SemioVizFamily}\{\meta{name}\}\{\meta{code}\} in the
  namespace package that owns it; inside \meta{code}, \Key{\#1} is the option list. Six obligations
  come with it.%
}{%
  Eine Familie ist ein mit \Cs{SemioVizFamily}\{\meta{Name}\}\{\meta{Code}\} registrierter Renderer im
  Namensraumpaket, dem sie gehört; in \meta{Code} ist \Key{\#1} die Optionsliste. Damit kommen sechs
  Pflichten.%
}

\SemioTableLong[text-size=8pt]{\ApiText{Obligations of a family}{Pflichten einer Familie}}{0.22,0.78}{\ApiText{Obligation}{Pflicht} & \ApiMeaning}{%
  \SemioTableRow{\ApiText{Document the keys}{Schlüssel dokumentieren} & \ApiText{A \Key{\%region \E{🔖}Keys} block at the top of the family lists every option with its type, its default and its meaning. There are no undocumented options.}{Ein \Key{\%region \E{🔖}Keys}-Block am Kopf der Familie führt jede Option mit Typ, Vorgabe und Bedeutung auf. Es gibt keine undokumentierten Optionen.}}
  \SemioTableRow{\ApiText{Declare them in the schema}{Im Schema deklarieren} & \ApiText{The same names go into \Key{x-semio-family-options.<family>.options} with a description in both languages; the catalogue check gates on it.}{Dieselben Namen kommen in \Key{x-semio-family-options.<Familie>.options}, mit einer Beschreibung in beiden Sprachen; die Katalogprüfung hängt daran.}}
  \SemioTableRow{\ApiText{Consume the data}{Die Daten verwenden} & \ApiText{Two kinds of one family with different options must render differently, and their probe projections must differ. This is verified.}{Zwei Arten einer Familie mit verschiedenen Optionen müssen verschieden zeichnen, und ihre Prüfprojektionen müssen sich unterscheiden. Das wird geprüft.}}
  \SemioTableRow{\ApiText{Draw through the kernel}{Über den Kern zeichnen} & \ApiText{Positions come from scales and the coordinate system, geometry from marks and shape generators. No renderer writes a raw drawing primitive.}{Positionen kommen von Skalen und dem Koordinatensystem, Geometrie von Marken und Formgeneratoren. Kein Renderer schreibt einen rohen Zeichenbefehl.}}
  \SemioTableRow{\ApiText{Take colours from the theme}{Farben vom Thema nehmen} & \ApiText{Only through the theme kernel's colour, shade and chrome accessors; never a literal colour and never a private palette.}{Nur über die Farb-, Schattierungs- und Chrome-Zugriffe des Themenkerns; nie eine wörtliche Farbe und nie eine private Palette.}}
  \SemioTableRow{\ApiText{Emit its geometry}{Ihre Geometrie ausgeben} & \ApiText{Call \Cs{semio\_viz\_probe\_geometry:nn} so the rendered primitives are testable numerically, without rasterising a page.}{\Cs{semio\_viz\_probe\_geometry:nn} aufrufen, damit die gezeichneten Grundbausteine numerisch prüfbar sind, ohne eine Seite zu rastern.}}
}

\section{\ApiText{How to add a test}{Wie man einen Test hinzufügt}}

\ApiText{%
  A test case is a directory \Key{\E{🧪}tests/<kebab-case>/} holding exactly three things: the
  language-neutral \Key{\E{🥒}.feature}, the TypeScript adapter \Key{\E{🟦}.ts} and \Key{\E{🧫}fixtures/}. Any
  other file in that directory is a contract breach. The case slug is lower-case kebab, with no emoji
  and no underscore.%
}{%
  Ein Testfall ist ein Verzeichnis \Key{\E{🧪}tests/<kebab-case>/} mit genau drei Dingen: der
  sprachneutralen \Key{\E{🥒}.feature}, dem TypeScript-Adapter \Key{\E{🟦}.ts} und \Key{\E{🧫}fixtures/}. Jede
  andere Datei dort ist ein Vertragsbruch. Der Fall-Slug ist Kleinschreibung mit Bindestrichen, ohne
  Emoji und ohne Unterstrich.%
}

\SemioTableLong[text-size=8pt]{\ApiText{Feature tags}{Feature-Marken}}{0.28,0.72}{\ApiText{Tag}{Marke} & \ApiMeaning}{%
  \SemioTableRow{\Key{@capability-<kebab>} & \ApiText{The one capability the feature specifies. It must be listed on the oracle entry it names.}{Die eine Fähigkeit, die das Feature festlegt. Sie muss beim genannten Orakeleintrag aufgeführt sein.}}
  \SemioTableRow{\Key{@oracle-<id>} & \ApiText{A third-party reference registered in \Key{\E{🔮}oracle/\E{🔣}.json} --- or \Key{@no-oracle-<id>} for a justified decision to specify the numbers instead.}{Eine in \Key{\E{🔮}oracle/\E{🔣}.json} registrierte Fremdreferenz --- oder \Key{@no-oracle-<id>} für eine begründete Entscheidung, die Zahlen stattdessen festzuschreiben.}}
  \SemioTableRow{\Key{@comparison-viz-probe-v1} & \ApiText{The numeric comparison profile of the print owner.}{Das numerische Vergleichsprofil dieses Produkts.}}
  \SemioTableRow{\Key{@id-<kebab>} & \ApiText{Per scenario; must equal the adapter's scenario key.}{Je Szenario; muss dem Szenarioschlüssel des Adapters entsprechen.}}
  \SemioTableRow{\Key{@level-…} & \ApiText{\Key{fundamental} (no TeX), \Key{quick} (kernel probes on small vectors), \Key{long} (family probes and the gallery build), \Key{exhaustive} (every kind in both themes and both languages).}{\Key{fundamental} (kein TeX), \Key{quick} (Kernprüfungen auf kleinen Vektoren), \Key{long} (Familienprüfungen und der Galeriebau), \Key{exhaustive} (jede Art in beiden Themen und beiden Sprachen).}}
  \SemioTableRow{\Key{@mode-…} & \ApiText{\Key{differential}, \Key{conformance}, \Key{property} or \Key{error}.}{\Key{differential}, \Key{conformance}, \Key{property} oder \Key{error}.}}
}

\ApiText{%
  The vectors live in the Gherkin data table, so the feature stays the single source of the numbers.
  The adapter's \Key{subject} compiles a probe document with the repository's own tectonic and
  returns the parsed projection; its \Key{oracle} computes the same projection with the referenced
  package. Round both sides onto the same decimal grid --- the oracle in TypeScript and the probe
  through \Cs{SemioVizProbePrecision} --- or the last digit of TeX's fixed-point arithmetic decides
  the test.%
}{%
  Die Vektoren stehen in der Gherkin-Datentabelle, das Feature bleibt so die einzige Quelle der
  Zahlen. Das \Key{subject} des Adapters übersetzt ein Prüfdokument mit dem repositorieeigenen
  Tectonic und liefert die geparste Projektion; sein \Key{oracle} berechnet dieselbe Projektion mit
  dem genannten Paket. Runden Sie beide Seiten auf dasselbe Dezimalraster --- das Orakel in
  TypeScript und die Prüfung über \Cs{SemioVizProbePrecision} --- sonst entscheidet die letzte Stelle
  der Festkommaarithmetik von TeX über den Test.%
}

\SemioTableLong[text-size=8pt]{\ApiText{The probe protocol}{Das Prüfprotokoll}}{0.34,0.66}{\ApiText{Macro}{Makro} & \ApiMeaning}{%
  \SemioTableRow{\Cs{SemioVizProbeBegin}\{\meta{case}\}\{\meta{scenario}\} & \ApiText{Opens the record stream and fixes the identity of what follows.}{Öffnet den Datenstrom und legt die Identität des Folgenden fest.}}
  \SemioTableRow{\Cs{SemioVizProbePage} & \ApiText{Typesets one line naming case and scenario. Mandatory in a probe that typesets nothing else: a TeX run with no page writes no files at all, and the stream is lost with them.}{Setzt eine Zeile mit Fall und Szenario. Pflicht in einer Prüfung, die sonst nichts setzt: Ein TeX-Lauf ohne Seite schreibt überhaupt keine Dateien, und der Datenstrom geht mit ihnen verloren.}}
  \SemioTableRow{\Cs{SemioVizProbeValues}\{\meta{key}\}\{\meta{clist}\} & \ApiText{One record; each item is emitted bare if it is a JSON number and quoted otherwise.}{Ein Datensatz; jeder Eintrag wird blank ausgegeben, wenn er eine JSON-Zahl ist, sonst in Anführungszeichen.}}
  \SemioTableRow{\Cs{SemioVizProbePoints}\{\meta{key}\}\{\meta{points}\} & \ApiText{One record with the flattened coordinates.}{Ein Datensatz mit den flach ausgeschriebenen Koordinaten.}}
  \SemioTableRow{\Cs{SemioVizProbeEval}\{\meta{key}\}\{\meta{expressions}\} & \ApiText{One record, each expression evaluated on the emission grid.}{Ein Datensatz, jeder Ausdruck auf dem Ausgaberaster ausgewertet.}}
  \SemioTableRow{\Cs{SemioVizProbeString}\{\meta{key}\}\{\meta{text}\} & \ApiText{One record with a single string value, commas included.}{Ein Datensatz mit einem einzelnen Textwert, Kommata eingeschlossen.}}
  \SemioTableRow{\Cs{SemioVizProbePrecision}\{\meta{decimals}\} & \ApiText{Decimal places every evaluated number is emitted on; the default is 6, a negative value emits the raw result.}{Nachkommastellen, auf denen jede ausgewertete Zahl ausgegeben wird; die Vorgabe ist 6, ein negativer Wert gibt das rohe Ergebnis aus.}}
  \SemioTableRow{\Cs{SemioVizProbeOn}, \Cs{SemioVizProbeOff} & \ApiText{Toggle geometry emission from marks and families.}{Schalten die Geometrieausgabe von Marken und Familien um.}}
  \SemioTableRow{\Cs{SemioVizProbeEnd} & \ApiText{Closes the stream early; it also runs at the end of the document.}{Schließt den Datenstrom vorzeitig; er läuft auch am Dokumentende.}}
}

\ApiText{%
  The output is \Key{<jobname>.probe.jsonl}, one JSON object per line with the fields \Key{case},
  \Key{scenario}, \Key{key} and \Key{values}. Emit one record per key, values in a fixed order; the
  harness groups by key and concatenates in emission order.%
}{%
  Die Ausgabe ist \Key{<jobname>.probe.jsonl}, je Zeile ein JSON-Objekt mit den Feldern \Key{case},
  \Key{scenario}, \Key{key} und \Key{values}. Geben Sie je Schlüssel einen Datensatz aus, die Werte in
  fester Reihenfolge; die Prüfumgebung gruppiert nach Schlüssel und hängt in Ausgabereihenfolge
  aneinander.%
}

\chapter{\ApiText{Numerics and limits}{Numerik und Grenzen}}

\section{\ApiText{What l3fp is, and what it costs}{Was l3fp ist und was es kostet}}

\ApiText{%
  Every number in this library is computed in l3fp, the LaTeX3 floating-point unit. It is a
  \emph{decimal} arithmetic with sixteen significant digits, evaluated by macro expansion. That has
  three consequences a caller should know about.%
}{%
  Jede Zahl dieser Bibliothek wird in l3fp berechnet, der Fließkommaeinheit von LaTeX3. Es ist eine
  \emph{dezimale} Arithmetik mit sechzehn signifikanten Stellen, ausgewertet durch Makroexpansion.
  Das hat drei Folgen, die ein Aufrufer kennen sollte.%
}

\ApiText{%
  \textbf{It is not IEEE-754.} A value whose decimal expansion sits exactly on a rounding boundary of
  a non-tie double can differ from JavaScript in the last digit. l3fp also breaks ties to even where
  JavaScript breaks them away from zero, so every place that mirrors a rounding step computes
  \Key{floor(x + 0.5)} explicitly. The test vectors avoid the ambiguous values instead of pretending
  it does not exist.%
}{%
  \textbf{Es ist nicht IEEE-754.} Ein Wert, dessen Dezimalentwicklung genau auf einer Rundungsgrenze
  eines nicht mittigen Doubles liegt, kann sich in der letzten Stelle von JavaScript unterscheiden.
  l3fp rundet Gleichstände zur geraden Zahl, JavaScript von der Null weg, deshalb rechnet jede Stelle,
  die einen Rundungsschritt nachbildet, ausdrücklich \Key{floor(x + 0.5)}. Die Testvektoren meiden die
  mehrdeutigen Werte, statt so zu tun, als gäbe es das Problem nicht.%
}

\ApiText{%
  \textbf{It is slow.} One evaluated expression costs on the order of a tenth of a millisecond on the
  machine these numbers were measured on. An algorithm that is linear in the data is instant; one
  that is quadratic is not. Every expensive algorithm therefore exposes the parameter that decides
  its cost, and the catalogue presets stay well inside a few seconds.%
}{%
  \textbf{Es ist langsam.} Ein ausgewerteter Ausdruck kostet auf der Maschine, auf der diese Zahlen
  gemessen wurden, in der Größenordnung einer Zehntelmillisekunde. Ein Algorithmus, der linear in den
  Daten ist, läuft sofort; ein quadratischer nicht. Jeder teure Algorithmus legt deshalb den Parameter
  offen, der seinen Aufwand bestimmt, und die Katalogvoreinstellungen bleiben deutlich unter wenigen
  Sekunden.%
}

\ApiText{%
  \textbf{It is exact where exactness was the point.} The force layout evaluates its many-body force
  once per unordered pair rather than through a Barnes--Hut approximation, and the geographic
  projections resample at d3's own tolerance. Both decisions were taken so that the library can be
  compared against a third-party reference number for number; the cost of that choice is in the table
  below, not hidden.%
}{%
  \textbf{Es ist exakt, wo Exaktheit der Zweck war.} Das Kräftelayout wertet seine Vielkörperkraft
  einmal je ungeordnetem Paar aus statt über eine Barnes-Hut-Näherung, und die geographischen
  Projektionen tasten mit d3s eigener Toleranz nach. Beide Entscheidungen wurden getroffen, damit die
  Bibliothek Zahl für Zahl gegen eine Fremdreferenz geprüft werden kann; der Preis dafür steht in der
  Tabelle unten und ist nicht versteckt.%
}

\section{\ApiText{Measured limits}{Gemessene Grenzen}}

\SemioTableLong[text-size=8pt]{\ApiText{Measured cost of the expensive algorithms}{Gemessener Aufwand der teuren Algorithmen}}{0.20,0.16,0.20,0.44}{\ApiText{Algorithm}{Algorithmus} & \ApiText{Input}{Eingabe} & \ApiText{Wall clock}{Laufzeit} & \ApiText{Complexity and ceiling}{Aufwand und Obergrenze}}{%
  \SemioTableRow{\Key{tree}, \Key{cluster} & \ApiText{211 nodes}{211 Knoten} & \Key{0.3 s} & \ApiText{Linear in the nodes plus the contour walk.}{Linear in den Knoten plus der Konturdurchlauf.}}
  \SemioTableRow{\Key{treemap} & \ApiText{211 nodes}{211 Knoten} & \Key{0.5 s} & \ApiText{Linear, plus the tiling per parent.}{Linear, plus die Kachelung je Elternknoten.}}
  \SemioTableRow{\Key{partition} & \ApiText{211 nodes}{211 Knoten} & \Key{0.7 s} & \ApiText{Linear.}{Linear.}}
  \SemioTableRow{\Key{pack} & \ApiText{211 nodes}{211 Knoten} & \Key{18.7 s} & \ApiText{The outlier, and the cost is l3fp rather than the algorithm: the front chain plus an expected-linear enclosing circle per parent, run twice by d3's default-radius branch. Real figures are far smaller --- a 13-node pack fixture compiles in 4.7 s.}{Der Ausreißer, und der Aufwand liegt an l3fp, nicht am Algorithmus: die Frontkette plus ein erwartungsgemäß linearer Umkreis je Elternknoten, von d3s Vorgaberadius-Zweig zweimal durchlaufen. Wirkliche Abbildungen sind viel kleiner --- eine Packung mit 13 Knoten übersetzt in 4,7 s.}}
  \SemioTableRow{\Key{force} & \ApiText{100 nodes, 300 iterations}{100 Knoten, 300 Iterationen} & \Key{5 min 30 s} & \ApiText{Proportional to iterations times (nodes squared plus edges), four evaluations per pair.}{Proportional zu Iterationen mal (Knoten zum Quadrat plus Kanten), vier Auswertungen je Paar.}}
  \SemioTableRow{\Key{force} & \ApiText{100 nodes, 30 iterations}{100 Knoten, 30 Iterationen} & \Key{2 min 12 s} & \ApiText{The catalogue presets sit at 60 to 160 iterations on graphs of ten nodes, which is seconds.}{Die Katalogvoreinstellungen liegen bei 60 bis 160 Iterationen auf Graphen mit zehn Knoten, das sind Sekunden.}}
  \SemioTableRow{\Key{force} & \ApiText{8 nodes, 20 iterations}{8 Knoten, 20 Iterationen} & \Key{1.6 s} & \ApiText{The size a figure should use.}{Die Größe, die eine Abbildung nutzen sollte.}}
  \SemioTableRow{\Key{delaunay}, \Key{voronoi} & \ApiText{12 points}{12 Punkte} & \Key{4 s} & \ApiText{Quadratic in the insertion scan. 24 points already cost 25 s; the hard ceiling is 120 points, and demo sets should stay at or below 40.}{Quadratisch im Einfügedurchlauf. 24 Punkte kosten bereits 25 s; die harte Obergrenze liegt bei 120 Punkten, Demomengen sollten bei höchstens 40 bleiben.}}
  \SemioTableRow{\Key{contour} & \ApiText{6 by 5 grid}{Gitter 6 mal 5} & \ApiText{instant}{sofort} & \ApiText{Linear in the cells; the hard ceiling is 4096 cells.}{Linear in den Zellen; die harte Obergrenze liegt bei 4096 Zellen.}}
  \SemioTableRow{\Key{density} & \ApiText{grid times points}{Gitter mal Punkte} & --- & \ApiText{Linear in cells times points; keep that product at or below 20000.}{Linear in Zellen mal Punkte; halten Sie dieses Produkt bei höchstens 20000.}}
  \SemioTableRow{\Key{projection} & \ApiText{4-vertex ring}{Ring mit 4 Ecken} & \ApiText{instant}{sofort} & \ApiText{Resampling at the default precision expands it to about six vertices at world scale. A 200-vertex coastline at \Key{precision=0} is the practical ceiling per figure; a full 10 degree graticule at d3's default precision is 7876 numbers and exceeds the 30 s per-adapter test budget.}{Das Nachtasten mit der Vorgabegenauigkeit erweitert ihn im Weltmaßstab auf etwa sechs Ecken. Eine Küstenlinie mit 200 Ecken bei \Key{precision=0} ist die praktische Obergrenze je Abbildung; ein volles 10-Grad-Gradnetz mit d3s Vorgabegenauigkeit sind 7876 Zahlen und sprengt das Prüfbudget von 30 s je Adapter.}}
}

\section{\ApiText{Traps worth knowing}{Fallen, die man kennen sollte}}

\SemioTableLong[text-size=8pt]{\ApiText{expl3 and l3fp traps}{Fallen in expl3 und l3fp}}{0.30,0.70}{\ApiText{Trap}{Falle} & \ApiMeaning}{%
  \SemioTableRow{\ApiText{No digits in expl3 variable names}{Keine Ziffern in expl3-Variablennamen} & \ApiText{A digit ends a control sequence name, so a variable written \Key{x0} silently becomes \Key{x} followed by a typeset zero. The failure surfaces far from its cause. The library uses \Key{xlo}, \Key{xhi}, \Key{ylo}, \Key{yhi} instead.}{Eine Ziffer beendet einen Kontrollsequenznamen, eine Variable \Key{x0} wird also stillschweigend zu \Key{x}, gefolgt von einer gesetzten Null. Der Fehler zeigt sich weit von seiner Ursache. Die Bibliothek nutzt stattdessen \Key{xlo}, \Key{xhi}, \Key{ylo}, \Key{yhi}.}}
  \SemioTableRow{\ApiText{The power operator binds tighter than unary minus}{Der Potenzoperator bindet stärker als das unäre Minus} & \ApiText{An expression that squares a possibly negative value must parenthesise it.}{Ein Ausdruck, der einen möglicherweise negativen Wert quadriert, muss ihn klammern.}}
  \SemioTableRow{\ApiText{The conditional operator evaluates both branches}{Der Bedingungsoperator wertet beide Zweige aus} & \ApiText{An index guard written as a conditional does not guard anything; the out-of-range branch is evaluated too.}{Eine als Bedingung geschriebene Indexabsicherung sichert nichts ab; auch der Zweig außerhalb des Bereichs wird ausgewertet.}}
  \SemioTableRow{\ApiText{Colons and percent signs are catcode-sensitive}{Doppelpunkte und Prozentzeichen hängen am Catcode} & \ApiText{A colon is a letter inside expl3 and punctuation in a document, so any key value split on punctuation must be stringified first. The same applies to a percent sign inside a format specifier.}{Ein Doppelpunkt ist in expl3 ein Buchstabe und im Dokument ein Satzzeichen, jeder an einem Satzzeichen zerlegte Schlüsselwert muss also zuerst in einen String verwandelt werden. Dasselbe gilt für ein Prozentzeichen in einer Formatangabe.}}
  \SemioTableRow{\ApiText{Determinism}{Determinismus} & \ApiText{Every randomised layout is seeded and every figure redraws identically; \Key{seed} is an option, not an accident.}{Jedes zufallsgestützte Layout hat einen Startwert und jede Abbildung wird identisch neu gezeichnet; \Key{seed} ist eine Option, kein Zufall.}}
}

\end{document}
